\documentclass[a4paper,reqno,11pt]{amsart}
\usepackage{amsmath,amssymb,mathtools,enumitem,microtype,tikz-cd,booktabs,tabularx,placeins}
\usepackage[a4paper,margin=30mm]{geometry}
\usepackage[colorlinks=true,linkcolor=blue,citecolor=blue,urlcolor=blue]{hyperref}

\setcounter{tocdepth}{2}
\usetikzlibrary{arrows.meta,positioning}
\makeatletter
\renewcommand{\subsection}{\@startsection{subsection}{2}%
  {\z@}{1.5ex plus .5ex minus .2ex}{.7ex plus .2ex}{\normalfont\bfseries}}
\def\l@subsection{\@tocline{2}{0pt}{2pc}{5pc}{}}
\makeatother
\renewcommand{\thepart}{\Roman{part}}

\newtheorem{theorem}{Theorem}[section]
\newtheorem{proposition}[theorem]{Proposition}
\newtheorem{lemma}[theorem]{Lemma}
\newtheorem{corollary}[theorem]{Corollary}
\theoremstyle{definition}
\newtheorem{definition}[theorem]{Definition}
\newtheorem{construction}[theorem]{Construction}
\newtheorem{setting}[theorem]{Setting}
\newtheorem{example}[theorem]{Example}
\newtheorem{remark}[theorem]{Remark}
\newtheorem*{theoremA}{Theorem A}
\newtheorem*{theoremB}{Theorem B}
\newtheorem*{aiuse}{Use of AI}

\DeclareMathOperator{\modu}{\mathsf{mod}}
\DeclareMathOperator{\ind}{\mathsf{ind}}
\DeclareMathOperator{\proj}{\mathsf{proj}}
\DeclareMathOperator{\add}{\mathsf{add}}
\DeclareMathOperator{\rad}{rad}
\DeclareMathOperator{\soc}{soc}
\DeclareMathOperator{\End}{End}
\DeclareMathOperator{\Hom}{Hom}
\DeclareMathOperator{\Ext}{Ext}
\DeclareMathOperator{\Tor}{Tor}
\DeclareMathOperator{\HH}{HH}
\DeclareMathOperator{\gr}{gr}
\DeclareMathOperator{\Irr}{Irr}
\DeclareMathOperator{\im}{im}
\DeclareMathOperator{\ann}{ann}
\DeclareMathOperator{\id}{id}
\DeclareMathOperator{\Aut}{Aut}
\DeclareMathOperator{\charac}{char}
\DeclareMathOperator{\Span}{span}
\newcommand{\op}{\mathrm{op}}
\newcommand{\JJ}{\mathcal J}
\newcommand{\C}{\mathcal C}
\newcommand{\D}{\mathcal D}
\newcommand{\B}{\mathcal B}
\newcommand{\ray}{\operatorname{ray}}
\newcommand{\Z}{\mathbb Z}
\newcommand{\N}{\mathbb N}
\newcommand{\F}{\mathbb F}
\newcommand{\sixlifttree}{%
\begin{tikzcd}[ampersand replacement=\&,row sep=large,column sep=large]
v_1\arrow[r,"\gamma"]\&v_2\arrow[r,"\alpha"]\&v_3\\
v_4\arrow[r,"\alpha"]\&v_5\arrow[u,"\lambda"]\&v_6\arrow[l,"\beta"']
\end{tikzcd}}

\title{Multiplicative bases and deformations of representation-finite algebras}
\author[H.\ Enomoto]{Haruhisa Enomoto}
\address{Parakeet Inc., Japan}
\email{haruhisa.enomoto.math@gmail.com}
\subjclass[2020]{16G10, 16G20, 16G60}
\keywords{multiplicative basis, representation-finite algebra,
Auslander--Reiten theory, mesh category, deformation}
\date{October 16, 2026}

\begin{document}

\begin{abstract}
Bautista, Gabriel, Roiter, and Salmer\'on proved that every
finite-dimensional representation-finite algebra over an algebraically
closed field admits a multiplicative basis and, outside characteristic two, is
standard, meaning that its category of indecomposable modules is presented by
the Auslander--Reiten quiver with mesh relations. We give a new,
self-contained proof of both theorems by viewing the algebra as a filtered
deformation of its standard form, which we construct from the radical layers
of the module category and equip with a multiplicative basis. Outside
characteristic two, this deformation is trivial, which yields both theorems
at once. In characteristic two, we embed the algebra into an associative
family of deformations of the standard form obtained by modifying the
multiplication table only on pairs of basis elements whose product in the
standard form is zero. After rescaling, the original algebra is isomorphic to
a member of this family in which each modified product is either zero or a
single basis element, so the multiplicative basis persists.
\end{abstract}

\maketitle
\tableofcontents
\clearpage

\section{Introduction}
\label{sec:introduction}

A finite-dimensional algebra is \emph{representation-finite} if it has only
finitely many isomorphism classes of indecomposable finite-dimensional
modules. Although this condition is module-theoretic, it imposes strong
combinatorial restrictions on the multiplication of the algebra itself: over
an algebraically closed field, every representation-finite algebra admits a
\emph{multiplicative basis}, a basis in which the product of any two elements
is either zero or another basis element (Theorem~A below). Consequently, the
entire algebra multiplication is encoded by a finite combinatorial table
without scalar structure constants.

A companion structural theorem governs the category of indecomposable
modules. A basic representation-finite algebra is called \emph{standard} if
one can choose representatives of the irreducible morphisms so that every
morphism between indecomposable modules is a linear combination of paths in
these representatives, and all relations among such paths follow from the
quadratic mesh relations induced by almost split sequences. Outside
characteristic two, every basic representation-finite algebra over an
algebraically closed field is standard (Theorem~B below). This structure
theory was established in the landmark work of Bautista, Gabriel, Roiter,
and Salmer\'on \cite{BGRS}, abbreviated BGRS below, building on the
construction of the standard form by Bretscher and Gabriel
\cite{BretscherGabriel}.

Because of the intricacy of the original combinatorial case analysis, finding
a more transparent proof has remained a longstanding problem. Bongartz
described the BGRS paper as ``beautiful but hard to read''
\cite[Section~1]{BongartzSurvey} and surveyed subsequent simplifications and
extensions; more recently, Ringel remarked that ``any publication which
provides help, provides improvement would be strongly appreciated''
\cite[Section~3]{RingelVossieck}.

In this paper we give an alternative, self-contained proof of both classical
theorems. The central idea is to pass first to the radical grading of the
module category to construct the standard form of the algebra, and then to
analyze the original algebra as a filtered deformation of its standard form.
All arguments required beyond basic Auslander--Reiten theory and elementary
homological algebra are developed in full.

\subsection{The classical results revisited}
\label{subsec:intro-results}

We begin with the existence theorem for multiplicative bases compatible with
primitive idempotents and powers of the Jacobson radical.

\begin{definition}
\label{def:normed-multiplicative-basis}
A basis $\B$ of a finite-dimensional algebra $A$ is \emph{multiplicative} if
$bc\in\B\cup\{0\}$ for all $b,c\in\B$. It is \emph{normed} if it contains
a complete set of primitive orthogonal idempotents and
$\B\cap(\rad A)^n$ is a basis of $(\rad A)^n$ for every $n\geq1$.
\end{definition}

Bautista, Gabriel, Roiter, and Salmer\'on proved that every
representation-finite algebra over an algebraically closed field admits a
normed multiplicative basis \cite[Sections~1.2 and~1.12]{BGRS}.

\begin{theoremA}[= Theorem~\ref{thm:final-basis}]
Every finite-dimensional representation-finite algebra over an algebraically
closed field has a normed multiplicative basis.
\end{theoremA}

Using matrix units to account for the multiplicities of indecomposable
projective summands, Lemma~\ref{lem:nonbasic-basis} reduces Theorem~A
immediately to basic algebras. We therefore focus on basic algebras in what
follows.

Let $A$ be a basic finite-dimensional representation-finite algebra over an
algebraically closed field $k$, let $\ind(\modu A)$ be a skeletal category
of indecomposable finite-dimensional right $A$-modules, and let $\Gamma_A$
be its Auslander--Reiten quiver. Writing compositions of morphisms in path
order from left to right, the components of an almost split sequence
\[
 0\longrightarrow\tau X\xrightarrow{(a_i)_i}
 \bigoplus_iY_i\xrightarrow{(b_i)_i}X\longrightarrow0
\]
satisfy $\sum_i a_i b_i=0$. Let $\JJ$ denote the categorical Jacobson radical
of $\ind(\modu A)$. Choosing a basis of each space $\JJ(X,Y)/\JJ^2(X,Y)$ of
irreducible morphisms—which is at most one-dimensional when $A$ is
representation-finite (Lemma~\ref{lem:square-free})—and expressing the
components $a_i$ and $b_i$ modulo $\JJ^2$ in these bases yields a quadratic
mesh relation at each nonprojective vertex $X$. The \emph{mesh category}
$k(\Gamma_A)$ has the vertices of $\Gamma_A$ as objects and paths in
$\Gamma_A$ modulo these quadratic mesh relations as morphisms. Up to
isomorphism fixing the objects, $k(\Gamma_A)$ is independent of the chosen
bases, and its mesh coefficients can all be rescaled to one
(Lemma~\ref{lem:unit-mesh-coefficients}).

\begin{definition}
\label{def:standardness}
A basic finite-dimensional representation-finite algebra $A$ over an
algebraically closed field $k$ is \emph{standard} if there is a $k$-linear equivalence
\[
 k(\Gamma_A)\simeq\ind(\modu A)
\]
sending each vertex of $\Gamma_A$ to the corresponding indecomposable module.
\end{definition}

\begin{theoremB}[= Corollary~\ref{cor:away-from-two-basis}]
Every basic finite-dimensional representation-finite algebra over an
algebraically closed field of characteristic different from two is standard.
\end{theoremB}

In characteristic two, standardness can fail
\cite[Section~3.7]{BongartzSurvey}, even though Theorem~A remains valid.
To treat all characteristics uniformly, we construct in
Section~\ref{sec:graded-category} a basic representation-finite algebra
$A_0$ whose category of indecomposable modules is equivalent to
$k(\Gamma_A)$: this is the \emph{standard form} of $A$, introduced
classically by Bretscher and Gabriel \cite{BretscherGabriel}. Specifically,
the powers $\JJ^n$ of the categorical radical filter the morphism spaces
$\Hom_A(P_y,P_x)$ between indecomposable projective modules by allowing
factorizations through arbitrary indecomposable modules, not merely
projective ones. This categorical filtration can be strictly finer than the
ordinary radical filtration of $A$, and its associated graded algebra is
$A_0$. Thus $A$ is naturally a positive filtered deformation of $A_0$, and
$A$ is standard if and only if $A\cong A_0$
(Corollary~\ref{cor:graded-standard-form}). Outside characteristic two, we
show in Section~\ref{sec:mesh-deformations} that any higher-degree error terms
in the mesh relations can be eliminated degree by degree by adjusting the
chosen irreducible morphisms; hence the filtered deformation is trivial and
Theorem~B follows.

\subsection{The construction and its relation to earlier proofs}
\label{subsec:intro-comparison}

Bretscher and Gabriel constructed the standard form as the opposite
endomorphism algebra of the projective vertices inside the mesh category
\cite[Section~3.1]{BretscherGabriel}. We recover this description directly
by showing that the associated graded category $\gr_\JJ\ind(\modu A)$ is
isomorphic to $k(\Gamma_A)$ and that evaluation at the projective objects
gives an equivalence $\add k(\Gamma_A)\simeq\modu A_0$
(Propositions~\ref{prop:graded-realization}
and~\ref{prop:projective-evaluation}; compare Iyama's graded realization
theorems for $\tau$-categories \cite[Theorems~5.2 and~9.2]{IyamaI}). In
Section~\ref{sec:scalar-normalization}, we construct a normed multiplicative
basis of $A_0$ by lifting the mesh category to the universal cover of the
Bongartz--Gabriel mesh complex \cite[Section~4]{BongartzGabriel}, extending
the basis one projective vertex at a time as in
\cite[Section~2.6]{BretscherGabriel}, and descending along the free deck
group.

In characteristic two, $A$ need not be isomorphic to $A_0$. Instead, we keep
the multiplicative basis of $A_0$ and study deformations of $A_0$ obtained
by modifying the multiplication table only on pairs of basis elements whose
product in $A_0$ is zero, setting each such modified product equal to a
scalar multiple of another basis element. Using a Bockstein connecting map
from coefficients modulo four, we show that cocycles of this form span the
positive-degree Hochschild cohomology spaces governing filtered deformations
of $A_0$ (Corollary~\ref{cor:zero-supported-coverage}). The local
factorization restrictions developed in
Section~\ref{sec:local-ray-geometry} then imply that these zero-supported
cocycles fit together into a global associative family of algebras over a
polynomial ring (Proposition~\ref{prop:zero-product-family}). By finite
filtered normalization and diagonal rescaling along the weight lattice
(Theorem~\ref{thm:finite-family-normalization} and
Corollary~\ref{cor:binary-normalization}), $A$ is isomorphic to a fibre of
this family whose parameters all take values in $\{0,1\}$. In any such
binary fibre, the basis of $A_0$ remains a normed multiplicative basis.

Bautista, Gabriel, Roiter, and Salmer\'on likewise obtained the
nonstandard characteristic-two algebras by replacing specified zero products
in the standard multiplication table with nonzero basis elements
\cite[Sections~9.2--9.7]{BGRS}. In our approach, the admissible
modifications are organized cohomologically by Hochschild two-cocycles and
their underlying derivation-lifting equations. Along the way, we establish
two general deformation-theoretic tools of independent interest: a module-lifting
theorem for positive filtered deformations based on internal-degree bounds
for $\Ext$ groups (Lemma~\ref{lem:positive-ext-vanishing} and
Proposition~\ref{prop:positive-module-lifting}), and an elementary
construction of associative multi-parameter families from compatible
multiplication tables (Lemma~\ref{lem:two-associative-tables}).
Table~\ref{tab:method-comparison} summarizes the correspondence between the
classical steps and the present proof.

\begin{table}[htbp]
\centering
\small
\renewcommand{\arraystretch}{1.15}
\begin{tabularx}{\textwidth}{@{}>{\raggedright\arraybackslash}p{.19\textwidth}
 >{\raggedright\arraybackslash}X >{\raggedright\arraybackslash}X@{}}
\toprule
Task & Classical constructions & Present organization \\
\midrule
Construct the standard form &
Take the endomorphism algebra of the projective vertices of the mesh
category \cite{BretscherGabriel}; alternatively, build an algebra from the
multiplication table of basis elements up to scalars \cite{BGRS}. &
Take the projective corner of the radical grading of $\ind(\modu A)$,
then normalize the scalar structure constants on a universal cover. \\
\addlinespace
Prove standardness ($\charac k\ne2$) &
Compare paths with equal values \cite{BGRS}, or correct the mesh
relations on stable components step by step \cite{Riedtmann1984}. &
Eliminate higher-degree error terms in the mesh relations degree by degree
by modifying the chosen irreducible morphisms. \\
\addlinespace
Treat characteristic two &
Replace specified zero products by nonzero basis elements in explicit normal
forms \cite{BGRS}. &
Represent positive-degree first-order deformations by zero-supported
cocycles via $\Z/4$-lifts of derivations, and realize $A$ as a binary fibre
of the resulting associative family. \\
\addlinespace
Restrict local configurations &
Exhibit representation-infinite subcategories and analyze contours locally
\cite{BGRS}; see also the shorter mild-contour proofs in
\cite{BongartzMild}. &
Bound trace spaces on $\tau$-periodic components of $\Gamma_A$, and restrict
factorizations of basis elements in $A_0$ via bimodule retractions onto
representation-infinite incidence and bound-quiver algebras. \\
\bottomrule
\end{tabularx}
\caption{Comparison between the classical proofs and the present approach.}
\label{tab:method-comparison}
\end{table}

The combinatorial case analysis that remains is concentrated in two places:
Appendices~\ref{app:periodic-mesh} and~\ref{app:residual-mesh} analyze the
linearized mesh derivative on $\tau$-periodic components of $\Gamma_A$,
while Section~\ref{sec:local-ray-geometry} and
Appendix~\ref{app:local-rays} restrict how basis elements of $A_0$ can
factor through one another. In the periodic mesh calculations we bound spaces
of paths modulo cyclic rotation—recovering special cases of Schedler's
description of zeroth Hochschild homology of preprojective algebras
\cite[Theorem~13.1.1]{Schedler}—as well as twisted rotation under tree
automorphisms. The degree-by-degree correction of irreducible morphisms
follows Riedtmann's method for stable blocks \cite{Riedtmann1984}.

\subsection{Reading routes}
\label{subsec:reading-routes}

Readers interested only in the standard form or in the theorems outside
characteristic two need only the first part of the paper. The table below
indicates the sections required for each main result.

\begin{center}
\small
\renewcommand{\arraystretch}{1.15}
\begin{tabularx}{\textwidth}{@{}>{\raggedright\arraybackslash}p{.26\textwidth}
 >{\raggedright\arraybackslash}X@{}}
\toprule
Reading goal & Route and conclusion \\
\midrule
Construct the standard form &
Sections~\ref{sec:graded-category}--\ref{sec:scalar-normalization},
culminating in Corollary~\ref{cor:graded-standard-form}. \\
\addlinespace
Prove Theorems~A and~B outside characteristic two &
Continue through Section~\ref{sec:mesh-deformations} to
Corollary~\ref{cor:away-from-two-basis} (with Lemma~\ref{lem:nonbasic-basis}
covering nonbasic algebras). Appendix~\ref{app:periodic-mesh} proves
Propositions~\ref{prop:periodic-tree-model}
and~\ref{prop:mesh-correction-spaces}, whose statements suffice for the main
text. \\
\addlinespace
Prove Theorem~A in characteristic two &
After Sections~\ref{sec:graded-category}--\ref{sec:scalar-normalization},
read Subsections~\ref{subsec:mesh-derivative}
and~\ref{subsec:brick-hypothesis} together with
Sections~\ref{sec:positive-deformations}--\ref{sec:binary-families}, ending
at Theorems~\ref{thm:final-basis} and~\ref{thm:final-standardness}.
Section~\ref{sec:deformation-refinements} establishes further properties of
module categories and polynomial parameter spaces. \\
\bottomrule
\end{tabularx}
\end{center}

\medskip
\phantomsection\label{subsec:intro-organization}
\noindent\textbf{Organization.}
Part~\ref{part:graded-form}
(Sections~\ref{sec:graded-category}--\ref{sec:mesh-deformations})
constructs the standard form $A_0$, equips it with a normed multiplicative
basis, and completes the proof of Theorems~A and~B outside characteristic
two. Part~\ref{part:mesh-deformations}
(Sections~\ref{sec:positive-deformations}--\ref{sec:weights-and-parameters})
turns to characteristic two: it relates positive filtered deformations of
$A_0$ to the residual mesh cokernels on periodic components of $\Gamma_A$
and reduces Theorem~A to constructing zero-supported Hochschild cocycles
(Proposition~\ref{prop:family-criterion}).
Part~\ref{part:characteristic-two}
(Sections~\ref{sec:local-ray-geometry}--\ref{sec:deformation-refinements})
constructs those cocycles by lifting derivation equations from $\F_2$ to
$\Z$ and $\Z/4$ via induction on the number of basis elements, builds the
resulting associative family of deformations to complete the proof of
Theorem~A, and develops additional consequences for module categories and
mesh parameters in Section~\ref{sec:deformation-refinements}.
Appendix~\ref{app:periodic-mesh} contains the periodic mesh calculations
used in Part~\ref{part:graded-form}, while
Appendices~\ref{app:residual-mesh} and~\ref{app:local-rays} contain the
characteristic-two mesh calculations and local ray exclusions.

\medskip
\noindent\textbf{Conventions and notation.}
Throughout the paper, $k$ denotes an algebraically closed field unless
explicitly stated otherwise. All algebras and modules are finite-dimensional,
and modules are right modules. We assume familiarity with standard
Auslander--Reiten theory, Morita equivalence \cite{ARS}, and basic
homological algebra. We write $D_k=\Hom_k(-,k)$ for standard vector-space
duality and $\tau X$ for the starting term of the almost split sequence
ending at a nonprojective indecomposable module $X$.

For a basic algebra $A$ with a chosen complete set of primitive orthogonal
idempotents $\{e_x\}$, we set $A(x,y)=e_xAe_y$ and $P_x=e_xA$. Products of
paths and morphisms are written in path order from left to right, whereas
$\circ$ always denotes ordinary function composition: thus for morphisms
$f:X\to Y$ and $g:Y\to Z$, their path-order product is $fg=g\circ f$. An
arrow $a:x\to y$ of the quiver of $A$ lies in $A(x,y)=e_xAe_y$. Under the
canonical identification $A(x,y)\cong\Hom_A(P_y,P_x)$ sending $a$ to left
multiplication by $a$, endpoints are reversed, so the algebra product $ab$ of
$a\in A(x,y)$ and $b\in A(y,z)$ corresponds to the composite
$a\circ b:P_z\to P_x$. Because endomorphism algebras are formed using the
path-order product $fg=g\circ f$, we have $A\cong\End_A(A)^{\op}$ and,
similarly, $A_0=\End_{\C}(\bigoplus_x P_x)^{\op}$ in
Construction~\ref{con:projective-algebra}.

For an additive category $\C$, we write $\add\C$ for its closure under
finite direct sums and direct summands. Any $k$-linear category $C$ with
finitely many objects and finite-dimensional morphism spaces is identified
with its category algebra $\bigoplus_{i,j} C(i,j)$, where the product of
$f\in C(i,j)$ and $g\in C(j,l)$ is $fg=g\circ f$ and composites of
noncomposable morphisms are zero; thus $C(i,j)$ is the $(i,j)$-corner of the
category algebra. We write $\JJ$ for the categorical Jacobson radical of a
category of indecomposable modules and $\rad A$ for the Jacobson radical of
an algebra $A$.

\begin{aiuse}
The proof was developed by GPT-6-Astra under the author's direction.
The manuscript was written with GPT-6-Astra, followed by Claude Opus~5.5
and Gemini~4 Argon. GPT-6-Astra and Claude Opus~5.5 were used for
formalization. The author directed the research and revisions and is
responsible for the mathematical content and the final manuscript.
\end{aiuse}

\clearpage
\part{The standard form and the proof outside characteristic two}
\label{part:graded-form}

In this part we construct the standard form $A_0$ of a basic
representation-finite algebra $A$ in arbitrary characteristic, equip $A_0$
with a normed multiplicative basis, and prove Theorems~A and~B when
$\charac k\ne2$.

\section{Reduction to basic algebras and graded reconstruction}
\label{sec:graded-category}

We begin by reducing the existence of a normed multiplicative basis to the
case of basic algebras. We then pass from the category of indecomposable
modules to its radical grading and reconstruct the standard form $A_0$ from
the graded projective objects. Both constructions, together with the
normalization of $A_0$ in the next two sections, apply in arbitrary
characteristic.

\subsection{Multiplicative bases and matrix units}
\label{subsec:nonbasic-basis}

Using matrix units to account for the multiplicities of indecomposable
projective summands, a normed multiplicative basis of a basic algebra lifts
immediately to any Morita-equivalent finite-dimensional algebra.

\begin{lemma}
\label{lem:nonbasic-basis}
Let $A$ be a finite-dimensional algebra over an arbitrary field, and let
$B_0$ be its basic algebra. If $B_0$ has a normed multiplicative basis, then
so does $A$.
\end{lemma}
\begin{proof}
Let $\{e_i\}$ be the complete set of primitive orthogonal idempotents
contained in the given normed multiplicative basis of $B_0$, and let $n_i$
be the multiplicity of the corresponding indecomposable projective module in
the right regular module $A_A$. Every element $b$ of the basis of $B_0$ lies
in a unique corner $e_i B_0 e_j$: indeed, in the decomposition
$b=\sum_{i,j} e_i b e_j$, each nonzero summand belongs to the basis by
multiplicativity, and since basis elements in distinct corners are linearly
independent, only one summand can be nonzero.

Under Morita equivalence, the regular module $A_A$ corresponds to the
$B_0$-module $\bigoplus_i (e_i B_0)^{n_i}$. Taking the opposite of its
endomorphism algebra under our left-to-right composition convention
identifies $A$ with the block matrix algebra whose $(i,j)$-block is
$\operatorname{Mat}_{n_i\times n_j}(e_i B_0 e_j)$. For each basis element
$b\in e_i B_0 e_j$, let $E_{i\alpha,j\beta}(b)$ ($1\leq\alpha\leq n_i$,
$1\leq\beta\leq n_j$) denote the matrix having $b$ in entry $(\alpha,\beta)$
of the $(i,j)$-block and zero elsewhere. These matrices form a basis of $A$
and multiply according to the rule
\[
 E_{i\alpha,j\beta}(b)E_{l\gamma,m\delta}(c)
 =\begin{cases}
 E_{i\alpha,m\delta}(bc),&\text{if }(j,\beta)=(l,\gamma),\\
 0,&\text{otherwise}.
 \end{cases}
\]
Since the given basis of $B_0$ is multiplicative, each product is either
zero or another basis element. Moreover, the diagonal elements
$E_{i\alpha,i\alpha}(e_i)$ form a complete set of primitive orthogonal
idempotents of $A$, and for each $n\geq1$ the $(i,j)$-block of $(\rad A)^n$
is $\operatorname{Mat}_{n_i\times n_j}(e_i(\rad B_0)^n e_j)$. Thus the
constructed basis intersects each radical power $(\rad A)^n$ in a basis.
\end{proof}

From now on we may therefore assume that $A$ is basic, invoking
Lemma~\ref{lem:nonbasic-basis} at the conclusion of each characteristic
branch to recover the general case.

\subsection{Radical layers and mesh relations}

We filter the morphisms in $\ind(\modu A)$ by powers of the categorical
radical, that is, by sums of composites of nonisomorphisms between
indecomposable modules. Because intermediate modules in such a composite need
not be projective, the induced filtration on morphisms between indecomposable
projective modules is in general finer than the ordinary radical filtration
of $A$.

Let $A$ be a basic representation-finite algebra over $k$, and let $\JJ$
denote the categorical Jacobson radical of $\ind(\modu A)$. Since $A$ is
representation-finite, the direct sum of a representative set of
indecomposable modules is a finite-dimensional additive generator; its
endomorphism algebra is finite-dimensional, and $\JJ$ consists of the matrix
entries of its Jacobson radical, so $\JJ$ is nilpotent. We define the
\emph{radical grading} of $\ind(\modu A)$ to be the graded category $\C$
with the same objects as $\ind(\modu A)$ and morphism spaces
\[
 \C(X,Y)=\bigoplus_{n\geq0}\JJ^n(X,Y)/\JJ^{n+1}(X,Y),
\]
writing $\C_n(X,Y)=\JJ^n(X,Y)/\JJ^{n+1}(X,Y)$ for the homogeneous component of
degree $n$. For a morphism $f\in \JJ^n(X,Y)\setminus \JJ^{n+1}(X,Y)$, its
\emph{initial form} is its residue class in $\C_n(X,Y)$. Composition in $\C$
is induced by composition in $\ind(\modu A)$. Since the endomorphism ring of
each indecomposable module is local with residue field $k$, we have
$\C_0(X,X)=k\cdot 1_X$ and $\C_0(X,Y)=0$ for $X\ne Y$.

Let us recall how the Auslander--Reiten quiver $\Gamma_A$ and its mesh
relations arise in this setting. The vertices of $\Gamma_A$ are the objects
of $\ind(\modu A)$, and the arrows $X\to Y$ index a chosen basis of the
space $\Irr(X,Y)=\JJ(X,Y)/\JJ^2(X,Y)$ of irreducible morphisms. For each
nonprojective indecomposable module $Y$, choose an almost split sequence
\[
 0\longrightarrow\tau Y\longrightarrow E_Y
   \longrightarrow Y\longrightarrow0.
\]
Each indecomposable direct summand $X$ of $E_Y$ contributes a component map
$X\to Y$ and a component map $\tau Y\to X$; modulo $\JJ^2$, the former
components form a basis of $\Irr(X,Y)$ and the latter a basis of
$\Irr(\tau Y,X)$ indexed by the multiplicity of $X$ in $E_Y$. The next
well-known lemma ensures that each multiplicity is at most one, so that these
components differ from the chosen arrows only by nonzero scalars.

\begin{lemma}
\label{lem:square-free}
Let $A$ be a basic representation-finite algebra over an algebraically
closed field. Its Auslander--Reiten quiver $\Gamma_A$ has no loops and at
most one arrow between each ordered pair of vertices.
\end{lemma}
\begin{proof}
Consider any irreducible morphism $f:X\to Y$ between indecomposable modules
and factor $f$ through its image. By irreducibility, either the canonical
surjection $X\twoheadrightarrow\im f$ is a split monomorphism or the
inclusion $\im f\hookrightarrow Y$ is a split epimorphism; thus one of the
two factors is an isomorphism, and $f$ is either a proper monomorphism or a
proper epimorphism. In particular, $\dim X\ne\dim Y$, so $\Gamma_A$ has no
loops.

Suppose for a contradiction that there are at least two arrows $X\to Y$ in
$\Gamma_A$. If $\dim X>\dim Y$, any irreducible morphism $X\to Y$ is an
epimorphism, so $Y$ cannot be projective (otherwise the epimorphism would
split). The middle term of the almost split sequence ending at $Y$ then
contains $X^{\oplus 2}$ as a direct summand, which gives
\[
 \dim\tau Y\geq2\dim X-\dim Y>\dim X,
\]
and the first map of the sequence provides at least two arrows
$\tau Y\to X$. Iterating this step produces indecomposable modules of
strictly increasing dimensions, contradicting the assumption that $A$ is
representation-finite. If $\dim X<\dim Y$, the dual argument applied to the
almost split sequence starting at $X$ yields the same contradiction.
\end{proof}

Thanks to Lemma~\ref{lem:square-free}, we can formulate the mesh category of
$\Gamma_A$ using scalar coefficients rather than matrices.

\begin{definition}
\label{def:translation-quiver}
A \emph{translation quiver} is a quiver $Q$ equipped with a bijection $\tau$
from a subset of vertices, called \emph{nonprojective}, onto a subset of
vertices, called \emph{noninjective}, together with a bijection $\sigma$
from the set of arrows ending at nonprojective vertices onto the set of
arrows starting at noninjective vertices, such that $\sigma a:\tau x\to y$
for every arrow $a:y\to x$ with $x$ nonprojective. We call $\sigma a$ the
\emph{mate} of $a$. Vertices outside the domain (respectively, image) of
$\tau$ are called \emph{projective} (respectively, \emph{injective}). When
both $x$ and $y$ are nonprojective, we write
$\tau a=\sigma^2a:\tau y\to\tau x$. The translation quiver $Q$ is
\emph{stable} if every vertex is both nonprojective and noninjective.
\end{definition}

\begin{definition}
\label{def:scalar-mesh-category}
Let $Q$ be a translation quiver with at most one arrow between each ordered
pair of vertices. Given a family $c=(c_{x,a})$ of nonzero scalars
$c_{x,a}\in k^*$ indexed by nonprojective vertices $x$ and arrows $a:y\to x$
ending at $x$, the associated \emph{scalar mesh category} is the path-length
graded quotient category
\[
 C(Q,c)=kQ/\langle r_x:x\text{ nonprojective}\rangle,
 \qquad r_x=\sum_{a:y\to x}c_{x,a}(\sigma a)a,
\]
where $kQ$ denotes the $k$-linear path category of $Q$.
\end{definition}

Multiplying a mesh relation $r_x$ by a nonzero scalar leaves the relation
ideal unchanged, while rescaling the arrows of $Q$ rescales the coefficients
$c_{x,a}$ and induces a graded isomorphism between the corresponding scalar
mesh categories.

For $Q=\Gamma_A$, choose a basis element in each one-dimensional space
$\Irr(X,Y)$ to serve as the unique arrow $a:X\to Y$. For each nonprojective
indecomposable module $Y$, the components of the first and second maps of
the chosen almost split sequence ending at $Y$ take the form
$u_a(\sigma a)$ and $v_a a$ modulo $\JJ^2$, with $u_a,v_a\in k^*$. Since the
composite of the two maps in the almost split sequence is zero, passing to
degree two in $\C$ yields the mesh relation
\begin{equation}
 r_Y=\sum_{a:X\to Y}c_{Y,a}(\sigma a)a,
 \qquad c_{Y,a}=u_av_a.
 \label{eq:scalar-mesh}
\end{equation}
We write $k(\Gamma_A)$ for the scalar mesh category $C(\Gamma_A,c)$ defined
by these coefficients.

Up to isomorphism fixing the objects, $k(\Gamma_A)$ is independent of both
the chosen arrow representatives and the chosen almost split sequences.
Indeed, any isomorphism between two almost split sequences ending at $Y$ acts
by scalars $\xi,\eta\in k^*$ on the indecomposable end terms $\tau Y$ and
$Y$. Because the indecomposable summands of the middle term are pairwise
nonisomorphic by Lemma~\ref{lem:square-free}, the automorphism of the middle
term acts by scalars $\lambda_a\in k^*$ on those summands modulo radical
maps, whose composites with the maps of the sequence vanish modulo $\JJ^2$.
The coefficients therefore transform according to
\[
 u'_a=\lambda_a u_a\xi^{-1},\qquad
 v'_a=\eta v_a\lambda_a^{-1},\qquad
 c'_{Y,a}=\eta\xi^{-1}c_{Y,a},
\]
which simply rescales the entire relation $r_Y$ by $\eta\xi^{-1}$. Moreover,
Lemma~\ref{lem:unit-mesh-coefficients} will show that the arrows can always
be rescaled so that $c_{Y,a}=1$ for all $Y$ and $a$.

In Sections~\ref{sec:scalar-normalization} and~\ref{sec:mesh-deformations}
we will restrict a scalar mesh category $C(Q,c)$ to a subset $S$ of vertices
that is closed under $\tau$ and $\tau^{-1}$ wherever defined (such as the
set of $\tau$-periodic vertices). Let $Q_S$ be the full translation
subquiver on $S$. Quotienting $C(Q,c)$ by the ideal generated by the
identity morphisms $1_x$ for $x\notin S$ and restricting to the objects in
$S$ gives a canonical isomorphism
\begin{equation}
 \left(C(Q,c)/\langle 1_x:x\notin S\rangle\right)\big|_S
       \cong C(Q_S,c|_S).
 \label{eq:mesh-object-quotient}
\end{equation}
Indeed, in the path category $kQ$, killing $1_x$ for $x\notin S$ annihilates
precisely those paths that pass through a vertex outside $S$. Under this
quotient, a mesh relation $r_x$ with $x\notin S$ maps to zero, while a mesh
relation $r_x$ with $x\in S$ retains precisely the terms $(\sigma a)a$ whose
middle vertex lies in $S$, which is the mesh relation of $Q_S$ at $x$. In
particular, if $C(Q,c)$ is finite-dimensional, then so is $C(Q_S,c|_S)$.

\subsection{Exactness on radical layers}

We next show that almost split sequences remain exact after passing to the
associated graded category $\C$. This exactness is the key to identifying
$\C$ with the mesh category $k(\Gamma_A)$ in
Proposition~\ref{prop:graded-realization}(3) and to reconstructing the module
category $\modu A_0$ in Proposition~\ref{prop:projective-evaluation}. Our
argument adapts Iyama's complementary-monomorphism technique for
$\tau$-categories \cite[Sections~3.3, 3.6 and~4.2]{IyamaI} directly to module
categories.

\begin{lemma}
\label{lem:strict-radical-layers}
Let $A$ be an Artin algebra, let $\JJ$ be the categorical radical of its
category of finite-length modules, and let $a:U\to Y$ be a radical
monomorphism. Suppose that $a+h$ is a monomorphism for every
$h\in \JJ^2(U,Y)$. Then for every finite-length module $W$ and every
$n\geq1$,
\[
             a\circ z\in \JJ^n(W,Y)\quad\Longleftrightarrow\quad
             z\in \JJ^{n-1}(W,U).
\]
Furthermore, any finite direct sum of left almost split monomorphisms
satisfies the hypothesis that all such $\JJ^2$-perturbations remain
monomorphisms.
\end{lemma}
\begin{proof}
\emph{Step 1: Almost split kernels and the perturbation property.}
For any finite direct sum $Y$ of indecomposable modules, let $p:E_Y\to Y$ be
the direct sum of the minimal right almost split morphisms into its
indecomposable summands (using the radical inclusion $\rad P\hookrightarrow P$
at projective summands). By the lifting property of right almost split maps,
every radical morphism into $Y$ factors through $p$. For $n\geq0$, writing a
morphism in $\JJ^{n+1}(W,Y)$ as a sum of composites of $n+1$ radical morphisms
and factoring the last radical morphism through $p$ shows that
\begin{equation}
 p\circ \JJ^n(W,E_Y)=\JJ^{n+1}(W,Y)\qquad(n\geq0).
 \label{eq:right-radical-surjection}
\end{equation}
At each nonprojective indecomposable summand of $Y$, the kernel of $p$ is
the left almost split morphism of the corresponding almost split sequence,
while at each projective summand the inclusion $\rad P\hookrightarrow P$ has
zero kernel. Consequently, the kernel inclusion $i:T\to E_Y$ of $p$ is a
finite direct sum of left almost split monomorphisms.

More generally, let $i=\bigoplus_j i_j:T\to E$ be any finite direct sum of
left almost split monomorphisms $i_j:T_j\to E_j$. Any morphism in
$\JJ^2(T,Z)$ is a sum of composites $v\circ u$ with $u$ and $v$ radical. Since
the restriction of $u$ to each $T_j$ is radical, it factors through $i_j$ as
$u_j\circ i_j$, and $v\circ u_j$ remains radical. Thus
\[
 \JJ^2(T,Z)=\{q\circ i:q\in \JJ(E,Z)\}.
\]
Taking $Z=E$, any $h\in \JJ^2(T,E)$ has the form $h=q\circ i$ with
$q\in \JJ(E,E)$. Since $1+q$ is an automorphism of $E$, the perturbation
$i+h=(1+q)\circ i$ is again a monomorphism.

\emph{Step 2: Splitting off the given monomorphism.}
Returning to the given monomorphism $a:U\to Y$, since $a$ is radical and
$p:E_Y\to Y$ is right almost split, we can write $a=p\circ u$ for some
$u:U\to E_Y$. By performing row and column operations to split off any
invertible components of $u$ between indecomposable summands, we obtain
decompositions $U=U''\oplus U'$ and $E_Y=U''\oplus E'$ under which
$u=1_{U''}\oplus u'$ with $u'\in \JJ(U',E')$. Since $p$ is radical, the
restriction $a|_{U'}=p|_{E'}\circ u'$ lies in $\JJ^2(U',Y)$. Setting
$h=-a|_{U'}\circ\operatorname{pr}_{U'}\in \JJ^2(U,Y)$, we find that $a+h$
vanishes on $U'$. Since $a+h$ is a monomorphism by hypothesis, we must have
$U'=0$, so $u$ is a split monomorphism. Replacing $p$ and $i$ by
$p\circ\alpha^{-1}$ and $\alpha\circ i$ for a suitable automorphism $\alpha$
of $E_Y$ preserves \eqref{eq:right-radical-surjection} and the factorization
property of $i$, and allows us to write
\[
 E_Y=V\oplus U,\qquad p=(f,a),\qquad i=\binom b c.
\]

\emph{Step 3: The complementary monomorphism.}
We claim that the first component $b:T\to V$ of $i$ is also a radical
monomorphism satisfying the $\JJ^2$-perturbation hypothesis. First, $b$ and $c$
are radical because $i$ is radical. If $x\in\ker b$, then $i(x)=(0,c(x))$
and $0=p(i(x))=a(c(x))$; since $a$ is injective, $c(x)=0$ and hence
$i(x)=0$, so $x=0$. Next, given any $h\in \JJ^2(T,V)$, Step~1 lets us write
$h=q\circ i$ for a radical map $q=(q_V,q_U):V\oplus U\to V$. The upper
triangular matrix
\[
 L=\begin{pmatrix}1+q_V&q_U\\0&1\end{pmatrix}
\]
is an automorphism of $V\oplus U$, and the kernel inclusion of
$p\circ L^{-1}$ is $L\circ i=\binom{b+h}{c}$. On the other hand, the
restriction of $p\circ L^{-1}$ to $U$ is
$a-f\circ(1+q_V)^{-1}\circ q_U$, which differs from $a$ by a morphism in
$\JJ^2(U,Y)$ and is therefore a monomorphism by hypothesis. As just shown for
$p$ and $b$, the injectivity of the restriction of $p\circ L^{-1}$ to $U$
implies that the first component $b+h$ of its kernel inclusion is a
monomorphism.

\emph{Step 4: Induction on the radical degree.}
We now prove by induction on $n\geq1$, simultaneously for all radical
monomorphisms $a$ satisfying the perturbation hypothesis, that
$a\circ z\in \JJ^n(W,Y)$ if and only if $z\in \JJ^{n-1}(W,U)$. Since $a$ is
radical, the ``if'' direction is immediate, as is the base case $n=1$. For
 the induction step, suppose $a\circ z\in \JJ^{n+1}(W,Y)$ with $n\geq1$. By
\eqref{eq:right-radical-surjection}, there exist $v\in \JJ^n(W,V)$ and
$u\in \JJ^n(W,U)$ such that $a\circ z=p\circ\binom{v}{u}=f\circ v+a\circ u$.
Thus $\binom{v}{u-z}$ lies in the kernel of $p$, so
$\binom{v}{u-z}=\binom{b}{c}\circ s$ for some $s:W\to T$. In particular,
$b\circ s=v\in \JJ^n(W,V)$. Applying the induction hypothesis to the
complementary monomorphism $b$ yields $s\in \JJ^{n-1}(W,T)$. Since $c$ is
radical, it follows that $u-z=c\circ s\in \JJ^n(W,U)$, and hence
$z\in \JJ^n(W,U)$, completing the induction.
\end{proof}

Applying Lemma~\ref{lem:strict-radical-layers} and its dual to the maps in
each almost split sequence yields the graded realization theorem below
(compare Iyama \cite[Theorems~5.2 and~9.2]{IyamaI} in the setting of
$\tau$-categories).

\begin{proposition}
\label{prop:graded-realization}
Let $\C=\gr_\JJ\ind(\modu A)$. Then the following statements hold:
\begin{enumerate}[label={\upshape(\arabic*)}]
\item Let $0\to\tau Y\to E_Y\to Y\to0$ be an almost split sequence in
$\modu A$, and let $\tau Y\xrightarrow{\bar\imath}E_Y\xrightarrow{\bar p}Y$
be the sequence of initial forms of its morphisms in $\add\C$. We call
$\bar p$ the \emph{right mesh map} at $Y$. For every object $W$ of $\C$, the
induced sequences
\[
 0\to\C(W,\tau Y)\to\C(W,E_Y)\to\C(W,Y),\qquad
 0\to\C(Y,W)\to\C(E_Y,W)\to\C(\tau Y,W)
\]
are exact; in particular, $\bar\imath$ is a monomorphism and $\bar p$ is an
epimorphism in $\add\C$.
\item For each indecomposable projective module $P$, the initial forms of
the components of the inclusion $\rad P\hookrightarrow P$ define a
monomorphism in $\add\C$. Dually, for each indecomposable injective module
$I$, the initial forms of the components of $I\twoheadrightarrow I/\soc I$
define an epimorphism in $\add\C$.
\item The chosen arrow representatives in $\C_1$ induce an isomorphism of
graded categories $k(\Gamma_A)\xrightarrow{\sim}\C$, where $k(\Gamma_A)$ has
the scalar mesh relations \eqref{eq:scalar-mesh}.
\end{enumerate}
\end{proposition}
\begin{proof}
Let $i:\tau Y\to E_Y$ and $p:E_Y\to Y$ be the maps in the almost split
sequence ending at $Y$. By Lemma~\ref{lem:strict-radical-layers}, $i$ induces
an injection $\C_{n-1}(W,\tau Y)\hookrightarrow\C_n(W,E_Y)$ for every
$n\geq1$. Moreover, if $w\in \JJ^n(W,E_Y)$ ($n\geq1$) satisfies
$p\circ w\in \JJ^{n+2}(W,Y)$, then \eqref{eq:right-radical-surjection} yields
some $w'\in \JJ^{n+1}(W,E_Y)$ with $p\circ w'=p\circ w$. Since
$p\circ(w-w')=0$, we have $w-w'=i\circ z$ for some $z:W\to\tau Y$, and
$i\circ z\in \JJ^n(W,E_Y)$ implies $z\in \JJ^{n-1}(W,\tau Y)$ by
Lemma~\ref{lem:strict-radical-layers}. In degree $n=0$, any
$w\in\Hom_A(W,E_Y)$ with $p\circ w\in \JJ^2(W,Y)$ differs by an element of
$\JJ(W,E_Y)$ from a morphism killed by $p$, which factors through the radical
map $i$ and therefore also lies in $\JJ(W,E_Y)$; thus $\bar p$ is injective on
$\C_0(W,E_Y)$. Since $p\circ i=0$, this proves exactness of the first
sequence in (1).

For an indecomposable projective module $P$, let $p:\rad P\hookrightarrow P$
be the inclusion. If $z\in \JJ^n(W,\rad P)$ ($n\geq0$) satisfies
$p\circ z\in \JJ^{n+2}(W,P)$, then \eqref{eq:right-radical-surjection} gives
$z'\in \JJ^{n+1}(W,\rad P)$ with $p\circ z'=p\circ z$, and injectivity of $p$
forces $z=z'\in \JJ^{n+1}(W,\rad P)$. Thus the initial form of $p$ is a
monomorphism in $\add\C$. Applying the dual arguments in the opposite module
category establishes exactness of the second sequence in (1) and the
statement for injective modules in (2).

Finally, sending each arrow of $\Gamma_A$ to its chosen class in $\C_1$
defines a graded functor $k\Gamma_A\to\C$. By \eqref{eq:scalar-mesh}, the
degree-two initial form of the relation $p\circ i=0$ is precisely the mesh
relation $r_Y$, so this functor descends to $k(\Gamma_A)\to\C$. It is full
because each radical power $\JJ^n$ is generated by products of $n$ morphisms
in $\JJ$. To see that it is faithful, we show by induction on $n\geq0$ that any
homogeneous combination $F$ of paths of length $n$ from $W$ to $Y$ whose
image vanishes in $\C_n(W,Y)$ belongs to the mesh ideal; the cases $n=0,1$
hold by definition of $\Gamma_A$. For $n\geq2$, group $F$ by its final arrow
as $F=\sum_a F_a a$, where $a:X_a\to Y$. If $Y$ is nonprojective, the
components of $\bar p:E_Y\to Y$ are $v_a a$, so the tuple
$(F_a v_a^{-1})_a$ maps to zero under $\bar p:\C_{n-1}(W,E_Y)\to\C_n(W,Y)$.
By the exactness proved in (1), there is a path combination $H$ of length
$n-2$ from $W$ to $\tau Y$ whose image under $\bar\imath$ equals the image of
$(F_a v_a^{-1})_a$ in $\C_{n-1}(W,E_Y)$. Since the $a$-component of
$\bar\imath$ is $u_a(\sigma a)$ and $u_a v_a=c_{Y,a}$, each difference
$F_a - H c_{Y,a}(\sigma a)$ has length $n-1$ and maps to zero in $\C$. By
the induction hypothesis, each such difference lies in the mesh ideal, and
hence so does
\[
 F = H r_Y + \sum_a \bigl(F_a - H c_{Y,a}(\sigma a)\bigr)a.
\]
If $Y$ is projective, then $\bar p$ is a monomorphism by (2), so each $F_a$
maps to zero in $\C_{n-1}(W,X_a)$ and lies in the mesh ideal by induction.
\end{proof}

We also record the cokernels of the right mesh maps in $\C$, which provide
minimal projective resolutions for the simple functors on $\C$.

\begin{lemma}
\label{lem:mesh-simple-tops}
For every nonprojective object $X$ of $\C$, the image of the right mesh map
$\bar p:\C(-,E_X)\to\C(-,X)$ is the radical $\bigoplus_{n\geq1}\C_n(-,X)$ of
the representable functor $\C(-,X)$, and its cokernel is the simple functor
$\C_0(-,X)$. The same conclusion holds at each projective object $P$ for the
map $\C(-,\rad P)\to\C(-,P)$ of Proposition~\ref{prop:graded-realization}(2).
\end{lemma}
\begin{proof}
Let $p:E_X\to X$ be the minimal right almost split map in $\modu A$ (or the
inclusion $\rad X\hookrightarrow X$ when $X$ is projective). By
\eqref{eq:right-radical-surjection}, $p\circ \JJ^n(W,E_X)=\JJ^{n+1}(W,X)$ for
every $n\geq0$ and every object $W$. Passing to the associated graded spaces
shows that the image of $\bar p$ on $\C(W,E_X)$ is precisely the
positive-degree ideal $\bigoplus_{n\geq1}\C_n(W,X)$. Since this ideal is
nilpotent and the degree-zero quotient has $\C_0(X,X)=k$ and $\C_0(W,X)=0$
for $W\ne X$, the positive-degree part is the radical of $\C(-,X)$ and its
cokernel is the simple functor $\C_0(-,X)$.
\end{proof}

\subsection{Evaluation on projective objects}

Following Bretscher and Gabriel \cite[Section~3.1]{BretscherGabriel}, we now
define the standard form $A_0$ as the opposite endomorphism algebra of the
projective objects in $\C$, and we show using the graded exactness properties
above that evaluating objects of $\add\C$ on the projectives gives an
equivalence $\add\C\simeq\modu A_0$.

\begin{construction}
\label{con:projective-algebra}
Let $\D$ be the full subcategory of $\C$ on the objects $P_x=e_x A$
corresponding to the indecomposable projective $A$-modules. We define the
graded algebra $A_0=\bigoplus_{x,y}A_0(x,y)$ by
\[
 A_0(x,y)=\D(P_y,P_x)
       =\bigoplus_{n\geq0}\JJ^n(P_y,P_x)/\JJ^{n+1}(P_y,P_x),
\]
with multiplication $bc=b\circ c\in\D(P_z,P_x)$ for $b\in\D(P_y,P_x)$ and
$c\in\D(P_z,P_y)$ (so that $c$ is applied first, then $b$). Equivalently,
under our path-order convention for endomorphism algebras,
$A_0=\End_{\C}(\bigoplus_x P_x)^{\op}$. The \emph{projective evaluation
functor} sends each object $X\in\add\C$ to the right $A_0$-module
$\bigoplus_x\C(P_x,X)$, on which $b\in A_0(x,y)$ acts on $m\in\C(P_x,X)$ by
$mb=m\circ b\in\C(P_y,X)$.
\end{construction}

To prove that projective evaluation is an equivalence of categories, we use a
standard homological realization criterion from the Auslander correspondence
(compare \cite[Theorem~3.3]{IyamaRealization}).

\begin{lemma}
\label{lem:projective-evaluation-criterion}
Let $\Lambda$ be a finite-dimensional algebra over a field, let $e\in\Lambda$
be an idempotent, and set $\Lambda_e=e\Lambda e$. Suppose that
\[
 \Hom_\Lambda(T,\Lambda)=\Ext^1_\Lambda(T,\Lambda)=0
 \qquad\text{for every simple $\Lambda$-module $T$ with }Te=0.
\]
Then the functor $P\mapsto Pe$ is fully faithful on $\proj\Lambda$. If, in
addition, $\operatorname{gl.dim}\Lambda\leq2$ and every finite-dimensional
injective $\Lambda_e$-module is isomorphic to $Qe$ for some
$Q\in\proj\Lambda$, then $P\mapsto Pe$ is an equivalence of categories
$\proj\Lambda\xrightarrow{\sim}\modu\Lambda_e$, and in particular
$\Lambda_e$ is representation-finite.
\end{lemma}
\begin{proof}
By induction on composition length using the long exact sequence for $\Ext$,
the assumed vanishings hold for every finite-dimensional $\Lambda$-module $T$
satisfying $Te=0$. For any right $\Lambda$-module $M$, the canonical
multiplication map $(Me)\otimes_{\Lambda_e}e\Lambda\to M$ becomes an
isomorphism upon multiplying on the right by $e$, so its kernel and cokernel
are annihilated by $e$. For any $Q\in\proj\Lambda$, applying
$\Hom_\Lambda(-,Q)$ and using the tensor--Hom adjunction therefore yields
canonical isomorphisms
\[
 \Hom_\Lambda(M,Q)\xrightarrow{\sim}
 \Hom_\Lambda((Me)\otimes_{\Lambda_e}e\Lambda,Q)
 \xrightarrow{\sim}\Hom_{\Lambda_e}(Me,Qe)
\]
whose composite sends a homomorphism $f:M\to Q$ to its restriction
$f|_{Me}:Me\to Qe$. Taking $M\in\proj\Lambda$ shows that $P\mapsto Pe$ is
fully faithful on $\proj\Lambda$.

Now suppose $\operatorname{gl.dim}\Lambda\leq2$ and every injective
$\Lambda_e$-module lies in the essential image of $\proj\Lambda$. Given any
$X\in\modu\Lambda_e$, choose the beginning of an injective coresolution
$0\to X\to I^0\to I^1$ in $\modu\Lambda_e$. By full faithfulness and the
assumption on injectives, we may lift $I^0\to I^1$ to a homomorphism
$f:Q^0\to Q^1$ in $\proj\Lambda$. Because $\operatorname{gl.dim}\Lambda\leq2$,
the kernel $\ker f$ is a second syzygy over $\Lambda$ and is therefore
projective. Since the functor $M\mapsto Me$ is exact, we obtain
$X\cong(\ker f)e$ with $\ker f\in\proj\Lambda$, proving essential
surjectivity. Finally, since $\proj\Lambda$ has only finitely many
indecomposable objects up to isomorphism, the equivalence implies that
$\Lambda_e$ is representation-finite.
\end{proof}

We now verify the hypotheses of Lemma~\ref{lem:projective-evaluation-criterion}
for $\add\C$ and $A_0$.

\begin{proposition}
\label{prop:projective-evaluation}
Let $A$ be a basic representation-finite $k$-algebra, and set
$\C=\gr_\JJ\ind(\modu A)$ and $A_0$ as in
Construction~\ref{con:projective-algebra}.
\begin{enumerate}[label={\upshape(\arabic*)}]
\item The projective evaluation functor induces an equivalence of categories
$\add\C\xrightarrow{\sim}\modu A_0$.
\item The algebra $A_0$ is basic and representation-finite.
\item If $X_0\in\modu A_0$ denotes the evaluation of an object $X\in\C$, then
\[
                         \dim_k X_0e_x=\dim_k Xe_x
\]
for every primitive idempotent $e_x$.
\end{enumerate}
\end{proposition}
\begin{proof}
Because the radical filtration on $\Hom_A(X,Y)$ is finite, passing to the
associated graded vector space preserves total dimension:
\begin{equation}
 \dim_k\C(X,Y)=\sum_{n\geq0}\dim_k \JJ^n(X,Y)/\JJ^{n+1}(X,Y)
             =\dim_k\Hom_A(X,Y).
 \label{eq:graded-hom-dimensions}
\end{equation}
In particular, $\dim_k X_0 e_x=\dim_k\C(P_x,X)=\dim_k\Hom_A(P_x,X)=\dim_k Xe_x$,
which proves (3).

To prove (1) and (2), let $G=\bigoplus_{X\in\operatorname{ob}\C}X$ and set
$\Lambda=\End_{\add\C}(G)^{\op}$. The Yoneda functor $X\mapsto R_X=\C(G,X)$
is an equivalence $\add\C\xrightarrow{\sim}\proj\Lambda$, and as $X$ ranges
over $\operatorname{ob}\C$, the modules $R_X$ form a complete set of
pairwise nonisomorphic indecomposable projective $\Lambda$-modules with
simple tops $L_X=\C_0(G,X)$. Let $e\in\Lambda$ be the sum of the idempotents
corresponding to the projective objects $P_x$. Then $e\Lambda e=A_0$, and
under the Yoneda equivalence, projective evaluation on $\add\C$ coincides
with the functor $\operatorname{ev}(M)=Me$ on $\proj\Lambda$.

\emph{Global dimension and full faithfulness.}
For each nonprojective object $X\in\C$, Proposition~\ref{prop:graded-realization}(1)
and Lemma~\ref{lem:mesh-simple-tops} yield a projective resolution in
$\modu\Lambda$:
\[
 0\longrightarrow R_{\tau X}\longrightarrow
 \bigoplus_{Y\to X}R_Y\longrightarrow R_X\longrightarrow L_X
 \longrightarrow0.
\]
For each projective object $P_x$, the same results yield a projective
resolution $0\to R_{\rad P_x}\to R_{P_x}\to L_{P_x}\to0$ of length at most
one. Consequently, every simple $\Lambda$-module has projective dimension at
most two, so $\operatorname{gl.dim}\Lambda\leq2$. Moreover, when $X$ is
nonprojective, applying $\Hom_\Lambda(-,\Lambda)$ to the resolution of $L_X$
gives by Yoneda's lemma the sequence
$0\to\C(X,G)\to\C(E_X,G)\to\C(\tau X,G)$, which is exact by
Proposition~\ref{prop:graded-realization}(1). Thus
$\Hom_\Lambda(L_X,\Lambda)=\Ext^1_\Lambda(L_X,\Lambda)=0$ for every simple
module $L_X$ annihilated by $e$, and Lemma~\ref{lem:projective-evaluation-criterion}
shows that $\operatorname{ev}$ is fully faithful on $\proj\Lambda$.

\emph{Essential surjectivity.}
Let $S_x=P_x/\rad P_x$ and $I_x=D_k(Ae_x)$ denote the simple and
indecomposable injective $A$-modules at vertex $x$. By
\eqref{eq:graded-hom-dimensions}, the evaluation $S_x^0=\operatorname{ev}(R_{S_x})$
satisfies $\dim_k S_x^0 e_y=\dim_k\Hom_A(P_y,S_x)=\delta_{xy}$, so $S_x^0$ is
the simple $A_0$-module at $x$. Set $N_x=\operatorname{ev}(R_{I_x})$. Since
$\operatorname{ev}$ is fully faithful on $\proj\Lambda$, we have
\[
 \dim_k\Hom_{A_0}(S_y^0,N_x)=\dim_k\C(S_y,I_x)=\dim_k\Hom_A(S_y,I_x)=\delta_{xy},
\]
which shows that $\soc N_x\cong S_x^0$. On the other hand,
\eqref{eq:graded-hom-dimensions} gives
\[
 \begin{aligned}
 \dim_k N_xe_y&=\dim_k\C(P_y,I_x)=\dim_k I_xe_y=\dim_k e_yAe_x\\
 &=\dim_k\C(P_x,P_y)=\dim_k e_yA_0e_x=\dim_k D_k(A_0 e_x)e_y.
 \end{aligned}
\]
Since $N_x$ has simple socle $S_x^0$, it embeds into the injective hull
$D_k(A_0 e_x)$, and equality of dimensions implies $N_x\cong D_k(A_0 e_x)$.
Thus every indecomposable injective $A_0$-module lies in the image of
$\operatorname{ev}$. Lemma~\ref{lem:projective-evaluation-criterion} now
implies that $\add\C\xrightarrow{\sim}\modu A_0$ is an equivalence. Since
the objects $P_x$ are pairwise nonisomorphic in $\C$, the algebra $A_0$ is
basic as well as representation-finite.
\end{proof}

\begin{corollary}
\label{cor:same-ar-quiver}
The Auslander--Reiten quivers of $A_0$ and $A$ are isomorphic as translation
quivers: $\Gamma_{A_0}\cong\Gamma_A$.
\end{corollary}
\begin{proof}
By Proposition~\ref{prop:projective-evaluation}(1), we may compute
$\Gamma_{A_0}$ directly inside $\add\C$. Since $\C$ is generated in degree
one, the radical of $\add\C$ is its positive-degree ideal and the radical
square is the ideal of degree at least two; hence the arrows of $\Gamma_A$
form bases of the spaces of irreducible morphisms in $\add\C$. By
Lemma~\ref{lem:mesh-simple-tops}, each right mesh map $\bar p:E_X\to X$ is
right almost split in $\add\C$. Because the indecomposable summands of $E_X$
are pairwise nonisomorphic by Lemma~\ref{lem:square-free}, any endomorphism
of $E_X$ is diagonal modulo the radical, and if it commutes with $\bar p$,
the linear independence of the arrow components modulo the radical square
forces its diagonal entries to be $1$. Thus every such endomorphism is an
automorphism, which means that $\bar p$ is right minimal. The mesh sequences
in $\add\C$ are therefore the almost split sequences of $\modu A_0$, so the
translation of $\Gamma_{A_0}$ agrees with that of $\Gamma_A$.
\end{proof}

We conclude this section with a simple example illustrating how the
categorical filtration $\JJ^n(P_y,P_x)$ used to define $A_0$ can shift the
degrees of generators of $\rad A$.

\begin{example}
\label{ex:dual-number-filtration}
Let $A=k[t]/(t^2)$ be the algebra of dual numbers. Up to isomorphism, its
indecomposable modules are $P=A$ and $S=A/(t)$, connected by the canonical
projection $\pi:P\twoheadrightarrow S$ and the inclusion $\iota:S\hookrightarrow P$
sending $1+(t)$ to $t$. In path order, the multiplication-by-$t$
endomorphism of $P$ factors as
\[
 P\xrightarrow{\pi}S\xrightarrow{\iota}P,
 \qquad \pi\iota=t,\qquad \iota\pi=0.
\]
Consequently, $\JJ$ is spanned by $\pi$, $\iota$, and $t$, while $\JJ^2$ is
spanned by $t$ and $\JJ^3=0$. On the projective endomorphism ring $A=\End_A(P)^{\op}$,
the categorical radical filtration is therefore
\[
 A\supset \JJ(P,P)=kt=\JJ^2(P,P)\supset \JJ^3(P,P)=0,
\]
whereas $(\rad A)^2=0$. The associated graded algebra $A_0$ is again
isomorphic to $k[t]/(t^2)$, but its radical generator $t$ sits in
categorical degree two, reflecting its factorization through the
nonprojective simple module $S$.
\end{example}

\section{Rays of the algebra and of its standard form}
\label{sec:projective-rays}

To compare the original algebra $A$ with its standard form $A_0$, we study
the orbits of elements of $A(x,y)$ under the left and right actions of the
unit groups of the local endomorphism rings $R_x=e_xAe_x$ and $R_y=e_yAe_y$.
These orbits are the \emph{rays} of $A$ (Definition~\ref{def:rays}). In this
section we show that the nonzero rays of $A$ correspond bijectively via
initial forms to the nonzero homogeneous components $A_0(x,y)_d$, each of
which is one-dimensional, and that the homogeneous lines of $A_0$ form an
associative category with zero under ray multiplication.

\subsection{Endpoint subbimodules}

For each primitive idempotent $e_x$ of $A$, set $R_x=e_xAe_x$. The classical
subbimodule structure of $A(x,y)$ over $R_x$ and $R_y$ follows from the
finiteness of the ideal lattice of a representation-finite algebra (see Jans
\cite[Corollary~1.3 and Theorem~2.1]{Jans} and Bongartz
\cite[Section~2.3]{BongartzSurvey}).

\begin{lemma}
\label{lem:endpoint-chains}
Let $A$ be a basic representation-finite algebra over $k$.
\begin{enumerate}[label={\upshape(\arabic*)}]
\item The algebra $A$ has only finitely many two-sided ideals, and every
idempotent corner $eAe$ is representation-finite (these two properties hold
without assuming $A$ basic).
\item For each primitive idempotent $e_x$, there is an isomorphism
$R_x\cong k[t]/(t^{N_x})$ for some $N_x\geq1$.
\item For each pair of vertices $x,y$, the $R_x$--$R_y$-subbimodules of
$A(x,y)$ form a finite chain with one-dimensional successive quotients. In
particular, for every nonzero element $f\in A(x,y)$, either
$R_xfR_y=R_xf$ or $R_xfR_y=fR_y$.
\end{enumerate}
\end{lemma}
\begin{proof}
Each two-sided ideal $I$ of $A$ is the annihilator of the cyclic right module
$A/I$. Because $A$ has only finitely many isomorphism classes of modules of
dimension at most $\dim_k A$, it has only finitely many two-sided ideals.
For any idempotent $e\in A$ and $C=eAe$, the induction functor
$V\mapsto V\otimes_C eA$ is fully faithful because the unit of its
adjunction with the restriction functor $(-)e$ is the canonical isomorphism
\[
 V\xrightarrow{\sim}(V\otimes_C eA)e\cong V\otimes_C eAe.
\]
Since induction preserves finite dimension, indecomposability, and
isomorphism classes, $C$ is also representation-finite. This proves (1).

For a primitive idempotent $e=e_x$, every ideal $I$ of the local algebra
$C=eAe$ equals $e(AIA)e$, so $C$ has only finitely many ideals. Every linear
subspace of $\rad C/(\rad C)^2$ is the image of an ideal containing
$(\rad C)^2$; since $k$ is infinite, a space of dimension at least two would
contain infinitely many lines and thus produce infinitely many ideals of $C$.
Hence $\dim_k\rad C/(\rad C)^2\leq1$. By Nakayama's lemma, any lift
$t\in\rad C$ of a spanning element of $\rad C/(\rad C)^2$ generates all
powers of $\rad C$. Since $\rad C$ is nilpotent, $t$ generates $C$ as a
$k$-algebra, giving $C\cong k[t]/(t^N)$ for some $N\geq1$, which proves (2).

Similarly, any $R_x$--$R_y$-subbimodule $N\subseteq A(x,y)$ satisfies
$e_x(ANA)e_y=N$, so $A(x,y)$ has only finitely many
$R_x$--$R_y$-subbimodules. We may view $A(x,y)$ as a module over the
commutative local algebra $T=R_x\otimes_k R_y^{\op}$, which has residue
field $k$. If $A(x,y)$ contained two incomparable $T$-submodules, their sum
modulo their intersection would contain a direct sum of two nonzero
submodules and hence a two-dimensional semisimple submodule $k^{\oplus 2}$
in its socle; pulling back the infinitely many lines in $k^{\oplus 2}$ would
yield infinitely many $T$-submodules of $A(x,y)$, a contradiction. Thus the
$R_x$--$R_y$-subbimodules of $A(x,y)$ form a chain, and its successive
quotients are simple $T$-modules, hence one-dimensional.

Finally, for any nonzero $f\in A(x,y)$, the cyclic $T$-submodule
$Tf=R_xfR_y$ is isomorphic to $T/\ann_T(f)$. Because the ideals of the
quotient ring $T/\ann_T(f)$ also form a chain, its radical modulo its
radical square has dimension at most one and is therefore spanned by the
image of the radical generator of either $R_x$ or $R_y$. By Nakayama's
lemma, the corresponding endpoint algebra surjects onto $T/\ann_T(f)$, which
means that $R_xfR_y=R_xf$ or $R_xfR_y=fR_y$.
\end{proof}

\begin{definition}[{\cite[Section~1]{BGRS}}]
\label{def:rays}
For a nonzero element $f\in A(x,y)$, the \emph{ray} of $f$ is the orbit
\[
 [f]=\{ufv:u\in R_x^*,\ v\in R_y^*\}.
\]
We set $U(f)=R_xfR_y$ and
\[
 R(f)=(\rad R_x)fR_y+R_xf(\rad R_y).
\]
In each corner $A(x,y)$ we also include a zero ray $0$. Since $U(f)$ depends
only on the ray $r=[f]$, we often write $U_r=U(f)$.
\end{definition}

By Lemma~\ref{lem:endpoint-chains}(3), the $R_x$--$R_y$-subbimodule $U(f)$ is
uniserial with unique maximal proper subbimodule $R(f)$, and $U(f)$ is cyclic
over either $R_x$ or $R_y$. By Nakayama's lemma, the generators of $U(f)$ over
that endpoint algebra are precisely the elements of $U(f)\setminus R(f)$, and
any two such generators differ by a unit of that endpoint algebra. Thus
\begin{equation}
                        [f]=U(f)\setminus R(f).
\label{eq:ray-layer}
\end{equation}

\begin{corollary}
\label{cor:ray-transversal}
Let $A$ be a basic representation-finite algebra over $k$. Choose one
representative from each nonzero ray of $A$, taking $e_x$ as the
representative of $[e_x]$. Then these representatives form a basis of $A$
whose intersection with every two-sided ideal $I$ of $A$ is a basis of $I$.
\end{corollary}
\begin{proof}
In each corner $A(x,y)$, the chosen representatives project to bases of the
one-dimensional successive quotients $U(f)/R(f)$ of the finite chain of
$R_x$--$R_y$-subbimodules. They therefore form a basis of $A(x,y)$ adapted
to every subbimodule in the chain. Since $I\cap A(x,y)=e_xIe_y$ is an
$R_x$--$R_y$-subbimodule, it belongs to this chain and is spanned by the
representatives it contains.
\end{proof}

To define multiplication of rays independently of the choice of
representatives, one must quotient out contributions coming from the radical
of the middle vertex.

\begin{definition}
\label{def:ray-product}
For composable nonzero elements $f\in A(x,y)$ and $h\in A(y,z)$, define the
subbimodule
\[
                  M(f,h)=R_x f(\rad R_y)hR_z.
\]
The \emph{ray product} of $[f]$ and $[h]$ is defined by
\begin{equation}
 [f][h]=
 \begin{cases}
 [fh],&\text{if }fh\notin M(f,h),\\
 0,&\text{if }fh\in M(f,h),
 \end{cases}
 \label{eq:ray-product}
\end{equation}
and any product involving a zero ray is defined to be zero.
\end{definition}

As the following example illustrates, a product $fh$ of specific
representatives can be nonzero in $A$ even when its ray product $[f][h]$ is
zero because $fh\in M(f,h)$.

\begin{example}
\label{ex:zero-ray-span}
Let $B$ be the bound quiver algebra given by
\[
\begin{tikzcd}
b\arrow[r,"\alpha"]&c\arrow[r,"\beta"]\arrow[loop above,"q"]&a
\end{tikzcd}
\qquad q^2=0,\quad\alpha\beta=0.
\]
This is a representation-finite string algebra with basis
\[
 e_b,e_c,e_a,\quad\alpha,q,\beta,\quad\alpha q,q\beta,
       \quad u=\alpha q\beta.
\]
Each basis vector spans a distinct ray. Here
$U_{\alpha}=\Span_k\{\alpha,\alpha q\}$ and
$U_{\beta}=\Span_k\{\beta,q\beta\}$, while $M(\alpha,\beta)=ku$. Notice that
$\alpha$ and $\alpha(1+q)$ represent the same ray $[\alpha]$, yet
$\alpha\beta=0$ whereas $\alpha(1+q)\beta=u\ne0$. Because both products lie
in $M(\alpha,\beta)=ku$, Definition~\ref{def:ray-product} assigns
$[\alpha][\beta]=0$.
\end{example}

\begin{lemma}
\label{lem:ray-product-independent}
The ray product \eqref{eq:ray-product} is independent of the choice of
representatives. In particular, if $[f][h]\ne0$, then $f'h'\ne0$ for all
$f'\in[f]$ and $h'\in[h]$.
\end{lemma}
\begin{proof}
Every unit of $R_y$ has the form $\lambda(1+n)$ with $\lambda\in k^*$ and
$n\in\rad R_y$. Changing the representatives $f$ and $h$ within their rays
therefore replaces $fh$, up to action by outer units in $R_x^*\times R_z^*$,
by $fh+fnh$ with $fnh\in M(f,h)$, while the subbimodule $M(f,h)$ itself is
unchanged. If $fh\in M(f,h)$, then $fh+fnh\in M(f,h)$ as well. If
$fh\notin M(f,h)$, then $U(fh)\not\subseteq M(f,h)$; since the subbimodules
of $A(x,z)$ form a chain and $R(fh)$ is the unique maximal proper
subbimodule of $U(fh)$, it follows that $M(f,h)\subseteq R(fh)$. Hence
$fh+fnh\in U(fh)\setminus R(fh)=[fh]$, so the perturbed product is nonzero
and defines the same ray $[fh]$.
\end{proof}

We call the set of rays of $A$ equipped with the operation
\eqref{eq:ray-product} the \emph{ray table} $\ray A$. At this stage we do
not yet claim that $\ray A$ is associative for an arbitrary $A$; we first
establish associativity for the graded standard form $A_0$, and later
transfer it to $A$ in Proposition~\ref{prop:dimension-induction}(2).

\subsection{Homogeneous rays of the standard form}

Consider the categorical radical filtration $F^n A(x,y)=\JJ^n(P_y,P_x)$ used
in Construction~\ref{con:projective-algebra} to define $A_0=\gr_F A$. Since
each subspace $F^n A(x,y)$ is an $R_x$--$R_y$-subbimodule, multiplication by
units in $R_x^*\times R_y^*$ preserves the filtration. Every nonzero ray $r$
of $A$ therefore has a well-defined categorical degree $\deg_\JJ(r)$, the
unique integer $d$ such that $r\subset F^d A\setminus F^{d+1}A$ (compare
\cite[Section~1.1]{BretscherGabriel}).

\begin{lemma}
\label{lem:ray-degree}
Let $A$ be a basic representation-finite algebra over $k$. Each homogeneous
component $A_0(x,y)_d$ has dimension at most one, and passing to initial
forms induces a bijection between the nonzero rays of $A$ in $A(x,y)$ and
the nonzero homogeneous spaces $A_0(x,y)_d$. Moreover, if two rays $r,r'$ in
$A(x,y)$ satisfy $U_{r'}\subsetneq U_r$, then $\deg_\JJ(r')>\deg_\JJ(r)$.
\end{lemma}
\begin{proof}
Let $A(x,y)=U_0\supsetneq U_1\supsetneq\cdots\supsetneq U_l=0$ be the
complete chain of $R_x$--$R_y$-subbimodules of $A(x,y)$. Its successive
quotients $U_i/U_{i+1}$ are one-dimensional, and
$U_{i+1}=(\rad R_x)U_i+U_i(\rad R_y)$. Since every $F^d A(x,y)$ is an
$R_x$--$R_y$-subbimodule, it must equal one of the $U_i$. Because
$\rad R_x\subseteq \JJ(P_x,P_x)$ and $\rad R_y\subseteq \JJ(P_y,P_y)$, we have
\[
 F^d A(x,y)=U_i\quad\Longrightarrow\quad
 U_{i+1}\subseteq F^{d+1}A(x,y)\subseteq U_i.
\]
As $d$ increases from $0$, the filtration starts at $U_0=A(x,y)$, ends at
$0$, and at each step either remains at $U_i$ or drops by a single step to
$U_{i+1}$. Consequently, each one-dimensional quotient $U_i/U_{i+1}$ appears
as $A_0(x,y)_d=F^d A(x,y)/F^{d+1}A(x,y)$ for a unique degree $d$. Since the
nonzero rays of $A(x,y)$ are precisely the layers $U_i\setminus U_{i+1}$,
all claims follow.
\end{proof}

We now examine the intrinsic rays of the representation-finite graded
algebra $A_0$ itself.

\begin{proposition}
\label{prop:projective-ray-product}
Let $P=\ray A_0$ be the ray table of $A_0$. Then:
\begin{enumerate}[label={\upshape(\arabic*)}]
\item each nonzero ray of $A_0$ contains the nonzero vectors of a unique
homogeneous line $A_0(x,y)_d$, giving a bijection between the nonzero rays in
$P(x,y)$ and the nonzero homogeneous lines $A_0(x,y)_d$;
\item for any two nonzero rays of $A_0$, the product of their homogeneous
lines in $A_0$ is nonzero if and only if their ray product in $P$ is
nonzero, in which case the product line represents that ray product;
\item $P$ is an associative category with zero.
\end{enumerate}
\end{proposition}
\begin{proof}
By Proposition~\ref{prop:projective-evaluation}(2), $A_0$ is basic and
representation-finite, so Lemma~\ref{lem:endpoint-chains} applies to $A_0$.
For each $a\in k^*$, the grading automorphism of $A_0$ that scales degree
$d$ by $a^d$ preserves every endpoint subbimodule of $A_0(x,y)$, because the
subbimodules of $A_0(x,y)$ form a chain with at most one member of each
dimension. Since $k$ is infinite, invariance under these scaling
automorphisms implies that every endpoint subbimodule is homogeneous. In
particular, its one-dimensional successive quotients are the nonzero
homogeneous lines $A_0(x,y)_d$, and because the radical of each endpoint
algebra of $A_0$ lies in positive degrees, each nonzero ray $U_i\setminus U_{i+1}$
contains the nonzero vectors of a unique homogeneous line $A_0(x,y)_d$. This proves (1).

For homogeneous elements $f\in A_0(x,y)_p$ and $g\in A_0(y,z)_q$, the
product $fg$ is homogeneous of degree $p+q$, whereas every element of
$M(f,g)=R_x f(\rad R_y)g R_z$ is supported in degrees strictly greater than
$p+q$. Thus $fg\in M(f,g)$ if and only if $fg=0$, which proves (2).
Finally, (3) follows immediately from (2) and the associativity of
multiplication in $A_0$.
\end{proof}

In Section~\ref{sec:scalar-normalization} we will show that $A_0$ is
isomorphic to the linearization of $P=\ray A_0$, and in
Theorem~\ref{thm:final-standardness}(3) we will also apply our constructions
to algebras defined directly from combinatorial ray categories. We record the
precise terminology for both settings here.

\begin{definition}
\label{def:finite-ray-category}
We adopt the following conventions:
\begin{enumerate}[label={\upshape(\arabic*)}]
\item When the ray table $\ray A$ of a basic algebra $A$ is associative, we
call $\ray A$ the \emph{ray category} of $A$: its objects are the vertices
indexing the primitive idempotents of $A$, and its composition is the ray
product of Definition~\ref{def:ray-product}.
\item Given a finite skeletal category with zero $Q$, its
\emph{linearization} $kQ$ is the $k$-algebra with basis consisting of the
nonzero morphisms of $Q$ and multiplication extending the composition in $Q$
(with zero compositions in $Q$ giving $0$ in $kQ$). Note that here $Q$ is
already a category with zero, rather than a quiver whose path algebra is
being formed.
\item A finite skeletal category with zero $Q$ is called a \emph{finite ray
category} if the identity morphisms of $Q$ form a complete set of primitive
orthogonal idempotents in $kQ$, the ray product \eqref{eq:ray-product} on
$kQ$ is independent of the choice of representatives, and the map
$r\mapsto[r]$ identifies $Q$ with $\ray(kQ)$.
\end{enumerate}
\end{definition}

When $\ray A$ is associative, we write $k\ray A$ for its linearization. For
$P=\ray A_0$, choose a nonzero homogeneous representative $b_r\in A_0$ for
each nonzero ray $r\in P$, taking $b_{1_x}=e_x$ for the identity rays. By
Proposition~\ref{prop:projective-ray-product}, there exist nonzero structure
constants $\kappa(r,s)\in k^*$ such that
\begin{equation}
 b_rb_s=\begin{cases}\kappa(r,s)b_{rs},&\text{if }rs\ne0,\\0,&\text{if }rs=0,\end{cases}
 \label{eq:projective-scalars}
\end{equation}
and associativity of $A_0$ is equivalent to the two-cocycle identity
$\kappa(r,s)\kappa(rs,t)=\kappa(s,t)\kappa(r,st)$ for all composable triples
with $rst\ne0$.

\section{Scalar normalization on a universal cover}
\label{sec:scalar-normalization}

Starting from the homogeneous ray representatives $b_r\in A_0$ and their
structure constants $\kappa(r,s)\in k^*$ in \eqref{eq:projective-scalars}, we
wish to find nonzero scalars $a_r\in k^*$ such that the rescaled elements
$a_r b_r$ multiply with structure constants equal to $1$, that is,
\[
                   a_ra_s\kappa(r,s)=a_{rs}\qquad\text{whenever }rs\ne0.
\]
Because distinct pairs of rays can have the same nonzero product, these
equations are coupled across all corners of $A_0$. Following Bongartz and
Gabriel \cite[Section~4]{BongartzGabriel} and Bretscher and Gabriel
\cite[Sections~2.6 and~3.1--3.2]{BretscherGabriel}, we solve them by
constructing a two-dimensional cell complex $K$ from the Auslander--Reiten
quiver $\Gamma_A$ and its meshes, and lifting the mesh category $\C$ to the
universal cover $\widetilde K$. On the universal cover, the category of
lifted projectives is directed and thin, allowing us to construct a
multiplicative basis one vertex at a time. Because every multiplicative
scalar weight on the cover is a ratio of object values, and the fundamental
group $G=\pi_1(K)$ is free (Lemma~\ref{lem:mesh-walk-free}), we can then
rescale the lifted basis to make it $G$-invariant and descend it to $A_0$.

\subsection{A cell complex of meshes and its periodic components}
\label{subsec:mesh-walk-cover}

Since a basic algebra is a direct product of connected blocks, we may work
with one connected block at a time and assume throughout this section that
$A$ is connected; then $\Gamma_A$ is connected by
\cite[Theorem~VII.2.1]{ARS}. Write $\Gamma=\Gamma_A$, and identify $\C$ with
$\ind(\modu A_0)$ via Proposition~\ref{prop:projective-evaluation}. To form
the cell complex $K$ of \cite[Section~4.1]{BongartzGabriel}, start with the
underlying undirected graph of $\Gamma$, adjoin a translation edge
$t_X:\tau X\to X$ for every nonprojective vertex $X$, and for every arrow
$a:Y\to X$ with $X$ nonprojective, attach a two-dimensional triangle cell
whose boundary is the closed walk $(\sigma a)a t_X^{-1}$ from $\tau X$
through $Y$ and $X$ back to $\tau X$. Here a \emph{walk} in $K$ means a path
of edges traversed in either orientation; each triangle cell identifies the
two-edge walk $(\sigma a)a$ up to homotopy with the translation edge $t_X$
(rather than equating the individual summands of the mesh relation in the
linear category).

Let $\Gamma_{\mathrm{per}}$ denote the full subquiver of $\Gamma$ on the
vertices that are periodic under $\tau$, and call its connected components
the \emph{periodic components} of $\Gamma$. Because $\tau$ and $\tau^{-1}$
are everywhere defined on $\Gamma_{\mathrm{per}}$ and the mate bijection
$\sigma$ restricts to arrows between periodic vertices, $\Gamma_{\mathrm{per}}$
is a stable translation quiver. By \eqref{eq:mesh-object-quotient}, its
scalar mesh category is the object quotient of $\C$ obtained by killing all
nonperiodic vertices, and in particular is finite-dimensional. Riedtmann's
structure theorem \cite[Sections~1.5--1.6 and~2.1]{Riedtmann1980} describes
such stable components as quotients of repetition quivers $\Z\Delta$, where
for an oriented tree $\Delta$ the quiver $\Z\Delta$ has vertices
$(x,j)\in\Delta_0\times\Z$, arrows $(x,j)\to(y,j)$ and $(y,j)\to(x,j+1)$ for
each arrow $x\to y$ of $\Delta$, and translation $\tau(x,j)=(x,j-1)$. We
prove the following self-contained form in Appendix~\ref{app:periodic-trees}.

\begin{proposition}
\label{prop:periodic-tree-model}
Let $Q$ be a finite stable translation quiver with at most one arrow between
each ordered pair of vertices. Fix a coefficient system $c$ and suppose that
the scalar mesh category $C(Q,c)$ of Definition~\ref{def:scalar-mesh-category}
is finite-dimensional.
\begin{enumerate}[label={\upshape(\arabic*)}]
\item Every connected component of $Q$ containing at least one arrow is
isomorphic to $\Z\Delta/\Theta$, where $\Delta$ is a finite simply-laced
Dynkin tree and $\Theta$ is an infinite cyclic group of automorphisms of
$\Z\Delta$ acting freely on vertices and commuting with $\tau$.
\item If $\Delta$ has type $D_n$, $E_6$, $E_7$, $E_8$, or $A_{2n+1}$, root
$\Delta$ at its unique branch vertex (for types $D$ and $E$) or at its middle
vertex (for type $A_{2n+1}$), and orient all edges of $\Delta$ away from the
root. Then $\Theta$ is generated by $\tau^s\rho$ for some integer $s\geq1$,
where $\rho$ is the automorphism $(v,j)\mapsto(\rho v,j)$ of $\Z\Delta$
induced by a root-fixing automorphism $\rho$ of $\Delta$. If $\Delta=A_{2n}$,
then $\Theta$ is generated by $\eta^r$ for some integer $r\geq1$, where
$\eta$ is an automorphism of $\Z\Delta$ satisfying $\eta^2=\tau^{-1}$ that
interchanges the two halves of $\Delta$.
\end{enumerate}
\end{proposition}

Geometrically, the mesh triangles lying over each edge of $\Delta$ form an
infinite strip in $\Z\Delta$, so the subcomplex of $K$ on a periodic
component is the mapping torus of the tree $\Delta$ under the monodromy
automorphism generating $\Theta$. Using this description, we can deformation-retract
each periodic component onto a circle and $K$ itself onto a finite graph.

\begin{lemma}[{\cite[Theorem~4.2 and its corollary]{BongartzGabriel}}]
\label{lem:mesh-walk-free}
Let $K$ be the mesh cell complex of the connected Auslander--Reiten quiver
$\Gamma_A$ defined in Subsection~\ref{subsec:mesh-walk-cover}. Then $K$ is
homotopy equivalent to a finite graph, and consequently its fundamental
group $G=\pi_1(K)$ is free.
\end{lemma}
\begin{proof}
The partial mate bijection $\sigma$ partitions the arrows of $\Gamma_A$ into
finite chains and finite cycles. Since $\sigma^2=\tau$, an arrow lies in a
cycle of $\sigma$ if and only if both of its endpoints are $\tau$-periodic;
thus the cyclic $\sigma$-orbits are precisely the arrows of
$\Gamma_{\mathrm{per}}$. At the end of any finite $\sigma$-chain of length at
least two, the terminal arrow belongs to only one mesh triangle of $K$. We
may therefore collapse that free arrow together with its unique triangle, and
repeat the collapse inductively along each finite $\sigma$-chain until only a
single quiver edge per chain and the translation edges $t_X$ remain outside
the periodic subcomplexes.

Next, consider a periodic component $\Z\Delta/\Theta$, modeled as in
Proposition~\ref{prop:periodic-tree-model}. When $\Delta$ has a central root
vertex fixed by $\rho$, we inductively prune the outermost leaves of $\Delta$
in complete $\rho$-orbits. Specifically, for a leaf $v$ of $\Delta$ with
parent $u$, the arrows between their $\tau$-orbits in $\Z\Delta$ are
\[
 a_j:u_j\longrightarrow v_j,\qquad
 b_j:v_j\longrightarrow u_{j+1},
\]
and the incident mesh triangles have boundaries
$a_jb_jt_{u_{j+1}}^{-1}$ and $b_ja_{j+1}t_{v_{j+1}}^{-1}$. Because $v$ is a
leaf, the translation edge $t_{v_{j+1}}$ belongs only to the second triangle;
collapsing $t_{v_{j+1}}$ across that triangle leaves the edge $b_j$ incident
only to the first triangle, and collapsing $b_j$ across the first triangle
leaves only the edges $a_j:u_j\to v_j$, which then contract onto the parent
vertices $u_j$. Performing these elementary collapses simultaneously across
a $\rho$-orbit of leaves is compatible with the $\Theta$-action, and
repeating the pruning down to the root contracts the periodic component onto
the circle formed by the translation edges over the root.

When $\Delta=A_{2n}$, its center is the middle edge rather than a vertex.
Pruning all other edges as above reduces the component to the periodic strip
over the central edge. Subdividing this strip along its midline into two
half-strips gives a deformation retraction onto the midline that is
equivariant with respect to both $\tau$ and the involution $\eta$ exchanging
the two endpoints of the central edge; hence the quotient again contracts
onto a circle. Likewise, any periodic component consisting of a single
$\tau$-orbit with no arrows is already a circle of translation edges $t_X$.

Finally, after contracting all periodic subcomplexes to circles and
collapsing all nonperiodic triangles as in the first paragraph, only
one-dimensional cells remain, so $K$ deformation-retracts onto a finite
graph, whose fundamental group is free.
\end{proof}

Reversing the elementary triangle collapses of Lemma~\ref{lem:mesh-walk-free}
also shows that any closed edge weight on the periodic subcomplexes extends
to all of $K$. Here an \emph{edge weight} on $K$ with values in an abelian
group assigns opposite values to the two orientations of each edge, and it
is \emph{closed} if the sum of its values around the boundary of every
triangle cell is zero.

\begin{lemma}
\label{lem:periodic-weight-extension}
Every closed edge weight on the union of the periodic subcomplexes of $K$
extends to a closed edge weight on $K$.
\end{lemma}
\begin{proof}
As in the first paragraph of the proof of Lemma~\ref{lem:mesh-walk-free}, we
may collapse all triangles outside the periodic subcomplexes one by one along
finite $\sigma$-chains. Assign weight zero to all nonperiodic edges that
remain after these collapses, and then restore the collapsed triangles in
reverse order: since each restored triangle has a unique free edge that was
removed at that step, its weight is uniquely determined by requiring the
boundary sum around the triangle to vanish. This extends the given periodic
weight to a closed edge weight on all of $K$.
\end{proof}

As an immediate topological consequence of Lemma~\ref{lem:mesh-walk-free}, we
can normalize all scalar coefficients in the mesh relations of $k(\Gamma_A)$
to $1$.

\begin{lemma}
\label{lem:unit-mesh-coefficients}
In every characteristic, any system of nonzero mesh coefficients $c_{X,a}$ on
$\Gamma_A$ can be normalized to $1$ by rescaling the arrows and the mesh
relations.
\end{lemma}
\begin{proof}
Working on one connected component of $\Gamma_A$ at a time, view the family
of nonzero scalars $c_{X,a}\in k^*$ as a multiplicative cellular $2$-cochain
on $K$ assigning $c_{X,a}$ to the triangle cell with boundary
$(\sigma a)a t_X^{-1}$. A multiplicative $1$-cochain $\phi$ on $K$ assigns a
nonzero scalar $\phi(e)\in k^*$ to each edge $e$ of $K$, and its coboundary
on the triangle $(X,a)$ is
\[
 (d\phi)(X,a)=\phi(\sigma a)\phi(a)\phi(t_X)^{-1}.
\]
Rescaling each arrow $a$ of $\Gamma_A$ by $\phi(a)$ and dividing each mesh
relation $r_X$ by $\phi(t_X)$ multiplies the coefficient $c_{X,a}$ by
$(d\phi)(X,a)$ without changing the mesh ideal. Since $K$ has no $3$-cells
and is homotopy equivalent to a graph by Lemma~\ref{lem:mesh-walk-free}, we
have $H^2(K;k^*)=1$, so there exists a $1$-cochain $\phi$ with
$d\phi=c^{-1}$. This rescaling normalizes all mesh coefficients to $1$ and
induces an object-fixing isomorphism on $k(\Gamma_A)$.
\end{proof}

\subsection{The universal cover and the lifted mesh category}
\label{subsec:universal-cover}

By Proposition~\ref{prop:graded-realization}(3), we identify $\C$ with the
scalar mesh category $k(\Gamma_A)$ equipped with the coefficients
\eqref{eq:scalar-mesh}. Let $\widetilde K$ be the universal cover of $K$,
whose vertices are homotopy classes of walks in $K$ starting at a fixed
basepoint, and on which the fundamental group $G=\pi_1(K)$ acts freely by
deck transformations $\widetilde Y\mapsto g\widetilde Y$. Lifting the arrows
and translation edges of $\Gamma$ to $\widetilde K$ produces a translation
quiver $\widetilde\Gamma$. Because each two-arrow path $(\sigma a)a$ in a
mesh at $X$ is homotopic in $K$ to the translation edge $t_X$, all terms of
a mesh relation $r_X$ lift to paths in $\widetilde\Gamma$ sharing the same
source $\tau\widetilde X$ and target $\widetilde X$. Imposing the lifted
relations $r_{\widetilde X}$ with the same scalar coefficients $c_{X,a}$
defines the lifted scalar mesh category $\widetilde\C$.

Any \emph{contextual mesh relation}—that is, a relation of the form
$u r_Y v$ with $u$ and $v$ composable paths in $\Gamma$—lifts uniquely to
$\widetilde\Gamma$ once a lift of its starting vertex is chosen, and all of
its terms share the same lifted endpoints. Consequently, both the path space
and the mesh relation ideal decompose as direct sums over lifted endpoints,
so pushing paths down to $\Gamma$ induces a canonical linear isomorphism for
any chosen lifts $\widetilde X$ of $X$ and $\widetilde Y$ of $Y$:
\begin{equation}
 \bigoplus_{g\in G}\widetilde\C(\widetilde X,g\widetilde Y)
                 \xrightarrow{\sim}\C(X,Y).
 \label{eq:walk-hom-sum}
\end{equation}
In particular, the uniform bound on the length of nonzero paths in the
finite-dimensional category $\C$ also holds in $\widetilde\C$. Note also
that the lifted quiver $\widetilde\Gamma$ is connected: in a module category
every almost split sequence has a nonzero middle term, so every translation
edge $t_X$ in a walk on $\widetilde K$ can be replaced by a two-arrow path
$(\sigma a)a$ in $\widetilde\Gamma$.

Assigning length $1$ to each quiver arrow and length $2$ to each translation
edge $t_X$ preserves length across every triangle cell of $K$. Because closed
walks in the simply connected cover $\widetilde K$ are null-homotopic, this
edge length integrates on $\widetilde K$ to an integer-valued \emph{height
function} on the vertices of $\widetilde\Gamma$ that increases by $1$ along
every arrow. In particular, $\widetilde\Gamma$ is directed, every object of
$\widetilde\C$ has trivial endomorphism ring $k$, and only finitely many
vertices can be reached by directed paths between any two vertices of
prescribed heights.

\subsection{Finite subcategories and evaluation of lifted objects}
\label{subsec:lifted-projectives}

Let $\mathcal L$ denote the opposite of the full subcategory of
$\widetilde\C$ on the lifted projective vertices. Taking the opposite
reverses morphism directions so that they match those of
$A_0(x,y)=\C(P_y,P_x)$: choosing one reference lift $\widetilde x$ for each
projective vertex $P_x$, the covering decomposition \eqref{eq:walk-hom-sum}
gives a $G$-grading on $A_0$,
\[
 A_0(x,y)=\bigoplus_{g\in G}A_0(x,y)_g,\qquad
 A_0(x,y)_g=\mathcal L(\widetilde x,g\widetilde y)=\widetilde\C(P_{g\widetilde y},P_{\widetilde x}),
\]
which we refer to as the \emph{deck grading} to distinguish it from the
categorical integer grading of Section~\ref{sec:graded-category}. Since each
homogeneous component $A_0(x,y)_d$ of the categorical grading is at most
one-dimensional by Lemma~\ref{lem:ray-degree} and the two gradings are
compatible, each nonzero homogeneous ray line $A_0(x,y)_d$ is concentrated in
a single deck degree $g\in G$ and therefore lifts uniquely to a nonzero
morphism space $\mathcal L(\widetilde x,g\widetilde y)$ once the source lift
$\widetilde x$ is fixed.

More generally, both $\mathcal L$ and the finite categories constructed from
integer weights in Section~\ref{sec:local-ray-geometry} are instances of the
following standard covering construction for group-graded algebras.

\begin{definition}
\label{def:graded-covering-category}
Let $B$ be an algebra graded by a group $\mathcal G$, with a complete set of
primitive orthogonal idempotents $\{e_x\}$ homogeneous of identity degree.
The \emph{graded covering category} of $B$ has objects $(x,g)$ with
$g\in\mathcal G$, morphism spaces
\[
 \Hom\bigl((x,g),(y,h)\bigr)=B(x,y)_{g^{-1}h},
\]
and composition induced by the multiplication of $B$ (for an additively
written group $\mathcal G$, the degree $g^{-1}h$ is written $h-g$).
\end{definition}

Because the deck group $G$ acts freely on $\widetilde K$ and the chosen
lifts $\widetilde x$ form a transversal for the $G$-orbits of lifted
projectives, sending $g\widetilde x$ to $(x,g)$ identifies $\mathcal L$ with
the graded covering category of $A_0$ for its deck grading:
\[
 \mathcal L(g\widetilde x,h\widetilde y)
 =\mathcal L(\widetilde x,g^{-1}h\widetilde y)
 =A_0(x,y)_{g^{-1}h}.
\]

To transfer finite representation type from $A_0$ to finite full
subcategories of $\mathcal L$ (and later to subquotients in
Sections~\ref{sec:local-ray-geometry} and~\ref{sec:integral-partial-lifts}),
we use the standard induction--restriction argument for split bimodule
retractions (compare the cleaving-functor method in
\cite[Sections~3.2 and~3.8]{BGRS}).

\begin{lemma}
\label{lem:split-representation-transfer}
Let $C\subseteq B$ be a unital subalgebra of a finite-dimensional algebra $B$
over any field, let $q:C\twoheadrightarrow D$ be a quotient algebra, and
suppose that $q$ extends to a $C$-bimodule homomorphism $\pi:B\to D$. If $B$
is representation-finite, then so is $D$.
\end{lemma}
\begin{proof}
View any right $D$-module $M$ as a right $C$-module via $q$. The canonical
$C$-module homomorphism $M\to(M\otimes_C B)|_C$ sending $m$ to $m\otimes 1$
admits the $C$-linear retraction $m\otimes b\mapsto m\pi(b)$. Thus $M$ is a
direct summand of the restriction to $C$ of the $B$-module $M\otimes_C B$.
Since $B$ has only finitely many indecomposable modules up to isomorphism,
only finitely many indecomposable $C$-modules can arise as direct summands of
their restrictions, and since inflation along $q:C\twoheadrightarrow D$ is
fully faithful, $D$ is representation-finite.
\end{proof}

\begin{lemma}
\label{lem:homogeneous-projection}
Let $B$ be a representation-finite algebra graded by a group $\mathcal G$,
with primitive idempotents in the identity degree. Then every finite full
subcategory of the graded covering category of $B$ is representation-finite.
\end{lemma}
\begin{proof}
Let $H$ be the full subcategory on finitely many objects $(x_1,g_1),\ldots,(x_m,g_m)$,
so that $H(i,j)=B(x_i,x_j)_{g_i^{-1}g_j}$, and embed $H$ into the ungraded
category $B'$ on the same object indices with morphism spaces
$B'(i,j)=B(x_i,x_j)$. Because $B'$ is Morita equivalent to the idempotent
corner of $B$ supported on the distinct vertices among $x_1,\ldots,x_m$, it
is representation-finite by Lemma~\ref{lem:endpoint-chains}(1). Moreover, the
linear projection $\pi:B'\to H$ that projects each corner $B(x_i,x_j)$ onto
its homogeneous component of degree $g_i^{-1}g_j$ is an $H$-bimodule
retraction: for $p\in H(i,j)$, $b\in B'(j,l)$, and $q\in H(l,r)$, the product
with the homogeneous component of $b$ of degree $g_j^{-1}g_l$ is the unique
component of $pbq$ of degree $(g_i^{-1}g_j)(g_j^{-1}g_l)(g_l^{-1}g_r)=g_i^{-1}g_r$,
so $\pi(pbq)=p\pi(b)q$. Applying Lemma~\ref{lem:split-representation-transfer}
with $C=D=H$ shows that $H$ is representation-finite.
\end{proof}

\begin{lemma}
\label{lem:lifted-projective-evaluation}
The category $\mathcal L$ is connected and thin (meaning that
$\dim_k\mathcal L(u,v)\leq1$ for all objects $u,v$). For each object
$\widetilde x$ of $\mathcal L$, write $P_{\widetilde x}$ for the
corresponding lifted projective vertex in $\widetilde\C$, and for each object
$\widetilde X\in\widetilde\C$ define the $\mathcal L$-module
$V_{\widetilde X}$ by
\[
 V_{\widetilde X}(\widetilde x)=\widetilde\C(P_{\widetilde x},\widetilde X).
\]
Then each $V_{\widetilde X}$ is a finite-dimensional indecomposable
$\mathcal L$-module, and the evaluation functor
$\widetilde X\mapsto V_{\widetilde X}$ on $\widetilde\C$ is fully faithful.
\end{lemma}
\begin{proof}
Because the negative of height strictly increases along nonidentity morphisms
of $\mathcal L$, every object of $\mathcal L$ has endomorphism ring $k$, and
in any two-object full subcategory $\{u,v\}$ of $\mathcal L$ with
$\mathcal L(u,v)\ne0$ we have $\mathcal L(v,u)=0$. Every linear subspace of
$\mathcal L(u,v)$ is then a two-sided ideal of the two-object subcategory.
Since this subcategory is representation-finite by
Lemma~\ref{lem:homogeneous-projection}, it has only finitely many ideals by
Lemma~\ref{lem:endpoint-chains}(1), which forces $\dim_k\mathcal L(u,v)\leq1$.
Thus $\mathcal L$ is thin.

For any $\widetilde X\in\widetilde\C$, the uniform bound on nonzero path
lengths in $\widetilde\C$ implies that $V_{\widetilde X}$ is supported on
finitely many objects of $\mathcal L$. If $X\in\C$ denotes the projection of
$\widetilde X$, the covering isomorphism \eqref{eq:walk-hom-sum} gives, for
each reference projective lift $\widetilde x$,
\[
 \C(P_x,X)=\bigoplus_{g\in G}\widetilde\C(P_{g\widetilde x},\widetilde X)
 =\bigoplus_{g\in G}V_{\widetilde X}(g\widetilde x).
\]
Summing over $x$ identifies the evaluated $A_0$-module $\operatorname{ev}(X)$
of Proposition~\ref{prop:projective-evaluation} with the direct sum of the
spaces $V_{\widetilde X}(\widetilde z)$ over all objects
$\widetilde z\in\mathcal L$. Any homomorphism
$\varphi:V_{\widetilde X}\to V_{\widetilde Y}$ of $\mathcal L$-modules
acts componentwise on these spaces and commutes with the action of each
deck-homogeneous element of $A_0$, so it defines an $A_0$-module
homomorphism $\operatorname{ev}(X)\to\operatorname{ev}(Y)$ preserving each
individual summand indexed by $\widetilde z$. By
Proposition~\ref{prop:projective-evaluation}(1), this $A_0$-module map has
the form $\operatorname{ev}(f)$ for a unique $f\in\C(X,Y)$. Using
\eqref{eq:walk-hom-sum}, decompose $f=\sum_{g\in G}\bar f_g$, where
$\bar f_g\in\C(X,Y)$ is the projection of $f_g\in\widetilde\C(\widetilde X,g\widetilde Y)$.
Because composition with $f_g$ maps $V_{\widetilde X}(\widetilde z)$ into
$\widetilde\C(P_{\widetilde z},g\widetilde Y)=V_{\widetilde Y}(g^{-1}\widetilde z)$
and $G$ acts freely, the component $\operatorname{ev}(\bar f_g)$ for $g\ne1$
maps the summand at $\widetilde z$ into the distinct summand at
$g^{-1}\widetilde z$. Since $\varphi=\operatorname{ev}(f)$ preserves each
summand $\widetilde z$, we must have $\operatorname{ev}(\bar f_g)=0$ for all
$g\ne1$; faithfulness of $\operatorname{ev}$ and \eqref{eq:walk-hom-sum}
then give $f_g=0$ for $g\ne1$, so $f=\bar f_1$ and $\varphi$ is induced by
$f_1\in\widetilde\C(\widetilde X,\widetilde Y)$. This proves fullness, and
the same decomposition shows faithfulness. In particular,
$\End_{\mathcal L}(V_{\widetilde X})\cong\widetilde\C(\widetilde X,\widetilde X)=k$,
so each $V_{\widetilde X}$ is indecomposable.

Finally, since each $V_{\widetilde X}$ is indecomposable, its support lies in
a single connected component of $\mathcal L$. For every arrow
$\widetilde X\to\widetilde Y$ of $\widetilde\Gamma$, the induced nonzero
homomorphism $V_{\widetilde X}\to V_{\widetilde Y}$ forces the supports of
$V_{\widetilde X}$ and $V_{\widetilde Y}$ to lie in the same connected
component of $\mathcal L$. Because $\widetilde\Gamma$ is connected and every
object $\widetilde x\in\mathcal L$ belongs to the support of
$V_{P_{\widetilde x}}$, it follows that $\mathcal L$ is connected.
\end{proof}

\subsection{Extending a basis by one object}

We will construct a multiplicative basis of $\mathcal L$ inductively along an
exhaustion by finite full subcategories, adding one source or sink at a time.
When we adjoin a source $s$ to a finite directed thin category, the objects
$y$ receiving a nonzero morphism from $s$ form a finite poset, and
normalizing the new products amounts to trivializing a multiplicative
$1$-cocycle on the incidence category of that poset (compare Bretscher and
Gabriel \cite[Section~2.6]{BretscherGabriel}). Recall that the
\emph{incidence category} of a finite poset $R$ has the elements of $R$ as
objects and a single basis morphism $x\to y$ whenever $x\leq y$, with the
composite of any two comparable morphisms equal to the basis morphism between
their endpoints; an \emph{additive weight} on $R$ with values in an abelian
group $M$ is a function $w(x,y)\in M$ on comparable pairs $x\leq y$
satisfying $w(x,y)+w(y,z)=w(x,z)$ for $x\leq y\leq z$.

\begin{lemma}
\label{lem:poset-weights}
If the incidence algebra of a finite poset $R$ is representation-finite, then
every additive weight $w$ on $R$ with values in an abelian group $M$ has the
form $w(x,y)=\zeta_y-\zeta_x$ for some family of object values
$\zeta_x\in M$ (written multiplicatively as $w(x,y)=\zeta_y/\zeta_x$ when
$M=k^*$).
\end{lemma}
\begin{proof}
We proceed by induction on $|R|$. Choose a minimal element $s\in R$ and set
$U=\{y\in R:s<y\}\subseteq R\setminus\{s\}$. We claim that distinct
connected components of the subposet $U$ lie in distinct connected components
of $R\setminus\{s\}$. Suppose to the contrary that two connected components
$U_1,U_2$ of $U$ lie in the same connected component $C$ of
$R\setminus\{s\}$. For each scalar $\lambda\in k^*$, define a representation
of $R$ of dimension vector $(1,\ldots,1)$ by placing $k$ at every vertex, the
identity map $1_k$ on every relation inside $R\setminus\{s\}$ and from $s$ to
$U\setminus U_1$, and multiplication by $\lambda$ on every relation from $s$
to $U_1$. Any isomorphism between two such representations must act by a
constant scalar on the connected component $C$ and therefore cannot change
the ratio of the maps from $s$ into $U_1$ and $U_2$. Varying $\lambda\in k^*$
thus yields infinitely many pairwise nonisomorphic indecomposable
representations, contradicting the hypothesis that $R$ is
representation-finite.

Because the incidence algebra of $R\setminus\{s\}$ is an idempotent corner of
that of $R$, it is representation-finite by Lemma~\ref{lem:endpoint-chains}(1).
By the induction hypothesis, after subtracting a coboundary $\zeta_y-\zeta_x$
on $R\setminus\{s\}$, we may assume that $w(x,y)=0$ for all $x\leq y$ in
$R\setminus\{s\}$. For any $y\leq z$ in $U$, the additivity condition
$w(s,y)+w(y,z)=w(s,z)$ then gives $w(s,y)=w(s,z)$, so $w(s,-)$ is constant on
each connected component of $U$. Since distinct connected components of $U$
lie in distinct connected components of $R\setminus\{s\}$, we may add a
constant to $\zeta$ on each connected component of $R\setminus\{s\}$ to make
$w(s,y)=\zeta_y-\zeta_s$ for all $y\in U$ without altering the zero weights
on $R\setminus\{s\}$.
\end{proof}

To apply Lemma~\ref{lem:poset-weights} when adjoining a source $s$, we note
that the functor $y\mapsto T(s,y)$ on a finite directed thin category $T$
defines a thin module whose faithful support is a poset. We will reuse this
observation in Sections~\ref{sec:weights-and-parameters}
and~\ref{sec:local-ray-geometry}.

\begin{lemma}
\label{lem:thin-module-incidence}
Let $T$ be a finite directed category, and let $M$ be a covariant or
contravariant $k$-linear functor from $T$ to vector spaces satisfying
$\dim_k M(y)\leq1$ for every object $y$. Let $R$ be the category obtained
from $T$ by passing to the quotient by the annihilator ideal $\ann_T(M)$ and
discarding the objects $y$ with $M(y)=0$. Then $R$ is isomorphic to the
incidence category of a finite poset, and if $T$ is representation-finite,
then so is $R$.
\end{lemma}
\begin{proof}
Fix a nonzero basis vector in each one-dimensional space $M(y)$ for $y\in R$.
Because $R$ acts faithfully on $M$, every morphism space $R(y,z)$ acts
faithfully on the one-dimensional spaces $M(y)$ and $M(z)$ and therefore has
dimension at most one, spanned when nonzero by the unique morphism sending
the chosen basis vector to the chosen basis vector. Since $T$ is directed,
the relation $R(y,z)\ne0$ defines a partial order on the objects of $R$, and
these normalized morphisms identify $R$ with the incidence category of that
poset. Finally, if $T$ is representation-finite, then the full subcategory of
$T$ on the objects of $R$ is an idempotent corner of $T$ and hence
representation-finite by Lemma~\ref{lem:endpoint-chains}(1); since $R$ is a
quotient of that corner, $R$ is representation-finite as well.
\end{proof}

\begin{lemma}
\label{lem:source-basis-extension}
Let $T$ be a finite directed thin representation-finite category, let $s$ be
a source of $T$, and set $T'=T\setminus\{s\}$ and $M=\rad T(s,-)$. Let $R$ be
the representation-finite poset incidence category obtained from $T'$ and $M$
as in Lemma~\ref{lem:thin-module-incidence}.
\begin{enumerate}[label={\upshape(\arabic*)}]
\item Every multiplicative basis of $T'$ extends to a multiplicative basis of
$T$.
\item Let $\operatorname{Scal}(T)$ denote the abelian group of multiplicative
$k^*$-valued weights on the nonzero nonidentity morphism lines of $T$, that
is, assignments $d(r)\in k^*$ satisfying $d(rr')=d(r)d(r')$ whenever
$rr'\ne0$. Then restriction to $T'$ is a split surjection of abelian groups
\[
 \operatorname{Scal}(T)\longrightarrow \operatorname{Scal}(T')
\]
whose kernel is isomorphic to $(k^*)^{\pi_0(R)}$.
\end{enumerate}
Both assertions hold dually when $s$ is a sink of $T$.
\end{lemma}
\begin{proof}
Fix a multiplicative basis $\{b_{yz}\}$ of $T'$, and choose an arbitrary
nonzero vector $f_y\in T(s,y)$ for each object $y$ of $R$. Whenever a basis
element $b_{yz}\in T'(y,z)$ acts nontrivially on $M$, so that
$T(s,y)b_{yz}=T(s,z)\ne0$, we have $f_y b_{yz}=\mu_{yz}f_z$ for a nonzero
scalar $\mu_{yz}\in k^*$. Associativity of composition in $T$ implies that
$\mu_{yz}\mu_{zw}=\mu_{yw}$ for all comparable triples $y\leq z\leq w$ in
$R$. By Lemma~\ref{lem:poset-weights}, there exist scalars $\zeta_y\in k^*$
such that $\mu_{yz}=\zeta_z/\zeta_y$. Replacing each $f_y$ by $\zeta_y f_y$
gives $(\zeta_y f_y)b_{yz}=\zeta_z f_z$, extending the multiplicative basis
of $T'$ to $T$. This proves (1).

For (2), an element $d\in\operatorname{Scal}(T')$ extends to $d\in\operatorname{Scal}(T)$
if and only if its values $d(sy)\in k^*$ at the objects $y\in R$ satisfy
$d(sy)d(b_{yz})=d(sz)$ whenever $y\leq z$ in $R$. Choosing a base object
$y_D$ in each connected component $D$ of $R$ and a spanning tree of walks in
$D$ from $y_D$ to every $y\in D$, Lemma~\ref{lem:poset-weights} shows that we
obtain a well-defined homomorphic section
$\operatorname{Scal}(T')\to\operatorname{Scal}(T)$ by setting $d(sy_D)=1$ and
propagating $d(sy)$ along the walks in $D$ via products and inverses of the
values $d(b_{yz})$. Any two extensions of the same weight on $T'$ differ by a
constant scalar on each connected component $D\in\pi_0(R)$, so the kernel of
the restriction map is $(k^*)^{\pi_0(R)}$.
\end{proof}

We now apply Lemma~\ref{lem:source-basis-extension} along an exhaustion of
$\mathcal L$. Recall that a full subcategory of $\mathcal L$ is \emph{convex}
if it contains every object lying on a chain of nonzero morphisms between any
two of its objects.

\begin{lemma}
\label{lem:cover-exhaustion}
There exists an increasing sequence $\mathcal L_1\subset\mathcal L_2\subset\cdots$
of finite convex full subcategories of $\mathcal L$ whose union is
$\mathcal L$, such that each $\mathcal L_{m+1}$ is obtained from
$\mathcal L_m$ by adjoining a single source or sink. Every finite full
subcategory of $\mathcal L$ is directed, thin, and representation-finite.
Consequently, $\mathcal L$ admits a multiplicative basis, and the group
$\operatorname{Scal}(\mathcal L)$ of multiplicative $k^*$-valued weights on
$\mathcal L$ is an abelian divisible group.
\end{lemma}
\begin{proof}
Because $\mathcal L$ is connected and locally finite, it has countably many
objects, and because the integer height strictly decreases along every
nonidentity morphism of $\mathcal L$, there are only finitely many directed
paths between any two objects. The convex hull in $\mathcal L$ of any finite
set of objects is therefore finite. Enumerating the objects of $\mathcal L$
and taking the convex hulls of initial segments yields a cofinal sequence of
finite convex subcategories. Whenever $C\subsetneq C'$ is a proper inclusion
of finite convex subcategories, reachability along nonzero morphisms defines
a partial order on $C'$. If $C'\setminus C$ contains an element lying below
some element of $C$, any maximal such element of $C'\setminus C$ is a source
in the union of $C$ with that element, and adjoining it preserves convexity;
otherwise, any minimal element of $C'\setminus C$ retains convexity and
adjoins a sink. Refining the cofinal sequence in this way produces the
desired one-object exhaustion $\mathcal L_1\subset\mathcal L_2\subset\cdots$.

Every finite full subcategory of $\mathcal L$ is directed by height, thin by
Lemma~\ref{lem:lifted-projective-evaluation}, and representation-finite by
Lemma~\ref{lem:homogeneous-projection}. Applying
Lemma~\ref{lem:source-basis-extension}(1) at each step $\mathcal L_m\subset\mathcal L_{m+1}$
therefore builds a compatible sequence of multiplicative bases whose union is
a multiplicative basis of $\mathcal L$. Likewise, by
Lemma~\ref{lem:source-basis-extension}(2), the inverse system of weight groups
$\operatorname{Scal}(\mathcal L_m)$ consists of split surjections whose
kernels are products of copies of the divisible group $k^*$ (since $k$ is
algebraically closed). Hence the inverse limit
$\operatorname{Scal}(\mathcal L)=\varprojlim\operatorname{Scal}(\mathcal L_m)$
is isomorphic to a direct product of copies of $k^*$ and is therefore a
divisible abelian group.
\end{proof}

\subsection{Scalar automorphisms and descent}
\label{subsec:scalar-descent}

To make the multiplicative basis of $\mathcal L$ invariant under the action
of the free deck group $G=\pi_1(K)$, we first show that every multiplicative
scalar weight on $\mathcal L$ is a coboundary of object values.

\begin{lemma}
\label{lem:scalar-object-values}
Every multiplicative scalar weight $d\in\operatorname{Scal}(\mathcal L)$ is a
ratio of object values: there exist nonzero scalars $\zeta_x\in k^*$ for
$x\in\operatorname{ob}\mathcal L$ such that $d(r)=\zeta_y/\zeta_x$ for every
nonzero morphism line $r\in\mathcal L(x,y)$.
\end{lemma}
\begin{proof}
Each weight $d\in\operatorname{Scal}(\mathcal L)$ acts on $\mathcal L$ by an
object-fixing automorphism $\alpha_d$ that scales each nonzero morphism line
$r$ by $d(r)$. Precomposing $\mathcal L$-module actions with $\alpha_d^{-1}$
defines an action of the group $\operatorname{Scal}(\mathcal L)$ on the set
of isomorphism classes of finite-dimensional indecomposable
$\mathcal L$-modules. For any fixed finite subset
$S\subset\operatorname{ob}\mathcal L$, the full subcategory of $\mathcal L$
on $S$ is representation-finite by Lemma~\ref{lem:cover-exhaustion}, so there
are only finitely many isomorphism classes of indecomposable
$\mathcal L$-modules supported in $S$. Because $\operatorname{Scal}(\mathcal L)$
is a divisible abelian group by Lemma~\ref{lem:cover-exhaustion}, it has no
nontrivial finite quotient groups and must therefore act trivially on this
finite set of isomorphism classes. In particular, for every object
$\widetilde X\in\widetilde\C$, the evaluated module $V_{\widetilde X}$ of
Lemma~\ref{lem:lifted-projective-evaluation} satisfies
$V_{\widetilde X}\cong T(V_{\widetilde X})$, where $T$ denotes pullback along
$\alpha_d^{-1}$.

Choose an $\mathcal L$-module isomorphism
$\theta_{\widetilde X}:V_{\widetilde X}\xrightarrow{\sim}T(V_{\widetilde X})$
for each $\widetilde X\in\widetilde\C$, taking
$\theta_{P_{\widetilde x}}(v)=\alpha_d^{-1}(v)$ on the representable modules
$V_{P_{\widetilde x}}=\mathcal L(\widetilde x,-)$. By the full faithfulness
of $\widetilde X\mapsto V_{\widetilde X}$
(Lemma~\ref{lem:lifted-projective-evaluation}), transporting morphisms via
$f\mapsto\theta_{\widetilde Y}^{-1}\circ T(f)\circ\theta_{\widetilde X}$
defines an object-fixing automorphism $\Phi$ of $\widetilde\C$ whose
restriction to the opposite subcategory $\mathcal L$ of lifted projectives is
precisely $\alpha_d$.

Because the endpoints of each arrow $a:\widetilde U\to\widetilde V$ of
$\widetilde\Gamma$ differ in height by $1$, the space
$\widetilde\C(\widetilde U,\widetilde V)$ is one-dimensional and spanned by
$a$, so $\Phi(a)=w(a)a$ for some $w(a)\in k^*$. Since $\Phi$ preserves the
mesh relation $r_{\widetilde X}$ at each nonprojective vertex $\widetilde X$,
it scales $r_{\widetilde X}$ by some $\nu_{\widetilde X}\in k^*$, which
gives
\[
 w(\sigma a)w(a)=\nu_{\widetilde X}
\]
for every arrow $a$ ending at $\widetilde X$. Assigning $\nu_{\widetilde X}$
to each translation edge $t_{\widetilde X}$ makes $w$ a closed multiplicative
edge weight on the simply connected complex $\widetilde K$. Integrating $w$
along walks from a basepoint produces vertex scalars
$\zeta_{\widetilde X}\in k^*$ such that
$w(a)=\zeta_{\widetilde V}/\zeta_{\widetilde U}$ for every arrow
$a:\widetilde U\to\widetilde V$, and hence $\Phi$ acts on every morphism in
$\widetilde\C(\widetilde X,\widetilde Y)$ by multiplication by
$\zeta_{\widetilde Y}/\zeta_{\widetilde X}$. Since
$\mathcal L(\widetilde x,\widetilde y)=\widetilde\C(P_{\widetilde y},P_{\widetilde x})$
and $\Phi$ acts on $\mathcal L(\widetilde x,\widetilde y)$ by $d(\widetilde x,\widetilde y)$,
we obtain $d(\widetilde x,\widetilde y)=\zeta_{P_{\widetilde x}}/\zeta_{P_{\widetilde y}}$,
which has the required form on setting $\zeta_{\widetilde x}=\zeta_{P_{\widetilde x}}^{-1}$.
\end{proof}

We can now descend the multiplicative basis from $\mathcal L$ to $A_0$
(compare Bretscher and Gabriel \cite[Sections~3.1--3.2]{BretscherGabriel}).

\begin{proposition}
\label{prop:scalar-normalization}
Let $A_0$ be the standard form of $A$ constructed in
Construction~\ref{con:projective-algebra}, and let $P=\ray A_0$. Then there
is a graded algebra isomorphism $A_0\cong kP$ preserving the primitive
idempotents and the homogeneous ray lines. In particular, $A_0$ admits a
normed multiplicative basis.
\end{proposition}
\begin{proof}
Let $\{b_{xy}\}$ be the multiplicative basis of $\mathcal L$ provided by
Lemma~\ref{lem:cover-exhaustion}, where $b_{xy}\in\mathcal L(x,y)$ when
$\mathcal L(x,y)\ne0$. For each generator $g$ of the free deck group $G$, the
deck action on the basis defines a multiplicative weight $d_g\in\operatorname{Scal}(\mathcal L)$
by
\[
 g(b_{xy})=d_g(x,y)b_{gx,gy}.
\]
By Lemma~\ref{lem:scalar-object-values}, for each free generator $g$ there
exist scalars $a_g(x)\in k^*$ such that $d_g(x,y)=a_g(y)/a_g(x)$. Because
$G$ is a free group acting freely on $\operatorname{ob}\mathcal L$, each
$G$-orbit of objects is a Cayley tree for the free generators of $G$, so we
can solve the equations $f(gx)=a_g(x)f(x)$ simultaneously for all free
generators $g$ and all objects $x\in\operatorname{ob}\mathcal L$. Rescaling
the basis of $\mathcal L$ to
\[
 b'_{xy}=\frac{f(y)}{f(x)}b_{xy}
\]
preserves multiplicativity and makes the basis $G$-invariant:
$g(b'_{xy})=b'_{gx,gy}$.

As shown in Subsection~\ref{subsec:lifted-projectives}, each nonzero
homogeneous ray line of $A_0$ lies in a single deck degree and therefore
lifts uniquely to a nonzero morphism space $\mathcal L(x,y)$ once a lift $x$
of its source is chosen; by $G$-invariance, the image of $b'_{xy}$ in $A_0$
is independent of the choice of lift $x$. Since nonzero products of
homogeneous ray lines in $A_0$ lift to nonzero compositions in $\mathcal L$,
the descended vectors form a homogeneous multiplicative basis of $A_0$
realizing the isomorphism $A_0\cong kP$, and Corollary~\ref{cor:ray-transversal}
shows that this basis is normed.
\end{proof}

\begin{corollary}
\label{cor:graded-standard-form}
For $P=\ray A_0$, the linearization $kP\cong A_0$ is representation-finite
and satisfies $\add k(\Gamma_A)\simeq\modu kP$. Moreover, $A$ is standard in
the sense of Definition~\ref{def:standardness} if and only if $A\cong A_0$.
\end{corollary}
\begin{proof}
The first statement combines Propositions~\ref{prop:graded-realization}(3),
\ref{prop:projective-evaluation}, and~\ref{prop:scalar-normalization}. If $A$
is standard, then the object-preserving equivalence
$k(\Gamma_A)\simeq\ind(\modu A)$ restricts on the projective vertices to an
isomorphism $A_0\cong A$. Conversely, any algebra isomorphism $A\cong A_0$
induces an equivalence $F:\ind(\modu A)\xrightarrow{\sim}\ind(\modu A_0)\simeq\C$,
possibly permuting the vertices of $\Gamma_A$. Since $\C\cong k(\Gamma_A)$ is
generated in degree one with homogeneous relations, its associated graded
category $\gr_\JJ\C$ is canonically isomorphic to $\C$ itself, and the induced
graded equivalence $\gr_\JJ F:\C\xrightarrow{\sim}\C$ effects the same vertex
permutation as $F$. Composing $F$ with $(\gr_\JJ F)^{-1}$ therefore produces an
equivalence $\ind(\modu A)\xrightarrow{\sim}\C\cong k(\Gamma_A)$ fixing every
vertex.
\end{proof}

\section{Standardness outside characteristic two}
\label{sec:mesh-deformations}

We now complete the proof of Theorems~A and~B when $\charac k\ne2$. By
Corollary~\ref{cor:graded-standard-form}, it suffices to prove that
$A\cong A_0$, or equivalently that $\ind(\modu A)\cong\C$. In the graded
category $\C=\gr_\JJ\ind(\modu A)$, the chosen arrow representatives satisfy
the quadratic mesh relations \eqref{eq:scalar-mesh} exactly, whereas any
lifts of those arrows to $\ind(\modu A)$ satisfy the mesh relations modulo
$\JJ^3$. If the linear derivative map sending degree-$m$ perturbations of the
arrows to their degree-$(m+1)$ mesh errors is surjective for every $m\geq2$,
we can inductively adjust the chosen irreducible morphisms degree by degree
until the mesh relations hold on the nose in $\ind(\modu A)$.

In Subsection~\ref{subsec:mesh-derivative} we show that, in any
characteristic, the cokernel of this linearized mesh derivative localizes to
the $\tau$-periodic components of $\Gamma_A$. In
Appendix~\ref{app:periodic-mesh} we prove that the localized cokernel
vanishes whenever $\charac k\ne2$, provided the mesh category satisfies a
natural necessary condition for finite representation type that we call the
\emph{brick hypothesis}: every nonempty object quotient contains an object
with endomorphism ring $k$. Subsection~\ref{subsec:brick-hypothesis} verifies
that $\C$ satisfies this hypothesis (which rules out the small-period
components of types $E_6$, $E_7$, and $E_8$ that would otherwise have nonzero
cokernels in characteristics $3$ and $5$), and
Subsection~\ref{subsec:successive-correction} carries out the successive
arrow corrections.

\subsection{The derivative and periodic vertices}
\label{subsec:mesh-derivative}

Throughout this subsection, let $Q$ be an arbitrary finite translation quiver
with at most one arrow between each ordered pair of vertices, and let
$C=C(Q,c)$ be the scalar mesh category of
Definition~\ref{def:scalar-mesh-category} for a fixed system of nonzero
coefficients $c=(c_{x,a})$ (we reserve the calligraphic letter $\C$ for the
graded module category $\gr_\JJ\ind(\modu A)$).

For each integer $m\geq2$, define the homogeneous spaces of degree-$m$ arrow
perturbations and degree-$(m+1)$ mesh errors by
\[
 U_m(C)=\bigoplus_{a:u\to v}C_m(u,v),\qquad
 V_{m+1}(C)=\bigoplus_{x\text{ nonprojective}}C_{m+1}(\tau x,x),
\]
where $C_m$ denotes the homogeneous component of path length $m$. Substituting
$a\mapsto a+f_a$ with $f_a\in C_m(u,v)$ into the quadratic mesh relations
$r_x=\sum_{a:y\to x}c_{x,a}(\sigma a)a$ and taking the linear term in $f$
yields the \emph{mesh derivative}
\begin{equation}
 \begin{split}
 D_m:U_m(C)&\longrightarrow V_{m+1}(C),\\
 (D_mf)_x&=\sum_{a:y\to x}c_{x,a}
             \bigl(f_{\sigma a}a+(\sigma a)f_a\bigr).
 \end{split}
 \label{eq:mesh-derivative}
\end{equation}
Thus $\operatorname{coker}D_m(C)$ measures the degree-$(m+1)$ mesh errors
that cannot be canceled by a degree-$m$ modification of the arrows.

Let $T\subseteq Q_0$ be the set of $\tau$-periodic vertices of $Q$, so that
the full subquiver $Q_T$ is a stable translation quiver, and let
$C_T=C(Q_T,c|_T)$ be its scalar mesh category, which by
\eqref{eq:mesh-object-quotient} is the object quotient of $C$ obtained by
killing the identity morphisms outside $T$ (compare Riedtmann's stable-block
construction \cite{Riedtmann1984}). The following localization proposition is
used both in Appendix~\ref{app:periodic-mesh} (for $\charac k\ne2$) and in
Section~\ref{sec:positive-deformations} (for $\charac k=2$).

\begin{proposition}
\label{prop:periodic-localization}
Let $C=C(Q,c)$ be a scalar mesh category on a finite translation quiver $Q$
with at most one arrow between each ordered pair of vertices, and let $C_T$
be its object quotient on the set $T$ of $\tau$-periodic vertices. For every
$m\geq2$, the quotient functor $C\twoheadrightarrow C_T$ induces an
isomorphism
\[
 \operatorname{coker}D_m(C)\xrightarrow{\sim}
 \operatorname{coker}D_m(C_T).
\]
\end{proposition}
\begin{proof}
Before imposing the mesh relations of $C$, let $\widetilde U_m$ and
$\widetilde V_{m+1}$ be the free vector spaces spanned by literal paths in
$Q$, with basis $(b,p)$ for $\widetilde U_m$ (where $b:u\to v$ is an arrow
and $p:u\to v$ is a path of length $m$) and basis $(x,w)$ for
$\widetilde V_{m+1}$ (where $x$ is a nonprojective vertex and
$w:\tau x\to x$ is a path of length $m+1$). Formula~\eqref{eq:mesh-derivative}
defines a lifted linear map $\widetilde D_m:\widetilde U_m\to\widetilde V_{m+1}$.

In the matrix of $\widetilde D_m$ with respect to these path bases, each
column $(b,p)$ has at most two nonzero entries: if $v$ is nonprojective, the
occurrence of $b$ as the second arrow in $r_v$ contributes
$c_{v,b}(\sigma b)p$ to the row $(v,(\sigma b)p)$; if $u$ is noninjective,
writing $b=\sigma b'$ for $b':v\to\tau^{-1}u$, the occurrence of $b$ as the
first arrow in $r_{\tau^{-1}u}$ contributes $c_{\tau^{-1}u,b'}pb'$ to the row
$(\tau^{-1}u,pb')$. Conversely, each row $(x,w)$ with $w:\tau x\to x$ of
length $m+1$ is incident to exactly two columns (counted with multiplicity),
corresponding to the first and last arrows of $w$: writing $w=bp$ with first
arrow $b=\sigma a:\tau x\to y$ gives the incidence with column $(a,p)$, while
writing $w=pa$ with last arrow $a:y\to x$ gives the incidence with column
$(\sigma a,p)$, which links $(x,pa)$ to $(y,(\sigma^2 a)p)$ whenever $y$ is
nonprojective and has no second entry when $y$ is projective.

Viewing the rows $(x,w)$ as vertices and the two-entry columns as edges (with
one-entry columns as half-edges at exposed endpoints), the connected
components of this incidence graph are finite chains and finite cycles. Every
chain component ends at a one-entry column, so Gaussian elimination from that
exposed end places the entire chain component in $\im\widetilde D_m$. Along a
two-entry column, moving from $(x,pa)$ to $(y,(\sigma^2 a)p)$ rotates the
last arrow $a$ of $w$ to the front and applies $\tau=\sigma^2$; iterating this
rotation $d=m+1$ times gives
\[
 \operatorname{rot}^d(x,w)=(\tau x,\tau w).
\]
Consequently, a row $(x,w)$ lies on a cycle component if and only if every
vertex visited by $w$ admits arbitrarily many positive and negative powers of
$\tau$, which in the finite quiver $Q$ means precisely that the path $w$ lies
entirely inside $Q_T$. Moreover, because a vertex is in $T$ if and only if
its $\tau$-translate is in $T$, no column of $\widetilde D_m$ links a path
contained in $Q_T$ to a path passing outside $T$. The free row space
$\widetilde V_{m+1}$ therefore decomposes as a direct sum of the subspace of
paths contained in $Q_T$ and the subspace of paths passing outside $T$; the
latter is a sum of chain components and lies in $\im\widetilde D_m$, so
projection onto paths in $Q_T$ induces an isomorphism
$\operatorname{coker}\widetilde D_m(Q)\xrightarrow{\sim}\operatorname{coker}\widetilde D_m(Q_T)$.

Finally, let $R_V\subseteq\widetilde V_{m+1}$ be the subspace spanned by the
contextual mesh relations in $\widetilde V_{m+1}$, so that
\[
 \operatorname{coker}D_m(C)=\widetilde V_{m+1}/\bigl(R_V+\im\widetilde D_m\bigr).
\]
Because the projection onto paths contained in $Q_T$ kills every contextual
mesh relation passing outside $T$ and sends the remaining terms to the
contextual mesh relations of $C_T$, it induces the stated isomorphism
$\operatorname{coker}D_m(C)\xrightarrow{\sim}\operatorname{coker}D_m(C_T)$.
\end{proof}

\subsection{The restriction imposed by finite representation type}
\label{subsec:brick-hypothesis}

We say that a $k$-linear category $C$ satisfies the \emph{brick hypothesis}
if every nonempty object quotient of $C$ contains a \emph{brick}, that is, an
object whose endomorphism ring is $k$. Here, for any subset $\Sigma$ of
objects of $C$, the object quotient $C/[\Sigma]$ denotes the quotient of $C$
by the ideal of morphisms factoring through finite direct sums of objects in
$\Sigma$, with the objects in $\Sigma$ removed. Although the periodic
quotient $\C_T$ of $\C=\gr_\JJ\ind(\modu A)$ need not itself be the category of
indecomposable modules of an algebra, both $\C$ and $\C_T$ inherit the brick
hypothesis from $\modu A_0$.

\begin{lemma}
\label{lem:deletion-brick}
Let $\mathcal D$ be a full additive subcategory of a finite-length
$k$-linear abelian category, closed under images and direct summands, with
finite-dimensional Hom spaces and Krull--Schmidt decomposition. For any subset
$\Sigma\subseteq\operatorname{ob}(\ind\mathcal D)$, every nonempty object
quotient $(\ind\mathcal D)/[\Sigma]$ contains an object whose endomorphism
ring is a division algebra (and hence is equal to $k$ when $k$ is
algebraically closed).
\end{lemma}
\begin{proof}
Choose an object $X\in\operatorname{ob}(\ind\mathcal D)\setminus\Sigma$ of
minimal composition length. Every radical endomorphism $f\in\rad\End(X)$ has
a proper image $\im f$ of strictly smaller composition length. By hypothesis,
$\im f$ belongs to $\mathcal D$, and each of its indecomposable summands has
length smaller than that of $X$ and therefore belongs to $\Sigma$. Thus every
radical endomorphism of $X$ vanishes in $(\ind\mathcal D)/[\Sigma]$. On the
other hand, $1_X$ cannot factor through a finite direct sum of objects in
$\Sigma$ by the Krull--Schmidt property, so the endomorphism ring of $X$ in
$(\ind\mathcal D)/[\Sigma]$ is the residue division ring of the local ring
$\End(X)$.
\end{proof}

By Proposition~\ref{prop:projective-evaluation}(1), $\add\C\simeq\modu A_0$
is a finite-length abelian category, so Lemma~\ref{lem:deletion-brick} shows
that $\C$ satisfies the brick hypothesis. Since every object quotient of
$\C_T$ is also an object quotient of $\C$, the periodic mesh category $\C_T$
satisfies the brick hypothesis as well. Using Proposition~\ref{prop:periodic-tree-model}
and Proposition~\ref{prop:periodic-localization}, we prove the following
vanishing theorem in Appendix~\ref{app:mesh-calculations}.

\begin{proposition}
\label{prop:mesh-correction-spaces}
Let $C=C(Q,c)$ be a finite-dimensional scalar mesh category on a finite
translation quiver $Q$ with at most one arrow between each ordered pair of
vertices. If $\charac k\ne2$ and $C$ satisfies the brick hypothesis, then the
mesh derivative $D_m(C)$ is surjective for every $m\geq2$.
\end{proposition}

\subsection{Successive correction of irreducible representatives}
\label{subsec:successive-correction}

When $D_m$ is surjective for all $m\geq2$, a standard degree-by-degree
perturbation argument (compare Riedtmann \cite[Section~3.3]{Riedtmann1984})
shows that any filtered category whose associated graded category is $C$ is
already isomorphic to $C$.

\begin{proposition}
\label{prop:filtered-normal-form}
Let $C=kQ/(r_1,\ldots,r_t)$ be a finite-dimensional path-length graded
algebra over an arbitrary field $k$, where $Q$ is a finite quiver and each
relation $r_i$ is homogeneous of degree $2$. For each $q\geq3$, let
$D_{q-1}$ denote the linear map sending a tuple $(f_a)_a$ of homogeneous
elements $f_a\in C_{q-1}$ parallel to the arrows $a$ of $Q$ to the tuple in
$C_q$ obtained by substituting $a\mapsto a+f_a$ into $(r_1,\ldots,r_t)$ and
retaining the linear terms in $f$. If $D_{q-1}$ is surjective for every
$q\geq3$, then any finite-dimensional filtered algebra or category $E$
equipped with a graded isomorphism $\gr E\cong C$ is isomorphic to $C$ via a
filtered isomorphism inducing the given identification on $\gr E$.
\end{proposition}
\begin{proof}
Let $F^\bullet E$ be the finite decreasing filtration of $E$. Lift the vertex
idempotents of $C$ to a complete set of orthogonal idempotents of $E$ (taking
the identity morphisms when $E$ is a category), and choose initial lifts in
$F^1 E$ of the arrows $a$ in the corresponding corners of $E$. Their
quadratic relation errors $r_i$ then lie in $F^3 E$.

Suppose by induction that we have chosen arrow lifts whose relation errors
lie in $F^q E$ for some $q\geq3$. Passing to $F^q E/F^{q+1}E\cong C_q$, the
leading classes of these relation errors lie in the image of $D_{q-1}$ by
hypothesis. Subtracting a parallel correction of degree $q-1$ from each arrow
lift cancels the degree-$q$ part of each relation error, while the terms
quadratic in the correction have degree at least $2(q-1)>q$; thus the new
relation errors lie in $F^{q+1}E$. Because the filtration $F^\bullet E$ is
finite, this process terminates after finitely many steps, yielding arrow
lifts in $E$ that satisfy the relations $r_1,\ldots,r_t$ on the nose. The
induced homomorphism $C\to E$ has associated graded map equal to the identity
on $C$ and is therefore an isomorphism.
\end{proof}

\begin{corollary}
\label{cor:away-from-two-basis}
Let $A$ be a basic finite-dimensional representation-finite algebra over an
algebraically closed field $k$ with $\charac k\ne2$.
\begin{enumerate}[label={\upshape(\arabic*)}]
\item There is an isomorphism $A\cong k\ray A$, and $A$ is standard.
\item The algebra $A$ admits a normed multiplicative basis consisting of one
representative from each nonzero ray of $A$, with the primitive idempotents
representing the identity rays.
\end{enumerate}
\end{corollary}
\begin{proof}
By Proposition~\ref{prop:graded-realization}(3) and
Lemma~\ref{lem:deletion-brick}, the graded category $\C=\gr_\JJ\ind(\modu A)$
is a finite-dimensional scalar mesh category $k(\Gamma_A)$ satisfying the
brick hypothesis. By Proposition~\ref{prop:mesh-correction-spaces}, its mesh
derivative $D_{q-1}(\C)$ is surjective for every $q\geq3$. Applying
Proposition~\ref{prop:filtered-normal-form} to $E=\ind(\modu A)$ with its
radical filtration $F^n E=\JJ^n$ yields a filtered category isomorphism
$\ind(\modu A)\cong\C\cong k(\Gamma_A)$ fixing the objects, which proves that
$A$ is standard. Restricting this isomorphism to the projective vertices and
taking opposite algebras gives $A\cong A_0$, and combining this with
Proposition~\ref{prop:scalar-normalization} gives
$A\cong A_0\cong k\ray A_0\cong k\ray A$. Finally, transporting the normed
multiplicative ray basis of $A_0$ across this isomorphism yields (2).
\end{proof}

Combining Corollary~\ref{cor:away-from-two-basis} with
Lemma~\ref{lem:nonbasic-basis} completes the proof of Theorem~A outside
characteristic two, as well as Theorem~B. All periodic mesh calculations
required for this branch are given in Appendix~\ref{app:periodic-mesh}; none
of the results of Parts~\ref{part:mesh-deformations}
and~\ref{part:characteristic-two} are needed when $\charac k\ne2$.


\clearpage
\part{Deformations in characteristic two}
\label{part:mesh-deformations}

In this part we assume $\charac k=2$. Viewing $A$ as a positive filtered
deformation of its standard form $A_0$, we show that first-order
deformations of $A_0$ are governed by residual mesh parameters on the
$\tau$-periodic components of $\Gamma_A$, and we reduce Theorem~A to the
construction of homogeneous Hochschild two-cocycles supported on zero
products (Proposition~\ref{prop:family-criterion}).

\section{Residual mesh parameters and deformation of module categories}
\label{sec:positive-deformations}

From now on we assume that $A$ is basic and $\charac k=2$, retaining the
graded category $\C=\gr_\JJ\ind(\modu A)$ and its standard form $A_0$ from
Part~\ref{part:graded-form}. In characteristic two, the mesh derivative
$D_m(\C)$ need not be surjective, so higher-degree error terms in the mesh
relations cannot in general be eliminated by adjusting the irreducible
morphisms. We begin by describing the cokernel of $D_m(\C)$ on the
$\tau$-periodic components of $\Gamma_A$. We then introduce positive
filtered deformations of $A_0$—of which $A$ itself is the primary
example—and establish internal-degree bounds on $\Ext$ groups that allow us
to lift graded $A_0$-modules and their morphisms to any positive filtered
deformation. In Section~\ref{sec:weights-and-parameters} we use this module
lifting to embed the positive-degree Hochschild cohomology of $A_0$ into the
cokernel of the mesh derivative.

\subsection{The remaining linear correction spaces}

By Proposition~\ref{prop:periodic-localization}, the cokernel of the mesh
derivative $D_m(\C)$ of \eqref{eq:mesh-derivative} is isomorphic to
$\operatorname{coker}D_m(\C_T)$ on the stable subquiver of $\tau$-periodic
vertices. Evaluating $\operatorname{coker}D_m$ on the periodic tree models
of Proposition~\ref{prop:periodic-tree-model} in characteristic two yields
the following classification, proved in
Appendix~\ref{app:characteristic-two-calculations} (where periodic
components of type $A$ are shown to have vanishing cokernel in
Appendix~\ref{app:char2-line-models}).

\begin{proposition}
\label{prop:residual-correction-spaces}
Let $C=C(Q,c)$ be a finite-dimensional scalar mesh category on a finite
translation quiver $Q$ with at most one arrow between each ordered pair of
vertices. Assume that $\charac k=2$ and that $C$ satisfies the brick
hypothesis of Subsection~\ref{subsec:brick-hypothesis}. Call a periodic
component of $Q$ containing at least one arrow \emph{residual} if
$\bigoplus_{m\geq2}\operatorname{coker}D_m(C)$ has a nonzero class supported
on that component, and let $I$ index the residual components of $Q$.
\begin{enumerate}[label={\upshape(\arabic*)}]
\item The $i$th residual component ($i\in I$) is isomorphic to
$\Z\Delta/\langle\tau^{s_i}\rho\rangle$, where $\Delta$ is a Dynkin tree of
type $D$ or $E$ and $s_i\geq1$, as in
Proposition~\ref{prop:periodic-tree-model}(2).
\item For each $i\in I$, there exist a vertex $v_i$ of the $i$th residual
component and a path $p_i:\tau v_i\to v_i$ of length $d_i=2s_i+2$ such that,
viewing $p_i$ as an element of $V_{d_i}(C)$ supported in the row of $v_i$
and writing $[p_i]$ for its class in $\operatorname{coker}D_{d_i-1}(C)$, we
have
\[
 \bigoplus_{m\geq2}\operatorname{coker}D_m(C)
                      =\bigoplus_{i\in I}k[p_i].
\]
In particular, $d_i=2s_i+2$ is the unique path degree in which the $i$th
residual component has a nonzero cokernel class.
\end{enumerate}
\end{proposition}

By Subsection~\ref{subsec:brick-hypothesis}, Proposition~\ref{prop:residual-correction-spaces}
applies directly to $C=\C$.

\subsection{Degrees of extensions between normalized modules}
\label{subsec:ext-degrees}

To relate filtered deformations of $A_0$ to the mesh cokernel classes
$[p_i]$, we will lift the graded indecomposable $A_0$-modules to any
positive filtered deformation $B$ of $A_0$.

\begin{definition}
\label{def:positive-deformation}
A \emph{positive filtered deformation} of a finite-dimensional nonnegatively
graded algebra $\Lambda=\bigoplus_{d\geq0}\Lambda_d$ is an algebra $B$
equipped with a finite decreasing multiplicative filtration $F^\bullet B$
and a specified graded algebra isomorphism $\gr_F B\cong\Lambda$.
\end{definition}

In particular, the categorical radical filtration
$F^n A(x,y)=\JJ^n(P_y,P_x)$ of Section~\ref{sec:projective-rays} makes the
original algebra $A$ a positive filtered deformation of its standard form
$A_0$.

In classical deformation theory \cite{Gerstenhaber,Lowen}, the obstructions
to lifting a graded $\Lambda$-module action to a positive filtered
deformation lie in $\Ext^2_\Lambda$ in positive internal degrees, while the
choices of module lifts and the obstructions to lifting morphisms lie in
$\Ext^1_\Lambda$ in positive internal degrees. Here, for graded right
$\Lambda$-modules $M$ and $N$, a homomorphism or extension has
\emph{internal degree} $t\in\Z$ if it is represented on a graded projective
resolution of $M$ by a cochain that shifts degrees by $+t$. Using the shift
convention $(N\langle a\rangle)_q=N_{q-a}$ (so that a degree-zero map
$M\langle a\rangle\to N\langle b\rangle$ has underlying internal degree
$a-b$), we have
\[
 \Ext^n_\Lambda(M,N)_t
   =\Ext^n_{\mathrm{gr}\text{-}\Lambda}(M,N\langle-t\rangle),
\]
where $\Ext_{\mathrm{gr}\text{-}\Lambda}$ is computed in the category of
graded right $\Lambda$-modules with degree-zero morphisms.

For $\Lambda=A_0$, enumerate the objects of $\C$ as $X_1,\ldots,X_h$, and
let $M_1,\ldots,M_h$ be the corresponding graded indecomposable right
$A_0$-modules obtained by projective evaluation, graded by
$(M_i)_d e_x=\C_d(P_x,X_i)$. We refer to $M_1,\ldots,M_h$ as the
\emph{degree-normalized} indecomposable $A_0$-modules; note that the
indecomposable projective $A_0$-modules $e_x A_0$ carry their natural
gradings under this normalization. Full faithfulness of projective
evaluation (Proposition~\ref{prop:projective-evaluation}(1)) immediately
gives
\begin{equation}
 \Hom_{A_0}(M_i,M_j)_r=0\quad(r<0),\qquad
 \Hom_{A_0}(M_i,M_j)_0=
 \begin{cases}k\cdot\id_{M_i},&\text{if }i=j,\\0,&\text{if }i\ne j.\end{cases}
 \label{eq:normalized-hom-degrees}
\end{equation}

\begin{lemma}
\label{lem:positive-ext-vanishing}
Let $A_0$ be the standard form of Construction~\ref{con:projective-algebra},
and let $M_1,\ldots,M_h$ be the degree-normalized indecomposable
$A_0$-modules defined above.
\begin{enumerate}[label={\upshape(\arabic*)}]
\item Every finite-dimensional graded right $A_0$-module is isomorphic to a
direct sum of degree shifts of the modules $M_i$.
\item For all $i,j$ and all $n\geq1$,
\[
                 \Ext^n_{A_0}(M_i,M_j)_t=0\qquad\text{for all }t>-2n.
\]
\end{enumerate}
\end{lemma}
\begin{proof}
Let $X$ be a finite-dimensional graded right $A_0$-module that is
indecomposable as a graded module. Decomposing $X$ as an ungraded
$A_0$-module into copies of the indecomposable modules $M_i$, let
$p_a:X\to M_{i_a}$ and $j_a:M_{i_a}\to X$ be the projection and inclusion
maps of the summands, so that $1_X=\sum_a p_a j_a$ in path order. Passing to
the homogeneous component of degree zero gives
\[
                   1_X=\sum_{a,r}(p_a)_r(j_a)_{-r}
\]
in the local ring $\End_{\mathrm{gr}\text{-}A_0}(X)$. Some single term
$(p_a)_r(j_a)_{-r}$ is therefore invertible, exhibiting $X$ as a graded
direct summand of $M_{i_a}\langle-r\rangle$. Since $M_{i_a}$ is already
indecomposable as an ungraded module, $X\cong M_{i_a}\langle-r\rangle$,
proving (1).

We prove (2) first for $n=1$. The vanishing is trivial if $M_i$ is
projective. If $X_i$ is nonprojective, Proposition~\ref{prop:graded-realization}(1)
and Lemma~\ref{lem:mesh-simple-tops} provide an exact sequence of
representable functors
\[
 0\longrightarrow\C(-,\tau X_i)\longrightarrow
 \bigoplus_{Y\to X_i}\C(-,Y)\longrightarrow\C(-,X_i)
\]
whose cokernel is the simple functor $\C_0(-,X_i)$. Restricting these
functors to the projective objects $P_x$ annihilates the simple cokernel
(since $X_i$ is nonprojective) and yields a short exact sequence of graded
$A_0$-modules with degree-one maps. Under the equivalence
$\add\C\simeq\modu A_0$, this short exact sequence is the almost split
sequence ending at $M_i$; inserting degree shifts so that its maps become
homogeneous of degree zero writes it as
\[
 0\longrightarrow M_{\tau i}\langle2\rangle\longrightarrow
 \bigoplus_{Y\to X_i}M_Y\langle1\rangle\longrightarrow M_i
 \longrightarrow0.
\]
Now suppose for a contradiction that
$0\to M_j\langle-t\rangle\to E\xrightarrow{v}M_i\to0$ is a nonsplit graded
extension with $t>-2$. It remains nonsplit as an ungraded sequence (since any
ungraded section would have a degree-zero component that is a graded
section). By the lifting property of the right almost split map, $v$ factors
through $\bigoplus_{Y\to X_i}M_Y\langle1\rangle\to M_i$, and taking the
degree-zero component produces a graded factorization. The induced graded
map on kernels $M_j\langle-t\rangle\to M_{\tau i}\langle2\rangle$ has
underlying internal degree $-t-2<0$ and therefore vanishes by
\eqref{eq:normalized-hom-degrees}. Consequently, the factorization descends
to a section of the almost split sequence ending at $M_i$, a contradiction.
Thus $\Ext^1_{A_0}(M_i,M_j)_t=0$ for all $t>-2$.

For $n\geq2$, observe that in a minimal graded projective presentation of a
nonprojective module $M_i$, every projective summand $P\langle a\rangle$ of
the projective cover has $a\geq1$ because its nonzero map to $M_i$ lies in
positive degree; similarly, every indecomposable summand $Z\langle b\rangle$
of the first syzygy $\Omega^1 M_i$ maps into $P\langle a\rangle$ via a
radical morphism of positive degree, so $b-a\geq1$ and hence $b\geq2$. By
induction on $r$, every indecomposable summand of the $r$th syzygy
$\Omega^r M_i$ has the form $Z\langle b\rangle$ with $b\geq 2r$. Dimension
shifting then gives
\[
 \Ext^n_{A_0}(M_i,M_j)_t
 =\Ext^1_{A_0}(\Omega^{n-1}M_i,M_j)_t=0
 \qquad\text{for all }t>-2n,
\]
because each indecomposable summand $Z\langle b\rangle$ of
$\Omega^{n-1}M_i$ has $b\geq 2(n-1)$ and thus contributes
$\Ext^1_{A_0}(Z,M_j)_{t+b}=0$ since $t+b>-2$.
\end{proof}

\subsection{Lifting actions and morphisms}

Before lifting modules to a positive filtered deformation $B$, we place $B$
in a split relative normal form over its degree-zero semisimple subalgebra.

For an algebra $\Lambda$ over a commutative ring $R$ and a unital subalgebra
$S\subseteq\Lambda$, the space of \emph{normalized $S$-relative Hochschild
$n$-cochains} ($n\geq1$) is
\begin{equation}
 C^n_S(\Lambda,\Lambda)
   =\Hom_{S\text{-}S}\bigl((\Lambda/S)^{\otimes_S n},\Lambda\bigr),
 \label{eq:normalized-relative-cochains}
\end{equation}
with $C^0_S(\Lambda,\Lambda)=\{a\in\Lambda:as=sa\text{ for all }s\in S\}$
in degree zero and the standard Hochschild differential (in particular,
$(df)(a,b)=af(b)-f(ab)+f(a)b$ for a $1$-cochain $f$). Viewed as an
$R$-multilinear map on $\Lambda^n$, a normalized $S$-relative $n$-cochain is
$S$-bilinear in the outer variables, $S$-balanced across each adjacent pair
of inputs, and vanishes whenever any input belongs to $S$. We will use these
cochains over fields and over $\Z/4$.

\begin{lemma}
\label{lem:split-positive-deformation}
Let $k$ be a field, let $\Lambda=\bigoplus_{d\geq0}\Lambda_d$ be a
finite-dimensional nonnegatively graded $k$-algebra whose degree-zero part
$S=\Lambda_0$ is split semisimple, and let $B$ be a positive filtered
deformation of $\Lambda$. Then $S$ lifts to a unital subalgebra of $B$, and
the filtration $F^\bullet B$ splits as a filtration of $S$-bimodules. Under
the resulting filtered $S$-bimodule identification $B\cong\Lambda$ (inducing
the given identification $\gr_F B=\Lambda$), the multiplication of $B$ takes
the form $\mu_0+\sum_{q>0}\mu_q$, where $\mu_0$ is the multiplication of
$\Lambda$ and each $\mu_q\in C^2_S(\Lambda,\Lambda)$ is a normalized
$S$-relative $2$-cochain raising internal degree by $q$.
\end{lemma}
\begin{proof}
Set $\mathfrak n=F^1 B$, so that $\mathfrak n$ is a nilpotent two-sided
ideal of $B$ with $B/\mathfrak n\cong S$. Because idempotents can be lifted
modulo any nilpotent ideal—iterating $a\mapsto a+(1-2a)(a^2-a)$ across
successive square-zero quotients—and orthogonal idempotents lift
successively in complementary corners, we may lift all diagonal matrix units
of the matrix blocks of $S$ to a complete set of orthogonal idempotents in
$B$.

Within a single matrix block $M_n(k)$ of $S$, let $p_1,\ldots,p_n\in B$ be
the lifted diagonal idempotents, and choose elements $u_i\in p_i B p_1$ and
$v_i\in p_1 B p_i$ lifting the elementary matrix units $E_{i1}$ and $E_{1i}$
(with $u_1=v_1=p_1$). Since $v_i u_i\equiv p_1\pmod{\mathfrak n}$, the
element $v_i u_i$ is invertible in $p_1 B p_1$; replacing $v_i$ by
$(v_i u_i)^{-1}v_i$ ensures that $v_i u_i=p_1$. Then $u_i v_i$ is an
idempotent in $p_i B p_i$ congruent to $p_i$ modulo $\mathfrak n$, so
$p_i-u_i v_i$ is an idempotent in $\mathfrak n$ and therefore vanishes. Thus
$u_i v_i=p_i$, and the products $u_i v_j$ satisfy the matrix-unit relations
of $M_n(k)$. Since idempotents in distinct matrix blocks are orthogonal,
these matrix units define a unital algebra embedding $S\hookrightarrow B$
lifting $B/\mathfrak n\cong S$.

Finally, because $S\otimes_k S^{\op}$ is a finite product of matrix algebras
over $k$, it is semisimple, so every short exact sequence of $S$-bimodules
splits. Choosing $S$-bimodule splittings $F^d B/F^{d+1}B\hookrightarrow F^d B$
(using the subalgebra $S$ itself in degree $d=0$) identifies $B$ with
$\Lambda$ as a filtered $S$-bimodule inducing the prescribed isomorphism
$\gr_F B\cong\Lambda$. Under this identification, the multiplication of $B$
is $S$-bilinear, restricts to the multiplication of $S$ on $S$, and raises
grading degree by at least $0$, with degree-zero component equal to $\mu_0$.
\end{proof}

\begin{proposition}
\label{prop:positive-module-lifting}
Let $k$ be a field, let $\Lambda$ be a finite-dimensional nonnegatively
graded $k$-algebra with split semisimple degree-zero part $S=\Lambda_0$, and
let $M_1,\ldots,M_h$ be finite-dimensional graded right $\Lambda$-modules
satisfying
\[
 \Ext^2_\Lambda(M_i,M_i)_q=0,\qquad
 \Ext^1_\Lambda(M_i,M_j)_q=0\qquad\text{for all }i,j\text{ and all }q>0.
\]
Let $B$ be a positive filtered deformation of $\Lambda$.
\begin{enumerate}[label={\upshape(\arabic*)}]
\item Each $M_i$ admits a filtered right $B$-module lift $M_i^B$ with
$\gr M_i^B=M_i$, unique up to a filtered $B$-module isomorphism inducing the
identity on $M_i$.
\item Every homogeneous $\Lambda$-module homomorphism $M_i\to M_j$ of degree
$r\geq0$ lifts to a filtered $B$-module homomorphism $M_i^B\to M_j^B$ in
filtration degree $r$ with that prescribed leading term. Consequently, the
associated graded Hom spaces in nonnegative degrees between the lifted
modules coincide with the original graded Hom spaces.
\item Suppose further that $\Hom_\Lambda(M_i,M_j)_r=0$ for all $r<0$ and
$\Hom_\Lambda(M_i,M_j)_0=\delta_{ij}k$. Then the full subcategory $E$ of
$\modu B$ on $M_1^B,\ldots,M_h^B$, filtered by map degree, has associated
graded category equal to the graded Hom category on $M_1,\ldots,M_h$; if the
latter is generated in degree one, then the map-degree filtration on $E$
coincides with its categorical radical filtration. Moreover, if
$M_i=e_x\Lambda$ for a primitive idempotent $e_x\in S$, we may take
$M_i^B=e_x B$.
\end{enumerate}
\end{proposition}
\begin{proof}
Identify $B$ as a filtered $S$-bimodule with $\Lambda$ via
Lemma~\ref{lem:split-positive-deformation}, so that the multiplication $a*b$
of $B$ has the form $\mu_0(a,b)+\sum_{q>0}\mu_q(a,b)$. Because $S$ is
semisimple, any filtered right $B$-module lift of $M_i$ can likewise be
identified with $M_i$ as a filtered right $S$-module, so that its $B$-action
takes the form $m\cdot a=\nu_0(m,a)+\sum_{q>0}\nu_q(m,a)$, where $\nu_0$ is
the $\Lambda$-action on $M_i$ and each $\nu_q:M_i\otimes_S(\Lambda/S)\to M_i$
is a right $S$-linear map raising internal degree by $q$. Recall that the
normalized $S$-relative bar resolution of a graded right $\Lambda$-module $M$
has degree-$n$ term
\[
           M\otimes_S(\Lambda/S)^{\otimes_S n}\otimes_S \Lambda,
\]
which is a graded projective right $\Lambda$-module because $S$ is
semisimple; applying $\Hom_\Lambda(-,N)$ computes $\Ext^n_\Lambda(M,N)$ via
normalized $S$-relative bar cochains.

Suppose by induction on $t\geq1$ that we have chosen $\nu_1,\ldots,\nu_{t-1}$
so that the associativity identity $(m\cdot a)\cdot b=m\cdot(a*b)$ holds in
internal degrees less than $t$. A direct check using the associativity of
$*$ and the lower-degree action identities shows that the degree-$t$ failure
of $(m\cdot a)\cdot b=m\cdot(a*b)$ is a normalized $S$-relative bar
$2$-cocycle of internal degree $t$ with values in $\Hom_k(M_i,M_i)$. Since
$\Ext^2_\Lambda(M_i,M_i)_t=0$, this $2$-cocycle is the bar coboundary of a
normalized $1$-cochain $\nu_t$ of internal degree $t$, which cancels the
degree-$t$ error. Because $M_i$ is finite-dimensional, this induction
terminates after finitely many steps and yields a filtered $B$-module action
on $M_i$. Similarly, if two such actions agree in degrees less than $t$,
their degree-$t$ difference is a bar $1$-cocycle of internal degree $t$; by
$\Ext^1_\Lambda(M_i,M_i)_t=0$, it is the coboundary of an $S$-linear
endomorphism $h_t$ of $M_i$ of degree $t$, so conjugating by the invertible
filtered automorphism $1+h_t$ brings the two actions into agreement through
degree $t$. This proves (1).

For (2), let $f_0:M_i\to M_j$ be homogeneous of degree $r\geq0$, and suppose
we have constructed $f=f_0+f_1+\cdots+f_{t-1}$ (with $f_s$ of degree $r+s$)
commuting with the $B$-actions in degrees less than $r+t$. The degree-$(r+t)$
error in $f(m\cdot a)=f(m)\cdot a$ is a bar $1$-cocycle in
$C^1_S(M_i,M_j)$ of internal degree $r+t>0$. Since
$\Ext^1_\Lambda(M_i,M_j)_{r+t}=0$, this error is the coboundary of a map
$f_t:M_i\to M_j$ of degree $r+t$, allowing us to continue the induction to
obtain a filtered $B$-linear lift of $f_0$. Since the leading homogeneous
term of any nonzero filtered $B$-module homomorphism $M_i^B\to M_j^B$ is a
nonzero homogeneous $\Lambda$-module homomorphism $M_i\to M_j$, this proves
(2).

Under the hypotheses of (3), there are no nonzero homomorphisms of negative
degree, so the associated graded category of $E$ is the full graded Hom
category on $M_1,\ldots,M_h$. The positive-degree filtration ideal $F^1 E$
is nilpotent and $E/F^1 E\cong\prod_i k$, so $F^1 E$ is the categorical
radical $\JJ(E)$ of $E$ (and in particular the lifted modules $M_i^B$ are
indecomposable and pairwise nonisomorphic). If $\gr_F E$ is generated in
degree one, then each $F^q E/F^{q+1}E$ is spanned by products of $q$
elements of $F^1 E/F^2 E$, so descending induction on $q$ gives
$F^q E=\JJ(E)^q$ for all $q\geq1$. Finally, when $M_i=e_x\Lambda$, the lifted
projective module $e_x B$ has associated graded module $e_x\Lambda$, so by
the uniqueness in (1) we may take $M_i^B=e_x B$.
\end{proof}

Applying Proposition~\ref{prop:positive-module-lifting} to $\Lambda=A_0$ and
its degree-normalized indecomposable modules $M_1,\ldots,M_h$ (whose Hom and
$\Ext$ groups satisfy the required bounds by \eqref{eq:normalized-hom-degrees}
and Lemma~\ref{lem:positive-ext-vanishing}(2)) shows that for every positive
filtered deformation $B$ of $A_0$, the lifted modules $M_i^B$ form a
filtered category with associated graded category $\C$, and its projective
corner recovers $\End_B(\bigoplus_x e_x B)^{\op}=B$.

In Section~\ref{sec:weights-and-parameters} and
Section~\ref{sec:integral-deformations} we also need the first-order version
of this lifting calculation over $R[\varepsilon]/(\varepsilon^2)$, where the
coefficient ring $R$ is either a field or $\Z/4$.

\begin{lemma}
\label{lem:square-zero-selected-lifts}
Let $R$ be a field or $\Z/4$, let $\Lambda$ be a finitely generated free
nonnegatively graded $R$-algebra with $\Lambda_0=S=R^n$, and let
$M_1,\ldots,M_h$ be finitely generated graded right $\Lambda$-modules that
are free over $R$. Suppose that each graded module $\Hom_\Lambda(M_i,M_j)$ is
free over $R$ and concentrated in nonnegative degrees, and that
\[
 \Ext^2_\Lambda(M_i,M_i)_t=0,\qquad
 \Ext^1_\Lambda(M_i,M_j)_t=0\qquad\text{for all }i,j\text{ and all }t>0.
\]
Let $\mu\in C^2_S(\Lambda,\Lambda)$ be a normalized $S$-relative Hochschild
$2$-cocycle of internal degree $q>0$, and form the graded first-order
deformation
$\Lambda_\mu=(\Lambda\otimes_R R[\varepsilon]/(\varepsilon^2),m+\varepsilon\mu)$
with $\deg\varepsilon=-q$. Then:
\begin{enumerate}[label={\upshape(\arabic*)}]
\item Each $M_i$ lifts to a graded right $\Lambda_\mu$-module $M_i^\mu$ free
over $R[\varepsilon]/(\varepsilon^2)$ (taking $M_i^\mu=e_x\Lambda_\mu$ when
$M_i=e_x\Lambda$), uniquely up to a degree-zero $\Lambda_\mu$-module
isomorphism reducing to the identity modulo $\varepsilon$.
\item Every homogeneous $\Lambda$-linear map $f:M_i\to M_j$ lifts to a
homogeneous $\Lambda_\mu$-linear map $M_i^\mu\to M_j^\mu$ of the same degree,
and for any homogeneous $R$-basis of $\Hom_\Lambda(M_i,M_j)$, any chosen
homogeneous lifts form an $R[\varepsilon]/(\varepsilon^2)$-basis of
$\Hom_{\Lambda_\mu}(M_i^\mu,M_j^\mu)$.
\item If $\Lambda$ and the $M_i$ carry a compatible grading by an abelian
group $L$ and $\mu$ is homogeneous of $L$-degree $\chi$, then assigning
$\varepsilon$ the $L$-degree $-\chi$ allows all module and morphism lifts to
be chosen homogeneous for the $L$-grading as well.
\end{enumerate}
\end{lemma}
\begin{proof}
Because the vertex components of $\Lambda$ and $M_i$ are direct summands of
finitely generated free $R$-modules, they are projective over $R$, so the
normalized $S$-relative bar resolution of $M_i$ consists of graded projective
$\Lambda$-modules and computes $\Ext^\bullet_\Lambda(M_i,M_j)$ over $R$ (including
when $R=\Z/4$). Writing a candidate lifted action on
$M_i\otimes_R R[\varepsilon]/(\varepsilon^2)$ as
$m\cdot a=ma+\varepsilon\nu_i(m,a)$ with $\nu_i\in C^1_S(M_i,M_i)_q$, the
associativity condition modulo $\varepsilon^2$ is
\[
 \nu_i(ma,b)-\nu_i(m,ab)+\nu_i(m,a)b=m\mu(a,b).
\]
Since $d\mu=0$, the right-hand side is a relative bar $2$-cocycle of internal
degree $q>0$, so the vanishing of $\Ext^2_\Lambda(M_i,M_i)_q$ guarantees the
existence of $\nu_i$. The difference of any two solutions is a relative bar
$1$-cocycle of degree $q$, which by $\Ext^1_\Lambda(M_i,M_i)_q=0$ is the
coboundary of a degree-$q$ endomorphism $h$ of $M_i$, so the two lifts are
identified by the automorphism $1+\varepsilon h$. This proves (1).

Similarly, a homogeneous map $f:M_i\to M_j$ of degree $r\geq0$ lifts to a
$\Lambda_\mu$-linear map $f+\varepsilon g$ if and only if
\[
 g(ma)-g(m)a=\nu_j(f(m),a)-f(\nu_i(m,a)).
\]
The right-hand side is a relative bar $1$-cocycle of internal degree
$r+q>0$, so the vanishing of $\Ext^1_\Lambda(M_i,M_j)_{r+q}$ provides the correction $g$.
If $F_1,\ldots,F_s$ are homogeneous lifts of a homogeneous $R$-basis
$f_1,\ldots,f_s$ of the entire module $\Hom_\Lambda(M_i,M_j)$, then
subtracting from any $\Phi\in\Hom_{\Lambda_\mu}(M_i^\mu,M_j^\mu)$ the $R$-linear
combination of $F_1,\ldots,F_s$ matching its reduction modulo $\varepsilon$
leaves an element of the form $\varepsilon H$ where $H\in\Hom_\Lambda(M_i,M_j)$
(when $\Phi$ is homogeneous of degree $r$, its reduction modulo $\varepsilon$
lies in degree $r$ and $H$ lies in degree $r+q$), which is therefore an
$R$-linear combination of $\varepsilon F_1,\ldots,\varepsilon F_s$. Conversely, if
$\sum_l(a_l+\varepsilon b_l)F_l=0$ with $a_l,b_l\in R$, reducing modulo
$\varepsilon$ gives $\sum_l a_l f_l=0$, so $a_l=0$; the remaining equation
$\varepsilon\sum_l b_l f_l=0$ in the free $R[\varepsilon]/(\varepsilon^2)$-module
$\Hom_R(M_i,M_j)\otimes_R R[\varepsilon]/(\varepsilon^2)$ then forces
$\sum_l b_l f_l=0$ and hence $b_l=0$. This proves (2). Finally, since the
bar differentials preserve the $L$-grading, projecting any solution onto its
homogeneous component of the required $L$-degree gives (3).
\end{proof}

\section{Parameter weights and the criterion for recovering the algebra}
\label{sec:weights-and-parameters}

We continue under the assumption $\charac k=2$, with $A_0=kP$ ($P=\ray A_0$)
and $\C=\gr_\JJ\ind(\modu A)\simeq\ind(\modu A_0)$. In this section we compute
the universal group of additive weights on the ray category $P$, show how
ray weights extend to mesh weights on $\C$, and use the first-order module
lifting of Section~\ref{sec:positive-deformations} to embed the
positive-degree Hochschild cohomology $\HH^2(A_0)_{>0}$ into the mesh
cokernel $\bigoplus_{m\geq2}\operatorname{coker}D_m(\C)$. Crucially, the
weight lattice $L=W(P)$ separates the residual cohomology classes attached to
distinct periodic components even when they share the same categorical
degree.

\subsection{Normalizing the scalar mesh coefficients}

By Lemma~\ref{lem:unit-mesh-coefficients}, we may assume from now on that all
scalar coefficients in the mesh relations of $\C$ are equal to $1$:
explicitly, rescaling each arrow $a$ of $\Gamma_A$ by $\mu_a\in k^*$ and
dividing the mesh relation at each nonprojective vertex $v$ by $\eta_v\in k^*$
normalizes the coefficients $c_{v,a}$ to $1$ whenever
\begin{equation}
 \mu_{\sigma a}\mu_a=\eta_v/c_{v,a}\qquad\text{for every arrow }a:u\to v.
 \label{eq:global-mesh-scales}
\end{equation}

\subsection{The integral group of ray weights}

We now establish an integral counterpart of the scalar weight trivialization
of Lemma~\ref{lem:scalar-object-values}. This result holds in arbitrary
characteristic (compare the fundamental-group and coefficient analysis in
Bretscher and Gabriel \cite[Sections~2.11--2.12]{BretscherGabriel}) and will
also be used in Part~\ref{part:characteristic-two} to lift modulo-two
weights on smaller ray tables.

\begin{definition}
\label{def:scalar-weight-group}
For a finite ray category $Q$, its \emph{scalar-weight group} is the abelian
group
\[
 W(Q)=\Z\langle[r]:r\ne0,1\rangle/
       \langle[rs]-[r]-[s]:rs\ne0\rangle,
\]
where the classes of identity morphisms are set to zero. Thus homomorphisms
$W(Q)\to M$ into an abelian group $M$ correspond bijectively to additive
weights on the nonzero morphisms of $Q$ (and, for $M=k^*$ written
multiplicatively, $\Hom(W(Q),k^*)$ is the group of multiplicative scalar
weights $a(rs)=a(r)a(s)$ for $rs\ne0$).
\end{definition}

\begin{proposition}
\label{prop:integral-ray-weights}
Let $A_0=kP$ ($P=\ray A_0$) for a connected basic representation-finite
algebra $A$ with $n$ primitive idempotents, let $G=\pi_1(K)$ be the free deck
group of Subsection~\ref{subsec:universal-cover}, and let $\mathcal L$ be the
category of lifted projectives. Let
$I_{\mathcal L}\subset\Z^{(\operatorname{ob}\mathcal L)}$ denote the subgroup
of formal integer combinations of objects of $\mathcal L$ whose coefficients
sum to zero.
\begin{enumerate}[label={\upshape(\arabic*)}]
\item Sending a lifted ray $r\in\mathcal L(x,y)$ to the difference
$e_y-e_x\in I_{\mathcal L}$ induces an isomorphism
$W(P)\cong(I_{\mathcal L})_G$ with the group of $G$-coinvariants of
$I_{\mathcal L}$. Under this identification, the coinvariant class
$\theta_g=[e_{gx}-e_x]$ is independent of $x\in\operatorname{ob}\mathcal L$,
and the map $g\mapsto\theta_g$ embeds the abelianization $G_{\mathrm{ab}}$
as a direct summand of $W(P)$ with free complement $\Z^{n-1}$:
\[
                       W(P)\cong\Z^{n-1}\oplus G_{\mathrm{ab}}.
\]
In particular, since $G=\pi_1(K)$ is finitely generated, $W(P)$ is a free
abelian group of finite rank.
\item Every additive weight on $\mathcal L$ with values in an abelian group
is a difference of object values.
\end{enumerate}
\end{proposition}
\begin{proof}
For any finite directed thin representation-finite category $T$, let $U(T)$
denote the abelian group generated by symbols $[r]$ for the nonzero
nonidentity morphism lines $r$ of $T$, subject to the relations
$[rr']=[r]+[r']$ whenever $rr'\ne0$. If $s$ is a source of $T$ and
$T'=T\setminus\{s\}$, let $R$ be the representation-finite poset incidence
category obtained from $T'$ and $\rad T(s,-)$ as in
Lemma~\ref{lem:source-basis-extension}. Because the tautological weight
$r\mapsto[r]\in U(T')$ is additive on $R$, Lemma~\ref{lem:poset-weights}
provides elements $u_y\in U(T')$ for $y\in R$, normalized to $u_{y_D}=0$ at a
chosen base object $y_D$ in each connected component $D\in\pi_0(R)$, such
that $[b_{yz}]=u_z-u_y$ for every nonzero product $(sy)b_{yz}=sz$. Since
the relations in $U(T)$ involving the new generators $[sy]$ are precisely
$[sy]+[b_{yz}]=[sz]$, substituting $[sy]=e_D+u_y$ gives a canonical direct-sum
decomposition
\[
 U(T)\cong U(T')\oplus\Z^{\pi_0(R)},
\]
and the dual decomposition holds when adjoining a sink. Passing to the direct
limit along the one-object exhaustion $\mathcal L_1\subset\mathcal L_2\subset\cdots$
of Lemma~\ref{lem:cover-exhaustion} shows that $U(\mathcal L)$ is a free
abelian group.

Now consider the $G$-equivariant homomorphism
$\partial:U(\mathcal L)\to I_{\mathcal L}$ sending $[r:x\to y]$ to
$e_y-e_x$. Because $\mathcal L$ is connected, $\partial$ is surjective. We
claim that $\partial$ is also injective. Since $U(\mathcal L)$ is free
abelian and $k^*$ is infinite (so for any nonzero integer $m$ there exists
$a\in k^*$ with $a^m\ne1$), any nonzero element
$z\in\ker\partial$ would be detected by some group homomorphism
$\varphi:U(\mathcal L)\to k^*$; however, by
Lemma~\ref{lem:scalar-object-values}, every multiplicative weight
$\varphi\in\Hom(U(\mathcal L),k^*)=\operatorname{Scal}(\mathcal L)$ has the
form $\varphi([r:x\to y])=\zeta_y/\zeta_x$ and therefore factors through
$\partial$, forcing $\varphi(z)=1$. Thus $\partial:U(\mathcal L)\xrightarrow{\sim}I_{\mathcal L}$
is a $G$-equivariant isomorphism, which immediately implies (2).

Because the nonzero rays of $P$ and their nonzero products are the $G$-orbits
of the nonzero morphism lines of $\mathcal L$ and their nonzero products,
$W(P)$ is the group of $G$-coinvariants $U(\mathcal L)_G\cong(I_{\mathcal L})_G$.
Fix reference lifts $x_0,x_1,\ldots,x_{n-1}\in\operatorname{ob}\mathcal L$ of
the $n$ projective vertices of $A_0$. In $(I_{\mathcal L})_G$, the relation
$e_{gy}-e_x=(e_{gy}-e_{gx})+(e_{gx}-e_x)\equiv(e_y-e_x)+(e_{gx}-e_x)$ shows
that $\theta_g=[e_{gx}-e_x]$ is independent of $x$ and satisfies
$\theta_{gh}=\theta_g+\theta_h$. Thus $(I_{\mathcal L})_G$ is generated by
the classes $z_v=[e_{x_v}-e_{x_0}]$ ($1\leq v\leq n-1$) and $\theta_g$
($g\in G$), and the homomorphism $(I_{\mathcal L})_G\to\Z^{n-1}\oplus G_{\mathrm{ab}}$
induced by sending the object $gx_v$ to $(e_v,[g])$ (with $e_0=0$) provides
the inverse isomorphism, proving (1).
\end{proof}

\subsection{Integer weights on rays and meshes}

We use integer weights on $P$ and on $\C$ to refine the categorical grading
and separate first-order deformations attached to different periodic
components.

\begin{definition}
\label{def:integer-weights}
Let $A_0=kP$ ($P=\ray A_0$) and $\C=\gr_\JJ\ind(\modu A)$.
\begin{enumerate}[label={\upshape(\arabic*)}]
\item An \emph{integer ray weight} is a function $h:P\setminus\{0\}\to\Z$
that vanishes on identity rays and satisfies $h(rs)=h(r)+h(s)$ whenever
$rs\ne0$. We write $H_{\Z}=\Hom_{\Z}(W(P),\Z)$ for the group of integer ray
weights, set $L=W(P)$ (called the \emph{weight lattice} of $A_0$), and write
$\ell(r)=[r]\in L$ for the weight class of a nonzero ray $r\in P$.
\item An \emph{integer mesh weight} on $\C$ is a pair $w=(w,\nu)$ assigning an
integer $w(a)\in\Z$ to each arrow $a$ of $\Gamma_A$ and an integer
$\nu_v\in\Z$ to each nonprojective vertex $v$ of $\Gamma_A$ such that
\[
                   w(\sigma a)+w(a)=\nu_v\qquad\text{for every arrow }a:u\to v.
\]
We write $H^{\mathrm{mesh}}_{\Z}$ for the abelian group of integer mesh
weights on $\C$.
\end{enumerate}
\end{definition}

Note that the categorical degree function $\deg_\JJ$ of Lemma~\ref{lem:ray-degree}
is an element of $H_{\Z}$. Because $L=W(P)$ is free abelian of finite rank by
Proposition~\ref{prop:integral-ray-weights}(1), biduality identifies $L$
with $\Hom_{\Z}(H_{\Z},\Z)$ via $\chi(h)=h(\chi)$ for $\chi\in L$ and
$h\in H_{\Z}$. The algebra $A_0=kP$ is naturally $L$-graded by placing each
ray $r\in P$ in degree $\ell(r)\in L$; a Hochschild $2$-cochain sending a
pair of rays $(r,s)$ to a multiple of $u$ shifts $L$-degree by
$\ell(u)-\ell(r)-\ell(s)\in L$.

\begin{lemma}
\label{lem:integer-weight-lifting}
For $A_0=kP$ and $\C=\gr_\JJ\ind(\modu A)$, assigning to each path in $\C$ the
sum of the weights of its arrows under a mesh weight $w\in H^{\mathrm{mesh}}_{\Z}$
induces a well-defined grading on $\C$ and restricts on the projective
corners to an integer ray weight in $H_{\Z}$. This restriction fits into a
short exact sequence of abelian groups
\[
 0\longrightarrow
 \{(\zeta_X)\in\Z^{\operatorname{ob}\C}:\zeta_P=0
                         \text{ for all projective }P\}
 \longrightarrow H^{\mathrm{mesh}}_{\Z}\longrightarrow H_{\Z}
 \longrightarrow0,
\]
where the left-hand map sends a tuple $(\zeta_X)$ to the coboundary mesh
weight $w(a)=\zeta_{\operatorname{target}(a)}-\zeta_{\operatorname{source}(a)}$
and $\nu_v=\zeta_v-\zeta_{\tau v}$.
\end{lemma}
\begin{proof}
For any $w\in H^{\mathrm{mesh}}_{\Z}$, every term $(\sigma a)a$ of the mesh
relation at a nonprojective vertex $v$ has total weight $\nu_v$, so the mesh
ideal of $\C$ is homogeneous with respect to path weight. Because each
nonzero homogeneous component $A_0(x,y)_d=\C_d(P_y,P_x)$ is one-dimensional
by Lemma~\ref{lem:ray-degree}, all paths representing a given ray $r\in P$
have the same $w$-weight, and Proposition~\ref{prop:projective-ray-product}(2)
shows that this weight is additive on nonzero ray products. Thus restriction
defines a homomorphism $H^{\mathrm{mesh}}_{\Z}\to H_{\Z}$.

To see that restriction is surjective, pull any $h\in H_{\Z}$ back to an
additive weight on the category $\mathcal L$ of lifted projectives. By
Proposition~\ref{prop:integral-ray-weights}(2), there exist integers $b_x$
for $x\in\operatorname{ob}\mathcal L$ such that $h(r)=b_y-b_x$ for every
nonzero morphism $r\in\mathcal L(x,y)$, and by $G$-invariance the difference
$\gamma(g)=b_x-b_{gx}$ is independent of $x\in\operatorname{ob}\mathcal L$
and homomorphic in $g\in G$. On the lifted projective vertices of
$\widetilde\C$, set $\zeta_{P_x}=-b_x$ (the sign change reflecting that
$\mathcal L$ is the opposite category), and extend $\zeta$ to all vertices of
$\widetilde\C$ by choosing arbitrary integers on one representative of each
$G$-orbit of nonprojective vertices and setting
$\zeta_{g\widetilde X}=\zeta_{\widetilde X}+\gamma(g)$. Then the differences
of $\zeta$ along lifted arrows and translation edges are $G$-invariant and
therefore descend to an integer mesh weight on $\C$ whose restriction to $P$
is $\zeta_{P_x}-\zeta_{P_y}=b_y-b_x=h(r)$.

Finally, if $w\in H^{\mathrm{mesh}}_{\Z}$ restricts to zero on $P$, pull $w$
back to a closed integer edge weight on the simply connected complex
$\widetilde K$ to write $w$ as the difference of vertex integers
$\zeta_{\widetilde X}$. Because $w$ vanishes on all morphisms between lifted
projectives and $\mathcal L$ is connected, $\zeta$ is constant on all lifted
projective vertices; subtracting this constant makes $\zeta_{P_x}=0$ for all
$x$. Since $\zeta_{g\widetilde X}-\zeta_{\widetilde X}$ is independent of
$\widetilde X$ and vanishes on lifted projectives, $\zeta$ is $G$-invariant
and descends to a vertex assignment $(\zeta_X)$ on $\C$ vanishing at the
projective vertices. Because every connected component of $\Gamma_A$ contains
at least one projective vertex, this coboundary map is injective.
\end{proof}

For each residual component $i\in I$ and its representative path
$p_i:\tau v_i\to v_i$ of length $d_i=2s_i+2$ from
Proposition~\ref{prop:residual-correction-spaces}(2), define the linear
functional $\chi_i$ on $H^{\mathrm{mesh}}_{\Z}$ by
\begin{equation}
 \chi_i(w)=w(p_i)-\nu_{v_i}.
 \label{eq:parameter-character}
\end{equation}
Because $p_i$ and every term of the mesh relation at $v_i$ share the same
endpoints $\tau v_i$ and $v_i$, any vertex coboundary $(\zeta_X)$ shifts
$w(p_i)$ and $\nu_{v_i}$ by the same amount $\zeta_{v_i}-\zeta_{\tau v_i}$.
By Lemma~\ref{lem:integer-weight-lifting}, $\chi_i(w)$ therefore depends only
on the restricted ray weight in $H_{\Z}$ and defines an element
$\chi_i\in\Hom_{\Z}(H_{\Z},\Z)=L$, which we call the \emph{residual
character} of the $i$th residual component.

More generally, each mesh weight $w\in H^{\mathrm{mesh}}_{\Z}$ grades
$V_{m+1}(\C)$ by assigning weight $w(p)-\nu_v$ to a path $p:\tau v\to v$ in
row $v$, and grades $U_m(\C)$ by assigning weight $w(p)-w(a)$ to a path $p$
parallel to an arrow $a$. The mesh derivative $D_m$ of
\eqref{eq:mesh-derivative} preserves this grading, so
$\operatorname{coker}D_m(\C)$ becomes an $L$-graded space in which the line
$k[p_i]$ is homogeneous of $L$-degree $\chi_i$ (and of positive internal
degree $\chi_i(\deg_\JJ)=d_i-2=2s_i>0$).

\begin{proposition}
\label{prop:primitive-parameter-characters}
Assume $\charac k=2$. In the weight lattice $L=W(P)$ of
Definition~\ref{def:integer-weights}, the residual characters
$\{\chi_i\}_{i\in I}$ defined by \eqref{eq:parameter-character} extend to a
$\Z$-basis of $L$.
\end{proposition}
\begin{proof}
For each $i\in I$, identify the $i$th residual component with
$\Z\Delta/\langle\tau^{s_i}\rho\rangle$ as in
Proposition~\ref{prop:residual-correction-spaces}(1), with $\Delta$ oriented
away from the root fixed by $\rho$. On the vertices $(x,j)$ of $\Z\Delta$,
set
\[
                   F_i(x,j)=\left\lfloor j/s_i\right\rfloor,
\]
and assign to every arrow and translation edge of $\Z\Delta$ the difference
of $F_i$ between its target and source. Because the deck generator
$(x,j)\mapsto(\rho x,j-s_i)$ subtracts $1$ from $F_i$ at every vertex, these
edge differences are invariant under $\langle\tau^{s_i}\rho\rangle$ and
descend to a closed integer edge weight on the $i$th residual component.
Assigning weight zero to all other periodic components of $\Gamma_A$ and
extending to a closed edge weight on the entire mesh complex $K$ via
Lemma~\ref{lem:periodic-weight-extension} produces an integer mesh weight
$w_i\in H^{\mathrm{mesh}}_{\Z}$ on $\C$; let $h_i\in H_{\Z}=\Hom_{\Z}(L,\Z)$
be its restriction to $P$.

Lift the path $p_i:\tau v_i\to v_i$ of length $2s_i+2$ to $\Z\Delta$ starting
at a vertex $(x,r)$ over $\tau v_i$. Its endpoint has the form
$(\rho^m x,r+1-ms_i)$ for some $m\in\Z$ and lies at the same distance from
the root of $\Delta$ as $x$. Since outward arrows in $\Z\Delta$ preserve the
second coordinate while inward arrows increase it by $1$, the path $p_i$ of
length $2s_i+2$ traverses an equal number $s_i+1$ of outward and inward
arrows and therefore ends at second coordinate $r+s_i+1$. Subtracting the
weight of the translation edge $t_{v_i}$ from $(x,r)$ to $(x,r+1)$ gives
\[
 h_i(\chi_i)=\chi_i(w_i)
 =\left\lfloor\frac{r+s_i+1}{s_i}\right\rfloor
       -\left\lfloor\frac{r+1}{s_i}\right\rfloor=1,
\]
whereas $h_i(\chi_j)=\chi_j(w_i)=0$ for $j\ne i$ because $w_i$ vanishes on
the $j$th periodic component. Consequently, the homomorphism
$(h_i)_{i\in I}:L\to\Z^I$ is a split surjection with section sending the
$i$th standard basis vector to $\chi_i$, so
$L=\bigoplus_{i\in I}\Z\chi_i\oplus\ker(h_i)_{i\in I}$. Since $\ker(h_i)_{i\in I}$
is a subgroup of the finitely generated free abelian group $L$, it is free
abelian, and any $\Z$-basis of it completes $\{\chi_i\}_{i\in I}$ to a
$\Z$-basis of $L$.
\end{proof}

\subsection{Hochschild cohomology and first-order mesh deformations}
\label{subsec:hochschild-restriction}

We now relate the Hochschild cohomology of $A_0$ to the mesh derivative
cokernel by lifting the degree-normalized indecomposable $A_0$-modules to
first-order deformations of $A_0$ over $k[\varepsilon]/(\varepsilon^2)$.
Because the degree-zero vertex subalgebra $S=\prod_x ke_x$ is separable over
$k$, the inclusion of the normalized $S$-relative cochain complex
$C^\bullet_S(A_0,A_0)$ of \eqref{eq:normalized-relative-cochains} into the
ordinary Hochschild cochain complex is a quasi-isomorphism
\cite[Theorem~1.2]{GerstenhaberSchack}. Both the internal categorical degree
and the $L$-grading on $A_0=kP$ induce gradings on $\HH^2(A_0)$; we write
$\HH^2(A_0)_q$ for the component of internal degree $q$,
$\HH^2(A_0)_{>0}=\bigoplus_{q>0}\HH^2(A_0)_q$, and $\HH^2(A_0)_\chi$ for the
component of $L$-weight $\chi\in L$. (In
Proposition~\ref{prop:full-hochschild-comparison} we will show that the
injection $\Psi$ below is in fact an isomorphism and $I_0=I$, though only
injectivity is needed to prove Theorem~A.)

\begin{proposition}
\label{prop:hochschild-mesh-restriction}
Let $A_0$ be the standard form of Construction~\ref{con:projective-algebra}
and $\C$ its graded mesh category. For every internal degree $q>0$, lifting
degree-normalized modules to first-order deformations induces an $L$-graded
injection
\[
 \Psi:\HH^2(A_0)_q\hookrightarrow\operatorname{coker}D_{q+1}(\C).
\]
In particular, when $\charac k=2$, there is a subset $I_0\subseteq I$ of the
residual components of Proposition~\ref{prop:residual-correction-spaces}
such that
\[
 \HH^2(A_0)_{>0}=\bigoplus_{i\in I_0}\HH^2(A_0)_{\chi_i},
 \qquad \dim_k\HH^2(A_0)_{\chi_i}=1\quad(i\in I_0).
\]
\end{proposition}
\begin{proof}
Represent a cohomology class in $\HH^2(A_0)_q$ by a normalized $S$-relative
$2$-cocycle $\mu\in C^2_S(A_0,A_0)_q$, and form the first-order deformation
$A_{0,\mu}=(A_0\otimes_k k[\varepsilon]/(\varepsilon^2),m+\varepsilon\mu)$
with $\deg\varepsilon=-q$. By Lemma~\ref{lem:positive-ext-vanishing}(2) and
Lemma~\ref{lem:square-zero-selected-lifts}, the degree-normalized
indecomposable $A_0$-modules $M_1,\ldots,M_h$ lift to graded
$A_{0,\mu}$-modules $M_i^\mu$ (with $M_i^\mu=e_x A_{0,\mu}$ at the projective
vertices), and each arrow $a$ of $\Gamma_A$ lifts to a degree-one morphism
$\tilde a$ in the Hom category $\C_\mu$ of the lifted modules. Because the
mesh relation $r_v=\sum_{a:u\to v}(\sigma a)a$ vanishes in $\C$, its evaluation
on the lifted arrows $\tilde a$ in $\C_\mu$ vanishes modulo $\varepsilon$; by
the freeness of the lifted Hom spaces over $k[\varepsilon]/(\varepsilon^2)$
(Lemma~\ref{lem:square-zero-selected-lifts}(2)), this mesh error has the form
$\varepsilon z_v$ for a unique element $z_v\in\C_{q+2}(\tau v,v)$ (since
$\deg\varepsilon=-q$ and $\deg r_v=2$). Modifying the arrow lifts by
$\tilde a\mapsto\tilde a+\varepsilon f_a$ with $f_a\in\C_{q+1}$ changes the
tuple $z=(z_v)_v$ by $D_{q+1}(f)$, while changing the module lifts $M_i^\mu$
by automorphisms reducing to the identity or replacing $\mu$ by a cohomologous
cocycle $\mu+dh$ (which corresponds to applying the algebra automorphism
$1+\varepsilon h$ of $A_{0,\mu}$) leaves the class $[z]\in\operatorname{coker}D_{q+1}(\C)$
unchanged. Thus $[\mu]\mapsto\Psi([\mu])=[z]$ is well-defined.

Because the first-order module-action and arrow-lifting equations in
Lemma~\ref{lem:square-zero-selected-lifts} are linear in $\mu$ and in the
first-order corrections, and any product of two $\varepsilon$-corrections
vanishes modulo $\varepsilon^2$, the map $\Psi$ is $k$-linear. To see that
$\Psi$ is injective, suppose $\Psi([\mu])=0$. Then we can adjust the arrow
lifts $\tilde a$ by $\varepsilon f_a$ so that every mesh relation $r_v(\tilde a)$
vanishes identically in $\C_\mu$. The assignment $a\mapsto\tilde a$ then
defines a graded $k[\varepsilon]/(\varepsilon^2)$-linear functor
$\C\otimes_k k[\varepsilon]/(\varepsilon^2)\to\C_\mu$ that reduces to the
identity modulo $\varepsilon$. By Lemma~\ref{lem:square-zero-selected-lifts}(2),
the morphism spaces of $\C_\mu$ are free over $k[\varepsilon]/(\varepsilon^2)$
with bases lifting those of $\C$, so by Nakayama's lemma over the local ring
$k[\varepsilon]/(\varepsilon^2)$ this functor is an isomorphism on every Hom
space. Restricting to the projective vertices $e_x A_{0,\mu}$ yields a
graded algebra isomorphism $A_0\otimes_k k[\varepsilon]/(\varepsilon^2)\xrightarrow{\sim}A_{0,\mu}$
reducing to the identity modulo $\varepsilon$, which means that $[\mu]=0$ in
$\HH^2(A_0)_q$.

Finally, choosing an $L$-homogeneous cocycle $\mu$ of weight $\chi\in L$ and
lifting the integer ray weights to mesh weights on $\C$ via
Lemma~\ref{lem:integer-weight-lifting}, Lemma~\ref{lem:square-zero-selected-lifts}(3)
shows that when $\varepsilon$ is assigned $L$-weight $-\chi$, the lifted
modules and arrows can be chosen $L$-homogeneous, so the resulting mesh
error $z_v$ has weight $\chi$ relative to row $v$. Thus $\Psi$ preserves
$L$-weights. Combining the injectivity of $\Psi$ with
Proposition~\ref{prop:residual-correction-spaces}(2) immediately gives the
stated decomposition of $\HH^2(A_0)_{>0}$ in characteristic two.
\end{proof}

\subsection{Finite normalization of an associative family}
\label{subsec:characteristic-two-construction}

We conclude Part~\ref{part:mesh-deformations} with a general deformation-theoretic
criterion for realizing a positive filtered deformation as a specialized
fibre of an associative polynomial family (compare Gerstenhaber
\cite{Gerstenhaber}). For an algebra family over $k[t_1,\ldots,t_n]$, we call
its specialization at a parameter vector $\lambda\in k^n$ a \emph{scalar
fibre}, and we call it a \emph{binary fibre} when every coordinate of
$\lambda$ belongs to $\{0,1\}$.

\begin{theorem}
\label{thm:finite-family-normalization}
Let $\Lambda=\bigoplus_{d=0}^N\Lambda_d$ be a finite-dimensional nonnegatively
graded algebra over a field $k$ whose degree-zero part $S=\Lambda_0$ is a
finite product of copies of $k$. Let $\mu_1,\ldots,\mu_n\in C^2_S(\Lambda,\Lambda)$
be normalized $S$-relative Hochschild $2$-cocycles of positive internal
degrees $q_1,\ldots,q_n>0$. Assume that their cohomology classes span
$\HH^2(\Lambda)_{>0}$ and that the deformed product
\[
                         m_t=m+\sum_{i=1}^n t_i\mu_i
\]
is associative over the polynomial ring $k[t_1,\ldots,t_n]$, where $m$ is
the multiplication of $\Lambda$. Then every positive filtered deformation of
$\Lambda$ is isomorphic to a scalar fibre $(\Lambda,m_\lambda)$ via a
filtered isomorphism that induces the identity on $\gr\Lambda=\Lambda$ and
fixes the lifted vertex idempotents.
\end{theorem}
\begin{proof}
By Lemma~\ref{lem:split-positive-deformation}, after lifting the vertex
idempotents and splitting the filtration as $S$-bimodules, the given positive
filtered deformation is defined on the underlying graded vector space
$\Lambda$ by an associative $S$-relative multiplication $\star$ such that
$\star-m$ raises internal degree by at least $1$. We construct a parameter
vector $\lambda\in k^n$ and a sequence of filtered automorphisms of $\Lambda$
by induction on the degree $q\geq1$, starting with $q=1$ and $\lambda=0$.

Suppose that $\star$ agrees with $m_\lambda$ in all internal degrees less
than $q$, where only those coordinates $\lambda_i$ with $q_i<q$ have been
assigned. Let $\varphi_q\in C^2_S(\Lambda,\Lambda)_q$ denote the homogeneous
component of $\star-m_\lambda$ of internal degree $q$. Subtracting the
associativity identity for $m_\lambda$ from that for $\star$ in internal
degree $q$, all cross-terms involving both $\varphi_q$ and lower-degree
positive components of $\star$ or $m_\lambda$ have internal degree strictly
greater than $q$ and therefore vanish, leaving the Hochschild cocycle
condition $d\varphi_q=0$. Because the classes $[\mu_i]$ span
$\HH^2(\Lambda)_{>0}$, there exist scalars $c_i\in k$ (for those indices $i$
with $q_i=q$) and a normalized $S$-relative $1$-cochain
$f_q\in C^1_S(\Lambda,\Lambda)_q$ such that
\[
                 \varphi_q=\sum_{q_i=q}c_i\mu_i+df_q.
\]
Setting $\lambda_i=c_i$ for all $i$ with $q_i=q$ keeps $m_\lambda$
associative by hypothesis, and pulling the multiplication $\star$ back along
the filtered $S$-bimodule automorphism $1-f_q$ subtracts
$(df_q)(a,b)=a f_q(b)-f_q(ab)+f_q(a)b$ from its degree-$q$ component while
leaving all components of degree less than $q$ unchanged. Thus the
transformed product $\star$ agrees with $m_\lambda$ in all internal degrees
at most $q$. Since $\Lambda$ is finite-dimensional, $C^2_S(\Lambda,\Lambda)_q=0$
for all $q>N$, so this process terminates after at most $N$ steps and yields
the desired filtered isomorphism with $(\Lambda,m_\lambda)$.
\end{proof}

Because the residual characters $\chi_i$ extend to a $\Z$-basis of the
weight lattice $L=W(P)$ by Proposition~\ref{prop:primitive-parameter-characters},
any nonzero coordinates of the parameter vector $\lambda$ can then be rescaled
to $1$ by a diagonal grading automorphism of $\Lambda$.

\begin{corollary}
\label{cor:binary-normalization}
In the setting of Theorem~\ref{thm:finite-family-normalization}, suppose
that $\Lambda$ also carries a compatible grading by a finite-rank free
abelian group $L$ with $S$ in degree zero, and that each cocycle $\mu_i$ is
homogeneous of $L$-degree $\chi_i\in L$. If $\{\chi_1,\ldots,\chi_n\}$
extends to a $\Z$-basis of $L$, then every positive filtered deformation of
$\Lambda$ is isomorphic to a binary fibre $(\Lambda,m_\lambda)$ with
$\lambda\in\{0,1\}^n$. More generally, any two parameter vectors in $k^n$
having the same zero coordinates define isomorphic fibres.
\end{corollary}
\begin{proof}
Every group homomorphism $a:L\to k^*$ defines a grading automorphism of
$(\Lambda,m)$ that multiplies elements of $L$-weight $\ell$ by $a(\ell)$.
Transporting the deformed product $m_\lambda$ along this automorphism
rescales the $i$th parameter according to
\begin{equation}
                       \lambda'_i=\lambda_i a(\chi_i).
 \label{eq:diagonal-algebra-rescaling}
\end{equation}
Because $\{\chi_1,\ldots,\chi_n\}$ extends to a $\Z$-basis of the free
abelian group $L$, we can prescribe $a(\chi_i)\in k^*$ arbitrarily for
$i=1,\ldots,n$; in particular, setting $a(\chi_i)=\lambda_i^{-1}$ whenever
$\lambda_i\ne0$ and $a(\chi_i)=1$ when $\lambda_i=0$ transforms $m_\lambda$
into a binary fibre.
\end{proof}

Specializing Theorem~\ref{thm:finite-family-normalization} and
Corollary~\ref{cor:binary-normalization} to $A_0=kP$ yields our reduction
criterion for Theorem~A in characteristic two.

\begin{proposition}
\label{prop:family-criterion}
Let $A$ be a basic representation-finite algebra over an algebraically closed
field $k$ of characteristic two, let $A_0=kP$ ($P=\ray A_0$), and let
$I_0\subseteq I$ index the nonzero spaces $\HH^2(A_0)_{\chi_i}$ in
Proposition~\ref{prop:hochschild-mesh-restriction}. Suppose that we have
constructed normalized $S$-relative $2$-cochains
$\mu_i\in C^2_S(A_0,A_0)$ ($i\in I_0$) satisfying:
\begin{enumerate}[label={\upshape(\arabic*)}]
\item each $\mu_i$ is homogeneous of $L$-weight $\chi_i$, and its cohomology
class in $\HH^2(A_0)_{\chi_i}$ is nonzero;
\item the deformed product $m_t=m+\sum_{i\in I_0}t_i\mu_i$ is associative
over $k[t_i:i\in I_0]$;
\item in every binary fibre $(A_0,m_\lambda)$ ($\lambda\in\{0,1\}^{I_0}$),
the basis of nonzero rays $P\setminus\{0\}$ is a normed multiplicative basis,
and in every scalar fibre $(\lambda\in k^{I_0})$ the map sending each
$r\in P\setminus\{0\}$ to its ray in $(A_0,m_\lambda)$ identifies the ray
category of $(A_0,m_\lambda)$ with $P$.
\end{enumerate}
Then $A$ is isomorphic to a binary fibre $(A_0,m_\lambda)$ via a filtered
isomorphism that transports $P\setminus\{0\}$ to a normed multiplicative ray
basis of $A$ and identifies $\ray A$ with $P$ via initial forms.
\end{proposition}
\begin{proof}
As noted after Definition~\ref{def:positive-deformation}, $A$ is a positive
filtered deformation of $A_0$. By Proposition~\ref{prop:hochschild-mesh-restriction},
the cohomology classes $[\mu_i]$ ($i\in I_0$) in (1) form a basis of
$\HH^2(A_0)_{>0}$. Theorem~\ref{thm:finite-family-normalization} therefore
provides a filtered isomorphism between $A$ and a scalar fibre
$(A_0,m_{\lambda'})$ inducing the identity on $\gr A=A_0$, and
Proposition~\ref{prop:primitive-parameter-characters} together with
Corollary~\ref{cor:binary-normalization} rescales $(A_0,m_{\lambda'})$ via a
diagonal grading automorphism on the ray basis $P\setminus\{0\}$ to a binary
fibre $(A_0,m_\lambda)$. Because this diagonal rescaling preserves each
homogeneous ray line in $A_0=\gr A$, the composite filtered isomorphism
$(A_0,m_\lambda)\xrightarrow{\sim}A$ sends the normed multiplicative basis
$P\setminus\{0\}$ of (3) to a normed multiplicative basis of $A$ consisting
of one representative from each nonzero ray, and identifies $\ray A$ with
$\ray(A_0,m_\lambda)=P$ via the initial-form bijection of
Lemma~\ref{lem:ray-degree}.
\end{proof}


\clearpage
\part{Construction of multiplicative bases in characteristic two}
\label{part:characteristic-two}

In this final part we construct the zero-supported Hochschild two-cocycles
required by Proposition~\ref{prop:family-criterion}. Since the multiplicative
basis $P=\ray A_0$ of $A_0$ has structure constants in $\{0,1\}$, it defines
models $A^{(2)}=\F_2P$ and $A^{(4)}=(\Z/4)P$ with
$A_0=k\otimes_{\F_2}A^{(2)}$. As we show in
Proposition~\ref{prop:integral-bockstein}, every cohomology class in
$\HH^2(A^{(2)})$ of positive internal degree can be represented via the
Bockstein connecting map by $(d\delta)/2$, where $\delta$ is a one-cochain
on $A^{(4)}$ lifting a derivation of $A^{(2)}$ and division by two identifies
$2(\Z/4)$ with $\F_2$. For the resulting two-cocycle $(d\delta)/2$ to vanish
on all nonzero basis products $rs\ne0$ in $P$, the lift $\delta$ must satisfy
the Leibniz equation
\[
 \delta(rs)=\delta(r)s+r\delta(s)\qquad\text{whenever }rs\ne0.
\]

These Leibniz conditions form a system of linear congruences for the
coefficients of $\delta$, which we decouple and solve one weight character at
a time. Section~\ref{sec:local-ray-geometry} and
Section~\ref{sec:primitive-ray-equations} establish local factorization
restrictions on rays and classify how a primitive character shifts endpoint
layers. Using these restrictions, Section~\ref{sec:integral-partial-lifts}
shows that the equations on two-sided rays reduce to the additive-weight
relations of a strictly smaller representation-finite ray algebra, while the
remaining one-sided equations organize into a bipartite graph whose
components are controlled by excluding four-point crowns. Induction on the
number of basis elements then yields integer lifts of all prescribed modulo-two
derivation coefficients. Finally, Section~\ref{sec:integral-deformations}
and Section~\ref{sec:binary-families} pass via the Bockstein map to
zero-supported cocycles over $k$, assemble them into a global associative
family of deformations, and recover the original algebra $A$ as a binary
fibre, completing the proof of Theorems~A and~B;
Section~\ref{sec:deformation-refinements} records further consequences for
module categories and parameter spaces.

\section{Local restrictions on zero products and returns}
\label{sec:local-ray-geometry}

Our first step in characteristic two is to constrain the products and factorizations of ray-basis elements in the standard algebra $A_0=kP$, including return paths that leave a vertex and come back to it with a nonzero value in the local ring. These local constraints play a role analogous to the analysis of endpoint actions in \cite[Sections~2.4--2.5]{BGRS}. Combined with the integral coefficient-lifting theorem of Section~\ref{sec:integral-partial-lifts} and the modulo-four comparison of Section~\ref{sec:integral-deformations}, they will allow us to construct cocycle representatives for every class in $\HH^2(A_0)_{>0}$ that vanish on all nonzero products of basis elements.

The local structural results in this section do not themselves depend on an inductive hypothesis. In Section~\ref{sec:integral-partial-lifts}, however, lifting their associated modulo-two coefficient equations to the integers will require an induction on the size of a weighted multiplication table. We record the definition of such tables and the main inductive statements before beginning the local analysis.

\begin{definition}
\label{def:weighted-table}
A \emph{weighted multiplication table} $T$ consists of a finite set of objects, identity elements $1_x$, and finitely many nonzero morphism elements (also called \emph{labels}), equipped with an associative partial composition under which the product of two composable elements is either an element of $T$ or zero, together with an additive weight $\ell\colon T\setminus\{0\}\to M$ valued in an abelian group $M$ and a homomorphism $d\colon M\to\Z$ satisfying the following conditions:
\begin{enumerate}[label={\upshape(\roman*)}]
\item $\ell(1_x)=0$ for every object $x$, and $d(\ell(r))>0$ for every nonidentity element $r\in T\setminus\{0\}$;
\item two nonzero elements of $T$ with the same source, target, and weight are equal.
\end{enumerate}
The \emph{linearization} $kT$ is the finite-dimensional $k$-linear category having the nonzero elements of $T$ as a basis, and its scalar weight group $W(T)$ is defined by the presentation of Definition~\ref{def:scalar-weight-group}.
\end{definition}

\begin{proposition}
\label{prop:table-weight-lifting}
Let $T$ be a weighted multiplication table whose linearization $kT$ is representation-finite over the algebraically closed field $k$. Every additive weight $T\to\Z/2$ lifts to an additive weight $T\to\Z$.
\end{proposition}

Proposition~\ref{prop:table-weight-lifting} will be proved in Subsection~\ref{subsec:main-theorems} by strong induction on the total number $n=\dim kT$ of nonzero elements of $T$ (including identities). Assuming the statement for all tables with fewer than $n$ nonzero elements, we first use it in characteristic two—applied strictly to the smaller centralizer table of Lemma~\ref{lem:centralizer-ray-algebra}—to construct integral solutions to the nonzero-product Leibniz equations for the standard form $A_0$ of any basic connected representation-finite algebra $A$ of dimension $n$ (Theorem~\ref{thm:integral-partial-lifting}). The resulting cocycles and their associative deformation families recover $A$ and identify its ray multiplication table with that of $A_0$ (Lemma~\ref{lem:recovery-with-lifting}). Specializing to $A=kT$ for a connected table $T$ of size $n$ then identifies $T$ with $\ray(kT)$ and equates $W(T)$ with the free abelian weight group of its standard form (Proposition~\ref{prop:integral-ray-weights}), so that every $\Z/2$-valued weight on $T$ lifts to $\Z$. Outside characteristic two, the corresponding recovery is already given by Corollary~\ref{cor:away-from-two-basis}, and disconnected tables reduce immediately to their connected components.

From this inductive step we also obtain the main structural consequences for representation-finite algebras in arbitrary characteristic, likewise proved in Subsection~\ref{subsec:main-theorems}.

\begin{proposition}
\label{prop:dimension-induction}
Let $A$ be a basic representation-finite algebra over $k$, in any characteristic, and let $A_0$ be the projective ray algebra of Construction~\ref{con:projective-algebra}. Then:
\begin{enumerate}[label={\upshape(\arabic*)}]
\item $A$ admits a normed multiplicative basis consisting of one representative of each nonzero ray;
\item the initial-form bijection of Lemma~\ref{lem:ray-degree} identifies the ray multiplication in $\ray A$ with that in $\ray A_0$;
\item the scalar weight group $W(\ray A)$ is free abelian.
\end{enumerate}
\end{proposition}

\begin{setting}
\label{set:projective-companion}
Throughout Part~\ref{part:characteristic-two} and Appendix~\ref{app:local-rays}, we let $A_0=kP$ denote the projective ray algebra of Construction~\ref{con:projective-algebra}, identified with $kP$ via Proposition~\ref{prop:scalar-normalization} and equipped with its normed homogeneous ray basis. For each vertex $x$, we write $R_x=k[q_x]/(q_x^{N_x})$ when $R_x\ne k$, where $q_x$ is the chosen homogeneous generator of $\rad R_x$, and we set $q_x=0$ and $N_x=1$ when $R_x=k$.
\end{setting}

We use the following terminology for the rays of $A_0=kP$. Note that an \emph{ordinary arrow} below means an indecomposable nonidentity ray of $A_0$ (whose vertices correspond to the indecomposable projectives of $A$), whereas the arrows of the Auslander--Reiten quiver $\Gamma_A$ represent irreducible morphisms among all indecomposable $A$-modules.

\begin{definition}
\label{def:generating-rays}
Let $A_0=kP$ satisfy Setting~\ref{set:projective-companion}.
\begin{enumerate}[label={\upshape(\arabic*)}]
\item A ray $r\colon x\to y$ is \emph{generating} if it generates the full corner bimodule $A_0(x,y)$ as an $(R_x,R_y)$-bimodule; otherwise $r$ is \emph{proper}.
\item Given a generating ray $f\colon x\to y$, we say that the corner bimodule $A_0(x,y)$ (or simply $f$) is \emph{cyclic at $x$} if $A_0(x,y)=R_xf$, and \emph{cyclic at $y$} if $A_0(x,y)=fR_y$. When $A_0(x,y)$ is cyclic at $x$, its rays are the elements $q_x^jf$, and we call the exponent $j$ the \emph{endpoint layer} of $q_x^jf$; the endpoint layer at $y$ is defined symmetrically when $A_0(x,y)$ is cyclic at $y$.
\item An \emph{ordinary arrow} is a nonidentity ray that cannot be factored as a product of two nonidentity rays, and an \emph{ordinary path} is a composable sequence of ordinary arrows. Within a corner or quotient category, these terms refer to factorizations inside that corner or quotient.
\item Let $x\ne y$, and let $f\colon x\to y$ and $h\colon y\to x$ be rays such that $fh=q_x^b\ne0$ for some integer $b\geq1$. We call $fh$ a \emph{return} at $x$ through $y$, $h$ a \emph{reverse} of $f$, and $b$ the \emph{exponent} of the return.
\end{enumerate}
\end{definition}

The \emph{endpoint socle} of a corner bimodule $A_0(x,y)$ is the subspace of vectors annihilated by both $\rad R_x$ on the left and $\rad R_y$ on the right. By Lemma~\ref{lem:endpoint-chains}, every nonzero corner bimodule $A_0(x,y)$ has a one-dimensional endpoint socle.

\begin{definition}[{\cite[Section~9.1]{BGRS}}]
\label{def:stem-product}
For a ray $r\colon x\to y$, let $U_r=R_xrR_y\subseteq A_0(x,y)$ denote the endpoint subbimodule it generates. For composable rays $r\colon x\to y$ and $s\colon y\to z$, their \emph{stem product} $r*s$ is the unique ray $t\colon x\to z$ satisfying $U_t=U_rU_s$ when $U_rU_s\ne0$, and $r*s=0$ when $U_rU_s=0$. Uniqueness is guaranteed by the linear ordering of endpoint subbimodules in Lemma~\ref{lem:endpoint-chains}.
\end{definition}

Under the bijection between rays and the endpoint subbimodules they generate, the stem product $*$ is composition in the stem category of Kupisch~\cite{Kupisch}; see \cite[Section~2.3]{BongartzSurvey} for the formulation used here. Because $U_rU_s=R_xrR_ysR_z$, the stem product captures the entire subspace spanned by products of representatives with arbitrary insertions from the intermediate local ring $R_y$. We will use it to determine which pairs of rays with $rs=0$ can support a nonzero deformation term.

\subsection{Homogeneous subcategories and a dimension criterion}
\label{subsec:homogeneous-local-categories}

To establish the local restrictions on $A_0$, we show that any violation would produce a finite representation-infinite category built from rays of $A_0$, contradicting the representation-finiteness of $A_0$. We construct these test categories as follows. Equip $A_0$ with its grading by the weight lattice $L$ of Definition~\ref{def:integer-weights} (written additively), placing each ray $r$ in degree $\ell(r)$ and the primitive idempotents in degree zero. A \emph{homogeneous category} is the full subcategory $H$ of the associated graded covering category (Definition~\ref{def:graded-covering-category}) on a finite set of distinct objects $(x_i,w_i)$ with $w_i\in L$. Writing subscripts for $L$-degrees, its morphism spaces are
\[
  H(i,j)=A_0(x_i,x_j)_{w_j-w_i}.
\]
Because the categorical degree $\deg_\JJ$ is one of the coordinates of the weight lattice $L$, each nonzero morphism space $H(i,j)$ is one-dimensional and $H$ is a directed category. By Lemma~\ref{lem:homogeneous-projection}, every such homogeneous category $H$ is representation-finite. Passing to the graded covering allows us to separate distinct rays between the same vertices—or multiple copies of the same vertex at different weights—into distinct objects and morphisms of $H$. Whenever the categorical $\Z$-grading $\deg_\JJ$ alone suffices to separate the rays under consideration, we may perform the same construction using the $\Z$-graded covering category instead.

A primary source of forbidden configurations is the factorization poset induced by the action of $H$ on a thin projective or injective $H$-module. Taking the quotient by the annihilator of the module in the next definition is essential, because a morphism that is nonzero in $H$ may still act as zero on the chosen module.

\begin{definition}
\label{def:faithful-support}
Let $H$ be a homogeneous category (with respect to either the weight lattice $L$ or the categorical grading $\deg_\JJ$). Fix an object $X\in H$, set $V=H(X,-)$, and let $\ann V$ be the ideal of morphisms in $H$ that act as zero on $V$. The \emph{faithful projective support} at $X$ is the poset consisting of the objects $Y\in H$ with $V(Y)\ne0$, ordered by $Y\leq Z$ whenever $(H/\ann V)(Y,Z)\ne0$. Equivalently, if we label the points of this poset by the nonzero rays $r$ starting at $X$, then $r\leq s$ if and only if $rv=s$ for some ray $v$ in $H$.

The \emph{faithful injective support} at $X$ is defined dually via the thin module $D_k H(-,X)$; in terms of nonzero rays ending at $X$, its order is given by $r\leq s$ if and only if $vr=s$ for some ray $v$ in $H$.
\end{definition}

When restricted to the objects on which $V$ is nonzero, the quotient $H/\ann V$ is isomorphic to the incidence category of the faithful projective support at $X$ by Lemma~\ref{lem:thin-module-incidence} (positivity of the grading rules out nonidentity isomorphisms). By the same lemma, this incidence category is representation-finite, and hence so is the incidence category of every full subposet (which is an idempotent corner of $H/\ann V$). The dual statements hold for the faithful injective support.

The following standard Tits-form dimension criterion rules out the small representation-infinite quivers and posets that arise in our local arguments, sparing us from having to construct explicit one-parameter families of representations for each orientation.

\begin{lemma}
\label{lem:dimension-infinity}
Let $B=kQ/(\rho_1,\ldots,\rho_t)$ be a finite-dimensional bound quiver algebra, where each relation $\rho_i$ is a linear combination of radical paths sharing a common source and target. For a nonzero dimension vector $d=(d_v)$, set
\[
  D=\sum_v d_v^2,\quad N=\sum_{a:v\to w} d_v d_w,\quad
  F=\sum_{i=1}^t d_{\mathrm{source}(\rho_i)}d_{\mathrm{target}(\rho_i)}.
\]
If $N-F\geq D$, then $B$ is representation-infinite.
\end{lemma}
\begin{proof}
The affine variety of $B$-modules of dimension vector $d$ is cut out inside $\mathbb{A}^N$ by at most $F$ polynomial equations, and it is nonempty because it contains the zero representation. Every irreducible component of this variety therefore has dimension at least $N-F$. On the other hand, the algebraic group $\prod_v\operatorname{GL}_{d_v}(k)$ has dimension $D$, and the diagonally embedded subgroup $k^\times$ acts trivially on every representation, so every isomorphism orbit has dimension at most $D-1$. Since $N-F>D-1$, the representation variety cannot be covered by finitely many orbits. By the Krull--Schmidt theorem, this contradicts finite representation type.
\end{proof}

For the relation-free quivers ($F=0$) used repeatedly below, the critical dimension vectors achieving $N=D$ are listed in the following table:
\begin{center}
\small
\begin{tabular}{@{}lll@{}}
\toprule
Underlying graph & Dimension vector $d$ & $N=D$ \\\midrule
Crown ($2n$ vertices) & $1$ at every vertex & $2n$ \\
Four-leaf star ($\widetilde D_4$) & $2$ at center, $1$ at each leaf & $8$ \\
Tree $\widetilde D_5$ (two trivalent vertices) & $2$ at trivalent vertices, $1$ at leaves & $12$ \\
Tree $\widetilde E_7$ (arms $3,3,1$) & $4$ at center; $(3,2,1)$, $(3,2,1)$, $(2)$ on arms & $48$ \\
\bottomrule
\end{tabular}
\end{center}
Because $N$ and $D$ depend only on the underlying undirected graph when $F=0$, these counts apply to any orientation of the listed trees. Here a \emph{crown} means a bipartite cycle of $n$ minimal and $n$ maximal poset elements with no additional comparability relations. In particular, neither a four-point crown nor a four-leaf star can appear as a full subposet of any faithful projective or injective support.

We also apply Lemma~\ref{lem:dimension-infinity} to any quotient $B$ of $A_0$ (or of an idempotent corner of $A_0$) satisfying $(\rad B)^2=0$. Grading $B$ by assigning degree~$1$ to its arrows and degree~$0$ to its vertices, we pass to the full subcategory of its graded covering category on the objects $v_0=(v,0)$ and $v_1=(v,1)$ as $v$ ranges over the vertices of $B$. The only nonidentity morphisms in this \emph{two-level category} go from $v_0$ to $w_1$, one for each arrow $v\to w$ of $B$. By Lemma~\ref{lem:homogeneous-projection}, the two-level category is representation-finite; consequently, its underlying bipartite graph cannot contain a four-cycle (for which the all-ones dimension vector on the cycle has $N=D=4$ and $F=0$).

To rule out more general configurations with relations, we pull representations back along a split bimodule inclusion using the following observation.

\begin{lemma}
\label{lem:representation-transfer}
Let $F\colon S\to B$ be a functor between finite categories, and let $B_F$ be the category having the same object set as $S$ and morphism spaces $B_F(x,y)=B(Fx,Fy)$. Suppose that the induced functor $S\to B_F$ is injective on morphism spaces and admits a retraction $\pi\colon B_F\to S$ as an $S$-bimodule morphism. If $B$ is representation-finite, then so is $S$.
\end{lemma}
\begin{proof}
When several objects of $S$ map to the same object of $B$, the category $B_F$ is a matrix enlargement of an idempotent corner of $B$, and is therefore representation-finite by Lemma~\ref{lem:endpoint-chains}(1). Applying Lemma~\ref{lem:split-representation-transfer} to the split inclusion $S\hookrightarrow B_F$ with $C=D=S$ then shows that $S$ is representation-finite.
\end{proof}

\subsection{Unequal actions and generation at the middle endpoint}

With respect to our normed homogeneous ray basis, every nonzero endpoint action on a corner bimodule can be written as a power of the radical generator at an endpoint where the bimodule is cyclic. We first analyze what happens when the radical generator at the opposite endpoint acts as a second or higher power of the cyclic generator. Together with Lemma~\ref{lem:opposite-controls}, this bound will force any zero ray product $rs=0$ with nonzero stem product $r*s\ne0$ to land in the one-dimensional endpoint socle (Corollary~\ref{cor:active-zero-pairs}).

\begin{lemma}
\label{lem:unequal-endpoint-action}
In Setting~\ref{set:projective-companion}, suppose that $A_0(x,y)=R_xf$ has dimension $l$, where $f\colon x\to y$ is a generating ray. If $fq_y=q_x^af\ne0$ for some integer $a\geq2$, then $l=a+1$, and $fq_y$ spans the endpoint socle of $A_0(x,y)$. The dual statement holds when $A_0(x,y)=fR_y$.
\end{lemma}
\begin{proof}
Since $q_x^af=fq_y\ne0$, we have $l\geq a+1$. Suppose for a contradiction that $l\geq a+2$, so that $q_x^{a+1}f\ne0$. Set $d=\deg_\JJ(q_x)$ and $e=\deg_\JJ(f)$; the relation $fq_y=q_x^af$ then gives $\deg_\JJ(q_y)=ad$. We first check that there is no homogeneous reverse ray $h\colon y\to x$ satisfying $fh=q_x^b$ with $1\leq b\leq a+1$. Indeed, if such an $h$ existed, then $f(hf)=(fh)f=q_x^bf\ne0$ because $b\leq a+1<l$; thus $hf$ would be a nonzero homogeneous element of $R_y$ of degree $bd$, which is possible only if $bd$ is a multiple of $\deg_\JJ(q_y)=ad$, forcing $b=a$. Even for $b=a$, however, the element $hq_xf\in R_y$ has degree $(a+1)d$, which is not divisible by $ad$; hence $hq_xf=0$, contradicting $f(hq_xf)=(fh)q_xf=q_x^{a+1}f\ne0$.

Now consider the faithful projective support at $(x,0)$ (with respect to the categorical grading $\deg_\JJ$) and restrict to the six points
\[
  X_i=(x,id)\quad (i\in\{0,1,a,a+1\}),\qquad
  Y_j=(y,e+jd)\quad (j\in\{1,a\}),
\]
which correspond to the nonzero rays $q_x^i$ and $q_x^jf$. One point precedes another in this poset if and only if the second ray is a right multiple of the first. Thus the points $X_0<X_1<X_a<X_{a+1}$ form a chain, and $X_i<Y_j$ holds whenever $i\leq j$. No reverse inequality $Y_j\leq X_i$ can hold, because it would require a ray $h\colon y\to x$ with $q_x^jfh=q_x^i$, and hence $fh=q_x^{i-j}$ with $1\leq i-j\leq a$, which we have just ruled out. This same exclusion prevents $X_i<Y_j$ when $i>j$. Finally, $Y_1$ and $Y_a$ are incomparable: a morphism $Y_1\to Y_a$ would have degree $(a-1)d$, which is not a multiple of $\deg_\JJ(q_y)=ad$, while a morphism $Y_a\to Y_1$ would have negative degree. Relabeling $X_0,X_1,X_a,X_{a+1}$ as $u_0,u_1,u_2,u_3$ and $Y_1,Y_a$ as $v_1,v_2$, we see that the induced full subposet has Hasse diagram
\[
\begin{tikzpicture}[x=1.3cm,y=.8cm,>=Latex,every node/.style={font=\small}]
\node (u0) at (0,0) {$u_0$}; \node (u1) at (1,1) {$u_1$};
\node (u2) at (1,2) {$u_2$}; \node (u3) at (2,3) {$u_3$};
\node (v2) at (0,3) {$v_2$}; \node (v1) at (2,2) {$v_1$};
\draw[->] (u0)--(u1); \draw[->] (u1)--(u2); \draw[->] (u2)--(u3);
\draw[->] (u2)--(v2); \draw[->] (u1)--(v1);
\end{tikzpicture}
\]
which is an orientation of the affine Dynkin tree $\widetilde D_5$. Assigning dimension~$2$ to the trivalent vertices $u_1,u_2$ and dimension~$1$ to the four leaves gives $N=D=12$ and $F=0$, contradicting finite representation type by Lemma~\ref{lem:dimension-infinity}.

We conclude that $l=a+1$. Consequently $q_x(fq_y)=q_x^{a+1}f=0$ and $(fq_y)q_y=q_x^{2a}f=0$ (since $2a\geq a+1=l$), so $fq_y$ is annihilated by both endpoint radicals and therefore spans the endpoint socle of $A_0(x,y)$.
\end{proof}

\begin{lemma}
\label{lem:opposite-controls}
In Setting~\ref{set:projective-companion}, let $f\colon b\to c$ and $g\colon c\to a$ be generating rays whose corner bimodules $A_0(b,c)$ and $A_0(c,a)$ are cyclic at $c$ but not at $b$ and $a$, respectively, and let $q=q_c$ generate $\rad R_c$.
\begin{enumerate}[label={\upshape(\arabic*)}]
\item Either $fq^2=0$ or $q^2g=0$.
\item The product subspace $fR_cg\subseteq A_0(b,a)$ has dimension at most one and is annihilated by $\rad R_b$ on the left and by $\rad R_a$ on the right.
\end{enumerate}
\end{lemma}
\begin{proof}
Because the local ring $R_c=A_0(c,c)$ is commutative and hence cyclic on both sides, the noncyclicity assumptions at $b$ and $a$ imply that $b\ne c$ and $a\ne c$ (though $b$ and $a$ may coincide). Moreover, a one-dimensional corner bimodule is automatically cyclic over both endpoint rings, so $\dim A_0(b,c)\geq2$ and $\dim A_0(c,a)\geq2$. Writing $A_0(b,c)=fR_c\cong k[q]/(q^n)$ with $n\geq2$, take any $t\in\rad R_b$ and write $tf=f\varphi(q)$ for a polynomial $\varphi(q)\in k[q]$. Nilpotence of $t$ forces $\varphi(0)=0$, and if the linear coefficient of $\varphi$ were nonzero, then the powers of $\varphi(q)$ would generate $k[q]/(q^n)$, making $A_0(b,c)=R_bf$ cyclic at $b$, contrary to assumption. Thus $(\rad R_b)f\subseteq fq^2R_c$. Furthermore, for any morphism $h\in A_0(c,b)$, the composition $fh$ lies in $\rad R_b$ because $b\ne c$; hence $f(hf)=(fh)f\in fq^2R_c$, and since $\ann_{R_c}(f)=(q^n)\subseteq(q^2)$, it follows that $hf\in(q^2)$. Applying the dual argument to $g$ yields the four inclusions
\begin{equation}
  \begin{gathered}
    (\rad R_b)f\subseteq fq^2R_c,\qquad
    g(\rad R_a)\subseteq R_cq^2g,\\
    A_0(c,b)f\subseteq(q^2),\qquad
    gA_0(a,c)\subseteq(q^2).
  \end{gathered}
  \label{eq:opposite-control-inclusions}
\end{equation}

To prove (1), suppose for a contradiction that both $fq^2\ne0$ and $q^2g\ne0$. We distinguish two cases according to whether $fg$ is nonzero.

\smallskip
\noindent\textbf{Case 1: $fg\ne0$.}
Let $S$ be the path category of the four-leaf star quiver
\[
\begin{tikzcd}[row sep=small,column sep=small]
  & u\arrow[d,"q_1"] & \\
l\arrow[r,"f"] & c\arrow[r,"g"]\arrow[d,"q_2"] & r \\
  & d &
\end{tikzcd}
\]
in which the two paths $l\xrightarrow{f}c\xrightarrow{g}r$ and $u\xrightarrow{q_1}c\xrightarrow{q_2}d$ cross at $c$ with no relations. We map the five vertices $l,c,r,u,d$ of $S$ to the objects of the $\Z$-graded homogeneous category $H$ of $A_0$ having underlying vertices $b,c,a,c,c$ and grading coordinates
\[
  -\deg_\JJ(f),\quad 0,\quad \deg_\JJ(g),\quad -\deg_\JJ(q),\quad \deg_\JJ(q),
\]
respectively, sending the arrows $f,g,q_1,q_2$ to the rays $f,g,q,q$. All four length-two paths in $S$ then map to the nonzero rays $fg,fq,qg,q^2$ in $H$, so the functor $S\to H$ identifies each nonzero morphism space of $S$ with the corresponding one-dimensional morphism space of $H$.

Define a linear map $\pi\colon H\to S$ on each morphism space to be the inverse of this identification whenever the morphism space in $S$ is nonzero, and zero otherwise. To see that $\pi$ is an $S$-bimodule retraction as required by Lemma~\ref{lem:representation-transfer}, it suffices to check that $\pi$ commutes with left and right multiplication by the four arrows of $S$. Whenever the product of two composable morphisms in $H$ lies in a morphism space where $S$ is nonzero, compatibility holds automatically because that homogeneous component is one-dimensional. Otherwise, we must verify that no morphism of $H$ outside the image of $S$ can compose with an arrow image to land inside the image of $S$. At the outgoing fork $c\to r,d$, a nonzero composite factoring through an ambient ray $h\colon r\to d$ or $h\colon d\to r$ would yield $gh=q$ or $qh=g$ (up to a nonzero scalar); dually, at the incoming fork $l,u\to c$, an ambient ray between $l$ and $u$ would yield $hf=q$ or $hq=f$. However, $gh=q$ and $hf=q$ are ruled out by~\eqref{eq:opposite-control-inclusions}, while $qh=g$ and $hq=f$ contradict the hypothesis that $g$ and $f$ are generating rays (Lemma~\ref{lem:generating-paths}). Since no product involving a morphism of positive degree can be an identity, $\pi$ is indeed an $S$-bimodule retraction. By Lemma~\ref{lem:dimension-infinity}, the four-leaf star $S$ is representation-infinite, contradicting Lemma~\ref{lem:representation-transfer}.

\smallskip
\noindent\textbf{Case 2: $fg=0$.}
When $fg=0$, the last two inclusions in~\eqref{eq:opposite-control-inclusions} sharpen to
\[
  A_0(c,b)f\subseteq(q^3),\qquad gA_0(a,c)\subseteq(q^3),
\]
because if $hf$ had a nonzero $q^2$-term, multiplying on the right by $g$ would give $(hf)g=h(fg)=0$ while $q^2g\ne0$, and dually for $gh$. We now let $S$ be the bound quiver category on seven vertices with diagram and single relation
\[
\begin{tikzcd}[row sep=large,column sep=large]
u' & u\arrow[l,"g_1"']\arrow[d,"q_1"] & \\
l\arrow[r,"f_1"] & c\arrow[r,"g_2"]\arrow[d,"q_2"] & r \\
l_d\arrow[r,"f_2"] & d &
\end{tikzcd}
\qquad f_1g_2=0.
\]
We map $l,c,r,u,d$ to the same graded objects as in Case~1, and we send $l_d$ and $u'$ to $(b,\deg_\JJ(q)-\deg_\JJ(f))$ and $(a,\deg_\JJ(g)-\deg_\JJ(q))$, respectively, mapping $f_1,f_2$ to $f$, $g_1,g_2$ to $g$, and $q_1,q_2$ to $q$. The relation $f_1g_2=0$ holds in $H$ because $fg=0$, and the three remaining length-two paths map to the nonzero rays $q^2,fq,qg$. Defining the projection $\pi\colon H\to S$ as before, we verify that no ambient morphism in $H$ can compose with an arrow of $S$ to produce a nonzero morphism in the image of $S$. The possible unwanted products and the reasons they vanish are summarized below:
\[
  \begin{array}{c|c|l}
  \text{target ray in }S & \text{ambient product in }H & \text{reason for exclusion} \\ \hline
  q\text{ or }q^2 & gh & gA_0(a,c)\subseteq(q^3) \\
  qg & gv\ (v\in\rad R_a) & \text{inclusion~\eqref{eq:opposite-control-inclusions} and }q^3g\notin k\cdot qg \\
  g & qh & g\text{ is a generating ray} \\
  q\text{ or }q^2 & hf & A_0(c,b)f\subseteq(q^3) \\
  fq & vf\ (v\in\rad R_b) & \text{inclusion~\eqref{eq:opposite-control-inclusions} and }fq^3\notin k\cdot fq \\
  f & hq & f\text{ is a generating ray}.
  \end{array}
\]
At $u$, the first three rows rule out factorizations between $g_1\colon u\to u'$ and the paths through $c$; at $c$, they rule out factorizations between the two outgoing branches $g_2$ and $q_2$. The last three rows provide the dual incoming exclusions at $c$ and $d$, while the remaining vertices have at most one incident branch in the relevant direction. Thus $\pi$ is a bimodule retraction (note that Lemma~\ref{lem:representation-transfer} permits distinct vertices of $S$ to share the same underlying graded object in $H$). Assigning dimension~$3$ to $c$, dimension~$2$ to $u$ and $d$, and dimension~$1$ to the four outer vertices $u',l,r,l_d$ gives $N=22$, $D=21$, and $F=1$ from the single relation $f_1g_2=0$. Since $N-F=D$, Lemma~\ref{lem:dimension-infinity} and Lemma~\ref{lem:representation-transfer} again yield a contradiction, completing the proof of~(1).

\smallskip
For~(2), consider the action of $\rad R_b$ on the left or $\rad R_a$ on the right of $fR_cg$. By~\eqref{eq:opposite-control-inclusions}, any such product is a linear combination of terms $fq^mg$ with $m\geq2$. Since either $fq^2=0$ or $q^2g=0$ by~(1), every such term vanishes. Thus $fR_cg$ is annihilated by both $\rad R_b$ and $\rad R_a$, so it lies inside the one-dimensional endpoint socle of $A_0(b,a)$.
\end{proof}

\subsection{Stem products and generating paths}

We next show that the stem product of Definition~\ref{def:stem-product} is determined purely by the ray category $P$ via an explicit local formula, and can be nonzero even when the ordinary ray product vanishes.

\begin{proposition}[{\cite[Section~9.1]{BGRS}}]
\label{prop:intrinsic-stem-product}
In Setting~\ref{set:projective-companion}, the stem product on $A_0=kP$ is determined by the ray category $P$: for any composable rays $r\colon x\to y$ and $s\colon y\to z$,
\[
  r*s=\begin{cases}
  rq_y^ns, & \text{where }n\geq0\text{ is the least integer such that }rq_y^ns\ne0,\\
  0, & \text{if }rq_y^ns=0\text{ for all }n\geq0,
  \end{cases}
\]
with $q_y^0=1_y$ and $q_y=0$ when $R_y=k$. The operation $*$ is associative and agrees with the ray product $rs$ whenever $rs\ne0$. Moreover, the same formula computes the product of endpoint subbimodules in any algebra having ray category $P$ and satisfying the endpoint-chain properties of Lemma~\ref{lem:endpoint-chains}.
\end{proposition}
\begin{proof}
Let $r\colon x\to y$ and $s\colon y\to z$ be composable rays in $P$, choose representatives $f\in r$ and $h\in s$ in an algebra with ray category $P$, and let $u$ generate $\rad R_y$ (with $u=0$ if $R_y=k$). For each $j\geq0$, set $V_j=R_xfu^jR_yhR_z$. We have $fu^j\ne0$ if and only if the ray product $rq_y^j$ is nonzero: indeed, if $fu^j\ne0$, then $fu^j$ generates the subbimodule $U(fu^j)$ and therefore lies outside its radical $R(fu^j)\supseteq M(f,u^j)$, which means $rq_y^j\ne0$ by~\eqref{eq:ray-product}; conversely, if $fu^j=0$, then $fu^j\in M(f,u^j)$ and $rq_y^j=0$. When $rq_y^j\ne0$, its ray product with $s$ vanishes if and only if $fu^jh\in V_{j+1}$, or equivalently $V_j=V_{j+1}$. Consequently, if $n\geq0$ is the smallest integer for which $rq_y^ns\ne0$, then $U_rU_s=V_0=V_n$, and hence
\begin{equation}
  U_rU_s=U_{rq_y^ns}.
  \label{eq:intrinsic-stem-product}
\end{equation}
If no such $n$ exists, the chain of equalities $V_j=V_{j+1}$ and the nilpotence of $u$ force $U_rU_s=0$. Thus $r*s$ is given by the stated formula in $P$. Associativity of $*$ follows immediately from the associativity of subbimodule multiplication, and when $rs\ne0$ the least index is $n=0$, giving $r*s=rs$.
\end{proof}

\begin{corollary}
\label{cor:active-zero-pairs}
In Setting~\ref{set:projective-companion}, let $r\colon x\to y$ and $s\colon y\to z$ be composable rays such that $rs=0$ and $r*s\ne0$. Then $r$ and $s$ are generating rays, $r*s=rq_ys$, and $U_{r*s}$ is the one-dimensional endpoint socle of $A_0(x,z)$.
\end{corollary}
\begin{proof}
If the corner bimodule $A_0(x,y)$ containing $r$ were cyclic at $x$, then $rq_y\in\rad A_0(x,y)$ would equal $q_x^ar$ for some $a\geq1$. For every $n\geq1$ we would then have $rq_y^ns=q_x^{an}(rs)=0$, giving $r*s=0$ by Proposition~\ref{prop:intrinsic-stem-product}, a contradiction. By the dual argument for $s$, neither $A_0(x,y)$ nor $A_0(y,z)$ is cyclic at its outer endpoint, so both are cyclic only at $y$. Write $r=fq_y^i$ and $s=q_y^jg$ where $f$ and $g$ are the generating rays of $A_0(x,y)$ and $A_0(y,z)$. By Lemma~\ref{lem:opposite-controls}, $fq_y^mg=0$ for all $m\geq2$, and the entire subspace $fR_yg$ lies in the endpoint socle of $A_0(x,z)$. Since $rs=fq_y^{i+j}g=0$ while $rq_y^ns=fq_y^{i+j+n}g\ne0$ for some $n\geq1$, the only possibility is $i=j=0$ and $n=1$. Thus $r=f$ and $s=g$ are generating rays, $r*s=rq_ys$, and $U_{r*s}$ is the endpoint socle of $A_0(x,z)$.
\end{proof}

\begin{lemma}
\label{lem:generating-paths}
In Setting~\ref{set:projective-companion}, let $f\colon x\to z$ be a generating ray whose corner bimodule $A_0(x,z)$ has dimension at least two. Then $f$ admits no factorization of the form $f=aq_yb$ through a local radical generator $q_y$. In particular, every ordinary path representing $f$ is vertex-simple (visits no vertex twice), and every nonempty proper subsegment of such a path represents a generating ray.
\end{lemma}
\begin{proof}
Suppose for a contradiction that $f=aq_yb$ for some rays $a\colon x\to y$ and $b\colon y\to z$. If $a=1_x$, then $f=q_xb\in(\rad R_x)A_0(x,z)$ lies in the radical of $A_0(x,z)$, contradicting the assumption that $f$ generates $A_0(x,z)$; the case $b=1_z$ is ruled out similarly. If the corner bimodule $A_0(x,y)$ is cyclic at $x$, then $aq_y\in(\rad R_x)A_0(x,y)$, which again places $f=aq_yb$ inside $(\rad R_x)A_0(x,z)$, and symmetrically if $A_0(y,z)$ is cyclic at $z$. In the remaining case, both $A_0(x,y)$ and $A_0(y,z)$ are cyclic only at $y$. Writing $a=a'q_y^i$ and $b=q_y^jb'$ with $a',b'$ generating, we find $f\in a'R_yb'$. By Lemma~\ref{lem:opposite-controls}(2), $a'R_yb'$ is contained in the one-dimensional endpoint socle of $A_0(x,z)$, whereas the generating ray $f$ of a corner bimodule of dimension at least two cannot lie in its endpoint socle. This contradiction rules out any factorization $f=aq_yb$. Finally, any repeated vertex along an ordinary path representing $f$, or any proper (non-generating) subsegment of that path, would produce a factorization of $f$ through a local radical generator $q_y$.
\end{proof}

Example~\ref{ex:zero-ray-span} illustrates the difference between the ordinary ray product and the stem product: there $[\alpha][\beta]=0$ in the ray category, whereas $U_\alpha U_\beta=ku$ is nonzero because inserting the middle radical generator $q$ yields $\alpha q\beta=u\ne0$, so that $[\alpha]*[\beta]=[u]$.

\subsection{Successor coefficients and primitive returns}
\label{sec:primitive-ray-equations}

We next study the coefficients of normalized Hochschild one-cochains $\delta$ that shift ray weights by a fixed character $\chi\in L$. The Leibniz condition $\delta(rs)=\delta(r)s+r\delta(s)$ on nonzero ray products $rs\ne0$ is equivalent to the vanishing of the coboundary $d\delta$ on all nonzero products of basis elements. In Section~\ref{sec:integral-partial-lifts} we will lift solutions of these nonzero-product Leibniz equations from $\Z/2$ to $\Z$. Applied to a derivation over $\F_2$, such an integral lift produces a two-cocycle $(d\delta)/2$ over $\F_2$ that vanishes on all nonzero ray products and represents the connecting homomorphism class of $\delta$ in $\HH^2(\F_2P)_\chi$ (Proposition~\ref{prop:integral-bockstein}).

\begin{definition}
\label{def:primitive-characters}
Evaluating characters in $L$ on the categorical degree $\deg_\JJ\in H_\Z$, we say that a character $\chi\in L$ is \emph{positive} if $\chi(\deg_\JJ)>0$, and \emph{primitive} if it cannot be written as $n\eta$ for an integer $n>1$ and a character $\eta\in L$.
\end{definition}

By Proposition~\ref{prop:primitive-parameter-characters} and the identity $\chi_i(\deg_\JJ)=2s_i$, every residual character $\chi_i$ is both positive and primitive.

\begin{setting}
\label{set:fixed-character}
For the remainder of this section and Section~\ref{sec:integral-partial-lifts}, we fix a positive primitive character $\chi\in L$. For each vertex $x$, we define
\[
  \widehat q_x=\begin{cases}
  q_x, & \text{if }q_x\ne0\text{ and }\ell(q_x)=\chi,\\
  0, & \text{otherwise}.
  \end{cases}
\]
Thus $q_x$ continues to denote the actual homogeneous generator of $\rad R_x$ from Setting~\ref{set:projective-companion}, whereas $\widehat q_x$ retains $q_x$ only when its weight equals the fixed character $\chi$.
\end{setting}

\begin{definition}
\label{def:successor-types}
Work with $A_0=kP$ and a positive primitive character $\chi$ as in Settings~\ref{set:projective-companion} and~\ref{set:fixed-character}.
\begin{enumerate}[label={\upshape(\arabic*)}]
\item A vertex $x$ is \emph{supporting} (for $\chi$) if $\widehat q_x\ne0$, that is, if $q_x\ne0$ and $\ell(q_x)=\chi$. More generally, $x$ has \emph{primitive local weight} if $q_x\ne0$ and $\ell(q_x)$ is primitive in $L$.
\item A nonidentity ray $r\colon x\to y$ is \emph{shiftable} if there is a ray $r^+\colon x\to y$ of weight $\ell(r)+\chi$, which we call the \emph{successor} of $r$.
\item The \emph{endpoint type} of a shiftable ray $r\colon x\to y$ is classified by which of the endpoint actions $\widehat q_xr$ and $r\widehat q_y$ are nonzero:
\[
\begin{array}{c|c|l}
\text{Name} & \text{Symbol} & \text{Defining condition} \\ \hline
\text{left-sided} & \mathsf L & \widehat q_xr=r^+\ne0,\ r\widehat q_y=0 \\
\text{right-sided} & \mathsf R & \widehat q_xr=0,\ r\widehat q_y=r^+\ne0 \\
\text{two-sided} & \mathsf B & \widehat q_xr=r\widehat q_y=r^+\ne0.
\end{array}
\]
We say that $r$ is \emph{one-sided} if its endpoint type is $\mathsf L$ or $\mathsf R$. We also say that the source $x$ \emph{controls} $r$ when $\widehat q_xr\ne0$, and that the target $y$ \emph{controls} $r$ when $r\widehat q_y\ne0$.
\end{enumerate}
\end{definition}

To see why every shiftable ray $r\colon x\to y$ falls into one of the three types in Definition~\ref{def:successor-types}(3), recall that the corner bimodule $A_0(x,y)$ is cyclic over at least one of its endpoint rings $R_x$ or $R_y$, so its rays are successive powers of that endpoint's radical generator applied to the generating ray of $A_0(x,y)$. The weight difference $\ell(r^+)-\ell(r)=\chi$ must therefore be a positive integer multiple of the weight of that radical generator. Because $\chi$ is primitive, the multiple is necessarily~$1$, which means that $r^+$ is obtained from $r$ by a single multiplication by $\widehat q_x$ on the left or by $\widehat q_y$ on the right (and $r^+$ is unique because each weight space in $A_0(x,y)$ is at most one-dimensional). Notice the distinction between cyclicity in Definition~\ref{def:generating-rays}, which is a property of the entire corner bimodule $A_0(x,y)$ over $R_x$ or $R_y$, and the endpoint type $\mathsf L,\mathsf R,\mathsf B$, which records whether the action of weight $\chi$ on a specific shiftable ray $r$ remains nonzero on the left, on the right, or on both sides.

For example, if $A_0(x,y)$ is a three-dimensional corner bimodule generated on the left at a supporting vertex $x$, its rays are $f$, $\widehat q_xf$, and $\widehat q_x^2f$, which lie in endpoint layers $0$, $1$, and $2$, respectively. The rays $f$ and $\widehat q_xf$ are shiftable, with successors $\widehat q_xf$ and $\widehat q_x^2f$, whereas the socle ray $\widehat q_x^2f$ is unshiftable. If the right action of $\widehat q_y$ satisfies $f\widehat q_y=\widehat q_xf$, then both $f$ and $\widehat q_xf$ have type $\mathsf B$; if instead $f\widehat q_y=0$, then both have type $\mathsf L$.

Any normalized one-cochain $\delta$ of weight $\chi$ vanishes on the primitive idempotents $1_x$ and on all unshiftable rays, and is determined on shiftable nonidentity rays $r$ by scalar coefficients $d_r$ via $\delta(r)=d_r r^+$. On a nonzero product $rs=t$ of nonidentity rays, the Leibniz equation $\delta(rs)=\delta(r)s+r\delta(s)$ compares $d_t t^+$ with $d_r r^+s+d_s rs^+$. We refer to the products $r^+s$ and $rs^+$ as the \emph{successor substitutions} of $rs$; each substitution either equals $t^+$ or vanishes.

Before tabulating these equations, we note a basic consequence of primitive local weight that will be used throughout our analysis.

\begin{lemma}
\label{lem:primitive-action-generation}
In Setting~\ref{set:projective-companion}, let $f\colon x\to y$ be a generating ray, and suppose that $\ell(q_x)$ is primitive. If $q_xf\ne0$, then $A_0(x,y)=R_xf$. If in addition $A_0(x,y)$ is cyclic at $y$, then $\ell(q_y)=\ell(q_x)$ and $q_xf=fq_y\ne0$. The dual assertions hold when $\ell(q_y)$ is primitive and $fq_y\ne0$.
\end{lemma}
\begin{proof}
If $A_0(x,y)$ is cyclic at $y$, then $q_xf\ne0$ must take the form $q_xf=fq_y^a$ for some integer $a\geq1$. Comparing weights gives $\ell(q_x)=a\ell(q_y)$, and the primitivity of $\ell(q_x)$ forces $a=1$. Iterating the relation $q_xf=fq_y$ then shows that $R_xf=fR_y=A_0(x,y)$. If $A_0(x,y)$ is not cyclic at $y$, then Lemma~\ref{lem:endpoint-chains}(3) immediately implies that $A_0(x,y)=R_xf$. The second assertion follows by the same weight comparison, and the dual statements follow by passing to the opposite algebra.
\end{proof}

\begin{lemma}
\label{lem:successor-substitutions}
In Settings~\ref{set:projective-companion} and~\ref{set:fixed-character}, let $rs=t$ be a nonzero product of nonidentity rays. If $t$ is unshiftable, then both successor substitutions $r^+s$ and $rs^+$ vanish (whenever defined), so $\delta(rs)=\delta(r)s+r\delta(s)$ imposes no condition on the coefficients $d_r,d_s$. If $t$ is shiftable, the possible endpoint types of $(r,s)$ and the resulting Leibniz equations for $d_t$ are given by the following table (where $*$ denotes an arbitrary shiftable or unshiftable ray):
\begin{equation}
  \label{eq:product-type-table}
  \begin{array}{c|c|c}
  \text{Type of }t=rs & \text{Types of }(r,s) & \text{Leibniz equation} \\ \hline
  \mathsf L & (\mathsf L,*) & d_t=d_r \\
  \mathsf L & (\mathsf B,\mathsf L) & d_t=d_r+d_s \\
  \mathsf R & (*,\mathsf R) & d_t=d_s \\
  \mathsf R & (\mathsf R,\mathsf B) & d_t=d_r+d_s \\
  \mathsf B & (\mathsf B,\mathsf B)\text{ or }(\mathsf L,\mathsf R) & d_t=d_r+d_s.
  \end{array}
\end{equation}
\end{lemma}
\begin{proof}
Let $rs=t\ne0$ with $r\colon x\to y$ and $s\colon y\to z$. If $t$ is unshiftable, then $A_0(x,z)$ has no ray of weight $\ell(t)+\chi$, so neither $r^+s$ nor $rs^+$ can be nonzero.

Assume now that $t$ is shiftable. Suppose first that $t$ has type $\mathsf L$, so that $t^+=\widehat q_xrs\ne0$ while $rs\widehat q_z=0$. Since $\widehat q_xrs\ne0$, we have $\widehat q_xr\ne0$, so $r$ is shiftable of type $\mathsf L$ or $\mathsf B$ and $r^+s=\widehat q_xrs=t^+$.
\begin{itemize}
\item If $r$ has type $\mathsf L$, then $r\widehat q_y=0$. If $s$ is shiftable, its successor $s^+$ equals either $\widehat q_ys$ or $s\widehat q_z$; in either case $rs^+$ equals $r\widehat q_ys=0$ or $rs\widehat q_z=0$. Thus only the first substitution $r^+s=t^+$ is nonzero, yielding $d_t=d_r$.
\item If $r$ has type $\mathsf B$, then $r\widehat q_y=\widehat q_xr$, so $r\widehat q_ys=\widehat q_xrs=t^+\ne0$. Hence $\widehat q_ys\ne0$, which means $s$ is shiftable with a nonzero left action, while $s\widehat q_z$ cannot equal $\widehat q_ys$ because $rs\widehat q_z=0\ne r\widehat q_ys$. Thus $s$ has type $\mathsf L$ and $rs^+=r\widehat q_ys=t^+$, giving $d_t=d_r+d_s$.
\end{itemize}
The two rows for $t$ of type $\mathsf R$ follow symmetrically by passing to the opposite algebra.

Finally, suppose that $t$ has type $\mathsf B$, so that $t^+=\widehat q_xrs=rs\widehat q_z\ne0$. Then $\widehat q_xr\ne0$ and $s\widehat q_z\ne0$, so $r\in\{\mathsf L,\mathsf B\}$ and $s\in\{\mathsf R,\mathsf B\}$. The mixed pairs $(\mathsf L,\mathsf B)$ and $(\mathsf B,\mathsf R)$ cannot occur: for $(\mathsf L,\mathsf B)$ we would have $t^+=rs\widehat q_z=r\widehat q_ys=0$, and for $(\mathsf B,\mathsf R)$ we would have $t^+=\widehat q_xrs=r\widehat q_ys=0$. In the remaining cases $(\mathsf L,\mathsf R)$ and $(\mathsf B,\mathsf B)$, both $r^+s=\widehat q_xrs$ and $rs^+=rs\widehat q_z$ equal $t^+$, which gives $d_t=d_r+d_s$.
\end{proof}

We now turn to factorizations of a primitive local generator $q_x$ through other vertices. Any factorization $q_x=fh$ through a vertex $y\ne x$ has $q_x$ as its output, and when $\ell(q_x)=\chi$ is shiftable (so $q_x^2\ne0$), Lemma~\ref{lem:successor-substitutions} yields a linear relation linking $d_f$, $d_h$, and $d_{q_x}$. We first describe the corner subcategories formed by the vertices through which such factorizations pass.

\begin{definition}
\label{def:primitive-return-corner}
Using the local radical generators $q_x$ of Setting~\ref{set:projective-companion} (independently of the fixed character $\chi$):
\begin{enumerate}[label={\upshape(\arabic*)}]
\item A \emph{primitive reverse} of a ray $f\colon x\to y$ is a ray $h\colon y\to x$ satisfying $fh=q_x\ne0$ and $hf=q_y\ne0$ (neither $f$ nor $h$ is required to be shiftable).
\item A \emph{return partner} of $x$ is a vertex $y\ne x$ for which there exists a generating ray $f\colon x\to y$ that admits a primitive reverse and satisfies $q_xf=fq_y\ne0$.
\item Writing $V_x$ for the set consisting of $x$ together with all of its return partners, the \emph{primitive return corner} $S_x$ is the full corner subcategory of $A_0$ on $V_x$. We write $y\in S_x$ as shorthand for $y\in V_x$.
\end{enumerate}
\end{definition}

\begin{lemma}
\label{lem:primitive-cyclic-corners}
Work in Setting~\ref{set:projective-companion}, with primitive reverses and return corners as in Definition~\ref{def:primitive-return-corner}.
\begin{enumerate}[label={\upshape(\arabic*)}]
\item Let $f\colon x\to y$ ($x\ne y$) be a generating ray whose corner bimodule $A_0(x,y)$ is cyclic at both endpoints with $q_xf=fq_y\ne0$, and suppose that $\ell(q_x)$ is primitive. If $fA_0(y,x)\ne0$, then the generating ray $h$ of $A_0(y,x)$ satisfies
\[
  fh=q_x,\qquad hf=q_y.
\]
\item Let $x$ be a vertex of primitive local weight. The generating rays $f\colon x\to y$ ($y\ne x$) that admit a primitive reverse and satisfy $q_xf=fq_y\ne0$ are linearly ordered by factorization, and the primitive return corner $S_x$ is a basic cyclic Nakayama algebra.
\item As $x$ ranges over all vertices of primitive local weight with $N_x\geq3$, the vertex sets $V_x$ are pairwise disjoint or equal, and every factorization $q_x=uv$ into nonidentity rays $u\colon x\to y$ and $v\colon y\to x$ has its intermediate vertex $y$ in $S_x$.
\end{enumerate}
\end{lemma}
\begin{proof}
\textit{(1)} Because $A_0(x,y)=fR_y$ is cyclic at $y$, multiplying the generating ray $h$ of $A_0(y,x)$ by powers of $q_y$ or $q_x$ simply multiplies $fh$ by powers of $q_x$. The assumption $fA_0(y,x)\ne0$ therefore ensures that $fh=q_x^a\ne0$ for some least return exponent $a\geq1$. If $a\geq2$, then the faithful projective support at $x$ contains the four-point subposet with lower elements $q_x,f$ and upper elements $q_x^a,q_xf$. All four comparisons from $\{q_x,f\}$ to $\{q_x^a,q_xf\}$ hold because $q_xf=fq_y$ and $fh=q_x^a$. On the other hand, $q_x$ and $f$ are incomparable (since $f$ is generating and $a\geq2$ is minimal), and by left cancellation $q_x^a$ and $q_xf$ are likewise incomparable, while neither upper element can precede a lower element by degree. Thus $\{q_x,f,q_x^a,q_xf\}$ forms a forbidden four-point crown, a contradiction. Hence $a=1$ and $fh=q_x$. Since $f(hf)=(fh)f=q_xf=fq_y\ne0$, we also have $hf=q_y$. (Note that $h$ had to be generating from the outset, as any proper endpoint factor of $h$ would strip a power of $q_x$ or $q_y$ from $fh=q_x$, factoring an identity morphism through a distinct vertex.)

\textit{(2)} By Definition~\ref{def:primitive-return-corner} and Lemma~\ref{lem:primitive-action-generation}, the targets of the specified generating rays $f\colon x\to y$ are precisely the return partners of $x$, and each associated corner bimodule $A_0(x,y)$ is cyclic at both endpoints. Given two such rays $f_i\colon x\to y_i$ and $f_j\colon x\to y_j$, their primitive reverses show that both $f_i$ and $f_j$ precede both $q_xf_i$ and $q_xf_j$ in the faithful projective support at $x$. By left cancellation of $q_x$, $f_i$ and $f_j$ are comparable if and only if $q_xf_i$ and $q_xf_j$ are comparable; moreover, no upper element $q_xf_i$ can precede a lower element $f_j$ because a generating ray $f_j$ cannot factor through $q_x$ (Lemma~\ref{lem:generating-paths}). If $f_i$ and $f_j$ were incomparable, the four rays $\{f_i,f_j,q_xf_i,q_xf_j\}$ would therefore form a four-point crown. Consequently, these outgoing generating rays form a chain under factorization. Since their categorical degrees lie strictly between $0$ and $\deg_\JJ(q_x)$ and no two distinct rays from $x$ to the same partner have the same weight, the return partners of $x$ may be listed in strictly increasing factorization order as $y_1,\ldots,y_t$.

For $1\leq i<j\leq t$, let $h_{ij}\colon y_i\to y_j$ be the ray satisfying $f_ih_{ij}=f_j$, and let $h_i\colon y_i\to x$ be the primitive reverse of $f_i$. Then
\[
  f_ih_{ij}=f_j,\qquad h_{ij}h_j=h_i,\qquad h_{ji}:=h_jf_i\colon y_j\to y_i.
\]
Cancellation and Lemma~\ref{lem:generating-paths} imply that all $h_{ij}$ and $h_{ji}$ are generating rays and are mutual primitive reverses. The full corner $S_x$ is therefore generated by the oriented cycle of ordinary arrows inside $S_x$,
\[
  x\to y_1\to\cdots\to y_t\to x,
\]
in which every simple path is a generating ray, one full turn around the cycle at any vertex is its local radical generator, and the only relations are monomial nilpotence relations. Thus $S_x$ is a cyclic Nakayama algebra.

\textit{(3)} Let $a$ be any vertex of primitive local weight $\xi=\ell(q_a)$, and let $b,c\in V_a\setminus\{a\}$ be two distinct return partners of $a$, reached by generating rays $u\colon a\to b$ and $v\colon a\to c$ with primitive reverses $\bar u,\bar v$. By the chain property in~(2), we may assume after swapping $b$ and $c$ if necessary that $v=ut$ for a ray $t\colon b\to c$. Setting $h=\bar v u\colon c\to b$, we have
\[
  ht=\bar v ut=\bar v v=q_c,\qquad u(th)=(ut)\bar v u=v\bar v u=q_a u\ne0.
\]
Since $q_a$, $q_b$, and $q_c$ all have weight $\xi$, the nonzero ray $th\in R_b$ has weight $\xi$ and therefore equals $q_b$. Neither $t$ nor $h$ can be a proper ray, since pulling out a local generator from $th=q_b$ or $ht=q_c$ would factor $1_b$ or $1_c$ through a distinct vertex. By associativity,
\[
  q_bt=(th)t=t(ht)=tq_c,\qquad q_ch=(ht)h=h(th)=hq_b.
\]
If either $N_b\geq3$ or $N_c\geq3$, then $q_b^2=tq_ch$ or $q_c^2=hq_bt$ is nonzero, which forces $q_bt=tq_c\ne0$ and $q_ch=hq_b\ne0$; hence $b$ and $c$ are mutual return partners whenever $\max(N_b,N_c)\geq3$. Likewise, if $N_a\geq3$, then $u\bar u q_a=q_a^2\ne0$ forces $q_b\bar u=\bar u q_a\ne0$, so $a\in V_b$; if both $N_a\geq3$ and $N_b\geq3$, applying the preceding partner comparison at $a$ and at $b$ shows that $V_a=V_b$.

Now suppose that $N_x\geq3$ and $N_{x'}\geq3$, and let $v\in V_x\cap V_{x'}$. Then $v$ has the same primitive local weight as $x$ and $x'$. If $N_v\geq3$, the equality just proved gives $V_x=V_v=V_{x'}$. If instead $N_v=2$, then since $N_x\geq3$ and $N_{x'}\geq3$, the reverse rays from $v$ to $x$ and to $x'$ show as above that $x,x'\in V_v$. Applying the partner comparison at $v$ to its two partners $x$ and $x'$ (which satisfy $N_x,N_{x'}\geq3$) shows that $x$ and $x'$ are mutual return partners, and hence $V_x=V_{x'}$. Thus the sets $V_x$ ($N_x\geq3$) are pairwise disjoint or equal.

Finally, suppose that $q_x=uv$ is a factorization of $q_x$ into nonidentity rays $u\colon x\to y$ and $v\colon y\to x$ at a vertex $x$ of primitive local weight $\xi$ with $N_x\geq3$ (so $q_x^2\ne0$). Then $u(vu)=(uv)u=q_xu\ne0$ (since $q_x^2=q_xuv\ne0$), so $vu=q_y^a\ne0$ for some $a\geq1$. Comparing weights gives $\xi=a\ell(q_y)$, so primitivity of $\xi$ forces $a=1$ and $vu=q_y$. Because $q_x^2=uvuv\ne0$, we have $q_xu=uq_y\ne0$ and $vq_x=q_yv\ne0$, and both $u$ and $v$ are generating rays because $uv=q_x$ cannot contain an extra local radical factor. Thus $y\in V_x$, so the factorization passes through $S_x$.
\end{proof}

\begin{lemma}
\label{lem:external-hom-chain}
In Setting~\ref{set:projective-companion}, let $x$ be a vertex of primitive local weight with $N_x\geq3$. Then the generating rays $f\colon x\to y$ satisfying $q_xf\ne0$ and $y\notin S_x$ are linearly ordered by factorization. The dual statement holds for incoming generating rays $g\colon y\to x$ with $gq_x\ne0$ and $y\notin S_x$.
\end{lemma}
\begin{proof}
Suppose for a contradiction that $f\colon x\to y$ and $h\colon x\to z$ are two incomparable generating rays with $q_xf\ne0$, $q_xh\ne0$, and $y,z\notin S_x$. In the faithful projective support at $x$, consider the five rays
\[
  1_x,\quad q_x,\quad q_x^2,\quad q_xf,\quad q_xh.
\]
Clearly $1_x<q_x$, and $q_x$ precedes each of $q_x^2$, $q_xf$, and $q_xh$. By left cancellation of $q_x$, the elements $q_xf$ and $q_xh$ are incomparable because $f$ and $h$ are incomparable. Neither $q_x^2\leq q_xf$ nor $q_xf\leq q_x^2$ can hold: the first inequality would cancel to $q_x\leq f$, contradicting Lemma~\ref{lem:generating-paths}, while the second would cancel to $fa=q_x$ for some ray $a\colon y\to x$, placing $y$ in $S_x$ by Lemma~\ref{lem:primitive-cyclic-corners}(3). The same holds for $q_xh$, and degree considerations rule out any remaining comparisons. The five points therefore form a four-leaf star with center $q_x$ and leaves $1_x,q_x^2,q_xf,q_xh$, which is impossible by Lemma~\ref{lem:dimension-infinity}.
\end{proof}

\subsection{First and last arrows of generating rays}

To construct our integral coefficient lift in Section~\ref{sec:integral-partial-lifts}, we will define the coefficient $d_f$ of a generating one-sided ray $f$ in terms of the first (or last) ordinary arrow along a path representing $f$. For this prescription to be well defined, all ordinary paths representing $f$ must share the same initial (or terminal) arrow. Constraints on multiple paths representing the same ray play a central role in \cite{BGRS,BongartzMild}; we establish the exact uniqueness property we need in the next proposition.

\begin{proposition}
\label{prop:controlled-contours}
In Settings~\ref{set:projective-companion} and~\ref{set:fixed-character}, all ordinary paths representing a generating left-sided ray begin with the same ordinary arrow. Dually, all ordinary paths representing a generating right-sided ray end with the same ordinary arrow.
\end{proposition}

The proof of Proposition~\ref{prop:controlled-contours} is given in Appendix~\ref{app:controlled-contours}, where we show that two distinct initial arrows for a generating left-sided ray would allow one to construct a bimodule retraction onto the incidence category of a representation-infinite tree.

The Leibniz equations of Lemma~\ref{lem:successor-substitutions} also constrain products $rs$ of generating rays whose output $rs$ is a \emph{proper} shiftable ray. A key source of such proper products is a path that leaves the primitive return corner $S_x$ and returns to $x$ with value $\widehat q_x^2$ at a vertex where $\widehat q_x^3\ne0$.

\begin{definition}
\label{def:nondeep-return}
In Settings~\ref{set:projective-companion} and~\ref{set:fixed-character}, a \emph{nondeep external return at $x$} is a pair of generating rays $e\colon x\to y$ of type $\mathsf L$ and $l\colon y\to x$ of type $\mathsf R$, with $y\notin S_x$, such that
\[
  el=\widehat q_x^2\quad\text{and}\quad\widehat q_x^3\ne0.
\]
\end{definition}

The adjective \emph{nondeep} reflects the condition $\widehat q_x^3\ne0$, which means that the return value $el=\widehat q_x^2$ is shiftable (with successor $\widehat q_x^3$) rather than lying in the socle of $R_x$.

\begin{figure}[htbp]
\centering
\begin{tikzpicture}[>=Latex,scale=.9,every node/.style={font=\small}]
\draw[rounded corners] (-.65,-.7) rectangle (5.15,1.65);
\node at (4.6,1.3) {$S_x$};
\node (x) at (0,0) {$x$};
\node (v1) at (1.5,1) {$v_1$};
\node (dots) at (3,1) {$\cdots$};
\node (vt) at (4.5,0) {$v_t$};
\draw[->] (x)--(v1); \draw[->] (v1)--(dots);
\draw[->] (dots)--(vt); \draw[->] (vt) to[bend left=14] (x);
\node at (2.3,.3) {one cycle gives $\widehat q_x$};
\node (y) at (0,-2.1) {$y\notin S_x$};
\draw[->] (x) to[bend right=28] node[left] {$e\ (\mathsf L)$} (y);
\draw[->] (y) to[bend right=28] node[right] {$l\ (\mathsf R)$} (x);
\node[align=left] at (3,-1.8) {$el=\widehat q_x^2$,\quad $le=0$\\$\widehat q_x^3\ne0$,\quad $\widehat q_x^4=0$};
\end{tikzpicture}
\caption{A primitive return corner $S_x$ together with a nondeep external return $(e,l)$ through $y\notin S_x$. The cycle arrows depict the Nakayama generator cycle of $S_x$, whereas $e$ and $l$ are generating rays that may be composite paths in the quiver of $A_0$. The relations $le=0$ and $\widehat q_x^4=0$ (so $N_x=4$) are established in Lemma~\ref{lem:nondeep-two-object-return}.}
\label{fig:primitive-corner-return}
\end{figure}

In the configuration of Figure~\ref{fig:primitive-corner-return}, both $e$ and $l$ are controlled at $x$: we have $\widehat q_xe=e^+\ne0$ and $l\widehat q_x=l^+\ne0$, whereas $e\widehat q_y=0$ and $\widehat q_yl=0$. Because $\widehat q_x=q_x\ne0$, Lemma~\ref{lem:nondeep-two-object-return} in Appendix~\ref{app:proper-products} shows that $le=0$ and $N_x=4$.

The next proposition, proved in Appendix~\ref{app:proper-products}, catalogs all proper shiftable products of generating rays and describes how nondeep external returns interact with other rays at $x$.

\begin{proposition}
\label{prop:proper-ray-products}
Work in Settings~\ref{set:projective-companion} and~\ref{set:fixed-character}, with endpoint layers as in Definition~\ref{def:generating-rays}, endpoint types as in Definition~\ref{def:successor-types}, and primitive return corners $S_x$ as in Definition~\ref{def:primitive-return-corner}.
\begin{enumerate}[label={\upshape(\arabic*)}]
\item If a generating ray $f\colon x\to y$ has type $\mathsf L$ (respectively, $\mathsf R$), then $\dim A_0(x,y)\leq3$, with equality implying $N_x=3$ (respectively, $N_y=3$). Moreover, whenever a nonzero product $rs$ and its first factor $r$ both have type $\mathsf L$, they have the same endpoint layer; the dual statement holds when $rs$ and $s$ both have type $\mathsf R$.
\item If $r\colon x\to y$ of type $\mathsf B$ and $s\colon y\to z$ of type $\mathsf L$ are generating rays whose product $rs$ is proper and shiftable, then $rs$ has endpoint layer~$1$, the source $x$ has local length $N_x=3$ ($\widehat q_x^2\ne0$ and $\widehat q_x^3=0$), and $r$ admits a primitive reverse (so $y\in S_x$). The dual statement holds for generating rays of types $(\mathsf R,\mathsf B)$.
\item Let $r\colon x\to y$ of type $\mathsf L$ and $s\colon y\to z$ of type $\mathsf R$ be generating rays whose product $rs$ is proper and shiftable. Then $rs$ has endpoint layer~$2$:
\begin{enumerate}[label={\upshape(\alph*)}]
\item If $x=z$, so that $rs=\widehat q_x^2\in R_x$, then $(r,s)$ is a nondeep external return at $x$ (Definition~\ref{def:nondeep-return}).
\item If $x\ne z$, write $rs=\widehat q_x^2f=f\widehat q_z^2$ where $f\colon x\to z$ is the generating ray of $A_0(x,z)$. Then at least one of the factorizations
\[
  r=fa,\quad as=\widehat q_z^2\qquad\text{or}\qquad s=bf,\quad rb=\widehat q_x^2
\]
holds, where $(a,s)$ (respectively, $(r,b)$) is a nondeep external return at $z$ (respectively, at $x$) and $a$ (respectively, $b$) is a generating one-sided ray.
\end{enumerate}
\item If $x$ admits a nondeep external return, then there is no generating two-sided ray $f\colon x\to z$ with $z\notin S_x$ and $\widehat q_xf\ne0$, and dually no generating two-sided ray $g\colon z\to x$ with $z\notin S_x$ and $g\widehat q_x\ne0$.
\end{enumerate}
\end{proposition}

With these local restrictions in hand, we are ready to lift the modulo-two coefficient equations of Lemma~\ref{lem:successor-substitutions} to the integers: the two-sided rays are governed by a smaller centralizer table to which the induction hypothesis applies, while Proposition~\ref{prop:controlled-contours} and Proposition~\ref{prop:proper-ray-products} allow us to extend the resulting integral lift consistently to all one-sided rays.

\section{Integral lifting via centralizers and four-point crowns}
\label{sec:integral-partial-lifts}

The local restrictions established in Section~\ref{sec:local-ray-geometry} govern how endpoint layers and endpoint types behave under nonzero compositions of rays. Continuing in Settings~\ref{set:projective-companion} and~\ref{set:fixed-character} for a fixed positive primitive character $\chi\in L$, we now use those restrictions to lift modulo-two solutions of the nonzero-product Leibniz equations to the integers. The equations involving only two-sided rays turn out to be the additive-weight relations of a strictly smaller weighted multiplication table obtained from a centralizer quotient, to which our inductive hypothesis applies. Once the two-sided coefficients are lifted, the remaining equations for generating one-sided rays reduce to a linear system on a bipartite graph, and the absence of four-point crowns ensures that the right-hand sides are constant on each connected component.

\subsection{The coefficient equations and their two-sided part}

For each shiftable nonidentity ray $r$, let $d_r$ denote its scalar coefficient in a normalized one-cochain $\delta(r)=d_r r^+$ of weight $\chi$, and set $d_r=0$ when $r$ is an identity or an unshiftable ray. For any nonzero product $rs=t$ of nonidentity rays whose output $t$ is shiftable, the Leibniz identity $\delta(rs)=\delta(r)s+r\delta(s)$ takes the form
\begin{equation}
  d_t=\epsilon_L d_r+\epsilon_R d_s,
  \qquad
  \epsilon_L=\begin{cases}1, & \text{if }r^+s=t^+,\\ 0, & \text{otherwise},\end{cases}
  \quad
  \epsilon_R=\begin{cases}1, & \text{if }rs^+=t^+,\\ 0, & \text{otherwise}.
  \end{cases}
  \label{eq:partial-leibniz-group}
\end{equation}
By Lemma~\ref{lem:successor-substitutions}, if $t=rs\ne0$ is unshiftable, then $\epsilon_L=\epsilon_R=0$ and the Leibniz condition on $(r,s)$ holds trivially as $0=0$. No condition is imposed on pairs $(r,s)$ with $rs=0$.

\begin{theorem}
\label{thm:integral-partial-lifting}
Let $A_0=kP$ and $\chi$ satisfy Settings~\ref{set:projective-companion} and~\ref{set:fixed-character}. Assume that Proposition~\ref{prop:table-weight-lifting} holds for every weighted multiplication table whose linearization is representation-finite of dimension strictly less than $\dim A_0$. Then every modulo-two solution $(\bar d_r)$ of the nonzero-product Leibniz equations~\eqref{eq:partial-leibniz-group} lifts to an integer solution $(d_r)$, and in particular reduces to a solution over $\Z/4$.
\end{theorem}

The proof of Theorem~\ref{thm:integral-partial-lifting} proceeds in three steps: first we isolate a smaller weighted multiplication table that governs the coefficients of all two-sided rays; next we express the coefficients of generating one-sided rays in terms of their initial or terminal one-sided ordinary arrows and solve the resulting bipartite system; finally, in Subsection~\ref{subsec:verify-nonzero-equations}, we verify that our formulas automatically satisfy every remaining equation in~\eqref{eq:partial-leibniz-group}.

To construct the algebra governing the two-sided rays, set $\widehat q=\sum_x\widehat q_x\in A_0$ and define
\[
  C_\chi=\operatorname{Cent}_{A_0}(\widehat q),\qquad
  \mathfrak a_\chi=\{c\in C_\chi:\widehat q c=0\},\qquad
  B_\chi=C_\chi/\mathfrak a_\chi.
\]

\begin{lemma}
\label{lem:centralizer-ray-algebra}
The quotient algebra $B_\chi$ is basic and representation-finite. The images of the supporting vertex idempotents $1_x$ ($\widehat q_x\ne0$) and the two-sided rays of $A_0$ form a multiplicative basis of $B_\chi$, defining a weighted multiplication table $Q_\chi$ with the weights and positive degree inherited from $P$ so that $B_\chi=kQ_\chi$. Moreover, if $B_\chi\ne0$, then $Q_\chi$ has strictly fewer nonzero elements than $P$.
\end{lemma}
\begin{proof}
For any ray $r\colon x\to y$, the commutator $\widehat q r-r\widehat q=\widehat q_xr-r\widehat q_y$ vanishes if and only if $r$ is an identity, an unshiftable ray, or a two-sided ray (for a one-sided ray, $\widehat q_xr-r\widehat q_y=\pm r^+\ne0$, and distinct rays in $A_0(x,y)$ have distinct weights). Thus $C_\chi$ has a basis consisting of all vertex idempotents $1_x$, all two-sided rays, and all unshiftable rays. Among these basis elements, multiplication by $\widehat q$ annihilates precisely the non-supporting idempotents ($\widehat q_x=0$) and the unshiftable rays, so they span the ideal $\mathfrak a_\chi$. Consequently, the supporting idempotents and the two-sided rays project to a basis of $B_\chi$.

Let $\pi\colon A_0\to B_\chi$ be the linear projection onto this basis that sends all one-sided rays, unshiftable rays, and non-supporting idempotents to zero. We claim that $\pi$ is a $C_\chi$-bimodule morphism extending the quotient map $C_\chi\to B_\chi$. Indeed, for any composable rays $r$ and $s$, if their product $rs$ is a retained nonidentity ray in $B_\chi$, then $rs$ has type $\mathsf B$; by the table~\eqref{eq:product-type-table} of Lemma~\ref{lem:successor-substitutions}, the pair $(r,s)$ must have types $(\mathsf B,\mathsf B)$ or $(\mathsf L,\mathsf R)$. If at least one factor belongs to $C_\chi$, the case $(\mathsf L,\mathsf R)$ cannot occur because $C_\chi$ contains no one-sided rays. Hence no discarded ray can multiply an element of $C_\chi$ to produce a retained ray in $B_\chi$. Conversely, if both factors are retained in $B_\chi$ but their product in $A_0$ is unshiftable (or zero), that product lies in $\mathfrak a_\chi$ and vanishes on both sides of the bimodule identity. Since identity factors manifestly respect $\pi$, it follows that $\pi$ is a $C_\chi$-bimodule retraction of the quotient $C_\chi\to B_\chi$ through $C_\chi\hookrightarrow A_0$. By Lemma~\ref{lem:split-representation-transfer}, $B_\chi$ is therefore representation-finite.

Because the retained basis elements of $B_\chi$ inherit associative multiplication, identities, additive weights, endpoint-and-weight separation, and positive categorical degree from $P$, they form a weighted multiplication table $Q_\chi$ in the sense of Definition~\ref{def:weighted-table} with $B_\chi=kQ_\chi$. Its positive-degree ideal is nilpotent with semisimple quotient $\prod_{\widehat q_x\ne0}k$, so $B_\chi$ is basic. Finally, if $B_\chi\ne0$, then $\widehat q\ne0$, and the highest nonzero power $\widehat q^d$ of the nilpotent element $\widehat q$ is unshiftable and hence lies in $\mathfrak a_\chi\setminus\{0\}$. Therefore
\[
  \dim B_\chi<\dim C_\chi\leq\dim A_0,
\]
so $Q_\chi$ has strictly fewer nonzero elements than $P$.
\end{proof}

Write $W_\chi=W(Q_\chi)$ for the scalar weight group of $Q_\chi$. By definition, $W_\chi$ is the abelian group generated by symbols $[r]$ for two-sided rays $r$, subject to the relations $[rs]=[r]+[s]$ whenever $r$, $s$, and $rs$ are all two-sided. Comparing with the $(\mathsf B,\mathsf B)$ row of~\eqref{eq:product-type-table}, we see that these relations are precisely the equations in~\eqref{eq:partial-leibniz-group} in which both inputs and the output are two-sided. Any modulo-two solution of~\eqref{eq:partial-leibniz-group} therefore restricts on two-sided rays to a homomorphism $W_\chi\to\Z/2$, which lifts to a homomorphism $W_\chi\to\Z$ by Lemma~\ref{lem:centralizer-ray-algebra} and the inductive hypothesis of Theorem~\ref{thm:integral-partial-lifting}.

\subsection{Coefficients of generating one-sided rays}

We next determine how the coefficients of generating one-sided rays depend on the coefficients of two-sided rays and one-sided ordinary arrows. The following lemma constrains the two-sided prefixes and suffixes that can appear along an ordinary path representing a one-sided ray.

\begin{lemma}
\label{lem:two-sided-factor-paths}
In Settings~\ref{set:projective-companion} and~\ref{set:fixed-character}, let $b$ be a two-sided ray. If $s$ is a ray of type $\mathsf L$ such that $bs$ is nonzero and shiftable, then every ordinary path representing $b$ consists entirely of two-sided ordinary arrows. Dually, if $r$ has type $\mathsf R$ and $rb$ is nonzero and shiftable, then every ordinary path representing $b$ consists entirely of two-sided ordinary arrows.
\end{lemma}
\begin{proof}
Suppose for a contradiction that $b$ admits an ordinary path containing at least one arrow that is not two-sided, and read the path from left to right up to the first arrow $a$ that is not of type $\mathsf B$. Because $\widehat q b\ne0$, every prefix of the path has a nonzero left action by $\widehat q$. Along the two-sided arrows preceding $a$, this left action commutes across each arrow to reach the source of $a$; hence $\widehat q a\ne0$, which means that $a$ has type $\mathsf L$. Consequently, the prefix $u$ ending at $a$ also has type $\mathsf L$ ($\widehat q u\ne0$ and $u\widehat q=0$). Since $b$ itself is two-sided, the remaining suffix $v$ of the path must be a nonidentity ray, and the factorization $b=uv$ has output of type $\mathsf B$ and first factor $u$ of type $\mathsf L$. By the table~\eqref{eq:product-type-table}, the second factor $v$ must have type $\mathsf R$.

Now consider the product $bs=uvs$. Because $bs$ is shiftable and $b$ is two-sided, we have
\[
  \widehat q bs=b\widehat q s=u(v\widehat q s)\ne0,
\]
so in particular $vs\ne0$ and $v\widehat q s=v^+s=vs^+\ne0$. However, $v$ has type $\mathsf R$ and $s$ has type $\mathsf L$, so both outer actions on $vs$ vanish: $\widehat q(vs)=(\widehat q v)s=0$ and $(vs)\widehat q=v(s\widehat q)=0$. Thus $vs$ cannot have a successor of weight $\ell(vs)+\chi$, contradicting $v\widehat q s\ne0$. The statement for $rb$ follows by passing to the opposite algebra.
\end{proof}

Because every nonidentity ray of $A_0$ has positive categorical degree, repeatedly factoring any decomposable ray into two rays of smaller degree shows that every nonidentity ray can be represented by an ordinary path. We now introduce formal symbols $X_a$ for each ordinary arrow $a$ of type $\mathsf L$ and $Y_b$ for each ordinary arrow $b$ of type $\mathsf R$, and work in the abelian group formed by the direct sum of $W_\chi$ with the free abelian group on these symbols.

\begin{lemma}
\label{lem:one-sided-path-coefficients}
Let $r$ be a generating ray of type $\mathsf L$, choose an ordinary path representing $r$, let $p$ be its maximal two-sided prefix, and let $a$ be the first arrow of the path that has type $\mathsf L$. Dually, for a generating ray $s$ of type $\mathsf R$, choose an ordinary path representing $s$, let $b$ be its last arrow of type $\mathsf R$, and let $p'$ be its maximal two-sided suffix. Then the formal expressions
\begin{equation}
  d_r=[p]+X_a,\qquad d_s=Y_b+[p']
  \label{eq:one-sided-path-coefficients}
\end{equation}
(where an empty prefix or suffix has class $0\in W_\chi$) do not depend on the choice of ordinary paths representing $r$ and $s$, and they satisfy every Leibniz equation in~\eqref{eq:partial-leibniz-group} whose output is a generating one-sided ray. In particular, the first $\mathsf L$-arrow $a$ along an ordinary path for $r$, and the last $\mathsf R$-arrow $b$ along an ordinary path for $s$, are independent of the chosen path.
\end{lemma}
\begin{proof}
We prove the statements for a generating left-sided ray $r$ by induction on its categorical degree $\deg_\JJ(r)$. Any ordinary path representing $r$ must contain at least one arrow that is not two-sided, since a shiftable product of two-sided rays is again two-sided. Because $\widehat q r\ne0$, every prefix of the path has a nonzero left action, and commuting $\widehat q$ across the initial two-sided arrows shows that the first non-two-sided arrow $a$ has $\widehat q a\ne0$ and hence is of type $\mathsf L$. Moreover, every prefix extending past $a$ also has type $\mathsf L$: if some longer proper prefix $u$ had type $\mathsf B$, then in the factorization $r=uv$ the suffix $v$ would have type $\mathsf L$ by~\eqref{eq:product-type-table}, and Lemma~\ref{lem:two-sided-factor-paths} applied to $uv=r$ would force every ordinary path for $u$ to consist solely of two-sided arrows, contradicting the presence of $a$ in $u$.

By Proposition~\ref{prop:controlled-contours}, all ordinary paths representing $r$ begin with the same initial ordinary arrow $u$. If $u$ has type $\mathsf L$, then $p$ is empty, $a=u$ is the first arrow of every path for $r$, and every path yields $d_r=X_u$. If instead $u$ has type $\mathsf B$, write $r=uv$. The suffix rays $v$ obtained from any two ordinary paths for $r$ share the same endpoints and the same weight $\ell(r)-\ell(u)$, so they are the same ray $v$. By Lemma~\ref{lem:generating-paths}, $v$ is a generating ray, and by~\eqref{eq:product-type-table} it has type $\mathsf L$. Since $\deg_\JJ(v)<\deg_\JJ(r)$, the induction hypothesis ensures that the first $\mathsf L$-arrow $a$ of $v$ and the sum $[p_v]+X_a$ are independent of the path chosen for $v$. Prepending $u$ gives $[p]=[u]+[p_v]$ in $W_\chi$, so $[p]+X_a$ is independent of the path for $r$. The statements for a generating right-sided ray $s$ follow by duality.

Finally, for any nonzero product $rs=t$ whose output $t$ is generating and one-sided (say of type $\mathsf L$), we concatenate an ordinary path for $r$ with an ordinary path for $s$ to obtain an ordinary path for $t$. If $r$ has type $\mathsf L$, then $r$ already contains the first $\mathsf L$-arrow $a$ of $t$ and shares the same two-sided prefix $p$, so $d_t=[p]+X_a=d_r$. If $r$ has type $\mathsf B$ and $s$ has type $\mathsf L$, then by Lemma~\ref{lem:two-sided-factor-paths} every arrow in the path of $r$ is two-sided, so the two-sided prefix of $t=rs$ is $rp$, which satisfies $[rp]=[r]+[p]$ in $W_\chi$; hence $d_t=[r]+[p]+X_a=d_r+d_s$. Both cases match~\eqref{eq:product-type-table}, and the right-sided case is dual.
\end{proof}

Consider now a nonzero product $rs=t$ in which $r$ has type $\mathsf L$, $s$ has type $\mathsf R$, and the output $t$ is a \emph{generating} two-sided ray. Concatenating ordinary paths for $r$ and $s$ yields an ordinary path for $t$. Writing $r=pa r'$ and $s=s'bp'$ where $p$ is the two-sided prefix of $r$, $a$ is its first $\mathsf L$-arrow, $b$ is the last $\mathsf R$-arrow of $s$, and $p'$ is its two-sided suffix, we have $t=pfp'$, where the middle subsegment $f=ar's'b$ is a generating ray by Lemma~\ref{lem:generating-paths} and is two-sided because the outer actions on $t$ commute across $p$ and $p'$. Since $[t]=[p]+[f]+[p']$ in $W_\chi$ while $d_r=[p]+X_a$ and $d_s=Y_b+[p']$, subtracting $[p]+[p']$ reduces the Leibniz equation $d_r+d_s=d_t$ to
\begin{equation}
  X_a+Y_b=[f],
  \label{eq:generating-crown-row}
\end{equation}
where $f$ is a generating two-sided ray represented by an ordinary path that begins with $a$ and ends with $b$.

\subsection{Four-point crowns and equal right-hand sides}

Because the same variable $X_a$ or $Y_b$ can appear in several equations of the form~\eqref{eq:generating-crown-row}, we must check that these equations are mutually consistent over $\Z$. Excluding four-point crowns shows that any two equations sharing a variable have the exact same right-hand side $[f]$.

\begin{lemma}
\label{lem:crown-output-comparison}
Any two two-sided rays that admit ordinary paths beginning with the same $\mathsf L$-arrow $a$ are comparable under right factorization. Dually, any two two-sided rays that admit ordinary paths ending with the same $\mathsf R$-arrow $b$ are comparable under left factorization. Consequently, any two equations of the form~\eqref{eq:generating-crown-row} (with generating two-sided outputs) that share a variable $X_a$ or $Y_b$ have the same output ray $f$.
\end{lemma}
\begin{proof}
Let $f$ and $t$ be two-sided rays admitting ordinary paths that start with the same $\mathsf L$-arrow $a\colon x\to y$. Write $f=av$ and $t=aw$. Because $f$ is two-sided, we have $\widehat q_xf=f\widehat q_z=a(v\widehat q_z)\ne0$, so $a$ precedes $\widehat q_xf$ in the faithful projective support at $x$, and likewise $a$ precedes $\widehat q_xt$. Since $\widehat q_x$ obviously precedes both $\widehat q_xf$ and $\widehat q_xt$, all four upward comparisons between the lower pair $\{\widehat q_x,a\}$ and the upper pair $\{\widehat q_xf,\widehat q_xt\}$ hold in the faithful projective support at $x$.

Notice that $\widehat q_x$ and $a$ are incomparable: $\widehat q_x\leq a$ would factor the ordinary arrow $a$ through $\widehat q_x$, while $a\leq\widehat q_x$ would give $ah=\widehat q_x$ for some ray $h\colon y\to x$, forcing $a(ha)=\widehat q_xa\ne0$ and hence $ha=\widehat q_y$ (by primitivity of $\chi$), which would make $a\widehat q_y=\widehat q_xa\ne0$ two-sided rather than left-sided. Nor can either upper element precede $a$ (which would again factor $a$ through $\widehat q_x$) or $\widehat q_x$ (by positive degree). Therefore, if $\widehat q_xf$ and $\widehat q_xt$ were incomparable, the four rays $\{\widehat q_x,a,\widehat q_xf,\widehat q_xt\}$ would form a forbidden four-point crown. Thus $\widehat q_xf$ and $\widehat q_xt$ are comparable, and left cancellation of $\widehat q_x$ shows that $f$ and $t$ are comparable under right factorization.

Suppose now that $f$ and $t$ are \emph{generating} two-sided outputs of two equations~\eqref{eq:generating-crown-row} that share the variable $X_a$. By what we have just shown, $f$ and $t$ are comparable; assume without loss of generality that $t=fh$ for a ray $h$. If $h$ were a nonidentity ray, then writing $f=av$ and $t=aw$ (where $v=r's'b$ and $w$ are generating rays of type $\mathsf R$ by Lemma~\ref{lem:generating-paths} and~\eqref{eq:product-type-table}) and canceling $a$ would give $vh=w$. Because $w=vh$ is shiftable, $v$ has type $\mathsf R$, and $t=fh$ is two-sided (so $h$ has type $\mathsf B$), the dual of Lemma~\ref{lem:two-sided-factor-paths} would imply that every ordinary path representing $h$ consists of two-sided arrows. Concatenating a path for $v$ with a path for $h$ would then produce an ordinary path for $w$ ending in a two-sided arrow, whereas $w$ also has an ordinary path ending in its terminal $\mathsf R$-arrow from~\eqref{eq:generating-crown-row}. This contradicts the dual of Proposition~\ref{prop:controlled-contours}, which requires all ordinary paths representing the generating right-sided ray $w$ to share the same final ordinary arrow. Thus $h$ is an identity and $f=t$. The case of a shared variable $Y_b$ follows dually.
\end{proof}

We also need to incorporate the equations coming from nondeep external returns (Definition~\ref{def:nondeep-return}).

\begin{lemma}
\label{lem:nondeep-boundary-arrows}
Let $el=\widehat q_x^2$ (with $\widehat q_x^3\ne0$) be a nondeep external return at $x$, where $e\colon x\to y$ and $l\colon y\to x$ are generating rays of types $\mathsf L$ and $\mathsf R$, respectively. Then every ordinary path representing $e$ begins with a single left-sided arrow $a$, and every ordinary path representing $l$ ends with a single right-sided arrow $b$ (so both the two-sided prefix of $e$ and the two-sided suffix of $l$ are empty).
\end{lemma}
\begin{proof}
The first arrow $u\colon x\to v$ of an ordinary path for $e$ satisfies $\widehat q_xu\ne0$ (since $\widehat q_xe\ne0$) and therefore has type $\mathsf L$ or $\mathsf B$. Suppose for a contradiction that $u$ has type $\mathsf B$, and write $e=uc$. Because $e$ and $e^+=\widehat q_xe$ are distinct rays in $A_0(x,y)$, we have $\dim A_0(x,y)\geq2$, so Lemma~\ref{lem:generating-paths} implies that $x\ne v$ and that both $u$ and $c$ are generating rays. Since $u$ is generating and two-sided with $\widehat q_xu\ne0$ and $\chi$ is primitive, Lemma~\ref{lem:primitive-action-generation} shows that $A_0(x,v)$ is cyclic at both endpoints. Furthermore, $u(cl)=el=\widehat q_x^2\ne0$, so $uA_0(v,x)\ne0$. Lemma~\ref{lem:primitive-cyclic-corners}(1) therefore provides a generating primitive reverse $g\colon v\to x$ satisfying $ug=\widehat q_x$ and $gu=\widehat q_v$.

Comparing $u(cl)=\widehat q_x^2=u(\widehat q_vg)\ne0$ shows that $cl$ and $\widehat q_vg$ have the same endpoints and weight, so $cl=\widehat q_vg$. Its successor $\widehat q_v^2g$ is nonzero because $u(\widehat q_v^2g)=\widehat q_x^3\ne0$. On the other hand, applying~\eqref{eq:product-type-table} to $e=uc$ shows that $c$ has type $\mathsf L$, while $l$ has type $\mathsf R$ by hypothesis. Thus $cl=\widehat q_vg$ is a proper shiftable product of generating rays of types $(\mathsf L,\mathsf R)$ whose endpoint layer is~$1$, contradicting Proposition~\ref{prop:proper-ray-products}(3) (which forces any such product to have endpoint layer~$2$).

Consequently, the initial arrow of every ordinary path for $e$ has type $\mathsf L$, and by Proposition~\ref{prop:controlled-contours} (or Lemma~\ref{lem:one-sided-path-coefficients}) this initial arrow $a$ is independent of the path. The statement for $l$ follows by passing to the opposite algebra.
\end{proof}

By Lemma~\ref{lem:nondeep-boundary-arrows}, if $(e,l)$ is a nondeep external return at $x$ with initial arrow $a$ of $e$ and terminal arrow $b$ of $l$, then $d_e=X_a$ and $d_l=Y_b$. Since $\widehat q_x$ is two-sided and $\widehat q_x^3\ne0$, we have $[\widehat q_x^2]=2[\widehat q_x]$ in $W_\chi$, so the Leibniz equation $d_e+d_l=d_{\widehat q_x^2}$ for $el=\widehat q_x^2$ becomes
\begin{equation}
  X_a+Y_b=2[\widehat q_x].
  \label{eq:nondeep-coefficient-row}
\end{equation}
Moreover, neither $X_a$ nor $Y_b$ can appear in any generating-output equation~\eqref{eq:generating-crown-row}. Indeed, if a generating two-sided ray $f$ had an ordinary path starting with $a$, Lemma~\ref{lem:crown-output-comparison} would make $f$ and $\widehat q_x^2$ comparable under right factorization; however, $\widehat q_x^2h=f$ would contradict the fact that $f$ is generating (Lemma~\ref{lem:generating-paths}), while $fh=\widehat q_x^2$ would give $f$ (and hence its initial arrow $a$) a nonzero return to $x$, forcing $a$ to be two-sided by Lemma~\ref{lem:primitive-action-generation} and Lemma~\ref{lem:primitive-cyclic-corners}(1), a contradiction. The dual argument excludes $Y_b$ from~\eqref{eq:generating-crown-row}. Finally, if two nondeep-return equations~\eqref{eq:nondeep-coefficient-row} share $X_a$, their base vertex $x$ is the source of $a$, and if they share $Y_b$, their base vertex $x$ is the target of $b$; in either case their right-hand sides $2[\widehat q_x]$ coincide.

We now form the bipartite graph whose vertices are the formal variables $X_a$ and $Y_b$ and whose edges correspond to the equations~\eqref{eq:generating-crown-row} and~\eqref{eq:nondeep-coefficient-row}, labeled by their right-hand sides in $W_\chi$. By Lemma~\ref{lem:crown-output-comparison} and the preceding paragraph, any two edges incident to a common vertex carry the same label $\omega\in W_\chi$, so the edge labels are constant on each connected component of the bipartite graph. On a connected component with constant edge label $\omega$, the system of equations $X_a+Y_b=\omega$ has the general solution
\[
  X_a=\omega+z,\qquad Y_b=-z,
\]
parametrized by a single free variable $z$ per connected component (with an isolated vertex likewise contributing one free variable). Letting $T$ denote the quotient of the direct sum of $W_\chi$ and the free abelian group on all symbols $X_a,Y_b$ by the relations~\eqref{eq:generating-crown-row} and~\eqref{eq:nondeep-coefficient-row}, and writing $m$ for the number of connected components of the bipartite graph, we obtain an isomorphism
\begin{equation}
  T\cong W_\chi\oplus\Z^m.
  \label{eq:crown-coefficient-group}
\end{equation}
In particular, even when $\omega=2[\widehat q_x]$ in~\eqref{eq:nondeep-coefficient-row}, the equation $X_a+Y_b=2[\widehat q_x]$ imposes no modulo-two divisibility obstruction on lifting $X_a$ and $Y_b$, because each variable appears with coefficient~$1$ and the right-hand side is constant across each connected component.

\subsection{Verification of all nonzero-product equations}
\label{subsec:verify-nonzero-equations}

We now complete the proof of Theorem~\ref{thm:integral-partial-lifting} by assigning coefficients in the group $T\cong W_\chi\oplus\Z^m$ of~\eqref{eq:crown-coefficient-group} to all rays of $A_0$ and verifying that they satisfy every nonzero-product Leibniz equation~\eqref{eq:partial-leibniz-group}.

\begin{proof}[Proof of Theorem~\ref{thm:integral-partial-lifting}]
We first define a universal solution with values in the group $T$ of~\eqref{eq:crown-coefficient-group}:
\begin{itemize}
\item For every two-sided ray $b$, set $d_b=[b]\in W_\chi\subseteq T$.
\item For every generating one-sided ray $r$ (of type $\mathsf L$) or $s$ (of type $\mathsf R$), define $d_r$ and $d_s$ in $T$ by the path formulas~\eqref{eq:one-sided-path-coefficients}.
\item For every proper shiftable one-sided ray, write it as $\widehat q_x^jr$ or $s\widehat q_z^j$ with $r,s$ generating and $j\geq1$, and set
\begin{equation}
  d_{\widehat q_x^jr}=j[\widehat q_x]+d_r,\qquad
  d_{s\widehat q_z^j}=d_s+j[\widehat q_z]
  \label{eq:proper-successor-coefficients}
\end{equation}
(note that $\widehat q_x^{j+1}r\ne0$ implies $\widehat q_x^2\ne0$, so $\widehat q_x$ is two-sided and $[\widehat q_x]\in W_\chi$ is defined, and dually for $[\widehat q_z]$).
\item For all vertex idempotents and unshiftable rays, set $d_r=0$.
\end{itemize}
Work in the integral ray algebra $\Z P$ (which is free over $\Z$ on the nonzero rays of $P$, with basis products having coefficients $0$ or $1$), let $\delta\colon\Z P\to(\Z P)\otimes_\Z T$ be the resulting normalized $\Z$-linear one-cochain $\delta(r)=r^+\otimes d_r$, and let $\mu=d\delta$ be its Hochschild coboundary with values in $(\Z P)\otimes_\Z T\cong\bigoplus_{r\in P\setminus\{0\}}T$. We verify that $\mu(r,s)=0$ in $(\Z P)\otimes_\Z T$ for every pair of composable rays $(r,s)$ with $rs\ne0$ by checking five exhaustive cases.

\smallskip
\noindent\textbf{Case 1: One of the factors is a local ray.}
When $r$ or $s$ is an identity, $\mu(r,s)=0$ by normalization. Suppose that the left factor is a nonidentity local ray $c=q_x^j$ ($j\geq1$) with $cs\ne0$. If the output $cs$ is unshiftable, then $\mu(c,s)=0$ by Lemma~\ref{lem:successor-substitutions}. Assume therefore that $cs$ is shiftable.

If $x$ is a supporting vertex ($q_x=\widehat q_x$), then $s$ cannot have type $\mathsf R$ (since $q_xs\ne0$), so by~\eqref{eq:product-type-table} the local ray $c=\widehat q_x^j$ has type $\mathsf B$ and $s$ has type $\mathsf B$ or $\mathsf L$. When $s$ has type $\mathsf B$, the relation $d_{cs}=d_c+d_s$ holds in $W_\chi$. When $s$ has type $\mathsf L$, write $s=\widehat q_x^if$ with $f$ generating ($A_0(x,y)$ is cyclic at $x$ by Lemma~\ref{lem:primitive-action-generation}). Note that $f$ also has type $\mathsf L$, since if $f\widehat q_y=\widehat q_xf$ were nonzero, then $s\widehat q_y=\widehat q_x^{i+1}f=\widehat q_xs\ne0$ would make $s$ two-sided. Because $c=\widehat q_x^j$ is a product of $j$ copies of $\widehat q_x$ in $W_\chi$, we have $d_c=j[\widehat q_x]$, and~\eqref{eq:proper-successor-coefficients} gives
\[
  d_{cs}=(j+i)[\widehat q_x]+d_f=j[\widehat q_x]+\bigl(i[\widehat q_x]+d_f\bigr)=d_c+d_s,
\]
matching~\eqref{eq:product-type-table}.

If $x$ is not a supporting vertex ($\ell(q_x)\ne\chi$), we claim that no nonzero product $cs=q_x^js$ can be shiftable. Indeed, if the corner bimodule $A_0(x,y)$ containing $s$ is cyclic at $x$, then adjacent rays in $A_0(x,y)$ differ in weight by $\ell(q_x)\ne\chi$, and since $\chi$ is primitive no multiple of $\ell(q_x)$ can equal $\chi$; hence $A_0(x,y)$ contains no shiftable rays at all. If instead $A_0(x,y)$ is cyclic only at $y$ with generating ray $f$, then $A_0(x,y)$ can contain shiftable rays only when $\ell(q_y)=\chi$, in which case $\ell(q_x)\ne\ell(q_y)$ forces any nonzero left action $q_xf=fq_y^n$ to have $n\geq2$. By Lemma~\ref{lem:unequal-endpoint-action} (applied in the opposite algebra), $q_xf$ then lies in the endpoint socle of $A_0(x,y)$, so $cs=q_x^js$ is again unshiftable (and $cs^+=0$). Thus $\mu(c,s)=0$ in all cases, and the dual argument gives $\mu(r,c)=0$ for any local ray $c$ on the right.

Moreover, if $cr=0$ or $rc=0$ for a local ray $c$, then $c*r=0$ and $r*c=0$ by Corollary~\ref{cor:active-zero-pairs} (since a nonidentity local ray $c=q_x^j$ is never the generating ray of a non-cyclic corner), so $\mu(c,r)=0$ and $\mu(r,c)=0$ even when the product vanishes. Expanding the cocycle identity $d\mu=0$ on the triples $(c,r,s)$, $(r,s,c)$, and $(r,c,s)$ for any local ray $c$ therefore yields the local-linearity relations in $(\Z P)\otimes_\Z T$:
\begin{equation}
  \mu(cr,s)=c\mu(r,s),\qquad
  \mu(r,sc)=\mu(r,s)c,\qquad
  \mu(rc,s)=\mu(r,cs).
  \label{eq:local-ring-relativity}
\end{equation}

\smallskip
\noindent\textbf{Case 2: Generating outputs or two-sided inputs.}
Every nonzero product whose output is generating and one-sided satisfies $\mu(r,s)=0$ by Lemma~\ref{lem:one-sided-path-coefficients}, and every nonzero product of types $(\mathsf L,\mathsf R)$ with a generating two-sided output is one of the defining relations~\eqref{eq:generating-crown-row} of $T$. Likewise, every product of two-sided rays (whether generating or proper) with a shiftable output has a two-sided output and holds by the defining relations of $W_\chi$.

\smallskip
\noindent\textbf{Case 3: Generating inputs of types $(\mathsf B,\mathsf L)$ or $(\mathsf R,\mathsf B)$ with a proper output.}
Suppose that $r\colon x\to y$ of type $\mathsf B$ and $s\colon y\to z$ of type $\mathsf L$ are generating rays whose product $rs$ is proper and shiftable. By Proposition~\ref{prop:proper-ray-products}(2), $rs=\widehat q_xf$ has endpoint layer~$1$ (where $f\colon x\to z$ is generating of type $\mathsf L$) and $r$ admits a primitive reverse $u\colon y\to x$ with $ru=\widehat q_x$ and $ur=\widehat q_y$. Since $r(uf)=(ru)f=\widehat q_xf=rs\ne0$, the product $uf$ is a nonzero ray in $A_0(y,z)$ with the same endpoints and weight as $s$, so $uf=s$ by weight uniqueness in $A_0(y,z)$. Moreover, $\widehat q_ys=\widehat q_yuf=u\widehat q_xf=urs\ne0$, so $u$ has type $\mathsf B$. Because $s=uf$ is a generating left-sided output with $u$ of type $\mathsf B$ and $f$ of type $\mathsf L$, Case~2 gives $d_s=[u]+d_f$. Since $[\widehat q_x]=[r]+[u]$ in $W_\chi$, formula~\eqref{eq:proper-successor-coefficients} yields
\[
  d_{rs}=[\widehat q_x]+d_f=[r]+[u]+d_f=d_r+d_s,
\]
so $\mu(r,s)=0$. The case of types $(\mathsf R,\mathsf B)$ follows dually.

\smallskip
\noindent\textbf{Case 4: Generating inputs of types $(\mathsf L,\mathsf R)$ with a proper output.}
Let $r\colon x\to y$ of type $\mathsf L$ and $s\colon y\to z$ of type $\mathsf R$ be generating rays whose product $rs$ is proper and shiftable. By Proposition~\ref{prop:proper-ray-products}(3), if $x=z$, then $(r,s)$ is a nondeep external return at $x$, so $\mu(r,s)=0$ is one of the defining relations~\eqref{eq:nondeep-coefficient-row} of $T$. If $x\ne z$, then Proposition~\ref{prop:proper-ray-products}(3)(b) gives either $r=fa$ and $as=\widehat q_z^2$ (with $f,a$ generating and $(a,s)$ a nondeep external return at $z$) or the dual factorization. In the first case, expanding $d\mu(f,a,s)=0$ gives
\[
  f\mu(a,s)-\mu(fa,s)+\mu(f,as)-\mu(f,a)s=0.
\]
Here $\mu(a,s)=0$ by the local-return case just noted, $\mu(f,as)=\mu(f,\widehat q_z^2)=0$ by Case~1, and $\mu(f,a)=0$ by Case~2 because $fa=r$ is a generating left-sided ray. Hence $\mu(r,s)=\mu(fa,s)=0$, and the dual factorization is handled symmetrically.

\smallskip
\noindent\textbf{Case 5: Proper inputs.}
Suppose finally that at least one of $r$ or $s$ is a proper ray and $rs$ is shiftable.
\begin{itemize}
\item If the output $rs$ has type $\mathsf B$, then $(r,s)$ has types $(\mathsf B,\mathsf B)$ (already covered in Case~2) or $(\mathsf L,\mathsf R)$. For $(\mathsf L,\mathsf R)$, write $r=\widehat q_x^ir_0$ and $s=s_0\widehat q_z^j$ with $r_0,s_0$ generating; the local-linearity relations~\eqref{eq:local-ring-relativity} give $\mu(r,s)=\widehat q_x^i\mu(r_0,s_0)\widehat q_z^j=0$ by Case~2 or Case~4 applied to $(r_0,s_0)$.
\item If the output $rs$ has type $\mathsf L$ and $r$ has type $\mathsf B$, write $r=\widehat q_x^ir_0$ and $s=\widehat q_y^js_0$ with $r_0,s_0$ generating; then~\eqref{eq:local-ring-relativity} and the two-sidedness of $r_0$ give $\mu(r,s)=\widehat q_x^i\mu(r_0\widehat q_y^j,s_0)=\widehat q_x^{i+j}\mu(r_0,s_0)=0$ by Case~2 or Case~3 applied to $(r_0,s_0)$.
\item If both $r$ and $rs$ have type $\mathsf L$, write $r=\widehat q_x^ir_0$ with $r_0$ generating. By Proposition~\ref{prop:proper-ray-products}(1), $rs$ also has endpoint layer $i$, which means $rs=\widehat q_x^i(r_0s)$ with $r_0s$ generating (and $s$ is generating by Lemma~\ref{lem:generating-paths}). Then~\eqref{eq:local-ring-relativity} gives $\mu(r,s)=\widehat q_x^i\mu(r_0,s)=0$ by Case~2.
\end{itemize}
The remaining cases with output of type $\mathsf R$ follow by duality. By Lemma~\ref{lem:successor-substitutions}, this exhausts all nonzero products with shiftable output, so our universal coefficients in $T$ satisfy every equation in~\eqref{eq:partial-leibniz-group}.

\smallskip
\noindent\textbf{Lifting a prescribed modulo-two solution.}
Now let $(\bar d_r)$ be any solution of~\eqref{eq:partial-leibniz-group} over $\Z/2$. Its restriction to two-sided rays defines a homomorphism $\bar w\colon W_\chi\to\Z/2$, which lifts by Lemma~\ref{lem:centralizer-ray-algebra} and our inductive hypothesis to a homomorphism $w\colon W_\chi\to\Z$.

Evaluating the modulo-two equations~\eqref{eq:product-type-table} along the chosen ordinary paths for generating one-sided rays $r$ and $s$ forces
\[
  \bar d_r=\bar d_p+\bar d_a,\qquad \bar d_s=\bar d_b+\bar d_{p'},
\]
in the notation of~\eqref{eq:one-sided-path-coefficients}, and evaluating them on local products forces $\bar d_{\widehat q_x^jr}=j\bar d_{\widehat q_x}+\bar d_r$ and $\bar d_{s\widehat q_z^j}=\bar d_s+j\bar d_{\widehat q_z}$ as in~\eqref{eq:proper-successor-coefficients}. Moreover, splitting a generating two-sided ray $f$ of~\eqref{eq:generating-crown-row} after its first $\mathsf L$-arrow shows that $\bar d_a+\bar d_b=\bar d_f=\bar w([f])$, while for any nondeep external return $el=\widehat q_x^2$ Lemma~\ref{lem:nondeep-boundary-arrows} gives $\bar d_a+\bar d_b=\bar d_e+\bar d_l=\bar d_{\widehat q_x^2}=2\bar d_{\widehat q_x}=0=\bar w(2[\widehat q_x])$. Thus the modulo-two arrow coefficients $(\bar d_a,\bar d_b)$ satisfy every edge relation of the bipartite graph.

On each connected component of the bipartite graph with constant edge label $\omega\in W_\chi$, the modulo-two values therefore take the form $\bar d_a=\bar w(\omega)+\bar z$ and $\bar d_b=-\bar z$ for a unique element $\bar z\in\Z/2$. Choosing an integer lift $z\in\Z$ of $\bar z$, we specialize the component's formal variables to $X_a=w(\omega)+z$ and $Y_b=-z$ (and for each isolated vertex we lift its modulo-two value directly to $\Z$). Together with $w\colon W_\chi\to\Z$, these choices define a group homomorphism $T\to\Z$ whose composition with the universal coefficients in $T$ yields an integer solution $(d_r)$ of~\eqref{eq:partial-leibniz-group} reducing modulo two to $(\bar d_r)$.
\end{proof}

\section{Hochschild representatives from the integers modulo four}
\label{sec:integral-deformations}

We now construct a cocycle generator for each one-dimensional weight space $\HH^2(A_0)_{\chi_i}$ ($i\in I_0$) that is supported strictly on zero ray products $(r,s)$ with $rs=0$ and takes values in their stem-product lines $U_{r*s}$. Because $A_0=k\otimes_{\F_2}\F_2P$ and Hochschild cohomology commutes with field extensions, it suffices to construct these cocycles over the prime field $\F_2$; their $k$-linear combinations will then span all of $\HH^2(A_0)_{>0}$.

The strategy is based on the Bockstein connecting homomorphism for the short exact coefficient sequence $0\to\F_2\xrightarrow{2}\mathcal O\to\F_2\to0$, where
\[
  \F_2=\Z/2,\qquad \mathcal O=\Z/4,
\]
and the inclusion $\F_2\hookrightarrow\mathcal O$ is given by $a\mapsto 2a$. Setting $A^{(2)}=\F_2P$ and $A^{(4)}=\mathcal OP$, we show in Proposition~\ref{prop:integral-bockstein} that every cohomology class in $\HH^2(A^{(2)})_{>0}$ lies in the image of the connecting map from $\HH^1(A^{(2)})_{>0}$, and is therefore represented by $(d\delta)/2$ for any $\mathcal O$-linear one-cochain $\delta$ lifting a derivation $\bar\delta$ of $A^{(2)}$. By Theorem~\ref{thm:integral-partial-lifting}, the lift $\delta$ can be chosen so that $d\delta$ vanishes on all nonzero ray products, which forces $(d\delta)/2$ to be supported exclusively on pairs of rays with zero product. Throughout this section, all cochains are normalized with respect to the vertex idempotents and are homogeneous with respect to both the categorical grading and the weight lattice $L$. For background on the Bockstein homomorphism in Hochschild cohomology, see \cite[Sections~4--5]{GerstenhaberSchack} and \cite[Section~2]{Linckelmann}.

\subsection{The model over the two-element field and its module lifts}

The all-ones mesh relations of Section~\ref{sec:weights-and-parameters} define a graded $\F_2$-linear category $\mathcal C^{(2)}$ satisfying $k\otimes_{\F_2}\mathcal C^{(2)}=\C$. Its full subcategory on the projective objects $P_x$ has opposite endomorphism algebra $\End(\bigoplus_x P_x)^{\op}=A^{(2)}$, since each nonzero homogeneous component of $A^{(2)}(x,y)$ is one-dimensional over $\F_2$ and every nonzero product of basis rays has coefficient~$1$. For each object $X_i$ of $\mathcal C^{(2)}$, define the graded right $A^{(2)}$-module $M_i^{(2)}$ by
\[
  M_i^{(2)}e_x=\mathcal C^{(2)}(P_x,X_i).
\]
Evaluation on the projective objects identifies $\mathcal C^{(2)}$ with the full graded homomorphism category of the modules $M_i^{(2)}$, because it does so after the faithfully flat base change $-\otimes_{\F_2}k$. Since tensoring a finitely generated projective resolution with $k$ commutes with graded $\Ext$, Lemma~\ref{lem:positive-ext-vanishing} implies that
\begin{equation}
  \Ext^n_{A^{(2)}}(M_i^{(2)},M_j^{(2)})_d=0\qquad(n\geq1,\ d\geq0).
  \label{eq:prime-field-ext-vanishing}
\end{equation}
When $X_i=P_x$ is a projective object, $M_i^{(2)}=e_xA^{(2)}$ is the corresponding indecomposable projective $A^{(2)}$-module.

\begin{proposition}
\label{prop:fixed-integral-companion}
The graded $A^{(2)}$-modules $M_i^{(2)}$ lift to graded $A^{(4)}$-modules $M_i^{(4)}$ that are free over $\mathcal O=\Z/4$, and their full graded homomorphism category $E_4$ has the following properties:
\begin{enumerate}[label={\upshape(\arabic*)}]
\item Each morphism space of $E_4$ is a free $\mathcal O$-module, and $E_4/2E_4=\mathcal C^{(2)}$.
\item At the projective objects $X_i=P_x$, we have $M_i^{(4)}=e_xA^{(4)}$, so the projective corner of $E_4$ has opposite endomorphism algebra $\End(\bigoplus_x e_xA^{(4)})^{\op}=A^{(4)}$.
\item The category $E_4$ admits a quadratic scalar mesh presentation over $\mathcal O$ whose coefficients are units reducing modulo~$2$ to the all-ones mesh presentation of $\mathcal C^{(2)}$.
\end{enumerate}
Moreover, all of these lifts may be chosen to preserve the simultaneous integer weights in $L$: lifting a basis of $H_\Z$ to mesh weights via Lemma~\ref{lem:integer-weight-lifting} equips the modules $M_i^{(2)}$ with compatible $L$-gradings, and all obstruction and lifting steps can be carried out within each weight component.
\end{proposition}
\begin{proof}
Choose arbitrary degree-preserving lifts of the action matrices of $A^{(2)}$ on $M_i^{(2)}$ to free graded $\mathcal O$-modules. Their associativity defect is divisible by~$2$; dividing it by~$2$ and reducing modulo~$2$ produces a degree-zero cocycle representing a class in $\Ext^2_{A^{(2)}}(M_i^{(2)},M_i^{(2)})_0$, which vanishes by~\eqref{eq:prime-field-ext-vanishing}. Subtracting $2$ times a bounding one-cochain kills the defect modulo~$4$ (since $(2\mathcal O)^2=0$), yielding an associative graded $A^{(4)}$-module structure $M_i^{(4)}$. The same bar-resolution calculation using the vanishing of $\Ext^1$ in nonnegative degrees shows that $M_i^{(4)}$ is unique up to isomorphism and that every homogeneous morphism of nonnegative degree between the modules $M_i^{(2)}$ lifts to $E_4$. At the projective objects $X_i=P_x$, we take $M_i^{(4)}=e_xA^{(4)}$ to be the regular projective $A^{(4)}$-modules, which gives~(2).

To verify~(1), fix any pair of objects $X_i,X_j$ and lift an $\F_2$-basis $\bar f_1,\ldots,\bar f_r$ of $\mathcal C^{(2)}(X_i,X_j)_d$ ($d\geq0$) to homogeneous $A^{(4)}$-module homomorphisms $f_1,\ldots,f_r\in E_4(X_i,X_j)_d$. Any $g\in E_4(X_i,X_j)_d$ that reduces to zero modulo~$2$ has matrix entries in $2\mathcal O$; dividing by~$2$ gives an $A^{(2)}$-module homomorphism of degree~$d$, which is an $\F_2$-linear combination of $\bar f_1,\ldots,\bar f_r$, so $f_1,\ldots,f_r$ span $E_4(X_i,X_j)_d$ over $\mathcal O$. Likewise, if $\sum_{m=1}^r a_mf_m=0$ with $a_m\in\mathcal O$, reducing modulo~$2$ shows that $a_m=2b_m$ for all $m$, and dividing by~$2$ and reducing modulo~$2$ again shows that every $b_m$ is even, so $a_m=0$ in $\mathcal O$. The same two-step reduction shows that $E_4(X_i,X_j)_d=0$ for $d<0$.

For~(3), Nakayama's lemma implies that the degree-one morphisms generate $E_4$ as an $\mathcal O$-linear category. Because $(E_4)_2$ is a free $\mathcal O$-module, the composition map from length-two paths onto $(E_4)_2$ splits, so its kernel is a free $\mathcal O$-module direct summand. Lifting the mesh relations of $\mathcal C^{(2)}$ to a basis of this kernel defines a quadratic $\mathcal O$-linear category $\widetilde E_4$ equipped with a graded surjection $\widetilde E_4\to E_4$ that reduces modulo~$2$ to the identity on $\mathcal C^{(2)}$. Each graded piece of $\widetilde E_4$ is finitely generated over $\mathcal O$, and in degrees above the top degree of $\mathcal C^{(2)}$ its reduction modulo~$2$ is zero, so $N=2N$ and hence $N=4N=0$; thus $\widetilde E_4$ is finite over $\mathcal O$. Since $E_4$ is free over $\mathcal O$, the surjection $\widetilde E_4\to E_4$ splits as a map of $\mathcal O$-modules, and its kernel vanishes by Nakayama's lemma because it reduces to zero modulo~$2$. Finally, because each arrow $\tau X\to Y$ in a mesh of $\mathcal C^{(2)}$ has a unique mate $Y\to X$ and the irreducible morphism spaces are one-dimensional, every length-two path between $\tau X$ and $X$ is a mesh path, so the lifted relations are quadratic scalar mesh relations with unit coefficients in $\mathcal O^\times$.
\end{proof}

\subsection{Comparison with the mesh derivative}

Let $Q_m=\operatorname{coker}D_m(E_4)$ denote the cokernel of the mesh derivative of $E_4$ in degree $m$, as in Lemma~\ref{lem:integral-mesh-cokernels} of Appendix~\ref{app:integral-corrections}. For a Hochschild two-cocycle of $A^{(2)}$ of internal degree $q>0$, the corresponding arrow-correction degree on the mesh category is $m=q+1$, and Lemma~\ref{lem:integral-mesh-cokernels} gives
\begin{equation}
  2Q_{q+1}=0.
  \label{eq:mod-four-mesh-annihilator}
\end{equation}
For any cochain $\bar\mu$ over $\F_2$, we write $2\widetilde\mu$ for twice an arbitrary coefficient lift $\widetilde\mu$ of $\bar\mu$ to $\mathcal O=\Z/4$; since $2\mathcal O\cong\F_2$, the cochain $2\widetilde\mu$ depends only on $\bar\mu$.

\begin{lemma}
\label{lem:mod-four-comparison}
Let $\bar\mu$ be a normalized Hochschild two-cocycle of $A^{(2)}$ of internal degree $q>0$. Then $2\widetilde\mu$ is a coboundary in the normalized Hochschild complex of $A^{(4)}$. The same statement holds within each weight component $\chi\in L$ with $\chi(\deg_\JJ)=q>0$.
\end{lemma}
\begin{proof}
Set $R=\mathcal O[\varepsilon]/(\varepsilon^2)$ with $\deg(\varepsilon)=-q$, and equip $B=A^{(4)}\otimes_{\mathcal O}R$ with the deformed multiplication $m+2\varepsilon\widetilde\mu$. Its associator is $2\varepsilon\,d\widetilde\mu=0$ because $d\bar\mu=0$ over $\F_2$, so $B$ is an associative graded $R$-algebra.

Over the reduced ring $\F_2[\varepsilon]/(\varepsilon^2)$, we may apply Lemma~\ref{lem:square-zero-selected-lifts} to the first-order deformation $m+\varepsilon\bar\mu$ of $A^{(2)}$ and the modules $M_i^{(2)}$, since the required $\Ext^2$ and $\Ext^1$ groups vanish by~\eqref{eq:prime-field-ext-vanishing}. Let $\bar\nu_i$ be the resulting first-order corrections to the module actions and let $\bar g_a$ be the corrections to the generating arrow morphisms. On the free $R$-modules $M_i^{(4)}\otimes_{\mathcal O}R$, we then define deformed actions $\nu_i^{(4)}+2\varepsilon\widetilde\nu_i$ and deformed arrow maps $a^{(4)}+2\varepsilon\widetilde g_a$. Because every correction term carries a factor of $2\varepsilon$ and $(2\varepsilon)^2=0$, the module-action and module-homomorphism conditions depend only on the reductions $\bar\nu_i$ and $\bar g_a$ modulo~$2$ and are therefore satisfied over $R$. At the projective objects $P_x$, we take the lifted modules to be the regular projective $B$-modules $e_xB$, whose action corrections match the deformed multiplication $m+2\varepsilon\widetilde\mu$ of $B$.

Now evaluate the quadratic scalar mesh relations of $E_4$ on the deformed arrow maps $a^{(4)}+2\varepsilon\widetilde g_a$. The constant terms vanish in $E_4$, and because the unit coefficients of the mesh relations in $E_4$ reduce to~$1$ modulo~$2$, the remaining residual error at each mesh ending at $X$ is $2\varepsilon\widetilde z_X$, where $\bar z_X$ is the mesh error of the modulo-two lift. Choose a homogeneous morphism $z_X\in E_4(\tau X,X)_{q+2}$ lifting $\bar z_X$; since the underlying representation spaces are free over $\mathcal O$, evaluating $2\varepsilon z_X$ produces precisely the residual error $2\varepsilon\widetilde z_X$. By~\eqref{eq:mod-four-mesh-annihilator}, the class of $2z=(2z_X)$ vanishes in $Q_{q+1}=\operatorname{coker}D_{q+1}(E_4)$, so there exist homogeneous $A^{(4)}$-module morphisms $f_a\in E_4(\operatorname{source}a,\operatorname{target}a)_{q+1}$ satisfying
\[
  D_{q+1}(f)=2z.
\]
Replacing the deformed arrow maps by $a^{(4)}+2\varepsilon\widetilde g_a-\varepsilon f_a$ preserves the $B$-module homomorphism property (since $f_a$ is $A^{(4)}$-linear and $\varepsilon^2=0$) and subtracts $\varepsilon D_{q+1}(f)=2\varepsilon z$ from the mesh errors, killing them completely. The corrected arrow maps therefore satisfy all relations of $E_4$ and define an $R$-linear functor from $E_4\otimes_{\mathcal O}R$ to the category of lifted $B$-modules.

Restricting this functor to the projective objects $P_x$ and passing to opposite endomorphism algebras via Proposition~\ref{prop:fixed-integral-companion}(2) yields an $R$-algebra homomorphism
\[
  A^{(4)}\otimes_{\mathcal O}R\longrightarrow B
\]
that reduces to the identity modulo $\varepsilon$ and fixes the vertex idempotents. On the underlying free $R$-module it has the form $\id+\varepsilon h$ for a normalized one-cochain $h$ of $A^{(4)}$, and multiplicativity $(\id+\varepsilon h)(ab)=(\id+\varepsilon h)(a)\cdot_B(\id+\varepsilon h)(b)$ expands to
\[
  h(ab)=h(a)b+ah(b)+2\widetilde\mu(a,b),
\]
so $2\widetilde\mu=-dh$ is a coboundary. Giving $\varepsilon$ weight $-\chi$ keeps the entire construction homogeneous of weight zero, so $h$ may be taken in weight $\chi$.
\end{proof}

\begin{proposition}
\label{prop:integral-bockstein}
The Bockstein connecting homomorphism associated to the short exact coefficient sequence $0\to\F_2\xrightarrow{2}\mathcal O\to\F_2\to0$ is surjective in every positive internal degree $q>0$:
\[
  \beta\colon\HH^1(A^{(2)})_q\twoheadrightarrow\HH^2(A^{(2)})_q.
\]
Explicitly, for any derivation $\bar\delta$ of $A^{(2)}$ and any normalized $\mathcal O$-linear one-cochain $\delta$ of $A^{(4)}$ lifting $\bar\delta$, the class $\beta([\bar\delta])$ is represented by $(d\delta)/2$, where division by~$2$ denotes the inverse of the additive isomorphism $\F_2\xrightarrow{\sim}2\mathcal O$. The same statement holds within each weight component $\chi\in L$ with $\chi(\deg_\JJ)>0$.
\end{proposition}
\begin{proof}
The short exact sequence of normalized Hochschild cochain complexes
\[
  0\longrightarrow C^*(A^{(2)})\xrightarrow{\,2\,}C^*(A^{(4)})\longrightarrow C^*(A^{(2)})\longrightarrow0
\]
induces a long exact sequence in cohomology. By Lemma~\ref{lem:mod-four-comparison}, multiplication by~$2$ from $\HH^2(A^{(2)})_q$ to $\HH^2(A^{(4)})_q$ is zero for every $q>0$, so the connecting homomorphism $\beta\colon\HH^1(A^{(2)})_q\to\HH^2(A^{(2)})_q$ is surjective. By the standard definition of the connecting homomorphism, $\beta([\bar\delta])$ is represented by $(d\delta)/2$ for any cochain lift $\delta\in C^1(A^{(4)})$, and all maps respect the weight grading by $L$.
\end{proof}

\subsection{A coefficient criterion for supported representatives}

Before combining Proposition~\ref{prop:integral-bockstein} with Theorem~\ref{thm:integral-partial-lifting}, we isolate the elementary homological principle that produces cocycle representatives supported on a prescribed subset of basis elements. It holds for any prime $p$; we will apply it with $p=2$.

\begin{lemma}
\label{lem:supported-bockstein}
Let $p$ be a prime, let $C^*$ be a cochain complex of free $\Z/p^2$-modules, and set $\overline C^*=C^*/pC^*$. Let $U\subseteq C^2$ be a submodule spanned by a subset of a fixed $\Z/p^2$-basis of $C^2$, and assume that:
\begin{enumerate}[label={\upshape(\arabic*)}]
\item the map $H^2(\overline C^*)\xrightarrow{p}H^2(C^*)$ induced by the inclusion $\Z/p\hookrightarrow\Z/p^2$ is zero;
\item every one-cocycle $\bar\delta\in Z^1(\overline C^*)$ lifts to a one-cochain $\delta\in C^1$ satisfying $d\delta\in U$.
\end{enumerate}
Then every cohomology class in $H^2(\overline C^*)$ is represented by a two-cocycle lying in the subspace $U/pU\subseteq\overline C^2$.
\end{lemma}
\begin{proof}
By hypothesis~(1) and the long exact cohomology sequence for $0\to\overline C^*\xrightarrow{p}C^*\to\overline C^*\to0$, the Bockstein connecting map $H^1(\overline C^*)\to H^2(\overline C^*)$ is surjective. Given any class in $H^2(\overline C^*)$, choose a preimage represented by $\bar\delta\in Z^1(\overline C^*)$ and lift $\bar\delta$ to $\delta\in C^1$ with $d\delta\in U$ via~(2). Since $d\bar\delta=0$ in $\overline C^2$, the coboundary $d\delta$ lies in $U\cap pC^2=pU$ (as $U$ is spanned by a subset of a basis of $C^2$). Dividing $d\delta$ by $p$ therefore yields a representative of the given cohomology class in $U/pU$.
\end{proof}

\subsection{Moving the representative onto zero products}

Combining Proposition~\ref{prop:integral-bockstein} with the integral lifting of the nonzero-product Leibniz equations in Theorem~\ref{thm:integral-partial-lifting} now yields a cocycle generator for each residual line $\HH^2(A_0)_{\chi_i}$ ($i\in I_0$) that vanishes on all nonzero ray products and takes values in the stem-product lines on zero ray products.

\begin{corollary}
\label{cor:zero-supported-coverage}
Assume that Proposition~\ref{prop:table-weight-lifting} holds for all weighted multiplication tables with fewer than $\dim A_0$ nonzero elements, as in Theorem~\ref{thm:integral-partial-lifting}. Then for each $i\in I_0$, the one-dimensional weight space $\HH^2(A_0)_{\chi_i}$ has a nonzero cocycle representative supported exclusively on composable pairs of rays $(r,s)$ satisfying $rs=0$ and $r*s\ne0$, where $*$ is the stem product of Definition~\ref{def:stem-product}, and its value on each such pair $(r,s)$ lies in the one-dimensional socle line $U_{r*s}$. Moreover, this cocycle may be chosen over $\F_2$ and extended by scalars to $k$.
\end{corollary}
\begin{proof}
Fix $i\in I_0$ and work first over $\F_2$ with $A^{(2)}=\F_2P$ and $A^{(4)}=\mathcal OP$. By Proposition~\ref{prop:primitive-parameter-characters}, the residual character $\chi_i\in L$ is part of a $\Z$-basis of $L$ and satisfies $\chi_i(\deg_\JJ)=2s_i>0$, so $\chi_i$ is positive and primitive as required by Setting~\ref{set:fixed-character}. Let $C^*$ be the weight-$\chi_i$ subcomplex of the normalized Hochschild complex of $A^{(4)}$, so that $\overline C^*$ is the weight-$\chi_i$ subcomplex for $A^{(2)}$, and let $U\subseteq C^2$ be the free $\mathcal O$-submodule of cochains that vanish on all pairs $(r,s)$ with $rs\ne0$ or $r*s=0$ and take values in $U_{r*s}$ when $rs=0$ and $r*s\ne0$.

Condition~(1) of Lemma~\ref{lem:supported-bockstein} holds by Lemma~\ref{lem:mod-four-comparison}. For condition~(2), any derivation $\bar\delta\in Z^1(\overline C^*)$ of weight $\chi_i$ acts on shiftable rays by $\bar\delta(r)=\bar d_r r^+$ and vanishes on identities and unshiftable rays, where the coefficients $(\bar d_r)$ satisfy the nonzero-product Leibniz equations~\eqref{eq:partial-leibniz-group} modulo~$2$. By Theorem~\ref{thm:integral-partial-lifting}, $(\bar d_r)$ lifts to a solution $(d_r)$ over $\mathcal O=\Z/4$, so the lifted one-cochain $\delta(r)=d_r r^+\in C^1$ satisfies $(d\delta)(r,s)=0$ whenever $rs\ne0$. On any pair $(r,s)$ with $rs=0$, we have $\delta(rs)=0$, and any outer successor substitution $r^+=\widehat q_xr$ or $s^+=s\widehat q_z$ contributes $\widehat q_x(rs)=0$ or $(rs)\widehat q_z=0$ to $(d\delta)(r,s)$. Thus the only possible nonzero contributions to $(d\delta)(r,s)=d_r r^+s+d_s rs^+$ come from middle actions $r^+=r\widehat q_y$ or $s^+=\widehat q_ys$, both of which evaluate to $r\widehat q_ys$. By Corollary~\ref{cor:active-zero-pairs} and Proposition~\ref{prop:intrinsic-stem-product}, $r\widehat q_ys$ vanishes if $r*s=0$ and equals $r*s\in U_{r*s}$ if $r*s\ne0$. Hence $d\delta\in U$, verifying condition~(2) of Lemma~\ref{lem:supported-bockstein}.

Lemma~\ref{lem:supported-bockstein} therefore produces a nonzero cocycle generator $(d\delta)/2\in U/2U$ of $\HH^2(A^{(2)})_{\chi_i}$, and tensoring with $k$ over $\F_2$ gives the desired generator of $\HH^2(A_0)_{\chi_i}$.
\end{proof}

\section{Associative families and completion of the proof}
\label{sec:binary-families}

Under the inductive hypothesis of Corollary~\ref{cor:zero-supported-coverage}, the cocycles supported on zero ray products span the positive-degree Hochschild cohomology $\HH^2(A_0)_{>0}$. In this section we organize all such cocycles into a single polynomial family of associative multiplications, verify that every specialization with parameters in $\{0,1\}$ has a normed multiplicative basis, and apply Proposition~\ref{prop:family-criterion} to recover the original algebra $A$ and close the induction of Proposition~\ref{prop:table-weight-lifting}.

\subsection{Indicator cocycles on zero products}
\label{subsec:binary-families}

The stem product $r*s$ specifies the only possible nonzero value for a deformation term on a pair of rays $(r,s)$ with $rs=0$. Imposing the Hochschild cocycle condition on all such pairs simultaneously turns out to require only pairwise equalities between their scalar coefficients, partitioning the active zero pairs into equivalence classes whose indicator sums define independent deformation parameters.

We first record this phenomenon as a general observation about two compatible associative multiplication tables on the same basis.

\begin{lemma}
\label{lem:two-associative-tables}
Let $S$ be a finite set with an absorbing element $0$, and let $m$ and $m_*$ be two associative multiplications on $S$ that share a common finite set of orthogonal idempotents summing to the identity of both linearizations. Suppose that $m_*(a,b)=m(a,b)$ whenever $m(a,b)\ne0$. Set
\[
  Z=\{(a,b)\in(S\setminus\{0\})^2:m(a,b)=0\text{ and }m_*(a,b)\ne0\},
\]
and for each $z=(a,b)\in Z$ let $\epsilon_z$ be the bilinear cochain on the free abelian group $\Z[S\setminus\{0\}]$ that sends $(a,b)$ to $m_*(a,b)$ and vanishes on all other basis pairs. Assume that all nested insertions vanish:
\[
  \epsilon_z(\epsilon_w(a,b),c)=0=\epsilon_z(a,\epsilon_w(b,c))
  \qquad\text{for all }z,w\in Z\text{ and }a,b,c\in S\setminus\{0\}.
\]
Then there is an equivalence relation $\sim$ on $Z$ such that, over any commutative ring $R$, the Hochschild two-cocycles of $(R[S\setminus\{0\}],m)$ of the form $\sum_{z\in Z}a_z\epsilon_z$ ($a_z\in R$) are precisely the $R$-linear combinations of the class sums
\[
  \mu_E=\sum_{z\in E}\epsilon_z\qquad(E\in Z/{\sim}).
\]
Moreover, the deformed multiplication $m+\sum_{E\in Z/{\sim}}T_E\mu_E$ is associative over the polynomial ring $R[T_E:E\in Z/{\sim}]$, and at every specialization with $T_E\in\{0,1\}$ the set $S\setminus\{0\}$ is a multiplicative basis.
\end{lemma}
\begin{proof}
Set $\epsilon=\sum_{z\in Z}\epsilon_z$, so that $m_*=m+\epsilon$. Because both $m$ and $m_*$ are associative and all quadratic terms $\epsilon(\epsilon(a,b),c)$ and $\epsilon(a,\epsilon(b,c))$ vanish by hypothesis, the Hochschild coboundary $d\epsilon$ with respect to $m$ vanishes identically over $\Z$.

Now consider a general $R$-linear combination $\mu=\sum_{z\in Z}a_z\epsilon_z$ and evaluate its coboundary on a triple of nonzero basis elements $a,b,c\in S\setminus\{0\}$ (writing $xy$ for $m(x,y)$):
\[
  (d\mu)(a,b,c)=a\mu(b,c)-\mu(ab,c)+\mu(a,bc)-\mu(a,b)c.
\]
If $abc\ne0$, then $ab\ne0$, $bc\ne0$, $(ab)c\ne0$, and $a(bc)\ne0$, so all four terms vanish because $\mu$ is supported on $Z$. When $abc=0$, the surviving terms depend on whether $ab$ and $bc$ are zero:
\[
  \begin{array}{c|c|c}
  ab & bc & (d\mu)(a,b,c) \\\hline
  0 & 0 & a\mu(b,c)-\mu(a,b)c \\
  0 & \ne0 & \mu(a,bc)-\mu(a,b)c \\
  \ne0 & 0 & a\mu(b,c)-\mu(ab,c) \\
  \ne0 & \ne0 & \mu(a,bc)-\mu(ab,c).
  \end{array}
\]
In every row, $(d\mu)(a,b,c)$ is the difference of at most one positive term and at most one negative term, each of which is a single coefficient $a_z$ times a nonzero element of $S$ (since $a\epsilon_{(b,c)}(b,c)$ lies in $S$). Specializing to $a_z=1$ for all $z\in Z$, the identity $d\epsilon=0$ over $\Z$ shows that on each triple $(a,b,c)$ either both terms vanish or both terms are $a_z$ and $a_{z'}$ times the exact same basis element of $S\setminus\{0\}$. Consequently, $(d\mu)(a,b,c)=0$ imposes either no condition or a simple equality $a_z=a_{z'}$. Generating an equivalence relation $\sim$ on $Z$ by these equalities proves that the cocycles of the form $\sum_{z\in Z}a_z\epsilon_z$ have $R$-basis $\{\mu_E:E\in Z/{\sim}\}$.

Finally, the associator of $m+\sum_E T_E\mu_E$ has zero linear part in the variables $T_E$ because each $\mu_E$ is a cocycle, and zero quadratic part because all nested insertions vanish. When every $T_E$ is specialized to $0$ or $1$, the deformed product of $a,b\in S\setminus\{0\}$ equals $m(a,b)$ if $(a,b)\notin Z$, and equals either $0$ or $m_*(a,b)$ if $(a,b)\in Z$; in every case it belongs to $S$.
\end{proof}

We now specialize Lemma~\ref{lem:two-associative-tables} to the ray category $P$ of $A_0=kP$. Let $Z$ be the set of composable pairs $z=(r,s)$ of nonidentity rays of $P$ satisfying $rs=0$ and $r*s\ne0$, and for each $z=(r,s)\in Z$ let $\epsilon_z$ denote the bilinear cochain that sends $(r,s)$ to $r*s$ and vanishes on all other pairs of rays (including identities).

\begin{proposition}
\label{prop:indicator-cocycles}
Let $Z$ and $\epsilon_z$ be as above for $A_0=kP$ in Setting~\ref{set:projective-companion}. Over any commutative coefficient ring $R$:
\begin{enumerate}[label={\upshape(\arabic*)}]
\item Every nontrivial cocycle condition on a linear combination $\sum_{z\in Z}a_z\epsilon_z$ ($a_z\in R$) takes the form $a_z=a_{z'}$.
\item If $\sim$ denotes the equivalence relation on $Z$ generated by these equalities, then the Hochschild two-cocycles of the form $\sum_{z\in Z}a_z\epsilon_z$ have $R$-basis
\begin{equation}
  \mu_E=\sum_{z\in E}\epsilon_z\qquad(E\in Z/{\sim}).
  \label{eq:indicator-cocycles}
\end{equation}
\item For any two cochains $\mu,\nu$ in the span of $\{\epsilon_z:z\in Z\}$ and any rays $a,b,c$,
\begin{equation}
  \mu(\nu(a,b),c)=0,\qquad \mu(a,\nu(b,c))=0.
  \label{eq:mixed-insertions-zero}
\end{equation}
\end{enumerate}
\end{proposition}
\begin{proof}
We first verify~(3). By Corollary~\ref{cor:active-zero-pairs}, any nonzero value $u=r*s$ of $\nu(a,b)$ spans the endpoint socle of its corner bimodule and is therefore annihilated by both endpoint radicals. For any composable ray $c$, if $uc\ne0$, then $(u,c)\notin Z$ and hence $\mu(u,c)=0$; if instead $uc=0$, then $uq_y^nc=0$ for all $n\geq1$ because $uq_y=0$, so $u*c=0$ by Proposition~\ref{prop:intrinsic-stem-product} and again $\mu(u,c)=0$. The identity $\mu(a,\nu(b,c))=0$ follows dually.

Now view the ordinary ray product $m(r,s)=rs$ and the stem product $m_*(r,s)=r*s$ as binary operations on the finite set $S$ of rays of $P$ (including $0$). By Proposition~\ref{prop:intrinsic-stem-product}, $m_*$ is associative and agrees with $m$ whenever $m(r,s)\ne0$, so $m_*=m+\sum_{z\in Z}\epsilon_z$. Since all nested insertions vanish by~(3), Lemma~\ref{lem:two-associative-tables} immediately yields~(1) and~(2).
\end{proof}

We call the cocycles $\mu_E$ of~\eqref{eq:indicator-cocycles} the \emph{indicator cocycles} on $A_0$. Using them to build an affine family of multiplications verifies the associativity and basis hypotheses of Proposition~\ref{prop:family-criterion}, after which Corollary~\ref{cor:zero-supported-coverage} will allow us to select a subfamily spanning $\HH^2(A_0)_{>0}$.

\begin{proposition}
\label{prop:zero-product-family}
Let $A_0=kP$ satisfy Setting~\ref{set:projective-companion}, let $E$ range over the equivalence classes in $Z/{\sim}$ defining the indicator cocycles $\mu_E$ in~\eqref{eq:indicator-cocycles}, and let $(T_E)_{E\in Z/{\sim}}$ be independent polynomial variables. On the free $k[T_E]$-module having the nonzero rays of $P$ as basis, define the multiplication
\begin{equation}
  r\cdot s=
  \begin{cases}
  rs, & \text{if }rs\ne0\text{ in }P,\\
  T_E(r*s), & \text{if }(r,s)\in E\in Z/{\sim},\\
  0, & \text{otherwise},
  \end{cases}
  \label{eq:affine-ray-family}
\end{equation}
where $r*s$ is the stem product of Definition~\ref{def:stem-product}. Then:
\begin{enumerate}[label={\upshape(\arabic*)}]
\item The multiplication~\eqref{eq:affine-ray-family} is associative over $k[T_E]$.
\item Every scalar fiber (obtained by specializing the variables $T_E$ to elements of $k$) is a positive filtered deformation of $A_0$.
\item Every scalar fiber has ray category $P$, with the nonzero rays of $P$ as representatives, and every binary fiber (with $T_E\in\{0,1\}$ for all $E$) has a normed multiplicative basis consisting of the nonzero rays of $P$.
\end{enumerate}
\end{proposition}
\begin{proof}
Assertion~(1) follows directly from Lemma~\ref{lem:two-associative-tables} and Proposition~\ref{prop:indicator-cocycles}. For~(2), the vertex idempotents are unchanged, and for every pair $(r,s)\in E$ with middle vertex $y$, Corollary~\ref{cor:active-zero-pairs} shows that $r*s=rq_ys$, so the indicator term $\epsilon_{(r,s)}$ has weight
\[
  c_E=\ell(r*s)-\ell(r)-\ell(s)=\ell(q_y),
\]
which has positive categorical degree $\deg_\JJ(q_y)>0$ (and is constant across each equivalence class $E$ because the defining equations $a_z=a_{z'}$ arise from homogeneous triples). Thus every scalar fiber is a positive filtered deformation of $A_0$.

For~(3), Corollary~\ref{cor:active-zero-pairs} shows that $(r,s)\in Z$ requires both $r$ and $s$ to be generating rays between distinct vertices. Consequently, all products involving a local ray $q_x^j$ or a proper ray remain unchanged in every scalar fiber. In particular, the local rings $R_x$ and their left and right actions on every corner bimodule coincide with those of $A_0$, so the basis elements represent the exact same $(R_x^\times\times R_y^\times)$-orbits and generate the same endpoint subbimodules $U_r$. Furthermore, for any composable rays $r\colon x\to y$ and $s\colon y\to z$ and any $n\in\rad R_y$, the product $rn$ is a linear combination of proper rays, so $(rn)\cdot s=(rn)s$; hence the middle-radical insertion subspace $M(r,s)=R_xr(\rad R_y)sR_z$ of Definition~\ref{def:ray-product} is identical in $A_0$ and in the fiber algebra. Because each deformation term $T_E(r*s)=T_E(rq_ys)$ lies inside $M(r,s)$, the deformed product $r\cdot s$ belongs to $M(r,s)$ if and only if $rs=0$ in $A_0$ (and when $rs\ne0$, $r\cdot s=rs$). Thus every scalar fiber has ray category $P$.

Finally, at any binary specialization $T_E\in\{0,1\}$, all products of basis elements are either basis elements or zero. Because the endpoint subbimodule chains are unchanged, every corner of a two-sided ideal is an endpoint subbimodule and is therefore spanned by a subset of the basis rays. Intersecting with powers of the radical shows that the basis is normed.
\end{proof}

\begin{example}
\label{ex:zero-product-fibre}
In the nine-dimensional algebra $B$ of Example~\ref{ex:zero-ray-span}, replacing the single zero product $\alpha\beta=0$ by $\alpha\mathbin{\cdot_t}\beta=tu$ (where $u=\alpha q\beta=[\alpha]*[\beta]$) while keeping all other basis products fixed yields the one-parameter family
\[
  B_t=kQ/(q^2,\ \alpha\beta-t\alpha q\beta).
\]
The same nine paths form a basis of $B_t$, and at $t=0$ and $t=1$ they form a normed multiplicative basis; assigning degree~$1$ to $q,\alpha,\beta$ makes $B_t$ a positive filtered deformation of $B$. Notice that the linear change of basis $\alpha'=\alpha-t\alpha q$ satisfies $\alpha'\beta=0$ and $\alpha'q=\alpha q$, giving an isomorphism $B_t\cong B$ for every $t\in k$. Equivalently, the indicator cocycle $\epsilon_{(\alpha,\beta)}$ is the coboundary $d\delta$ of the cochain $\delta(\alpha)=\alpha q$. Thus an individual indicator cocycle $\mu_E$ can be a coboundary; to obtain cocycles that represent nonzero classes in $\HH^2(A_0)_{>0}$, we must combine Proposition~\ref{prop:zero-product-family} with Corollary~\ref{cor:zero-supported-coverage}.
\end{example}

\begin{remark}
\label{rem:classical-mixtures}
In \cite[Sections~9.2--9.7]{BGRS}, Bongartz, Gabriel, Roiter, and Scherback construct associative algebras whose basis multiplications agree with the ray product except on specified local configurations where they equal the stem product, using these "mixtures" to classify representation-finite algebras in characteristic two. From the deformation-theoretic perspective, the indicator cocycles $\mu_E$ in~\eqref{eq:indicator-cocycles} provide canonical cocycle coordinates for these mixtures: if $V=\bigoplus_{E\in Z/{\sim}}k\mu_E$ denotes their span, then under the inductive hypothesis of Corollary~\ref{cor:zero-supported-coverage}, that corollary together with Proposition~\ref{prop:hochschild-mesh-restriction} yields the short exact sequence
\[
  0\longrightarrow V\cap B^2(A_0)\longrightarrow V\longrightarrow\HH^2(A_0)_{>0}\longrightarrow0,
\]
where $B^2(A_0)$ is the space of normalized Hochschild two-coboundaries.
\end{remark}

\subsection{Existence of multiplicative bases in every characteristic}
\label{subsec:main-theorems}

We now assemble the preceding results to prove the inductive step for Proposition~\ref{prop:table-weight-lifting} and deduce our main theorems. We first show that if Proposition~\ref{prop:table-weight-lifting} holds for all weighted multiplication tables of dimension strictly less than $\dim A$, then every basic connected representation-finite algebra $A$ of dimension $\dim A$ in characteristic two satisfies all three conclusions of Proposition~\ref{prop:dimension-induction}.

\begin{lemma}
\label{lem:recovery-with-lifting}
Assume that $\charac k=2$ and let $A$ be a basic connected representation-finite algebra over $k$. If Proposition~\ref{prop:table-weight-lifting} holds for all weighted multiplication tables with fewer than $\dim A$ nonzero elements, then all three conclusions of Proposition~\ref{prop:dimension-induction} hold for $A$.
\end{lemma}
\begin{proof}
Since $\dim A_0=\dim A$, the smaller-table hypothesis applies to the centralizer tables $Q_\chi$ in Theorem~\ref{thm:integral-partial-lifting}. For each $i\in I_0$, Corollary~\ref{cor:zero-supported-coverage} therefore provides a cocycle generator of the one-dimensional weight space $\HH^2(A_0)_{\chi_i}$ supported on active zero pairs $Z$. By Proposition~\ref{prop:indicator-cocycles}(2), this cocycle is a $k$-linear combination of indicator cocycles $\mu_E$ of weight $\chi_i$. Consequently, for each $i\in I_0$ we can choose a single equivalence class $E_i\in Z/{\sim}$ of weight $\chi_i$ whose indicator cocycle $\mu_i:=\mu_{E_i}$ represents a nonzero cohomology class in $\HH^2(A_0)_{\chi_i}$.

Because the characters $\chi_i$ ($i\in I_0$) are pairwise distinct, the classes $E_i$ are distinct. Specializing the family of Proposition~\ref{prop:zero-product-family} by setting $T_{E_i}=t_i$ for $i\in I_0$ and $T_E=0$ for all other classes $E\notin\{E_i:i\in I_0\}$ produces the polynomial family of multiplications
\begin{equation}
  m_t=m+\sum_{i\in I_0}t_i\mu_i
  \label{eq:selected-binary-family}
\end{equation}
over $k[t_i:i\in I_0]$. By Proposition~\ref{prop:zero-product-family}, $m_t$ is associative, every scalar fiber has ray category $P$, and every binary fiber ($t_i\in\{0,1\}$) has the nonzero rays of $P$ as a normed multiplicative basis. Together with our choice of $\mu_i$ as a basis of $\HH^2(A_0)_{>0}=\bigoplus_{i\in I_0}\HH^2(A_0)_{\chi_i}$, all hypotheses of Proposition~\ref{prop:family-criterion} are satisfied. Hence there exists a binary specialization $A_t$ ($t_i\in\{0,1\}$) and a filtered algebra isomorphism $A_t\xrightarrow{\sim}A$ whose associated graded isomorphism $\gr A_t\xrightarrow{\sim}\gr A$ rescales each homogeneous ray line of $A_0$. Transporting the normed multiplicative ray basis of $A_t$ across this isomorphism gives a normed multiplicative ray basis of $A$, proving conclusion~(1) of Proposition~\ref{prop:dimension-induction}.

Furthermore, because $\ray A_t=P=\ray A_0$ by Proposition~\ref{prop:zero-product-family}(3) and the filtered isomorphism $A_t\xrightarrow{\sim}A$ induces a diagonal rescaling on $\gr A_t=A_0=\gr A$, the initial-form bijection of Lemma~\ref{lem:ray-degree} identifies the ray multiplication in $\ray A$ with that in $\ray A_0$, proving conclusion~(2). Finally, $W(\ray A)\cong W(P)$ is free abelian by Proposition~\ref{prop:integral-ray-weights}, proving conclusion~(3).
\end{proof}

\begin{proof}[Proof of Proposition~\ref{prop:table-weight-lifting}]
We proceed by strong induction on the number $n=\dim kT$ of nonzero elements of $T$. The base case $n=0$ is vacuous, and if $T$ is disconnected, it decomposes as a disjoint union of smaller weighted multiplication tables whose linearizations are representation-finite direct factors of $kT$, so the claim follows componentwise.

Assume now that $T$ is connected with $n$ nonzero elements, and set $B=kT$. Because nonidentity elements of $T$ have strictly positive degree $d(\ell(r))>0$, they span a nilpotent ideal of $B$ whose quotient is $\prod_x k\cdot 1_x$; thus $B$ is a basic connected representation-finite algebra of dimension $n$. We claim that the basis elements of $T$ represent the intrinsic rays of $B$ and identify $T\cong\ray B$ as multiplication tables. By Lemma~\ref{lem:endpoint-chains}, each local ring $R_x=B(x,x)$ is a truncated polynomial ring and each corner bimodule $B(x,y)$ is cyclic over $R_x$ or $R_y$. Since the product of two elements of $T$ is either an element of $T$ or zero and nonidentity elements have positive degree, the radical $\rad R_x$ and all endpoint subbimodules of $B(x,y)$ are spanned by subsets of $T$. Choosing elements $q_x\in T$ and $f\in T$ that project to generators of $\rad R_x/(\rad R_x)^2$ and of the top of $B(x,y)$, we see that the powers $q_x^j$ and the products $q_x^if$ (or $fq_y^i$) are nonzero elements of $T$ of strictly increasing degree, and therefore form a basis for the successive endpoint layers of $R_x$ and $B(x,y)$. Thus every nonzero ray of $B$ contains a unique element of $T$. Moreover, whenever two elements $r,s\in T$ have a nonzero product $rs\in T$, its degree $d(\ell(rs))=d(\ell(r))+d(\ell(s))$ is strictly smaller than the degree of any middle-radical insertion $rq_y^ms$ ($m\geq1$); hence $rs\notin M(r,s)$, which means that the ray product in $\ray B$ coincides with the product in $T$. Thus $T\cong\ray B$.

If $\charac k\ne2$, Corollary~\ref{cor:away-from-two-basis} identifies $\ray B$ with $\ray B_0$, and if $\charac k=2$, the strong induction hypothesis and Lemma~\ref{lem:recovery-with-lifting} applied to $B$ likewise identify $\ray B$ with $\ray B_0$. By Proposition~\ref{prop:integral-ray-weights}, the scalar weight group $W(\ray B_0)$ is free abelian. Since $W(T)\cong W(\ray B)\cong W(\ray B_0)$ is free abelian, every homomorphism $W(T)\to\Z/2$ lifts to a homomorphism $W(T)\to\Z$, completing the induction.
\end{proof}

\begin{proof}[Proof of Proposition~\ref{prop:dimension-induction}]
Decomposing $A$ into connected components, we may assume that $A$ is connected. When $\charac k=2$, Proposition~\ref{prop:table-weight-lifting} provides the smaller-table hypothesis of Lemma~\ref{lem:recovery-with-lifting}, which gives~(1), (2), and~(3). When $\charac k\ne2$, Corollary~\ref{cor:away-from-two-basis} gives~(1) and~(2), and Proposition~\ref{prop:integral-ray-weights} gives~(3).
\end{proof}

Together with the matrix-unit extension for nonbasic algebras in Lemma~\ref{lem:nonbasic-basis}, Proposition~\ref{prop:dimension-induction} establishes Theorem~A from the introduction.

\begin{theorem}[{\cite[Sections~1.2 and~1.12]{BGRS}}]
\label{thm:final-basis}
Let $A$ be a finite-dimensional representation-finite algebra over an algebraically closed field $k$.
\begin{enumerate}[label={\upshape(\arabic*)}]
\item The algebra $A$ admits a normed multiplicative basis (Definition~\ref{def:normed-multiplicative-basis}).
\item If $A$ is basic, this basis can be chosen to consist of one representative from each nonzero ray of $A$, with the primitive idempotents representing the identity rays.
\item If $A$ is basic and $\charac k\ne2$, these representatives may be chosen so that they induce an algebra isomorphism $k\ray A\xrightarrow{\sim}A$.
\end{enumerate}
\end{theorem}
\begin{proof}
When $A$ is basic, Proposition~\ref{prop:dimension-induction}(1) gives~(1) and~(2), and Corollary~\ref{cor:away-from-two-basis} gives~(3). When $A$ is nonbasic, Lemma~\ref{lem:nonbasic-basis} extends the normed multiplicative basis from its basic corner to $A$.
\end{proof}

Our next theorem establishes Theorem~B from the introduction: it identifies the standard form of Bretscher and Gabriel~\cite[Section~3]{BretscherGabriel} with the linearization $k\ray A$ of the ray category (as in~\cite[Section~9]{BGRS}) and records the freeness of the scalar weight group.

\begin{theorem}
\label{thm:final-standardness}
Let $A$ be a basic representation-finite algebra over an algebraically closed field $k$, and form $k(\Gamma_A)$ using the mesh coefficients from the almost split sequences of $\modu A$. In every characteristic:
\begin{enumerate}[label={\upshape(\arabic*)}]
\item Taking initial forms identifies $\ray A$ with $\ray A_0$. In particular, $\ray A$ is an associative ray category, $A_0\cong k\ray A$ is representation-finite, and
\[
  \add k(\Gamma_A)\simeq\modu k\ray A.
\]
\item The algebra $A$ is standard if and only if $A\cong k\ray A$.
\item For every finite ray category $Q$ whose linearization $kQ$ is representation-finite, the scalar weight group $W(Q)$ is free abelian and
\[
  k(\Gamma_{kQ})\simeq\ind(\modu kQ).
\]
\end{enumerate}
\end{theorem}
\begin{proof}
Parts~(1) and~(2) follow immediately from Proposition~\ref{prop:dimension-induction}(2), Proposition~\ref{prop:scalar-normalization}, and Corollary~\ref{cor:graded-standard-form}.

For~(3), we first check that $B=kQ$ is basic. By Definition~\ref{def:finite-ray-category}, the identity elements $e_x=1_x$ form a complete set of orthogonal primitive idempotents in $B$. If two distinct vertices $x\ne y$ gave isomorphic projective modules $e_xB\cong e_yB$, there would exist elements $f\in e_xBe_y$ and $g\in e_yBe_x$ satisfying $fg=e_x$ and $gf=e_y$. Conjugation by $f$ and $g$ would identify the local rings $R_x$ and $R_y$ and their radicals, ensuring that $M(f,g)\subseteq\rad R_x$ and $M(g,f)\subseteq\rad R_y$. The ray products in $\ray B\cong Q$ would then satisfy $[f][g]=[e_x]$ and $[g][f]=[e_y]$, making $x$ and $y$ isomorphic in $Q$ and contradicting the skeletality of $Q$ (Definition~\ref{def:finite-ray-category}). Thus $kQ$ is basic. Applying~(1) and~(2) together with Proposition~\ref{prop:dimension-induction}(3) to each connected component of $kQ$ (and noting that $\ray(kQ)\cong Q$ by Definition~\ref{def:finite-ray-category}, so $(kQ)_0\cong kQ$) shows that $W(Q)$ is free abelian and $k(\Gamma_{kQ})\simeq\ind(\modu kQ)$.
\end{proof}

We emphasize that in characteristic two, the existence of a normed multiplicative ray basis in Theorem~\ref{thm:final-basis}(2) does \emph{not} imply that $A$ is isomorphic to $k\ray A$. Outside characteristic two, the products of the chosen ray representatives match composition in $\ray A$. In characteristic two, however, the deformation family~\eqref{eq:selected-binary-family} allows a pair of rays with zero ray product $rs=0$ to have the nonzero product $r\cdot s=r*s=rq_ys$ in $A$. Because $rq_ys$ lies in the middle-radical insertion subspace $M(r,s)$ of Definition~\ref{def:ray-product}, the ray product in $\ray A$ remains zero even though the product of the basis representatives in $A$ is nonzero. Thus the deformation changes the algebra structure from standard to nonstandard while preserving the ray category $\ray A$.

\section{Further consequences for deformations and mesh parameters}
\label{sec:deformation-refinements}

The module-lifting techniques developed in Part~\ref{part:mesh-deformations} yield further structural information about the module categories of positive filtered deformations of $A_0$. In this final section we show that the lifted modules of Proposition~\ref{prop:positive-module-lifting} exhaust all indecomposable modules over any positive deformation of $A_0$, establish a polynomial normal form for filtered mesh categories, and prove that the comparison map $\Psi$ of Proposition~\ref{prop:hochschild-mesh-restriction} is an isomorphism in all positive degrees. None of the results in this section is required for the proofs of Theorems~A and~B.

\subsection{All modules over a positive deformation}

When we apply Proposition~\ref{prop:positive-module-lifting} to the full collection of normalized graded indecomposable $A_0$-modules, the resulting lifted modules exhaust all indecomposable modules over the deformed algebra $B$. Compare the correspondence between modules over a nonstandard algebra and its standard form in \cite[Theorem~1.9 and Remark~1.10]{BautistaSalmeron}.

\begin{corollary}
\label{cor:positive-deformation-completeness}
Let $A_0$ be the projective ray algebra of Construction~\ref{con:projective-algebra}. Every positive filtered deformation $B$ of $A_0$ is representation-finite, the lifted modules $M_i^B$ of Proposition~\ref{prop:positive-module-lifting} form a complete set of pairwise nonisomorphic indecomposable $B$-modules, and $\gr_\JJ\ind(\modu B)\cong\C$.
\end{corollary}
\begin{proof}
Set $G=\bigoplus_i M_i^B$. At the projective objects $X_i=P_x$, we may take $M_i^B=e_xB$ to be the indecomposable projective right $B$-modules. Likewise, equipping the $k$-duals of the indecomposable projective left $B$-modules with their dual filtrations produces filtered injective right $B$-modules whose associated graded $A_0$-modules are degree shifts of the graded indecomposable injectives of $A_0$; after shifting their filtrations to match our degree normalization, the uniqueness statement in Proposition~\ref{prop:positive-module-lifting} identifies them with the corresponding lifted modules $M_i^B$. Consequently, $G$ contains all indecomposable projective and injective right $B$-modules and is therefore a generator-cogenerator for $\modu B$.

Let $\Lambda_B=\End_B(G)^{\op}$ be the opposite endomorphism algebra of $G$ (so that $\Hom_B(G,-)$ takes values in right $\Lambda_B$-modules). By Proposition~\ref{prop:positive-module-lifting}(3), the associated graded algebra $\gr\Lambda_B$ is the Auslander algebra of $A_0$, which has global dimension at most~$2$. We show by a standard filtered lifting argument that $\operatorname{gl.dim}\Lambda_B\leq2$ as well. Because $F^1\Lambda_B$ is a nilpotent ideal with semisimple quotient $\Lambda_B/F^1\Lambda_B=\prod_i k$, the simple right $\Lambda_B$-modules coincide with the degree-zero simple modules of $\gr\Lambda_B$. Lifting the primitive idempotents $e_i$ and equipping each indecomposable projective $e_i\Lambda_B$ with its induced filtration, we can lift any graded projective cover $\bar P\to\gr X$ of the associated graded module of a finitely generated filtered $\Lambda_B$-module $X$ to a filtered morphism $P\to X$ (where $P$ is a direct sum of filtration-shifted copies of $e_i\Lambda_B$) by lifting the images of the homogeneous generators in each idempotent component.

Because the filtration on $X$ is finite, successive approximation in increasing filtration degree shows that $P\to X$ is a strict filtered surjection ($F^dP\twoheadrightarrow F^dX$ for all $d$). Equipping $K=\ker(P\to X)$ with the induced filtration $F^dK=K\cap F^dP$, strictness implies that any element of $F^dP$ whose image lies in $F^{d+1}X$ can be modified by an element of $F^{d+1}P$ to lie in $F^dK$, so the induced sequence of associated graded modules
\[
  0\longrightarrow\gr K\longrightarrow\gr P\longrightarrow\gr X\longrightarrow0
\]
is exact. Applying this filtered projective cover construction twice to any simple $\Lambda_B$-module $S$ produces a filtered second syzygy $K_2$ whose associated graded module $\gr K_2$ is a graded second syzygy over $\gr\Lambda_B$, and hence is a graded projective $\gr\Lambda_B$-module because $\operatorname{gl.dim}\gr\Lambda_B\leq2$. Since $\gr\Lambda_B$ is positively graded, $\gr K_2$ is a direct sum of shifted copies of $e_i\gr\Lambda_B$; lifting its homogeneous generators yields a filtered morphism from a direct sum of shifted copies of $e_i\Lambda_B$ to $K_2$ that is an isomorphism on associated graded modules, and therefore an isomorphism of $\Lambda_B$-modules. Thus $\operatorname{proj.dim}_{\Lambda_B}S\leq2$ for every simple $\Lambda_B$-module $S$, so $\operatorname{gl.dim}\Lambda_B\leq2$.

Finally, by Auslander's theorem, for any finitely generated right $B$-module $X$, choose the beginning of an injective copresentation $0\to X\to I^0\to I^1$ with $I^0,I^1\in\add G$. Applying the left-exact functor $\Hom_B(G,-)$ exhibits $\Hom_B(G,X)$ as the kernel of a morphism between projective $\Lambda_B$-modules, hence as a second syzygy over $\Lambda_B$, which is projective because $\operatorname{gl.dim}\Lambda_B\leq2$. Since $G$ is a generator for $\modu B$, the functor $\Hom_B(G,-)$ restricts to an equivalence between $\add G$ and $\proj\Lambda_B$, forcing $X\in\add G$. Because the modules $M_i^B$ have local endomorphism rings and pairwise nonisomorphic associated graded modules, they form a complete set of pairwise nonisomorphic indecomposable $B$-modules, and Proposition~\ref{prop:positive-module-lifting}(3) gives $\gr_\JJ\ind(\modu B)\cong\C$.
\end{proof}

\subsection{Normal forms for filtered mesh categories}

Extending Proposition~\ref{prop:filtered-normal-form} from the case where the mesh derivative $D_{q-1}$ is surjective to the general case, we can eliminate all higher-order terms in a filtered deformation of a quadratic category except for a chosen set of representatives spanning $\operatorname{coker}D_{q-1}$. In characteristic two, Proposition~\ref{prop:residual-correction-spaces} provides canonical path representatives $(p_i)_{i\in I}$ for the mesh category $\C$.

\begin{proposition}
\label{prop:characteristic-two-normal-form}
Let $C=kQ/(r_v:v\in V)$ be a finite-dimensional graded category (or algebra) over an arbitrary field $k$, where $Q$ is a finite quiver and the $r_v$ are homogeneous quadratic relations, and let $D_{q-1}$ be the linearized substitution map of Proposition~\ref{prop:filtered-normal-form}. Fix a finite family $(p_i)_{i\in I}$ of homogeneous path tuples of degrees $d_i\geq3$ such that, for each $q\geq3$, the classes of the tuples $p_i$ with $d_i=q$ span $\operatorname{coker}D_{q-1}$. Choose $N$ such that $C_d=0$ for all $d\geq N$, and let $\mathfrak m$ denote the arrow ideal of $kQ$. Then every filtered category $E$ equipped with a finite filtration and an identification $\gr E=C$ admits a presentation
\begin{equation}
  E\cong kQ/\Bigl(\mathfrak m^{N+1},\ r_v-\sum_{i\in I}\lambda_i(p_i)_v:v\in V\Bigr)
  \label{eq:filtered-mesh-normal-form}
\end{equation}
for some scalars $\lambda_i\in k$, via a filtered isomorphism that induces the identity on $\gr E=C$.
\end{proposition}
\begin{proof}
We adjust the lifted arrows inductively in increasing degree $q\geq3$. Suppose that arrow lifts have been chosen so that all terms of degree strictly between $2$ and $q$ in $r_v$ belong to the span of the retained path tuples $(p_i)_v$ with $d_i<q$. Let $z_v\in C_q$ be the leading degree-$q$ component of the remaining error $r_v-\sum_{d_i<q}\lambda_i(p_i)_v$. By assumption, the class of $(z_v)_{v\in V}$ in $\operatorname{coker}D_{q-1}$ can be written as $\sum_{d_i=q}\lambda_i[p_i]+D_{q-1}(g)$ for some scalars $\lambda_i\in k$ ($d_i=q$) and a degree-$(q-1)$ arrow tuple $g$. Modifying the arrow lifts in degree~$q-1$ by $-g$ cancels $D_{q-1}(g)$ in degree~$q$ while leaving all terms of degree less than $q$ unchanged (since substituting a degree-$(q-1)$ correction into a quadratic relation shifts degree by $q-2$, and substituting it into a previously retained term of degree $d_i\geq3$ shifts degree to at least $d_i+q-2>q$). Since the filtration of $E$ is finite, this inductive adjustment terminates at $q=N$.

The resulting arrow lifts satisfy the relations of~\eqref{eq:filtered-mesh-normal-form} in $E$ and therefore induce a filtered surjection from the right-hand side of~\eqref{eq:filtered-mesh-normal-form} onto $E$. Passing to associated graded categories gives a composition of graded surjections $C\twoheadrightarrow\gr(\text{RHS})\twoheadrightarrow\gr E=C$ that is the identity on $C$, so both maps are isomorphisms of finite-dimensional categories.
\end{proof}

When $C=\C$ in characteristic two, Proposition~\ref{prop:residual-correction-spaces} shows that we may take $(p_i)_{i\in I}$ to be the single-path residual representatives supported on the residual periodic components of $\C$. Similar higher-degree perturbations of mesh relations define the deformed preprojective algebras of~\cite{BES}, and filtered reparametrizations of arrows appear in Kaplan's comparison of additive and multiplicative preprojective algebras~\cite[Lemma~3.8]{Kaplan}.

\subsection{Dimensions of perturbed mesh categories}

Viewing finite $k$-linear categories via their category algebras as in the introduction, the category algebra of $\C\simeq\ind(\modu A_0)$ has global dimension at most~$2$: projective resolutions of its simple modules of length at most~$2$ are obtained by applying representable functors to the almost split sequences and to the projective-injective sequences of Proposition~\ref{prop:graded-realization}(2). The next lemma shows that perturbing the homogeneous relations of any such category by higher-degree terms preserves both the dimensions of all morphism spaces and the projective resolutions of the simple modules.

\begin{lemma}
\label{lem:flat-mesh-perturbations}
Let $k$ be an arbitrary field, and let $C=kQ/(r_1,\ldots,r_t)$ be a finite-dimensional graded category presented by homogeneous relations $r_l$ of degree at least~$2$. Assume that for each vertex $i$, the sequence of right $C$-modules obtained by factoring each relation $r_l$ starting at $i$ through its initial arrow is a projective resolution of the simple right module $S_i$ of length at most~$2$.

Choose $N$ such that $C_d=0$ for all $d\geq N$, and for each $l=1,\ldots,t$ let $f_l\in kQ$ be a linear combination of paths having the same source and target as $r_l$ and degree strictly greater than $\deg(r_l)$. Set
\[
  B=kQ/(\mathfrak m^{N+1},\ r_1-f_1,\ldots,r_t-f_t).
\]
Then $\gr_{\mathfrak m}B\cong C$, every homogeneous path basis of $C$ lifts to a basis of $B$, and factoring the perturbed relations $r_l-f_l$ through their initial arrows yields projective resolutions of length at most~$2$ for the simple right $B$-modules.
\end{lemma}
\begin{proof}
Taking leading terms gives a graded surjection $C\twoheadrightarrow\gr_{\mathfrak m}B$. Since $C_d=0$ for $d\geq N$, it follows that $\mathfrak m^NB=\mathfrak m^{N+1}B$, and because $\mathfrak m^{N+1}=0$ in $B$, we already have $\mathfrak m^NB=0$. In particular, if any path of length $N+1$ is factored as $ap$ with $a$ its initial arrow and $p$ of length $N$, the suffix $p$ already vanishes in $B$.

With our left-to-right path convention, the indecomposable projective right $B$-module at vertex $i$ is the covariant representable module $P_i=B(i,-)$. Factoring each relation $r_l-f_l$ starting at $i$ through its initial arrow yields an exact sequence of right $B$-modules
\begin{equation}
  0\longrightarrow K_i\longrightarrow
  \bigoplus_{r_l:i\to j}P_j\longrightarrow
  \bigoplus_{a:i\to j}P_j\longrightarrow P_i\longrightarrow S_i
  \longrightarrow0,
  \label{eq:perturbed-simple-resolution}
\end{equation}
where $K_i$ denotes the kernel at the relation term. (Exactness at the arrow term $\bigoplus_{a:i\to j}P_j$ holds because any relation in $B$ starting at $i$ is a sum of expressions $v(r_l-f_l)w$ and paths of length at least $N+1$: when $v$ is nonempty, stripping the initial arrow of $v$ leaves $v'(r_l-f_l)w=0$ in $B$; when $v$ is empty, stripping the initial arrow of $r_l-f_l$ lands in the image of $\bigoplus_{r_l:i\to j}P_j$; and stripping the initial arrow of a path of length $N+1$ leaves a path of length $N$, which is zero in $B$.)

Let $H$ and $H_B$ be the matrices whose $(i,j)$-entries are $\dim C(i,j)$ and $\dim B(i,j)$, let $M$ and $R$ be the adjacency matrices counting arrows $i\to j$ and relations $r_l\colon i\to j$, set $Z=I-M+R$, and let $K$ be the matrix with $(i,j)$-entry $\dim(K_i e_j)$. Taking Euler characteristics in the assumed projective resolutions over $C$ and in~\eqref{eq:perturbed-simple-resolution} over $B$ gives the matrix identities
\[
  ZH=I,\qquad ZH_B=I+K.
\]
Hence $Z$ is invertible with $Z^{-1}=H$, and multiplying the second identity on the left by $H$ gives $H_B=H+HK$. Because $H$ and $K$ have nonnegative entries, we have $H_B\geq H$ entrywise, while the graded surjection $C\twoheadrightarrow\gr_{\mathfrak m}B$ gives $H_B\leq H$ entrywise. Therefore $H_B=H$, $K=0$, and the surjection $C\twoheadrightarrow\gr_{\mathfrak m}B$ is an isomorphism.
\end{proof}

Note that for the category algebra of $C=\C$ (where $C(i,j)$ sits in the $(i,j)$-corner), the indecomposable projective right modules are the covariant functors $\C(X,-)$ rather than the contravariant functors $\C(-,X)$ used in Section~\ref{sec:graded-category}. Consequently, the projective resolution of the simple right $\C$-module at a noninjective object $X$ is induced by the almost split sequence $0\to X\to E_{\tau^{-1}X}\to\tau^{-1}X\to0$ starting at $X$ (the leftmost map $\C(\tau^{-1}X,-)\to\C(E_{\tau^{-1}X},-)$ is injective because $E_{\tau^{-1}X}\to\tau^{-1}X$ is an epimorphism in $\modu A_0$), while at an injective object $X$ the epimorphism of Proposition~\ref{prop:graded-realization}(2) gives a projective resolution of length at most~$1$. Thus Lemma~\ref{lem:flat-mesh-perturbations} applies directly to $C=\C$.

\begin{corollary}
\label{cor:polynomial-mesh-flatness}
If the coefficients of the higher-degree perturbation terms $f_l$ in Lemma~\ref{lem:flat-mesh-perturbations} are replaced by independent polynomial variables over $k$, then any homogeneous path basis of $C$ lifts to a free basis of the perturbed category over the polynomial ring, and remains a basis after arbitrary base change to any commutative $k$-algebra.
\end{corollary}
\begin{proof}
Using the homogeneous relations $r_l$ of $C$, every path in $kQ$ can be rewritten modulo $r_l-f_l$ as a $k$-linear combination of the chosen basis paths plus paths of strictly higher degree; iterating this reduction terminates at degree $N+1$ and shows that the basis paths span the perturbed category over the polynomial ring without dividing by any polynomials. Extending scalars to the rational function field and applying Lemma~\ref{lem:flat-mesh-perturbations} shows that the basis paths are linearly independent over the function field, and hence over the polynomial ring. Since the perturbed category is free with this basis over the polynomial ring, the same holds after any commutative base change.
\end{proof}

\subsection{The full Hochschild comparison}

With Corollary~\ref{cor:polynomial-mesh-flatness} in hand, we can now prove that the injective comparison map $\Psi$ of Proposition~\ref{prop:hochschild-mesh-restriction} is also surjective, and hence an isomorphism in every positive internal degree.

\begin{proposition}
\label{prop:full-hochschild-comparison}
For every internal degree $q>0$, the comparison homomorphism
\[
  \Psi\colon\HH^2(A_0)_q\xrightarrow{\;\sim\;}\operatorname{coker}D_{q+1}(\C)
\]
of Proposition~\ref{prop:hochschild-mesh-restriction} is an isomorphism preserving weights in $L$. In particular, when $\charac k=2$, we have $I_0=I$.
\end{proposition}
\begin{proof}
Injectivity and weight compatibility were established in Proposition~\ref{prop:hochschild-mesh-restriction}. To prove surjectivity, fix a tuple of homogeneous morphisms $z_v\in\C_{q+2}(\tau v,v)$ representing a class $[z]\in\operatorname{coker}D_{q+1}(\C)$, and form the perturbed category $\mathcal C_\varepsilon$ over $k[\varepsilon]/(\varepsilon^2)$ defined by the relations $r_v-\varepsilon z_v$ (and $\mathfrak m^{N+1}=0$). By Corollary~\ref{cor:polynomial-mesh-flatness}, $\mathcal C_\varepsilon$ is a free deformation of $\C$ over $k[\varepsilon]/(\varepsilon^2)$. Its projective corner is therefore a first-order deformation $A_\varepsilon$ of $A_0$, and evaluating each object $X_i$ of $\mathcal C_\varepsilon$ on the projective objects yields a free lift $M_i^\varepsilon=\mathcal C_\varepsilon(P_-,X_i)$ of the normalized $A_0$-module $M_i$. Giving $\varepsilon$ internal degree $-q$ (and weight $-\chi$ when $z$ has weight $\chi$) makes $A_\varepsilon$ and $M_i^\varepsilon$ homogeneous of degree zero, so by the uniqueness in Lemma~\ref{lem:square-zero-selected-lifts} these modules are the unique graded lifts used to define $\Psi$. Moreover, the evaluation map from $\mathcal C_\varepsilon(X_i,X_j)$ to $\Hom_{A_\varepsilon}(M_i^\varepsilon,M_j^\varepsilon)$ is a map of free $k[\varepsilon]/(\varepsilon^2)$-modules of the same rank that reduces modulo $\varepsilon$ to the isomorphism $\C(X_i,X_j)\xrightarrow{\sim}\Hom_{A_0}(M_i,M_j)$, so it is an isomorphism. By construction of $\mathcal C_\varepsilon$, the lifted arrows satisfy $r_v=\varepsilon z_v$, which means that $\Psi$ sends the cohomology class of the corner deformation $A_\varepsilon$ to $[z]$. Thus $\Psi$ is surjective. In characteristic two, each residual line $k[p_i]$ ($i\in I$) of Proposition~\ref{prop:residual-correction-spaces} therefore lies in the image of $\Psi$ in weight $\chi_i$, forcing $\HH^2(A_0)_{\chi_i}\ne0$ for all $i\in I$ and hence $I_0=I$.
\end{proof}

\subsection{Rescaling the mesh parameters}

Using the all-ones mesh normalization of Section~\ref{sec:weights-and-parameters} and the residual paths $(p_i)_{i\in I}$ of Proposition~\ref{prop:residual-correction-spaces}, let $C(\lambda)$ denote the perturbed mesh category of~\eqref{eq:filtered-mesh-normal-form} associated to a parameter vector $\lambda=(\lambda_i)_{i\in I}\in k^I$. By Lemma~\ref{lem:flat-mesh-perturbations}, $C(\lambda)$ has the same dimension as $\C$ for every $\lambda\in k^I$. We next show that up to isomorphism, $C(\lambda)$ depends only on the support $\{i\in I:\lambda_i\ne0\}$ of the parameter vector $\lambda$.

\begin{lemma}
\label{lem:same-support-forms}
Assume that $\charac k=2$, and let $C(\lambda)$ and $C(\lambda')$ be two perturbed mesh categories of the form~\eqref{eq:filtered-mesh-normal-form} constructed from the all-ones mesh relations of $\C$ and the residual path representatives $(p_i)_{i\in I}$ of Proposition~\ref{prop:residual-correction-spaces}. If $\lambda_i=0\iff\lambda'_i=0$ for every $i\in I$, then there is an object-preserving isomorphism $C(\lambda)\cong C(\lambda')$.
\end{lemma}
\begin{proof}
For each $i\in I$, choose a nonzero scalar $a_i\in k^\times$. On the $i$-th residual periodic component of $\Gamma_A$, set the arrow scale $\mu_a=a_i$ for every internal arrow $a$ and the mesh scale $\eta_v=a_i^2$ for every internal mesh $v$; on all other periodic components, set $\mu_a=1$ and $\eta_v=1$. Because all mesh coefficients $c_{v,a}$ equal~$1$, these choices satisfy the compatibility equations $\mu_a\mu_{\sigma(a)}=\eta_v$ of~\eqref{eq:global-mesh-scales} on all periodic components. Since the cycles of the mate permutation $\sigma$ consist precisely of the arrows whose endpoints are periodic (as $\sigma^2$ acts on endpoints as $\tau^{\pm1}$), all remaining arrows belong to open $\sigma$-chains; choosing $\eta_v=1$ on all remaining meshes and picking an arbitrary nonzero scale on the initial arrow of each open $\sigma$-chain determines $\mu_a\in k^\times$ globally so that $\mu_a\mu_{\sigma(a)}=\eta_v$ holds everywhere.

The resulting diagonal rescaling of the arrows of $Q$ multiplies each quadratic mesh relation $r_v$ by $\eta_v$ and multiplies the residual path $p_i$ (which is a path of length $d_i\geq3$ internal to the $i$-th residual component) by $a_i^{d_i}$ while scaling its mesh $r_{v_i}$ by $a_i^2$. Dividing the $v$-th relation by $\eta_v$ therefore transforms the parameter $\lambda_i$ according to
\[
  \lambda_i\longmapsto \lambda_i a_i^{d_i-2}.
\]
Because $k$ is algebraically closed and $d_i-2\geq1$, whenever $\lambda_i$ and $\lambda'_i$ are both nonzero we may choose $a_i\in k^\times$ with $a_i^{d_i-2}=\lambda'_i/\lambda_i$, yielding an object-preserving isomorphism $C(\lambda)\xrightarrow{\sim}C(\lambda')$.
\end{proof}

\subsection{Polynomial comparison of the two parameter systems}

Finally, we relate the algebra deformation parameters $t=(t_i)_{i\in I}$ of our family~\eqref{eq:selected-binary-family} to the mesh perturbation parameters $\lambda=(\lambda_i)_{i\in I}$ of the lifted module category. Because the residual characters $\chi_i$ form part of a $\Z$-basis of the weight lattice $L$, weight homogeneity forces the relationship between the two parameter systems to be purely diagonal and linear.

\begin{lemma}
\label{lem:nonresonant-polynomials}
Let $L$ be an abelian group, and equip the polynomial ring $k[t_1,\ldots,t_s]$ with an $L$-grading by setting $\deg(t_i)=\chi_i\in L$. Assume that the weights $\chi_1,\ldots,\chi_s$ are pairwise distinct and satisfy the nonresonance condition that no $\chi_j$ can be written as $\sum_{i=1}^s n_i\chi_i$ with integers $n_i\geq0$ and $\sum_{i=1}^s n_i\geq2$. Then every homogeneous polynomial $F_j(t)\in k[t_1,\ldots,t_s]$ of weight $\chi_j$ with zero constant term has the linear form $F_j(t)=a_jt_j$ for some scalar $a_j\in k$.
\end{lemma}
\begin{proof}
The nonresonance hypothesis rules out every monomial of degree at least~$2$ from having weight $\chi_j$, and distinctness of the $\chi_i$ leaves $t_j$ as the only degree-one monomial of weight $\chi_j$.
\end{proof}

Notice that any subset of a $\Z$-basis of a free abelian group $L$ automatically satisfies the hypotheses of Lemma~\ref{lem:nonresonant-polynomials}.

\begin{proposition}
\label{prop:polynomial-mesh-comparison}
Let $m_t=m+\sum_{i\in I}t_i\mu_i$ be an associative polynomial family of multiplications on $A_0$ over $K=k[t_i:i\in I]$ satisfying the hypotheses of Proposition~\ref{prop:family-criterion} (such as the family~\eqref{eq:selected-binary-family}, where $I_0=I$ by Proposition~\ref{prop:full-hochschild-comparison}). Then the normalized graded $A_0$-modules $M_j$, their morphisms, and their mesh presentation lift over the polynomial ring $K$ so that the perturbed mesh relations take the normal form~\eqref{eq:filtered-mesh-normal-form} with linear parameter polynomials
\begin{equation}
  F_i(t)=a_it_i\qquad(a_i\in k^\times,\ i\in I).
  \label{eq:diagonal-marking-map}
\end{equation}
Moreover, after specializing $t$ to any point of $k^I$, the specialized modules form a complete set of pairwise nonisomorphic indecomposable modules over the fiber algebra $A_t$.
\end{proposition}
\begin{proof}
Set $K=k[t_i:i\in I]$, and lift a $\Z$-basis of $H_\Z$ to mesh weights via Lemma~\ref{lem:integer-weight-lifting} so that the normalized $A_0$-modules $M_j$ carry simultaneous $L$-gradings compatible with the $L$-grading of $A_0$. Assign each variable $t_i$ internal degree $-2s_i=-\chi_i(\deg_\JJ)$ and weight $-\chi_i\in L$, so that the deformed multiplication $m_t=m+\sum_{i\in I}t_i\mu_i$ on $A_K=A_0\otimes_k K$ is homogeneous of degree~$0$ and weight~$0$. In each internal degree and weight, choose once and for all a $k$-linear right inverse to the bar differentials on their images in the module bar complexes of the $M_j$, and choose a $k$-linear splitting of the mesh derivative $D_{q-1}(\C)$ with respect to the complement spanned by the residual paths $(p_i)_{i\in I}$. Tensoring these splittings with $K$ makes them available over $K$ without introducing denominators.

Keeping the underlying $K$-modules $M_j\otimes_k K$ and their vertex idempotent actions fixed, we lift the module actions of $A_K$ inductively in increasing internal degree $q>0$: when the action has been constructed up to degree $q-1$, its degree-$q$ associativity defect is a two-cocycle of internal degree $q>0$ in the module bar complex tensored with $K$, which is killed in a single step by applying our chosen right inverse (since $\Ext^2_{A_0}(M_j,M_j)_{>0}=0$ by Lemma~\ref{lem:positive-ext-vanishing}). Because the underlying vector spaces $M_j$ are concentrated in a finite range of internal degrees, this inductive correction terminates after finitely many steps, yielding polynomial action matrices over $K$. By the same token, using $\Ext^1_{A_0}(M_j,M_l)_{>0}=0$ and our chosen right inverses, every homogeneous $A_0$-module homomorphism $f\in\Hom_{A_0}(M_j,M_l)_r$ ($r\geq0$) lifts in finitely many degree-raising steps to a polynomial $A_K$-module homomorphism $\widetilde f$, with $\widetilde f=f$ at $t=0$.

Choosing a homogeneous $k$-basis $f_1,\ldots,f_h$ of each $\Hom_{A_0}(M_j,M_l)$, we see by standard degree induction on the leading term that their lifts $\widetilde f_1,\ldots,\widetilde f_h$ form a free $K$-basis of the lifted homomorphism module $\mathcal M(j,l):=\Hom_{A_K}(M_j\otimes_k K,M_l\otimes_k K)$, and that the same remains true after specializing $t$ to any point of $k^I$. In particular, at every scalar specialization, the specialized modules exhaust all indecomposable modules of the fiber algebra by Corollary~\ref{cor:positive-deformation-completeness}.

Now apply the degree-by-degree arrow adjustment of Proposition~\ref{prop:characteristic-two-normal-form} to the lifted arrow morphisms in $\mathcal M$, using our $K$-linear splittings of $D_{q-1}(\C)\otimes_k K$ to eliminate all higher-degree mesh errors except for $K$-multiples $F_i(t)p_i$ of the residual paths $p_i$. Because the procedure is triangular and terminates above the top degree of $\C$, the coefficients $F_i(t)$ are polynomials in $K$ satisfying $F_i(0)=0$, and the resulting presentation is isomorphic to $\mathcal M$ over $K$.

Finally, because every step preserves the $L$-grading, each polynomial $F_j(t)$ is homogeneous of weight $-\chi_j$ (canceling the weight difference $\chi_j$ between $p_j$ and its mesh relation $r_{v_j}$). Reversing the signs of the weights and applying Proposition~\ref{prop:primitive-parameter-characters} and Lemma~\ref{lem:nonresonant-polynomials} shows that $F_j(t)=a_jt_j$ for some scalar $a_j\in k$. Specializing $t_j=\varepsilon$ with $\varepsilon^2=0$ and $t_i=0$ for $i\ne j$ reduces the construction to the first-order lift of Lemma~\ref{lem:square-zero-selected-lifts} for $m+\varepsilon\mu_j$, where the residual mesh error is $\varepsilon a_j[p_j]=\varepsilon\Psi([\mu_j])$. Since $[\mu_j]\ne0$ in $\HH^2(A_0)_{\chi_j}$ and $\Psi$ is injective by Proposition~\ref{prop:hochschild-mesh-restriction}, we have $a_j\ne0$.
\end{proof}


\addtocontents{toc}{\protect\medskip}
\clearpage
\appendix
\section{Periodic meshes and correction vanishing outside characteristic two}
\label{app:periodic-mesh}

In this appendix we prove the two results on periodic mesh categories used in Part~\ref{part:graded-form}: the Dynkin-tree covering description of periodic components (Proposition~\ref{prop:periodic-tree-model}, used in Section~\ref{sec:scalar-normalization}) and the vanishing of the mesh correction cokernels $\operatorname{coker}D_m(\C)$ outside characteristic two (Proposition~\ref{prop:mesh-correction-spaces}, which completes the proof of Theorem~A when $\charac k\ne2$). The characteristic-two and modulo-four calculations in Appendix~\ref{app:residual-mesh} reuse the tree models and path identities established here, but are not needed outside characteristic two.

\subsection{Finite trees and periodic mesh categories}
\label{app:periodic-trees}

We first prove Proposition~\ref{prop:periodic-tree-model}, stated in Subsection~\ref{subsec:mesh-walk-cover}, by lifting sectional paths to construct a repetition quiver cover. This proof also fixes the explicit Dynkin tree models used in our subsequent cokernel calculations. Both the repetition-quiver covering and its construction via sectional paths are due to Riedtmann~\cite[Sections~1.5--1.6 and~2.1]{Riedtmann1980}.

\begin{proof}[Proof of Proposition~\ref{prop:periodic-tree-model}]
We divide the proof into four steps.

\medskip
\noindent\textbf{Step 1: Sectional paths and a finite rooted tree.}
A path in a periodic translation quiver $Q$ is called \emph{sectional} if it contains no consecutive pair of arrows of the form $(\sigma a)a$. Because every term in a mesh relation is of the form $(\sigma a)a$, a sectional path can never appear as a monomial in any contextual mesh relation $u r_v w$. Consequently, the indicator functional of a sectional path annihilates the mesh ideal, so every sectional path represents a nonzero morphism in the mesh category $C(Q,c)$. Since $C(Q,c)$ is finite-dimensional, there is an integer $N$ such that all paths of length at least $N$ vanish in $C(Q,c)$; hence every sectional path in $Q$ has length strictly less than $N$.

Fix a base vertex $x_0$ in the connected periodic quiver $Q$, and let $\Delta$ be the directed graph whose vertices are all sectional paths in $Q$ starting at $x_0$ (including the trivial path at $x_0$), with a directed edge $p\to q$ whenever $q$ is obtained by appending a single arrow to $p$. By definition, $\Delta$ is a finite tree rooted at the trivial path, oriented away from the root. Its repetition quiver $\Z\Delta$ has vertices $(p,j)$ for $p\in\Delta_0$ and $j\in\Z$, translation $\tau(p,j)=(p,j-1)$, and for each directed edge $p\to q$ in $\Delta$ (where $q=pa$ for an arrow $a$ of $Q$) a pair of arrows
\[
  (p,j)\longrightarrow(q,j),\qquad
  (q,j)\longrightarrow(p,j+1).
\]
We define a translation-quiver morphism $\pi\colon\Z\Delta\to Q$ on vertices by
\[
  \pi(p,j)=\tau^{-j}(\operatorname{target}p),
\]
and on the two repetition arrows associated to $q=pa$ by sending $(p,j)\to(q,j)$ to $\tau^{-j}a$ and $(q,j)\to(p,j+1)$ to $\tau^{-j}(\sigma^{-1}a)$.

At the root $p=x_0$, the children $q$ correspond to all outgoing arrows from $x_0$. At a nonroot path $p$ ending with an arrow $a$, the children $q$ correspond to all outgoing arrows from $\operatorname{target}p$ except $\sigma^{-1}a$, while the returning arrow $(p,j)\to(p',j+1)$ to the parent $p'$ of $p$ maps to $\tau^{-j}(\sigma^{-1}a)$. Thus $\pi$ restricts to a bijection on the outgoing arrows at every vertex of $\Z\Delta$, and compatibility with the mate permutation $\sigma$ implies that $\pi$ is also bijective on incoming arrows. In particular, the image of $\pi$ is a union of connected components, hence equals all of $Q$.

Because $\pi$ is a local isomorphism on incoming and outgoing stars, paths in $Q$ lift uniquely to $\Z\Delta$ once a starting vertex in $\Z\Delta$ is fixed. Pulling back the mesh coefficients of $Q$ along $\pi$ equips $\Z\Delta$ with a mesh category whose contextual mesh relations from a fixed starting vertex $(p,j)$ project bijectively onto those of $Q$ starting at $\pi(p,j)$. Taking the direct sum over all endpoints in $\Z\Delta$ therefore identifies the space of paths in $\Z\Delta$ starting at $(p,j)$, modulo lifted mesh relations, with the space of paths in $C(Q,c)$ starting at $\pi(p,j)$. In particular, all paths of length at least $N$ vanish in the lifted mesh category on $\Z\Delta$.

\medskip
\noindent\textbf{Step 2: The tree $\Delta$ is a Dynkin diagram.}
We next rescale the arrows of $\Z\Delta$ so that the lifted mesh relations take the standard $\pm1$ signs of the signed preprojective algebra $\Pi(\Delta)$. Along each tree edge $u\to v$ of $\Delta$, denote the repetition arrows by $a_j\colon(u,j)\to(v,j)$ and $b_j\colon(v,j)\to(u,j+1)$. Prescribing mesh coefficients $+1$ at $(v,j)$ and $-1$ at $(u,j)$ imposes equations of the form $\mu_{a,j}\mu_{b,j-1}=\alpha_{v,j}$ and $\mu_{a,j-1}\mu_{b,j-1}=\beta_{u,j}$ on the arrow scales $\mu$. Taking their ratio determines $\mu_{a,j}$ recursively for all $j\in\Z$ from $\mu_{a,0}=1$, and then determines every $\mu_{b,j}$ (since the index $j$ runs over $\Z$, there is no periodicity obstruction, and distinct edges of $\Delta$ involve disjoint sets of arrows).

Recall that the \emph{signed preprojective algebra} $\Pi(\Delta)$ is the quotient of the path algebra of the doubled quiver $\overline\Delta$ (obtained by adding a reverse arrow $a^*\colon v\to u$ for each tree edge $a\colon u\to v$) by the quadratic relations
\[
  \sum_{a:x\to y}aa^*-\sum_{a:y\to x}a^*a=0\qquad(x\in\Delta_0).
\]
Forgetting the integer coordinate $j$ in $\Z\Delta$ identifies the rescaled mesh relations on $\Z\Delta$ with the defining relations of $\Pi(\Delta)$, and every path in $\overline\Delta$ starting at $u$ lifts uniquely to a path in $\Z\Delta$ starting at $(u,0)$. Consequently, all paths of length at least $N$ vanish in $\Pi(\Delta)$, so $\Pi(\Delta)$ is finite-dimensional.

For each vertex $i\in\Delta_0$, set $P_i=e_i\Pi(\Delta)$ and $n_i=\dim_k P_i$, and let $S_i$ be the simple right $\Pi(\Delta)$-module at $i$. Factoring paths and relations starting at $i$ through their initial arrows produces the exact sequence of right $\Pi(\Delta)$-modules
\begin{equation}
  P_i\longrightarrow\bigoplus_{j\sim i}P_j\longrightarrow P_i\longrightarrow S_i\longrightarrow0,
  \label{eq:preprojective-dimension-sequence}
\end{equation}
where $j\sim i$ ranges over the neighbors of $i$ in $\Delta$. (Exactness at $\bigoplus_{j\sim i}P_j$ holds because any relation in the free path algebra starting at $i$ is a sum of contextual relations $v r_x w$: if $v$ is nonempty, stripping its first arrow leaves a relation in $\Pi(\Delta)$ that vanishes in $P_j$, while if $v$ is empty, stripping the first arrow of $r_i w$ gives an element in the image of the leftmost map $P_i\to\bigoplus_{j\sim i}P_j$.) Letting $K_i$ denote the kernel of the leftmost map in~\eqref{eq:preprojective-dimension-sequence}, taking dimensions yields
\[
  2n_i-\sum_{j\sim i}n_j=1+\dim_k K_i>0\qquad(i\in\Delta_0).
\]
For the symmetric matrix $M=2I-\operatorname{Adj}(\Delta)$ and any real vector $z=(z_i)$, expanding the quadratic form gives
\[
  z^{\mathsf T}Mz=\sum_{\{i,j\}\in\Delta_1}n_in_j\left(\frac{z_i}{n_i}-\frac{z_j}{n_j}\right)^2+\sum_{i\in\Delta_0}\frac{(Mn)_i}{n_i}z_i^2.
\]
Since $n_i\geq1$ and $(Mn)_i=1+\dim_k K_i>0$ for all $i$, the matrix $M$ is positive definite. This immediately rules out any vertex of degree at least~$4$ (which contains a $\widetilde D_4$ star, on which $z=(2;1,1,1,1)$ satisfies $z^{\mathsf T}Mz=0$) as well as two trivalent vertices (which contain a $\widetilde D_m$ subtree, on which $z$ equal to~$2$ along the connecting path and~$1$ at the four leaves satisfies $z^{\mathsf T}Mz=0$). Thus $\Delta$ is either a linear path $A_n$ or a star with three arms of lengths $p\leq q\leq r$. In the star case, the Cartan matrix of an arm of length $m$ has inverse entry $m/(m+1)$ at its inner endpoint, so taking the Schur complement at the central vertex gives
\[
  2-\frac{p}{p+1}-\frac{q}{q+1}-\frac{r}{r+1}>0\quad\Longleftrightarrow\quad
  \frac1{p+1}+\frac1{q+1}+\frac1{r+1}>1,
\]
which forces $p=1$ and either $q=1$ ($D_n$) or $q=2$ with $r\in\{2,3,4\}$ ($E_6,E_7,E_8$).

\medskip
\noindent\textbf{Step 3: The fibers of $\pi$ are deck-group orbits.}
Write $R=\Z\Delta$ and let $o=(x_0,0)$ be the root at coordinate~$0$. For any vertex $z\in R$ with $\pi(z)=\pi(o)$, we can construct a translation-quiver endomorphism $H_z\colon R\to R$ over $\pi$ by setting $H_z(o)=z$, lifting the rooted tree edges of $\Delta\times\{0\}$ uniquely from $z$, and extending to all of $R$ via $\tau$ and the mate permutation $\sigma$. Because $\Delta$ is a tree, this extension is well defined and locally bijective at every vertex, hence surjective onto $R$. Moreover, any translation-quiver morphism $R\to R$ over $\pi$ is uniquely determined by the image of a single vertex via unique lifting of undirected walks. Choosing $z'\in R$ with $H_z(z')=o$, we have $\pi(z')=\pi(o)$, and the compositions $H_z\circ H_{z'}$ and $H_{z'}\circ H_z$ fix $o$ and $z'$ respectively, so $H_z$ is an automorphism of $R$ over $\pi$.

Consequently, the deck group $\Theta$ of translation-quiver automorphisms of $R$ over $\pi$ acts freely on $R$ and transitively on the fiber $\pi^{-1}(\pi(o))$. For any two vertices $u,v\in R$ with $\pi(u)=\pi(v)$, projecting an undirected walk from $o$ to $u$ down to $Q$ and lifting its reverse from $v$ produces a vertex $z\in\pi^{-1}(\pi(o))$ such that $H_z(u)=v$. Thus $\Theta$ acts transitively on every fiber of $\pi$, giving $R/\Theta\cong Q$.

\medskip
\noindent\textbf{Step 4: The deck group $\Theta$ is cyclic.}
Every translation-commuting automorphism of $\Z\Delta$ has the form
\[
  (v,j)\longmapsto(\rho(v),j+n_v),
\]
where $\rho$ is an undirected graph automorphism of $\Delta$ and the integer shifts $n_v\in\Z$ satisfy $n_v=n_u$ along a directed edge $u\to v$ if $\rho$ preserves its orientation, and $n_v=n_u+1$ if $\rho$ reverses its orientation. Thus the automorphism is completely determined by $\rho$ and its shift at a single vertex, and changing the orientation of $\Delta$ simply shifts the repetition coordinates of the vertices of $\Z\Delta$.

When $\Delta$ is of type $D_n$, $E_6$, $E_7$, $E_8$, or $A_{2n-1}$, every graph automorphism of $\Delta$ fixes a canonical root vertex (the trivalent vertex or the central vertex of $A_{2n-1}$). Orienting $\Delta$ away from this root makes every $\rho\in\Aut(\Delta)$ orientation-preserving, so $n_v$ is constant across $v\in\Delta_0$ and the full group of translation-commuting automorphisms of $\Z\Delta$ is $\Z\times\Aut(\Delta)$ (with $\Z$ generated by $\tau$). Since $\Theta$ acts freely on vertices and any element in the kernel of the projection $\Theta\to\Z$ fixes the root $(x_0,0)$, the projection $\Theta\to\Z$ is injective; since $Q=\Z\Delta/\Theta$ is finite, its image is a nontrivial subgroup $s\Z$ ($s\geq1$). Hence $\Theta=\langle\tau^s\rho\rangle$ for some $s\geq1$ and $\rho\in\Aut(\Delta)$, proving the first case of Proposition~\ref{prop:periodic-tree-model}(2).

When $\Delta=A_{2n}$, label the two halves of the tree by $L_i$ and $R_i$ ($1\leq i\leq n$) from the center outward, orient the central edge as $L_1\to R_1$, and orient both halves outward ($L_i\to L_{i+1}$ and $R_i\to R_{i+1}$). The translation-quiver automorphism
\[
  \eta(L_i,j)=(R_i,j),\qquad \eta(R_i,j)=(L_i,j+1)
\]
satisfies $\eta^2=\tau^{-1}$, and the edge-shift equations show that every translation-commuting automorphism of $\Z A_{2n}$ is a power of $\eta$. Hence $\Theta=\langle\eta^r\rangle$ for some $r\geq1$ (where even $r$ gives a pure translation quotient and odd $r$ twists translation by the reflection of $A_{2n}$), completing the proof of Proposition~\ref{prop:periodic-tree-model}.
\end{proof}

For our later calculations on the odd quotient $\Z A_{2n}/\langle\eta^r\rangle$ ($r$ odd), it is convenient to identify $L_{i,j}$ with $(i,2j)$ and $R_{i,j}$ with $(i,2j+1)$. The quotient quiver then has vertices $(i,j)$ with $1\leq i\leq n$ and $j\in\Z/r\Z$, translation $\tau(i,j)=(i,j-2)$, and arrows
\[
  \varepsilon_j\colon(1,j)\to(1,j+1),\qquad
  a_{i,j}\colon(i,j)\to(i+1,j),\qquad
  b_{i,j}\colon(i+1,j)\to(i,j+2).
\]

\subsection{Mesh corrections outside characteristic two}
\label{app:mesh-calculations}

We now prove the vanishing assertion $\operatorname{coker}D_m(\C)=0$ ($m\geq2$) of Proposition~\ref{prop:mesh-correction-spaces} when $\charac k\ne2$. By Proposition~\ref{prop:periodic-localization}, $\operatorname{coker}D_m(\C)$ decomposes as the direct sum of the cokernels of the periodic components of $\Gamma_A$ (and a periodic component without arrows consists of isolated vertices with $\C_d=0$ for $d\geq1$). By Proposition~\ref{prop:periodic-tree-model}, each nontrivial periodic component is modeled on a Dynkin repetition quotient $\Z\Delta/\Theta$. We first normalize the mesh coefficients on these periodic models (diagonal rescalings of arrows and mesh relations conjugate $D_m$ by invertible diagonal matrices, preserving both $\operatorname{coker}D_m$ and the brick hypothesis). We then show that $\operatorname{coker}D_{d-1}$ is a quotient of a \emph{twisted trace space} of the signed preprojective algebra $\Pi(\Delta)$, and compute these twisted trace spaces for each Dynkin type. Note that our normalization, the twisted-trace surjection, and the nonvanishing results in Lemmas~\ref{lem:d4-four-return} and~\ref{lem:e7-sixteen-return} hold in arbitrary characteristic; we will point out explicitly where the invertibility of~$2$ is used.

\subsubsection{Scalar coefficients in the periodic models}

Consider first a rooted periodic model $Q(\Delta,s,\rho)=\Z\Delta/\langle\tau^s\rho\rangle$, where $\Delta$ is oriented away from a root fixed by $\rho$. The orbit graph $\Delta/\langle\rho\rangle$ is again a finite rooted tree. For each $\rho$-orbit $v$ of vertices in $\Delta$, let $m_v=s|v|$ denote the period of the corresponding $\tau$-cycle in $Q(\Delta,s,\rho)$. Along any parent-child edge orbit $u\to v$, the $\rho$-stabilizer of a child vertex fixes its unique parent in the tree, so $|u|$ divides $|v|$ and hence $m_u\mid m_v$.

\begin{lemma}
\label{lem:periodic-scalar-gauges}
By rescaling arrows and mesh relations in $C(Q,c)$, any system of nonzero mesh coefficients $c$ can be normalized as follows:
\begin{enumerate}[label={\upshape(\arabic*)}]
\item On $Q=\Z\Delta/\langle\tau^s\rho\rangle$, where $\Delta$ is a finite tree oriented away from a root fixed by $\rho$, the coefficients $c$ can be rescaled to the standard $\pm1$ coefficients of the signed preprojective algebra $\Pi(\Delta)$.
\item On the odd quotient $\Z A_{2n}/\langle\eta^r\rangle$ ($r$ odd), written in the coordinates at the end of Appendix~\ref{app:periodic-trees}, the coefficients $c$ can be rescaled so that every length-two mesh path has coefficient $+1$.
\end{enumerate}
\end{lemma}
\begin{proof}
For~(1), label the arrows in $Q$ lying over a parent-child edge orbit $u\to v$ by $j\in\Z/m_v\Z$ as
\[
  a_j\colon(u,j\bmod m_u)\to(v,j),\qquad
  b_j\colon(v,j)\to(u,j+1\bmod m_u).
\]
Let $\mu_{a,j},\mu_{b,j}\in k^\times$ be arrow scales and let $\eta_{v,j}\in k^\times$ scale the mesh relation at $(v,j)$. Transforming the given mesh coefficients into the target $\pm1$ coefficients requires solving equations of the form
\[
  \mu_{a,j}\mu_{b,j-1}=t_{v,j}\eta_{v,j},\qquad
  \mu_{a,j-1}\mu_{b,j-1}=t_{u,j}\eta_{u,j\bmod m_u},
\]
for prescribed nonzero scalars $t_{v,j},t_{u,j}\in k^\times$ (when $|\text{orbit }v|>|\text{orbit }u|$, several child indices $j$ share the same parent mesh $\eta_{u,j\bmod m_u}$). Dividing the two equations gives a first-order recurrence for $\mu_{a,j}$ around the cycle $j\in\Z/m_v\Z$, which has a solution if and only if the product of the ratios around the cycle equals~$1$:
\[
  \prod_{j\in\Z/m_v\Z}\eta_{v,j}=c_{uv}\left(\prod_{j\in\Z/m_u\Z}\eta_{u,j}\right)^{m_v/m_u},
\]
where $c_{uv}=\prod_{j}(t_{u,j}/t_{v,j})\in k^\times$. Because the orbit graph $\Delta/\langle\rho\rangle$ is a rooted tree, we may set $\eta_{\text{root},j}=1$ at the root and inductively choose $\eta_{v,j}\in k^\times$ down the tree so that the cycle-product condition holds on every edge orbit, after which the recurrences determine all arrow scales $\mu_{a,j},\mu_{b,j}\in k^\times$.

For~(2), on $\Z A_{2n}/\langle\eta^r\rangle$ ($r$ odd), we set the scales of the central arrows $\varepsilon_j$ equal to~$1$ and choose the level-$1$ mesh scales $\eta_{1,j}$ so that each $\varepsilon_{j-2}\varepsilon_{j-1}$ has coefficient $+1$. Along each arm edge between level $i$ and level $i+1$, the ratio recurrence shifts the index $j\in\Z/r\Z$ by~$2$; because $r$ is odd, shifting by~$2$ generates the entire cycle $\Z/r\Z$, so the compatibility condition simply prescribes the product $\prod_{j\in\Z/r\Z}\eta_{i+1,j}$ in terms of $\prod_{j\in\Z/r\Z}\eta_{i,j}$ and can be solved inductively for $i=1,\ldots,n-1$.
\end{proof}

\subsubsection{Rotation and twisted traces}

Let $\Pi=\Pi(\Delta)$ be the signed preprojective algebra of a rooted tree $\Delta$. Given an orientation-preserving automorphism $\theta\in\Aut(\Delta)$, we define the \emph{$\theta$-twisted trace space} of $\Pi$ to be the graded vector space quotient
\begin{equation}
  T_\theta(\Pi)=\left(\bigoplus_{x\in\Delta_0}\Pi(x,\theta x)\right)
  \Big/\Span\bigl\{uv-v\theta(u):u\in\Pi(x,y),\ v\in\Pi(y,\theta x)\bigr\}.
  \label{eq:twisted-trace-space}
\end{equation}
In $T_\theta(\Pi)$, the relation $[uv]=[v\theta(u)]$ allows one to cut a path $uv$ from $x$ to $\theta x$ and move its prefix $u$ to the end after applying $\theta$; in particular, $[w]=[\theta(w)]$ for every $w\in\Pi(x,\theta x)$. When $\theta=\id$, $T_{\id}(\Pi)=\Pi/[\Pi,\Pi]$ is the ordinary cyclic trace space of $\Pi$.

\begin{lemma}
\label{lem:twisted-trace-surjection}
Let $Q=Q(\Delta,s,\rho)=\Z\Delta/\langle\tau^s\rho\rangle$ be a rooted periodic mesh model equipped with the signed preprojective mesh coefficients of $\Pi=\Pi(\Delta)$, and let $D_{d-1}$ be the mesh derivative~\eqref{eq:mesh-derivative} in path degree $d\geq3$.
\begin{enumerate}[label={\upshape(\arabic*)}]
\item The cokernel $\operatorname{coker}D_{d-1}$ vanishes unless $d$ is even and $s$ divides $d/2-1$.
\item When $d$ is even and $s\mid(d/2-1)$, set $l=(1-d/2)/s$ and $\theta=\rho^l$. Then there is a canonical surjection
\[
  T_\theta(\Pi)_d\twoheadrightarrow\operatorname{coker}D_{d-1}.
\]
\end{enumerate}
\end{lemma}
\begin{proof}
A path of length $d$ in $Q$ from $\tau(x,j)=(x,j-1)$ to $(x,j)$ lifts to a path in $\Z\Delta$ from $(x,j-1)$ to $(\rho^l x,j-sl)$ for some integer $l\in\Z$. Because $x$ and $\rho^l x$ lie at the same distance from the root of $\Delta$, the lifted path takes equal numbers of steps away from the root (which preserve the second coordinate) and toward the root (which increment the second coordinate by~$1$). Hence $d$ must be even and $(j-1)+d/2=j-sl$, which gives $d/2-1=-sl$; this proves~(1) and shows that the lifted path projects in $\Pi(\Delta)$ to a path of length $d$ from $x$ to $\theta x=\rho^l x$.

In the signed mesh presentation on $Q$, modifying an arrow $a$ by a parallel path $w$ of length $d-1$ adds $c[wa]+c'[(\tau a)w]$ to $\im D_{d-1}$, where $c$ and $c'$ are the two mesh coefficients on the tree edge underlying $a$. Because $c$ and $c'$ have opposite signs $\{+1,-1\}$, this identifies $[wa]=[(\tau a)w]$ in $\operatorname{coker}D_{d-1}$. Iterating this single-arrow rotation $d$ times shows that $[w]=[\tau w]$, so the class in $\operatorname{coker}D_{d-1}$ of a lifted path depends only on its underlying path in $\Pi(\Delta)$ from $x$ to $\theta x$, independently of the repetition coordinate $j$ (and invariant under $\rho$ by the deck relation). Moreover, cutting a path $uv$ from $(x,j-1)$ to $(\theta x,j-sl)$ and rotating the prefix $u$ to the end via $[uv]=[v(\tau^{-1}u)]$ in $Q$ corresponds upstairs in $\Pi(\Delta)$ to $[uv]=[v\theta(u)]$. Since contextual mesh relations in $\Pi(\Delta)$ lift to contextual mesh relations in $Q$, the projection $T_\theta(\Pi)_d\to\operatorname{coker}D_{d-1}$ is well defined and surjective.
\end{proof}

To compute $T_\theta(\Pi)$ for star-shaped trees $\Delta$, we first express closed paths at the central vertex $o$ in terms of the length-two return loops along its arms, following~\cite[Sections~11.3 and~13.1]{Schedler} and incorporating the arm-permuting twist $\theta$.

\begin{lemma}
\label{lem:star-return-reduction}
Let $\Delta$ be a star tree rooted at its central vertex $o$, with arms of lengths $L_1,\ldots,L_t$, and for each $i=1,\ldots,t$ let $x_i\in\Pi(o,o)_2$ denote the length-two closed path that steps out along the $i$-th arm and immediately returns to $o$. Then the corner algebra $\Pi(o,o)$ is generated by $x_1,\ldots,x_t$ and is a quotient of
\begin{equation}
  k\langle x_1,\ldots,x_t\rangle\big/\Bigl(\sum_{i=1}^t x_i,\ x_1^{L_1+1},\ldots,x_t^{L_t+1}\Bigr).
  \label{eq:star-corner-presentation}
\end{equation}
\end{lemma}
\begin{proof}
Label the vertices along the $i$-th arm as $v_0=o,v_1,\ldots,v_{L_i}$. At any internal vertex $v_{j-1}$ ($j\geq2$), the mesh relation replaces the outward backtrack $v_{j-1}\to v_j\to v_{j-1}$ (up to sign) by the inward backtrack $v_{j-1}\to v_{j-2}\to v_{j-1}$, strictly decreasing the sum of the distances from $o$ along the path. Iterating this replacement reduces every closed path at $o$ to a polynomial in $x_1,\ldots,x_t$. The mesh relation at $o$ gives $\sum_{i=1}^t x_i=0$, while pushing $x_i^{L_i+1}$ out along the $i$-th arm reaches the leaf relation at $v_{L_i}$ and vanishes.
\end{proof}

\begin{lemma}
\label{lem:arm-trace-reduction}
Let $\Delta$ be a star tree with center $o$, and let $\theta\in\Aut(\Delta)$ be an automorphism fixing $o$ and permuting the arms. In every positive degree, $T_\theta(\Pi)_+$ is spanned by the classes of positive-degree words in the arm loops $x_1,\ldots,x_t$ at $o$. Moreover, if the $i$-th arm is fixed pointwise by $\theta$, then
\begin{equation}
  [x_i^q]=0\quad\text{in }T_\theta(\Pi)\qquad\text{for all }q\geq1.
  \label{eq:arm-power-trace-zero}
\end{equation}
\end{lemma}
\begin{proof}
If a path $w$ from $v$ to $\theta v$ passes through the center $o$, writing $w=uv$ with $u\colon v\to o$ and $v\colon o\to\theta v$ gives $[uv]=[v\theta(u)]$ in $T_\theta(\Pi)$, which is a closed path at $\theta o=o$ and hence a linear combination of words in $x_1,\ldots,x_t$ by Lemma~\ref{lem:star-return-reduction}.

If instead $w$ never visits $o$, then $w$ is supported entirely on a single arm of $\Delta$. Because $w$ goes from $v$ to $\theta v$, that arm must be fixed pointwise by $\theta$, so $w$ is a closed path and twisted rotation on that arm is ordinary cyclic rotation. Rotating $w$ so that it begins at a vertex $v_j$ of minimum depth along $w$, we see that its two incident steps at $v_j$ form an inward-pointing turn $v_{j+1}\to v_j\to v_{j+1}$; applying the mesh relation at $v_{j+1}$ replaces this turn by $v_{j+1}\to v_{j+2}\to v_{j+1}$ (or by~$0$ at the leaf), strictly increasing the sum of the vertex depths along the cyclic word. Since the depth sum of a word of fixed length is bounded above, repeating this step eventually pushes the word off the leaf and shows that $[w]=0$. Applying the exact same outward-pushing argument to the loop power $x_i^q$ on any $\theta$-fixed arm yields~\eqref{eq:arm-power-trace-zero}.
\end{proof}

\subsubsection{Types A and D}

Assume now that $\charac k\ne2$. We apply Lemma~\ref{lem:twisted-trace-surjection} and Lemma~\ref{lem:arm-trace-reduction} to rule out all nonzero classes in $\operatorname{coker}D_{d-1}$ ($d\geq3$) for the rooted models of types $A$ and $D$:
\begin{itemize}
\item \emph{Type $A$ with trivial twist ($\theta=\id$).} This covers all pure translation quotients of type $A$, including $\Z A_{2n}/\langle\eta^r\rangle$ for even $r$ (since $\eta^r=\tau^{-r/2}$). Rooting the path $A_n$ at one of its endpoints makes it a star with a single arm, so $\Pi(o,o)$ is generated by $x_1$ subject to $x_1=0$ by Lemma~\ref{lem:star-return-reduction}; hence $T_{\id}(\Pi)_+=0$.
\item \emph{Type $A_{2n-1}$ with reflection ($\theta=\rho\ne\id$).} Rooting $A_{2n-1}$ at its central vertex gives two arms of length $n-1$ with loops $x_1=x$ and $x_2=-x$ exchanged by $\theta(x)=-x$. For every $q\geq1$, twisted rotation of one factor $x$ gives $[x^q]=[x^{q-1}\theta(x)]=-[x^q]$, so $2[x^q]=0$ and hence $[x^q]=0$ because $\charac k\ne2$.
\item \emph{Type $D_n$ with $\theta=\id$ or the short-arm transposition.} Let $x$ and $y$ be the loops along the two short arms of length~$1$, and let $z=-x-y$ be the loop along the long arm of length $n-3$. Since $x^2=y^2=0$, the corner $\Pi(o,o)$ is spanned by alternating words in $x$ and $y$.
  \begin{itemize}
  \item When $\theta=\id$, any alternating word of odd length $2q+1\geq3$ has the same letters at its two ends and therefore rotates to a word containing $x^2=0$ or $y^2=0$, while $[x]=[y]=0$ by~\eqref{eq:arm-power-trace-zero}. In even word length $2q$, cyclic rotation identifies $[(xy)^q]=[(yx)^q]$, while~\eqref{eq:arm-power-trace-zero} applied to the fixed long arm gives $0=[z^{2q}]=[(x+y)^{2q}]=[(xy)^q]+[(yx)^q]=2[(xy)^q]$. Since $\charac k\ne2$, we obtain $[(xy)^q]=0$.
  \item When $\theta$ interchanges $x$ and $y$, any alternating word of even length $2q$ twisted-rotates to a word containing $x^2=0$ or $y^2=0$ (for instance, $[(xy)^q]=[(yx)^{q-1}y\theta(x)]=[(yx)^{q-1}y^2]=0$). In odd word length $2q+1$, the two alternating words $w=x(yx)^q$ and $w'=y(xy)^q$ satisfy $[w]=[\theta(w)]=[w']$, while the fixed long arm gives $0=[(-z)^{2q+1}]=[(x+y)^{2q+1}]=[w]+[w']=2[w]$, so again $[w]=0$.
  \end{itemize}
\item \emph{Type $D_4$ with a three-cycle $\theta$.} When $\Delta=D_4$, all three arms have length~$1$, so $x^2=y^2=(x+y)^2=0$, which gives $yx=-xy$ and kills all words of word degree~$3$ (path degree~$6$). In word degree~$2$ (path degree~$4$), if $\theta(x)=y$, then $[xy]=[y\theta(x)]=[y^2]=0$, and if $\theta^2(y)=x$, then $[yx]=[x\theta^2(y)]=[x^2]=0$. Thus $T_\theta(\Pi)_d=0$ for all $d\geq3$.
\end{itemize}

\subsubsection{The odd quotients of even A}

Consider next the odd quotient $Q=\Z A_{2n}/\langle\eta^r\rangle$ ($r$ odd), normalized by Lemma~\ref{lem:periodic-scalar-gauges}(2) so that every mesh path of length~$2$ has coefficient $+1$ in the coordinates $(i,j)$ ($1\leq i\leq n$, $j\in\Z/r\Z$) of Appendix~\ref{app:periodic-trees}. Any path from $\tau(i,j)=(i,j-2)$ to $(i,j)$ takes equal numbers of outward steps $a$ (which change $j$ by~$0$) and inward steps $b$ (which change $j$ by $+2$); since each $\varepsilon$ step changes $j$ by $+1$, the total change in the second coordinate $j\in\Z/r\Z$ equals the path length $d$. Thus a path of length $d$ from $\tau(i,j)$ to $(i,j)$ can exist only when $d\equiv2\pmod r$.

Because all mesh coefficients equal $+1$, the column of $D_{d-1}$ corresponding to any arrow $u$ imposes the relation $[wu]+[(\tau u)w]=0$ in $\operatorname{coker}D_{d-1}$, so rotating a path by one arrow negates its class in $\operatorname{coker}D_{d-1}$ while preserving its cyclic sequence of levels $i\in\{1,\ldots,n\}$. Given any path of length $d$ from $\tau v$ to $v$ that reaches a maximum level $i+1>1$, the step into level $i+1$ and the step out of it must be $a_i b_i$ (as there are no level-preserving arrows above level~$1$). Rotating the path if necessary so that $a_i b_i$ does not straddle the cut, we may apply the mesh relation at level $i$ to replace $a_i b_i$ by $-b_{i-1}a_{i-1}$ when $i>1$ (or by $-\varepsilon\varepsilon$ when $i=1$), strictly decreasing the sum of the levels along the path. Repeating this reduction expresses every class in $\operatorname{coker}D_{d-1}$ as a linear combination of the classes $w_j=[\varepsilon_{j-2}\varepsilon_{j-1}\cdots\varepsilon_{j+d-3}]$ of pure level-$1$ paths of length $d$ ending at $(1,j)$ ($j\in\Z/r\Z$).

Since $d\equiv2\pmod r$, the path of length $d-1$ formed by $d-1$ consecutive $\varepsilon$-arrows starting at $(1,j-1)$ has the same source and target as $\varepsilon_{j-1}\colon(1,j-1)\to(1,j)$. Evaluating $D_{d-1}$ on this arrow correction in the two root meshes containing $\varepsilon_{j-2}\varepsilon_{j-1}$ and $\varepsilon_{j-1}\varepsilon_j$ yields
\[
  w_{j-1}+w_j=0\qquad(j\in\Z/r\Z).
\]
Because $r$ is odd, going all the way around the $r$-cycle gives $w_j=(-1)^r w_j=-w_j$, so $2w_j=0$ and hence $w_j=0$ whenever $\charac k\ne2$ (and for $n=1$ the root mesh $\varepsilon_{j-2}\varepsilon_{j-1}=0$ already kills all paths of length~$d\geq2$).

\subsubsection{Type $E_6$ and the small-period exclusion}

For $\Delta=E_6$, the three arms have lengths $1,2,2$. Let $x$ be the loop along the short arm and $y,z$ the loops along the two arms of length~$2$, so that $x+y+z=0$ and $x^2=y^3=z^3=0$. Assuming $\charac k\ne2$, set $v=y-z=x+2y$, so that $2y=v-x$ and $2z=-v-x$. Adding and subtracting $8y^3=0$ and $8z^3=0$ transforms the presentation of $\Pi(o,o)$ into
\begin{equation}
  x^2=0,\qquad v^3=-xvx,\qquad v^2x+vxv+xv^2=0.
  \label{eq:e6-corner-relations}
\end{equation}
From~\eqref{eq:e6-corner-relations} one checks directly that $\Pi(o,o)$ is spanned in each word degree $q=d/2\geq1$ by the following words:
\[
\begin{array}{c|l}
\text{Word degree }q & \text{Spanning words} \\\hline
1 & x,\ v \\
2 & xv,\ vx,\ v^2 \\
3 & xvx,\ xv^2,\ vxv \\
4 & \psi_4=xvxv,\ \psi_4'=vxv^2 \\
5 & \psi_5=xvxv^2 \\
\geq6 & 0.
\end{array}
\]
Indeed, multiplying~\eqref{eq:e6-corner-relations} by $x$ gives $xv^2x=-\psi_4$, $vxvx=\psi_4$, and $xv^3=v^3x=0$; hence $\psi_4x=x(vxvx)=0$, $\psi_4v=\psi_5$, $\psi_4'x=-v\psi_4=-\psi_5$, $\psi_4'v=vxv^3=0$, $\psi_5v=xvxv^3=0$, and $\psi_5x=xvx(v^2x)=-xvx(vxv+xv^2)=-\psi_4xv=0$.

When $\theta=\rho$ is the nontrivial reflection of $E_6$, $\theta$ fixes $x$ and negates $v$. In $T_\rho(\Pi)$, every word containing an odd number of $v$'s satisfies $[w]=[\theta(w)]=-[w]$ and therefore vanishes since $\charac k\ne2$. Among words with an even number of $v$'s, $[x]=0$ by~\eqref{eq:arm-power-trace-zero}, $[v^2]=[v\theta(v)]=-[v^2]=0$, the last relation of~\eqref{eq:e6-corner-relations} gives $0=[v^2x]+[vxv]+[xv^2]=[xv^2]-[xv^2]+[xv^2]=[xv^2]$ in word degree~$3$, and $[\psi_4]=[(xv)\theta(xv)]=-[\psi_4]=0$ in word degree~$4$. Thus $T_\rho(\Pi)_+=0$.

When $\theta=\id$ (ordinary trace), all pure arm powers vanish by~\eqref{eq:arm-power-trace-zero}. Expanding $0=4[y^2]=[(v-x)^2]$ and $0=4[z^2]=[(v+x)^2]$ gives $[xv]=[vx]=[v^2]=0$ in word degree~$2$. In word degree~$3$, $[xvx]=[x^2v]=0$ and $[vxv]=[xv^2]$, while the last relation of~\eqref{eq:e6-corner-relations} gives $3[xv^2]=0$. In word degree~$4$, $[\psi_4]=-[xv^2x]=-[x^2v^2]=0$ and $[\psi_4']=[xv^3]=0$, and in word degree~$5$, $[\psi_5]=-[\psi_4'x]=-[x\psi_4']=-[\psi_5]=0$. Consequently, the only possible nonzero class in $T_{\id}(\Pi)_+$ occurs in word degree~$3$ (path degree $d=6$) when $\charac k=3$, where it is spanned by $[xv^2]$.

To show that even this degree-$6$ class in characteristic~$3$ cannot occur in a periodic component containing a brick, we use the following characteristic-free lemma, which will also be needed for $E_7$ and $E_8$.

\begin{lemma}
\label{lem:d4-four-return}
Over every field $k$, every vertex of $\Pi(D_4)$ admits a nonzero closed path of path degree~$4$. Consequently, the pure translation mesh quotients $\Z D_4/\langle\tau^s\rangle$ of periods $s=1$ and $s=2$ contain no bricks.
\end{lemma}
\begin{proof}
Orient the three arms of $D_4$ outward from the center $0$ to the leaves $1,2,3$ with arrows $a_i\colon0\to i$ and $b_i\colon i\to0$, and set $x_i=a_ib_i$. The defining relations of $\Pi(D_4)$ are $b_ia_i=0$ ($i=1,2,3$) and $x_1+x_2+x_3=0$. At the center~$0$, the paths of length~$4$ are the six products $x_ix_j$ ($i\ne j$), subject only to the relations $x_i(x_1+x_2+x_3)=0$ and $(x_1+x_2+x_3)x_j=0$; in the order $(12,13,21,23,31,32)$, their classes in $\Pi(D_4)_4$ are $(z,-z,-z,z,z,-z)$ with $z=x_1x_2\ne0$. Likewise, at each leaf $i$, the two length-$4$ closed paths $b_ia_jb_ja_i$ ($j\ne i$) are subject only to the single relation $\sum_{j\ne i}b_ia_jb_ja_i=0$, so each is nonzero in $\Pi(D_4)_4$.

In $\Z D_4$, each of these length-$4$ closed paths increases the repetition coordinate by~$2$, and therefore projects in $\Z D_4/\langle\tau^s\rangle$ for $s\in\{1,2\}$ to a nonzero radical endomorphism at every vertex, ruling out the existence of any brick.
\end{proof}

Returning to $E_6$ in path degree $d=6$, Lemma~\ref{lem:twisted-trace-surjection}(1) requires $s\mid(6/2-1)=2$, so $s\in\{1,2\}$. Deleting the outer vertices of the two long arms of $E_6$ at every repetition coordinate leaves a translation-invariant subquiver isomorphic to a $D_4$ quotient of period $s$, whose mesh category is the object quotient of the $E_6$ mesh category by~\eqref{eq:mesh-object-quotient}. For a pure translation quotient ($s=1$ or $2$), Lemma~\ref{lem:d4-four-return} shows that this $D_4$ quotient has no bricks; for the reflected quotient $\Z E_6/\langle\tau^s\rho\rangle$, $s=2$ gives an odd twist $l=-1$ where $T_\rho(\Pi)_6=0$, while $s=1$ gives a $D_4$ quotient with $(\tau\rho)^2=\tau^2$, on which the degree-$4$ closed paths of Lemma~\ref{lem:d4-four-return} again close up to form nonzero radical endomorphisms at every vertex. Thus no periodic $E_6$ component containing a brick has a nonzero correction class.

\subsubsection{The remaining exceptional traces}

We treat $\Delta=E_7$ and $\Delta=E_8$ simultaneously. Let $x,y,z$ denote the length-two loops at the center $o$ along the long, medium, and short arms, respectively (note that $x$ now denotes the long-arm loop, whereas in $E_6$ it denoted the short-arm loop). By Lemma~\ref{lem:star-return-reduction}, their relations in $\Pi(o,o)$ are
\[
  x+y+z=0,\qquad x^m=y^3=z^2=0,\qquad
  \text{where }m=\begin{cases}4 & \text{for }E_7,\\ 5 & \text{for }E_8.\end{cases}
\]
Eliminating $z=-x-y$, we see from Lemmas~\ref{lem:star-return-reduction} and~\ref{lem:arm-trace-reduction} that the positive-degree part of the ordinary trace space $T_{\id}(\Pi)_+$ is a quotient of the positive-degree part of $T_m$, where for any commutative coefficient ring $R$ we define the graded $R$-algebra $R_m$ and its $R$-module trace quotient $T_m$ by
\begin{equation}
  \begin{split}
  R_m &= R\langle x,y\rangle/\bigl(x^m,y^3,(x+y)^2\bigr),\\
  T_m &= R_m\big/\bigl([R_m,R_m]+\Span_R\{x^q,y^q,(x+y)^q:q\geq1\}\bigr),
  \end{split}
  \label{eq:exceptional-trace-spaces}
\end{equation}
with word degree in $x,y$ equal to half the path degree in $\Pi(\Delta)$. We first compute a spanning set and torsion relations for $(T_5)_+$ over $\Z$; this integral calculation reproduces the torsion groups of~\cite[Theorem~13.1.1]{Schedler} and will be used both here (for $\charac k\ne2$) and in Appendix~\ref{app:residual-mesh} (for $\charac k=2$ and $\mathcal O=\Z/4$).

\begin{lemma}
\label{lem:e78-traces}
Let $R_m$ and $T_m$ be defined as in~\eqref{eq:exceptional-trace-spaces} over an arbitrary commutative ring, and write $[w]$ for the class of a word $w$ in $T_m$.
\begin{enumerate}[label={\upshape(\arabic*)}]
\item Over any commutative ring, set $a=yx^3$, $b=yx^4$, and $t_j=[yx^j]$. The positive-degree part $(T_5)_+$ is spanned by the elements
\[
  t_1,\ t_2,\ t_3,\ t_4,\qquad [a^2],\qquad [ab^r]\quad(r\geq1),
\]
which satisfy the torsion relations
\begin{equation}
  \begin{gathered}
  2t_1=2t_3=0,\qquad 3t_2=0,\qquad 5t_4=0,\\
  2[a^2]=0,\qquad 3[ab]=0,\qquad 2[ab^r]=0\quad(r\geq2).
  \end{gathered}
  \label{eq:integral-exceptional-trace-relations}
\end{equation}
Moreover, for every nonempty word $W$ in the blocks $a$ and $b$,
\begin{equation}
  [a^2W]=0,\qquad [abW]=-[baW],\qquad [b^q]=0\quad(q\geq2),
  \label{eq:integral-exceptional-block-relations}
\end{equation}
and $T_4$ is the quotient of $T_5$ obtained by setting $x^4=0$.
\item If $2$ is invertible in the coefficient ring, then $(T_5)_+$ is spanned by
\[
  [yx^2],\qquad [yx^4],\qquad [(yx^3)(yx^4)],
\]
which are annihilated by $3$, $5$, and $3$, respectively, and $(T_4)_+$ is spanned by $[yx^2]$, annihilated by $3$.
\end{enumerate}
\end{lemma}
\begin{proof}
\textit{(1)} In $R_5$ over $\Z$, expanding $(x+y)^2=0$ and $y^3=0$ yields the four identities
\begin{equation}
\begin{aligned}
  y^2 &= -yx-xy-x^2,\\
  yxy &= xyx+x^3,\\
  yx^2y &= -yx^3-xyx^2-x^2yx-x^3y-2x^4,\\
  0 = F &:= yx^3yx+yx^4y+xyx^3y-x^2yx^3-x^3yx^2.
\end{aligned}
  \label{eq:e8-word-identities}
\end{equation}
(Indeed, substituting the first identity into $y^2\cdot y=0$ gives the second; expanding $y^2(xy)=(y^2x)y$ gives the third; and equating $(yxy)x^2y=xyx^3y+x^5y$ with $yx(yx^2y)=-yx^5+x^2yx^3+x^3yx^2+x^6-yx^3yx-yx^4y$ using $x^5=0$ gives the fourth.)

Any cyclic word in $x$ and $y$ can be described by the lengths of the $x$-gaps between consecutive occurrences of $y$. Whenever a gap has length $0$, $1$, or $2$, applying the first, second, or third identity in~\eqref{eq:e8-word-identities} rewrites the word in terms of words with strictly fewer occurrences of $y$. Because $x^5=0$, iterating this reduction expresses every class in $(T_5)_+$ as a linear combination of words with at most one $y$ (namely $t_j=[yx^j]$ for $1\leq j\leq4$, as $[x^q]=0$) and cyclic words in the two blocks $a=yx^3$ and $b=yx^4$.

Set $c=yx^2$. Note first that any word containing $c=yx^2$ flanked on both sides by $a$ or $b$ (so that the adjacent $x$-gaps have lengths $p,q\geq3$) has trace zero: rotating it to the form $yx^2yx^qUx^p$ and expanding $yx^2y$ via the third identity in~\eqref{eq:e8-word-identities} produces five terms, three of which contain $x^5=0$ and the other two of which begin with $x^2$ or $x^3$ and therefore contain $x^{p+2}=0$ or $x^{p+3}=0$ after cyclic rotation.

Using the relation $F=0$ from~\eqref{eq:e8-word-identities}, we now establish~\eqref{eq:integral-exceptional-block-relations}:
\begin{itemize}
\item For any nonempty word $W$ in $\{a,b\}$, taking the trace of $Fx^2W=0$ gives
\[
  [a^2W]+[bcW]+[xacW]-[x^3bW]=0.
\]
Here $[bcW]=0$ and $[xacW]=0$ by the flanked-$c$ observation (or because rotating the initial $x$ to the end of $W$ produces $x^5=0$), and rotating $x^3$ to the end of $[x^3bW]$ likewise produces $x^6=0$ or $x^7=0$; hence $[a^2W]=0$. When $W$ is empty, $0=[Fx^2]=2[a^2]+[bc]$, and multiplying the third identity of~\eqref{eq:e8-word-identities} by $x^4$ gives $[bc]=[yx^2yx^4]=0$, so $2[a^2]=0$.
\item Taking the trace of $Fx^3W=0$ gives
\[
  [abW]+[baW]+[xaaW]=0.
\]
When $W$ is nonempty, rotating the initial $x$ of $[xaaW]$ to the end either creates $x^5=0$ (if $W$ ends in $b$) or gives $[a^2W'b]=0$ (if $W=W'a$); thus $[abW]=-[baW]$. When $W$ is empty, all three terms equal $[ab]$, giving $3[ab]=0$.
\item Taking the trace of $Fx^4b^r=0$ ($r\geq0$) gives $[b^{r+2}]+[xab^{r+1}]=0$; rotating the initial $x$ of the second term to the end creates $x^5=0$, so $[b^q]=0$ for all $q\geq2$.
\end{itemize}
From~\eqref{eq:integral-exceptional-block-relations}, in any word in $\{a,b\}$ of block length at least~$3$ containing at least two $a$'s, anticommuting $a$ and $b$ brings two $a$'s adjacent and kills the trace via $[a^2W]=0$; in a word with a single $a$ and $r\geq2$ copies of $b$, anticommuting $a$ all the way through $b^r$ and cyclically rotating back gives $[ab^r]=(-1)^r[b^ra]=(-1)^r[ab^r]$ and also $[ab\cdot b^{r-1}]=-[bab^{r-1}]=-[ab^r]$, so $2[ab^r]=0$.

Finally, for $t_j=[yx^j]$ ($1\leq j\leq4$), taking the trace of the first identity in~\eqref{eq:e8-word-identities} gives $2t_1=0$; comparing the trace of the second identity with that of $y^2x$ gives $3t_2=0$; comparing the trace of the third identity with that of $y^2x^2$ gives $2t_3=0$; and taking the trace of the third identity multiplied by $x$ gives $t_4=-4t_4$, or $5t_4=0$. This proves~(1), and~(2) follows immediately on inverting~$2$.
\end{proof}

To exclude the remaining periods for $E_7$ and $E_8$ in characteristics~$3$ and~$5$, we establish a lower bound on the dimensions of path spaces in $\Pi(\Delta)$. For any finite graph $\Delta$ with adjacency matrix $M$, define the Chebyshev-type integer matrices $U_d$ ($d\geq0$) by the recurrence
\[
  U_0=I,\qquad U_1=M,\qquad U_{d+1}=MU_d-U_{d-1}.
\]
If $U_0,U_1,\ldots,U_N$ have nonnegative entries, then over every field $k$ we have the entrywise lower bound
\begin{equation}
  \label{eq:preprojective-lower-bound}
  \bigl(\dim_k e_i\Pi(\Delta)_N e_j\bigr)_{i,j}\geq U_N.
\end{equation}
Indeed, let $H(t)=\sum_{d\geq0}(\dim_k e_i\Pi(\Delta)_d e_j)t^d$ be the matrix Hilbert series of $\Pi(\Delta)$. Recording the graded dimensions of the kernels $K_i$ of the leftmost maps in the graded exact sequences~\eqref{eq:preprojective-dimension-sequence} in a power series $K(t)$ with nonnegative matrix coefficients gives
\[
  (I-tM+t^2I)H(t)=I+K(t)\quad\Longrightarrow\quad
  H(t)=\Bigl(\sum_{d\geq0}U_dt^d\Bigr)\bigl(I+K(t)\bigr),
\]
whose degree-$N$ coefficient is $U_N+\sum_{r=1}^N U_{N-r}K_r\geq U_N$ whenever $U_0,\ldots,U_N\geq0$.

Using~\eqref{eq:preprojective-lower-bound}, we verify that every vertex of $\Pi(E_7)$ has a nonzero closed path in the top socle degree $d=16$ (compare~\cite[Section~6.2]{Kaplan}).

\begin{lemma}
\label{lem:e7-sixteen-return}
Over every field $k$, every vertex of $\Pi(E_7)$ admits a nonzero closed path of path degree~$16$.
\end{lemma}
\begin{proof}
Ordering the vertices of $E_7$ according to its bipartition (with the three vertices at odd distance from the center followed by the four vertices at even distance), the off-diagonal block $B$ of the adjacency matrix $M$ and its products $S=BB^{\mathsf T}$ and $T=B^{\mathsf T}B$ are
\[
  B=\begin{pmatrix}1&1&1&0\\0&1&0&0\\0&0&1&1\end{pmatrix},\quad
  S=BB^{\mathsf T}=\begin{pmatrix}3&1&1\\1&1&0\\1&0&2\end{pmatrix},\quad
  T=B^{\mathsf T}B=\begin{pmatrix}1&1&1&0\\1&2&1&0\\1&1&2&1\\0&0&1&1\end{pmatrix}.
\]
Setting $f(X)=X^3-6X^2+9X-3I$ and $w=(-1,0,1,-1)^{\mathsf T}$, direct matrix multiplication gives $f(S)=0$, $f(T)=-ww^{\mathsf T}$, and $Tf(T)=0$. Writing $U_{2j}=\operatorname{diag}(F_j(S),F_j(T))$, the recurrence $U_{d+1}=MU_d-U_{d-1}$ yields
\[
  F_0=I,\quad F_1(X)=X-I,\quad F_2(X)=X^2-3X+I,
\]
\[
  \begin{array}{ll}
  F_3(S)=F_2(S)+I, & F_3(T)=F_2(T)+I-ww^{\mathsf T},\\
  F_4(S)=2(S-I), & F_4(T)=2(T-I)+ww^{\mathsf T}.
  \end{array}
\]
Explicit evaluation shows that all entries of $F_1,F_2,F_3,F_4$ on $S$ and $T$ are nonnegative:
\[
  F_2(S)=\begin{pmatrix}3&1&2\\1&0&1\\2&1&0\end{pmatrix},\qquad
  F_2(T)=\begin{pmatrix}1&1&1&1\\1&1&2&1\\1&2&2&0\\1&1&0&0\end{pmatrix},
\]
\[
  F_3(T)=\begin{pmatrix}1&1&2&0\\1&2&2&1\\2&2&2&1\\0&1&1&0\end{pmatrix},\qquad
  F_4(T)=\begin{pmatrix}1&2&1&1\\2&2&2&0\\1&2&3&1\\1&0&1&1\end{pmatrix}.
\]
Likewise, the upper-right blocks of the odd matrices $U_1,U_3,U_5,U_7$ are $B$, $(S-2I)B$, $(S^2-4S+3I)B$, and $(S-I)B$, which are nonnegative because
\[
  (S-2I)B=\begin{pmatrix}1&2&2&1\\1&0&1&0\\1&1&1&0\end{pmatrix},\qquad
  S^2-4S+3I=\begin{pmatrix}2&0&1\\0&1&1\\1&1&0\end{pmatrix}.
\]
Thus $U_0,\ldots,U_8$ are entrywise nonnegative. Since $U_9=MU_8-U_7=U_7$, the reversibility of the recurrence $U_{d-1}=MU_d-U_{d+1}$ implies by induction that $U_{8+j}=U_{8-j}$ for $0\leq j\leq8$. In particular, $U_0,\ldots,U_{16}$ are all entrywise nonnegative and $U_{16}=U_0=I$, so the claim follows from~\eqref{eq:preprojective-lower-bound}.
\end{proof}

We can now complete the proof of Proposition~\ref{prop:mesh-correction-spaces} for $E_7$ and $E_8$ outside characteristic two. By Lemma~\ref{lem:e78-traces}(2), the only possible nonzero classes in $(T_4)_+$ and $(T_5)_+$ when $2$ is invertible are $[yx^2]$ (word degree~$3$, path degree $d=6$, in characteristic~$3$), $[yx^4]$ (word degree~$5$, path degree $d=10$, for $E_8$ in characteristic~$5$), and $[(yx^3)(yx^4)]$ (word degree~$9$, path degree $d=18$, for $E_8$ in characteristic~$3$). By Lemma~\ref{lem:twisted-trace-surjection}(1), path degrees $d=6,10,18$ require the period $s$ to divide $d/2-1=2,4,8$, respectively.

However, periods $s=1$ and $s=2$ are ruled out by deleting the outer vertices to pass to $\Z D_4/\langle\tau^s\rangle$ and applying Lemma~\ref{lem:d4-four-return}. For periods $s=4$ and $s=8$, the degree-$16$ closed paths at every vertex of $\Pi(E_7)$ from Lemma~\ref{lem:e7-sixteen-return} advance the repetition coordinate by~$8$, and therefore project to nonzero radical endomorphisms at every vertex of $\Z E_7/\langle\tau^s\rangle$ (which is also an object quotient of $\Z E_8/\langle\tau^s\rangle$ obtained by deleting the terminal vertex of the long arm). Thus no periodic component of type $E_7$ or $E_8$ containing a brick can have period dividing~$8$. Since $E_7$ and $E_8$ have no nontrivial graph automorphisms, this completes the proof of Proposition~\ref{prop:mesh-correction-spaces}.

\clearpage
\section{Residual and integral mesh correction spaces}
\label{app:residual-mesh}

In this appendix we complete the periodic trace calculations required for the characteristic-two and integral branches of the main argument. Appendix~\ref{app:characteristic-two-calculations} establishes Proposition~\ref{prop:residual-correction-spaces} by bounding the positive-degree cokernels $\operatorname{coker}D_m(\C)$ ($m\geq2$) over a field of characteristic two, and Appendix~\ref{app:integral-corrections} lifts these calculations to $\mathcal O=\Z/4\Z$ to prove Lemma~\ref{lem:integral-mesh-cokernels} (used in Section~\ref{sec:integral-deformations}). Throughout, we retain the periodic tree models and trace reductions of Appendix~\ref{app:periodic-mesh}.

\subsection{Bounds for the characteristic-two correction spaces}
\label{app:characteristic-two-calculations}

Throughout this subsection we assume $\charac k=2$ and prove Proposition~\ref{prop:residual-correction-spaces}. We show that each periodic component of $\Gamma_A$ contributes at most one correction line across all degrees $m\geq2$, spanned by the class of a single row path. Defining $I$ to be the set of periodic components for which this cokernel line is actually nonzero then yields Proposition~\ref{prop:residual-correction-spaces} directly, without requiring a classification of which periodic components occur in $\Gamma_A$ or a separate nonvanishing proof for each model.

As in Appendix~\ref{app:periodic-mesh}, all vertex deletions below retain $\tau$-invariant subsets of vertices; by~\eqref{eq:mesh-object-quotient}, the resulting object quotients are the mesh categories of the retained subquivers, allowing us to restrict nonzero paths from smaller models to rule out small periods under the brick hypothesis.

\subsubsection{Trees of type A}
\label{app:char2-line-models}

In the preprojective category $\Pi(A_M)$, number the vertices $1,\ldots,M$ from left to right and denote the rightward and leftward steps by $s_+$ and $s_-$. The interior mesh relations commute $s_+s_-$ with $s_-s_+$, while the endpoint backtracks at~$1$ and~$M$ vanish. Sorting the steps of any path starting at vertex $i$ shows that the paths
\begin{equation}
  s_+^a s_-^b\qquad(0\leq a\leq M-i,\ 0\leq b\leq i-1)
  \label{eq:a-rectangle}
\end{equation}
span the paths from $i$, because any step pushing $a$ past $M-i$ or $b$ past $i-1$ can be commuted to an endpoint backtrack and vanishes. To see that the paths~\eqref{eq:a-rectangle} are linearly independent, consider the representation of $\Pi(A_M)$ with basis vectors $v_{a,b}$ ($0\leq a\leq M-i$, $0\leq b\leq i-1$) placed at vertex $i+a-b$, on which $s_+$ and $s_-$ increment $a$ and $b$ respectively (and act by zero at the boundary $a=M-i$ or $b=i-1$). Since these increments commute and satisfy the endpoint relations, the paths~\eqref{eq:a-rectangle} map $v_{0,0}$ to the distinct basis vectors $v_{a,b}$. In particular, all paths of degree at least $M$ in $\Pi(A_M)$ vanish, and for each vertex $i$ there is a unique nonzero path of top degree $M-1$, which terminates at the reflected vertex $M+1-i$.

We now examine the three periodic models of type~$A$:
\begin{itemize}
\item \emph{Pure translation quotients $\Z A_M/\langle\tau^s\rangle$.} The positive ordinary trace of $\Pi(A_M)$ vanishes even in characteristic two: any closed path at $i$ sorts to $(s_+s_-)^q$, and cyclic rotation combined with the interior commutation relation transports $(s_+s_-)^q$ step by step to an endpoint, where it vanishes (without dividing by $q$). Lemma~\ref{lem:twisted-trace-surjection}(1) therefore gives $\operatorname{coker}D_{d-1}=0$ for all $d\geq3$ and every period $s\geq1$.

\item \emph{Reflected odd trees $\Z A_{2m+1}/\langle\tau^s\rho\rangle$.} Here the brick hypothesis forces $s>m$: if $s\leq m$, retaining the central vertex and the first $s$ vertices of each arm yields the quotient $\Z A_{2s+1}/\langle\tau^s\rho\rangle$, in which the top-degree paths of length $2s$ in $\Pi(A_{2s+1})$ connect each vertex to its $\rho$-reflection while advancing the repetition coordinate by $s$, and hence project to nonzero radical endomorphisms at every vertex. With $s>m$ established, any row path from $\tau v$ to $v$ in $\Z A_{2m+1}/\langle\tau^s\rho\rangle$ has even path degree $d=2q\leq2m$ satisfying $s\mid(q-1)$ by Lemma~\ref{lem:twisted-trace-surjection}. Since $0\leq q-1<m<s$ for $q\geq1$, this divisibility forces $q=1$ ($d=2$), so no correction of path degree $d\geq3$ exists.

\item \emph{Folded even trees $\Z A_{2n}/\langle\eta^r\rangle$ ($r$ odd).} In the coordinates $(i,j)\in\{1,\ldots,n\}\times(\Z/r\Z)$ of Appendix~\ref{app:periodic-trees}, the arrows $\varepsilon_j\colon(1,j)\to(1,j+1)$, $a_{i,j}\colon(i,j)\to(i+1,j)$, and $b_{i,j}\colon(i+1,j)\to(i,j+2)$ have the property that any path starting and ending at the same level $i$ contains an equal number of $a$-steps and $b$-steps, so its total shift in the second coordinate $j$ equals its path degree $d$. If $r\leq2n-1$, retaining the first $(r+1)/2$ levels yields $\Z A_{r+1}/\langle\eta^r\rangle$; the top-degree paths of length $r$ in $\Pi(A_{r+1})$ connect each vertex to its reflection at the same level $i$ and shift $j$ by $r\equiv0\pmod r$, giving nonzero radical endomorphisms at every vertex and contradicting the brick hypothesis. Thus $r>2n-1$. Since any row path from $\tau(i,j)=(i,j-2)$ to $(i,j)$ starts and ends at level $i$, its path degree $d$ satisfies $d\equiv2\pmod r$; as $d\leq2n-1<r$, we must have $d=2$, ruling out any correction of degree $d\geq3$.
\end{itemize}

\subsubsection{The two short arms of type D}

For $\Delta=D_n$, set $N=n-2$. By Lemma~\ref{lem:star-return-reduction}, the central corner $e_{x_0}\Pi(D_n)e_{x_0}$ is a quotient of
\[
  R_N=k\langle x,y\rangle/\bigl(x^2,y^2,(x+y)^N\bigr),
\]
where $x$ and $y$ represent the length-two closed paths along the two short arms. In word degree $j<N$, $R_N$ is spanned by the two alternating words $w_j$ and $w'_j$ of length $j$ starting with $x$ and $y$, respectively; in word degree $N$, the relation $(x+y)^N=0$ becomes $w_N+w'_N=0$. Multiplying $w_N+w'_N=0$ on the left or right by $x$ and $y$ kills both alternating words of degree $N+1$, so $(R_N)_j=0$ for all $j>N$.

By Lemma~\ref{lem:arm-trace-reduction}, the positive twisted trace of $\Pi(D_n)$ is a quotient of the corresponding twisted trace of $R_N$:
\begin{itemize}
\item For the \emph{ordinary trace} ($\theta=\id$), the degree-one classes $[x]$ and $[y]$ vanish by~\eqref{eq:arm-power-trace-zero}. Any alternating word of odd length $j\geq3$ begins and ends with the same letter, so cyclic rotation produces $x^2$ or $y^2$ and the class vanishes. For even $j$, cyclic rotation identifies $[w_j]=[w'_j]$. Hence the positive ordinary trace is at most one-dimensional in each path degree $d=2j$ with $j$ even and $2\leq j\leq N$, and vanishes in all other positive degrees.
\item For the \emph{short-arm interchange} ($\theta=\iota$ swapping $x$ and $y$), twisted cyclic rotation $[uv]=[v\iota(u)]$ turns an alternating word of even length $j$ into a word containing $x^2$ or $y^2$, while for odd $j$ it identifies $[w_j]=[w'_j]$. Hence the positive $\iota$-twisted trace is at most one-dimensional in each path degree $d=2j$ with $j$ odd and $1\leq j\leq N$, and vanishes in all other positive degrees.
\end{itemize}
In either case, Lemma~\ref{lem:twisted-trace-surjection} shows that a nonzero cokernel in path degree $d=2j\geq3$ can occur only when
\begin{equation}
  2\leq j\leq N\qquad\text{and}\qquad s\mid(j-1).
  \label{eq:d-possible-correction}
\end{equation}
(These degree and parity restrictions match Riedtmann's characteristic-two analysis of type~$D$ in~\cite[Section~2, p.~487, and Sections~3.3, 3.9]{Riedtmann1984}.)

For $D_4$ ($N=2$), we also check a $3$-cycle $\rho$ permuting the three arms $x,y,x+y$. Here $R_2$ is spanned by $1,x,y,xy=yx$ with $(R_2)_j=0$ for $j\geq3$. In word degree~$1$, twisted cyclic rotation under $\rho$ identifies $[x]=[y]=[x+y]$, whence $[x]=0$; in word degree~$2$, $[xy]=[y\rho(x)]=[y^2]=0$ (and for $\rho^2$, $[yx]=[x\rho^2(y)]=[x^2]=0$). The only possible ordinary trace class has path degree $d=4$ ($j=2$), for which $s\mid(j-1)$ forces $s=1$; since $3\nmid s$, the induced twist $\theta=\rho^{\pm1}$ is nontrivial, so no correction class arises from a $3$-cycle on $D_4$.

To show that~\eqref{eq:d-possible-correction} admits at most one value of $j$ when $\theta\in\{\id,\iota\}$, we bound $N$ in terms of the period $s$ using the top-degree socle paths of $\Pi(D_n)$ (cf.\ \cite[Section~6.2]{Kaplan}).

\begin{lemma}
\label{lem:d-top-paths}
Let $\Pi(D_n)$ be the preprojective algebra of type $D_n$ over an arbitrary field, set $N=n-2$, and let $\nu\in\Aut(D_n)$ be the identity if $N$ is even and the short-arm interchange if $N$ is odd. For every vertex $x$ of $D_n$, there is a nonzero path of degree $2N$ in $\Pi(D_n)$ from $x$ to $\nu x$.
\end{lemma}
\begin{proof}
We apply the integer recurrence criterion~\eqref{eq:preprojective-lower-bound}. First consider the line $A_M$ on vertices $1,\ldots,M$ with adjacency matrix $\mathcal A$, and define $V_{-1}=0$, $V_0=I$, and $V_{d+1}=\mathcal AV_d-V_{d-1}$. On the infinite integer line $\Z$, the function
\[
  c_d(r)=\begin{cases}
    1,&\text{if }|r|\leq d\text{ and }r\equiv d\pmod2,\\
    0,&\text{otherwise}
  \end{cases}
  \qquad(c_{-1}=0)
\]
satisfies $c_{d+1}(r)=c_d(r-1)+c_d(r+1)-c_{d-1}(r)$. By the method of images, the reflected sum
\[
  (V_d)_{ij}=\sum_{h\in\Z}\bigl(c_d(i-j+2h(M+1))-c_d(i+j+2h(M+1))\bigr)\qquad(1\leq i,j\leq M)
\]
satisfies the same recurrence and initial conditions, and vanishes at the artificial boundaries $i=0$ (since $c_d(r)=c_d(-r)$) and $i=M+1$ (by shifting $h\mapsto h-1$). For $0\leq d\leq M-1$, only the $h=0$ term can contribute to the first sum and only the $h=0,-1$ terms can contribute to the second, subtracting the ranges $d\geq i+j$ and $d\geq2M+2-i-j$. Thus, for $0\leq d\leq M-1$, $(V_d)_{ij}=1$ if and only if
\[
  |i-j|\leq d\leq\min(i+j-2,2M-i-j)\qquad\text{and}\qquad d\equiv i-j\pmod2,
\]
and $(V_d)_{ij}=0$ otherwise. In particular, $V_0,\ldots,V_{M-1}$ are entrywise nonnegative, and $V_{M-1}$ is the reflection matrix $i\mapsto M+1-i$.

Now set $M=2N+1$ and shift the vertex labels of $A_{2N+1}$ to $-N,\ldots,N$. On the reflection-invariant subspace spanned by $f_0=e_0$ and $f_r=e_r+e_{-r}$ ($1\leq r\leq N$), the adjacency matrix $\mathcal A$ sends $f_0\mapsto f_1$, $f_1\mapsto 2f_0+f_2$, and $f_r\mapsto f_{r-1}+f_{r+1}$ for $r\geq2$ (omitting $f_{N+1}$). On $D_n$, label the two short-arm leaves by $a,b$ and the remaining vertices from the center outward by $1,\ldots,N$. On the subspace spanned by $g_0=e_a+e_b$ and $g_r=e_r$ ($1\leq r\leq N$), the adjacency matrix of $D_n$ sends $g_0\mapsto 2g_1$, $g_1\mapsto g_0+g_2$, and $g_r\mapsto g_{r-1}+g_{r+1}$ for $r\geq2$. Since these two $(N+1)\times(N+1)$ matrices are transposes of one another, the recurrence polynomials $U_d$ for $D_n$ and $V_d$ for $A_{2N+1}$ satisfy the corresponding transpose relations. Using the symmetries $(V_d)_{rt}=(V_d)_{tr}=(V_d)_{-r,-t}$, we obtain for $1\leq r,t\leq N$:
\[
  (U_d)_{rt}=(V_d)_{r,t}+(V_d)_{r,-t},\qquad
  (U_d)_{ar}=(U_d)_{br}=(V_d)_{r,0}=(V_d)_{0,r}.
\]
On the antisymmetric line $e_a-e_b$, the adjacency matrix of $D_n$ acts by zero, so $U_d(e_a-e_b)$ follows the recurrence $u_{-1}=0$, $u_0=1$, $u_{d+1}=-u_{d-1}$, giving $u_{2q}=(-1)^q$ and $u_{2q+1}=0$. Since the symmetric leaf coefficient $(V_d)_{00}$ equals~$1$ for even $d\in[0,2N]$ and~$0$ for odd $d$, writing $e_a=\frac12(g_0+e_a-e_b)$ and $e_b=\frac12(g_0-(e_a-e_b))$ yields the $\{a,b\}\times\{a,b\}$ block of $U_d$ for $0\leq d\leq2N$: it vanishes for odd $d$, and for $d=2q$ it equals
\[
  \frac12\begin{pmatrix}1+(-1)^q&1-(-1)^q\\1-(-1)^q&1+(-1)^q\end{pmatrix}.
\]
Thus $U_0,\ldots,U_{2N}$ are entrywise nonnegative integer matrices, and at $d=2N$ ($q=N$) the matrix $U_{2N}$ is the permutation matrix of $\nu$. The lemma now follows from~\eqref{eq:preprojective-lower-bound}.
\end{proof}

Returning to $\Z D_n/\langle\tau^s\rho\rangle$ with $\rho\in\{\id,\iota\}$, the brick hypothesis forces $N<2s$: otherwise, deleting the outer $N-2s$ vertices of the long arm leaves $\Z D_{2s+2}/\langle\tau^s\rho\rangle$, in which Lemma~\ref{lem:d-top-paths} (with $N'=2s$ even) provides nonzero paths of degree $4s$ fixing every vertex of $D_{2s+2}$ and advancing the repetition coordinate by $2s$, which project to nonzero radical endomorphisms at every vertex for either choice of $\rho$. Combining $N<2s$ with~\eqref{eq:d-possible-correction} yields $1\leq j-1\leq N-1<2s-1$ and $s\mid(j-1)$, which uniquely forces $j=s+1$ and $d=2s+2$ (occurring only when $s+1\leq N<2s$, with $s$ odd for the trivial twist $\theta=\id$ and $s$ even for $\theta=\iota$).

\subsubsection{Exceptional upper bounds}

For $\Delta=E_6,E_7,E_8$, we retain the presentation $R_m=k\langle x,y\rangle/(x^m,y^3,(x+y)^2)$ from Lemma~\ref{lem:e78-traces} for $m=3,4,5$, respectively. In $R_5$ over a field of characteristic two, set
\[
  \omega_3=yx^3y,\qquad \omega_4=yx^4y,\qquad \omega_5=\omega_4 x^3y,\qquad u=x^2yx^3+x^3yx^2.
\]
Reducing the word identities~\eqref{eq:e8-word-identities} modulo~$2$ gives
\begin{gather}
  \omega_3 x=\omega_4+x\omega_3+u,\quad \omega_3 y=x\omega_3+u,\quad
  \omega_4 y=\omega_4 x,\quad \omega_4 xy=0,\quad \omega_4 x^2y=\omega_4 x^3,
  \label{eq:e2-reduction-one}\\
  0=\omega_3 x^5=\sum_{r=0}^4 x^r\omega_4 x^{4-r},
  \label{eq:e2-reduction-two}
\end{gather}
where~\eqref{eq:e2-reduction-two} follows by iterating the first relation of~\eqref{eq:e2-reduction-one} five times and using $x^5=0$ to kill all resulting $u$-terms. Multiplying~\eqref{eq:e2-reduction-two} on the right by $y$ and applying~\eqref{eq:e2-reduction-one} gives $\omega_4 x^4y=x\omega_5+x^2\omega_4 x^3+x^4\omega_4 x$. Left-multiplying the last identity of~\eqref{eq:e8-word-identities} by $yx^4$ also gives $\omega_5 x=\omega_4 x^4y$, while expanding $y^2=xy+yx+x^2$ at the end of $\omega_5 y$ gives $\omega_5 y=0$; hence
\begin{equation}
  \omega_5 x=x\omega_5+x^2\omega_4 x^3+x^4\omega_4 x,\qquad \omega_5 y=0.
  \label{eq:e2-reduction-three}
\end{equation}
From~\eqref{eq:e8-word-identities} and~\eqref{eq:e2-reduction-one}--\eqref{eq:e2-reduction-three}, one checks immediately that the span of the words
\begin{equation}
  \label{eq:e2-spanning-list}
  \begin{gathered}
  x^i\ (0\leq i\leq4),\qquad
  x^iyx^j\ (0\leq i,j\leq4),\qquad
  x^i\omega_3\ (0\leq i\leq4),\\
  x^i\omega_4 x^j\ (0\leq i\leq4,\ 0\leq j\leq3),\qquad
  x^i\omega_5\ (0\leq i\leq4)
  \end{gathered}
\end{equation}
contains~$1$ and is closed under right multiplication by $x$ and $y$, so~\eqref{eq:e2-spanning-list} spans $R_5$ (with maximum word degree~$14$). Passing to the quotients $R_4=R_5/(x^4)$ and $R_3=R_5/(x^3)$ kills $\omega_4,\omega_5$ in $R_4$ (leaving $x^i$, $x^iyx^j$, $x^i\omega_3$ with $0\leq i,j\leq3$, of maximum word degree~$8$) and also kills $\omega_3$ in $R_3$ (leaving $x^i$, $x^iyx^j$ with $0\leq i,j\leq2$, of maximum word degree~$5$).

Applying Lemma~\ref{lem:e78-traces}(1) in characteristic two with $a=yx^3$ and $b=yx^4$, the relations $2=0$ and $3[ab]=0$ simplify the block identities to
\[
  [a^2W]=0\ (W\ne1),\qquad [abW]=[baW]\ (W\ne1),\qquad [ab]=0,\qquad [b^q]=0\ (q\geq2).
\]
Thus any word in $a,b$ of length at least~$3$ containing at least two $a$'s vanishes by commuting the $a$'s to be adjacent, while a word with a single $a$ reduces to $[ab^r]$, which vanishes for $r=1$ (as $[ab]=0$) and for $r\geq3$ (as $\deg(ab^r)=4+5r\geq19>14$). Among the single-$y$ words, $3[yx^2]=0$ and $5[yx^4]=0$ still force $[yx^2]=[yx^4]=0$ in characteristic two, leaving only the following spanning representatives for the positive ordinary trace:
\[
  \begin{array}{c|c|c}
  \text{Dynkin type}&\text{spanning trace representatives}&\text{path degrees }d\\ \hline
  E_6&[yx]&4\\
  E_7&[yx],\ [yx^3],\ [(yx^3)^2]&4,\ 8,\ 16\\
  E_8&[yx],\ [yx^3],\ [(yx^3)^2],\ [(yx^3)(yx^4)^2]&4,\ 8,\ 16,\ 28
  \end{array}
\]
Each listed path degree is at most one-dimensional, and all other positive ordinary trace degrees vanish.

For the nontrivial graph automorphism $\theta$ of $E_6$ (which swaps $x$ and $y$ in $R_3$), the spanning words of $R_3$ in word degrees $1$ through $5$ are:
\[
  \begin{array}{c|l}
  \text{word degree}& \text{spanning words in }R_3\\ \hline
  1&x,\ y\\
  2&x^2,\ xy,\ yx\\
  3&\xi_1=x^2y,\ \xi_2=xyx,\ \xi_3=yx^2\\
  4&\xi_4=x^2yx,\ \xi_5=xyx^2\\
  5&\xi_6=x^2yx^2
  \end{array}
\]
The $\theta$-twisted trace relations $[xw]=[wy]$ and $[yw]=[wx]$ identify the words in degree~$1$ ($[x]=[y]$) and in degree~$2$ ($[x^2]=[xy]=[yx]$, using $y^2=xy+yx+x^2$). In degrees~$3$ and~$4$, using the algebra expansions $xy^2=x^2y+xyx=\xi_1+\xi_2$, $x^2y^2=x^2yx=\xi_4$, and $y^2x^2=xyx^2=\xi_5$ (since $x^3=0$), five twisted-trace cuts kill all spanning classes:
\begin{align*}
  w=x^2\text{ in }[xw]=[wy]&\implies [\xi_1]=[x^2y]=[x^3]=0,\\
  w=x^2\text{ in }[yw]=[wx]&\implies [\xi_3]=[yx^2]=[x^3]=0,\\
  w=xy\text{ in }[xw]=[wy]&\implies [\xi_1]=[xy^2]=[\xi_1]+[\xi_2]\implies [\xi_2]=0,\\
  w=x^2y\text{ in }[xw]=[wy]&\implies 0=[x^3y]=[x^2y^2]=[\xi_4],\\
  w=yx^2\text{ in }[yw]=[wx]&\implies [\xi_5]=[y^2x^2]=[yx^3]=0.
\end{align*}
Degree~$5$ is spanned by $\xi_6$ alone, and the fixed long arm contributes no new classes. Thus the positive $\theta$-twisted trace of $\Pi(E_6)$ is at most one-dimensional in path degrees $d=2,4,10$ and vanishes in all other positive degrees.

Applying Lemma~\ref{lem:twisted-trace-surjection} to these candidate degrees:
\begin{itemize}
\item For an ordinary trace class ($\theta=\id$) in path degree $d\in\{4,8,16,28\}$, the period $s$ must divide $d/2-1\in\{1,3,7,13\}$. (In the reflected $E_6$ model $\Z E_6/\langle\tau^s\rho\rangle$, Lemma~\ref{lem:twisted-trace-surjection}(1) requires $(d/2-1)/s$ to be even for the untwisted trace; for the sole $E_6$ candidate $d=4$, $s\mid1$ forces $s=1$, for which $(d/2-1)/s=1$ is odd and hence twisted.)
\item For a $\theta$-twisted trace class on $E_6$ in path degree $d\in\{4,10\}$, Lemma~\ref{lem:twisted-trace-surjection}(2) requires $s\mid(d/2-1)$ with $(d/2-1)/s$ odd, which forces $s=1$ for $d=4$ and $s=4$ for $d=10$.
\end{itemize}
Finally, period $s=1$ is excluded in every case by deleting the vertices outside the central $D_4$ subquiver and applying Lemma~\ref{lem:d4-four-return} (noting that the degree-$4$ closed paths in $\Pi(D_4)$ advance the repetition coordinate by~$2$ and therefore also close in the reflected quotient $\Z D_4/\langle\tau\rho\rangle$ since $\rho^2=\id$). Consequently, the only remaining possibilities for a positive correction class of degree $d\geq3$ in types $E_6,E_7,E_8$ are:
\begin{equation}
  \label{eq:exceptional-correction-bounds}
  \begin{array}{c|c|c}
  \text{Dynkin type and action}&\text{period }s&\text{path degree }d\\ \hline
  E_6\text{ (with reflection }\rho\text{)}&4&10\\
  E_7&3,\ 7&8,\ 16\text{ respectively}\\
  E_8&3,\ 7,\ 13&8,\ 16,\ 28\text{ respectively}
  \end{array}
\end{equation}
In every case, each periodic component admits at most one correction degree $d=2s+2$, spanned by a single central row path. Taking $I$ to be the set of periodic components for which $\operatorname{coker}D_{2s_i+1}(\C)\ne0$ completes the proof of Proposition~\ref{prop:residual-correction-spaces}.

\subsection{Correction spaces over the integers modulo four}
\label{app:integral-corrections}

We now lift the preceding calculations to $\mathcal O=\Z/4\Z$ for the integral mesh category $E_4$ constructed in Proposition~\ref{prop:fixed-integral-companion}.

\begin{lemma}
\label{lem:integral-mesh-cokernels}
Let $\mathcal O=\Z/4\Z$, and let $E_4$ be the finite free mesh category over $\mathcal O$ of Proposition~\ref{prop:fixed-integral-companion}, with mod-$2$ reduction $\mathcal C^{(2)}$ over $\F_2$. Write $D_m(E_4)$ for its mesh derivative with respect to the unit coefficients of that presentation. With $I$ and $\chi_i$ as in Proposition~\ref{prop:residual-correction-spaces} and~\eqref{eq:parameter-character}, there is a degree- and weight-preserving isomorphism
\[
  \bigoplus_{m\geq2}\operatorname{coker}D_m(E_4)
  \cong\bigoplus_{i\in I}(\mathcal O/2\mathcal O)_{\chi_i},
\]
where the summand indexed by $i\in I$ lies in $\operatorname{coker}D_{d_i-1}(E_4)$ ($d_i=2s_i+2$) and has internal degree $d_i-2$.
\end{lemma}
\begin{proof}
Set $Q_m=\operatorname{coker}D_m(E_4)$. By Proposition~\ref{prop:fixed-integral-companion}, reducing the source, target, and matrix of $D_m(E_4)$ modulo~$2$ recovers $D_m(\mathcal C^{(2)})$, so right exactness of the tensor product gives degree- and weight-preserving identifications
\[
  Q_m/2Q_m\cong\operatorname{coker}D_m(\mathcal C^{(2)}),\qquad
  \operatorname{coker}D_m(\mathcal C^{(2)})\otimes_{\F_2}k\cong\operatorname{coker}D_m(\C).
\]
By Proposition~\ref{prop:residual-correction-spaces}, the retained row paths of path degree $m+1$ span $\operatorname{coker}D_m(\C)$ over $k$ and hence already form an $\F_2$-basis of $\operatorname{coker}D_m(\mathcal C^{(2)})$. Lifting these retained rows to $E_4$ and letting $U_m\subseteq Q_m$ denote the $\mathcal O$-submodule they generate, Nakayama's lemma (or directly $Q_m=U_m+2Q_m=U_m+2(U_m+2Q_m)=U_m$ since $4Q_m=0$) shows that $Q_m=U_m$, and in particular $Q_m=0$ in all degrees and weights without a retained row. Since the mod-$2$ reductions of the retained rows are linearly independent over $\F_2$, it suffices to show that twice each retained row lies in $\operatorname{im}D_m(E_4)$, i.e., vanishes in $Q_m$.

First, any row path that leaves its periodic component passes through a nonperiodic vertex and therefore lies in a finite triangular incidence chain as in the proof of Proposition~\ref{prop:periodic-localization}; because the mesh coefficients of $E_4$ are units in $\mathcal O$, eliminating from an exposed end expresses any such row path as a derivative over $\mathcal O$. Consequently, all terms in a contextual mesh relation that leave the periodic component vanish in $Q_m$. Likewise, the scalar normalizations of Lemma~\ref{lem:periodic-scalar-gauges}(1) hold over $\mathcal O$, since they divide only by unit mesh coefficients and rescale arrows and rows by units in $\mathcal O$. After this normalization to the signed preprojective coefficients $\pm1$, each arrow-correction column on a retained periodic component has opposite unit coefficients $\pm\alpha$ in its two row incidences, giving $\alpha[wa]-\alpha[(\tau a)w]=0$ and hence $[wa]=[(\tau a)w]$ in $Q_m$. Thus the class of a retained row in $Q_m$ factors through the positive $\theta$-twisted trace of the integral signed preprojective category $\Pi_{\mathcal O}(\Delta)$.

Moreover, the star and arm reductions of Lemmas~\ref{lem:star-return-reduction} and~\ref{lem:arm-trace-reduction} hold verbatim over $\mathcal O$, because every replacement step uses only coefficients $\pm1$:
\begin{itemize}
\item Replacing an outward backtrack $a_j b_j$ along an arm by $\pm b_{j-1}a_{j-1}$ decreases the total arm depth until any closed path at the center is a word in the length-two arm loops, whose signed sum at the center vanishes.
\item Along an arm of length $L$, replacing $b_j a_j$ at internal vertices by $\pm a_{j+1}b_{j+1}$ rewrites $(a_1 b_1)^q$ as $\pm a_1\cdots a_q b_q\cdots b_1$ for $q\leq L$ and kills $(a_1 b_1)^{L+1}$ via the leaf relation $b_L a_L=0$.
\item In the twisted trace along a $\theta$-fixed arm, rotating a closed word to a vertex of minimal depth and replacing the inward backtrack by the outward backtrack strictly increases the depth sum until the word reaches the leaf and vanishes, showing that positive powers of a fixed-arm loop have zero trace over $\mathcal O$.
\end{itemize}

We now check that $2$ annihilates the retained trace classes in each Dynkin type that can contribute to $I$:
\begin{itemize}
\item \emph{Type $D_n$.} In the central corner quotient over $\mathcal O$, the short-arm relations $x^2=y^2=0$ leave the alternating words $w_j$ and $w'_j$ of word degree $j=s+1$. For the trivial twist $\theta=\id$ ($s$ odd, $j$ even), cyclic rotation gives $[w_j]=[w'_j]$; for the short-arm interchange $\theta=\iota$ ($s$ even, $j$ odd), twisted cyclic rotation $[uv]=[v\iota(u)]$ likewise gives $[w_j]=[w'_j]$. Since the long-arm loop equals $-(x+y)$ in the central mesh relation and its $j$-th power has zero trace along the fixed long arm, we obtain in both cases
\[
  0=[(-x-y)^j]=(-1)^j([w_j]+[w'_j])=\pm2[w_j],
\]
so $2[w_j]=0$ over $\mathcal O$.

\item \emph{Reflected type $E_6$ ($s=4$, $d=10$, word degree~$5$).} In the central corner $R_3$ over $\mathcal O$, the relations $x^3=y^3=(x+y)^2=0$ give
\[
  y^2=-yx-xy-x^2,\qquad yxy=xyx,\qquad yx^2y=-xyx^2-x^2yx,
\]
where the second identity follows from $0=y^3=(-yx-xy-x^2)y=-yxy-xy^2-x^2y=-yxy-x(-yx-xy-x^2)-x^2y=-yxy+xyx+x^3=-yxy+xyx$ and the third follows by comparing $(yxy)y=(xyx)y$ with $yx(y^2)=yx(-yx-xy-x^2)$. These three identities eliminate any pair of $y$'s separated by $0$, $1$, or $2$ factors of $x$, so $R_3$ is spanned over $\mathcal O$ by $x^i$ and $x^iyx^j$ ($0\leq i,j\leq2$), and word degree~$5$ is spanned by $\xi_6=x^2yx^2$. Applying the $\theta$-twisted trace relation $[xw]=[wy]$ ($\theta(x)=y$) together with the third identity and $x^3=0$ yields
\[
  [\xi_6]=[x(xyx^2)]=[(xyx^2)y]=[x(yx^2y)]=[-x^2yx^2-x^3yx]=-[\xi_6],
\]
so $2[\xi_6]=0$ over $\mathcal O$.

\item \emph{Types $E_7$ and $E_8$.} Lemma~\ref{lem:e78-traces}(1) holds over any commutative ring, and in particular over $\mathcal O=\Z/4\Z$. Every spanning class in the positive ordinary trace is annihilated by $2$, $3$, or $5$; since $3$ and $5$ are units in $\Z/4\Z$, the entire positive trace is annihilated by~$2$.
\end{itemize}
Thus $2Q_m=0$ for all $m\geq2$, completing the proof.
\end{proof}

\clearpage
\section{Local restrictions on ray products}
\label{app:local-rays}

In this appendix we prove the two local structural results stated in Section~\ref{sec:local-ray-geometry}: the uniqueness of initial and terminal ordinary arrows for one-sided generating rays (Proposition~\ref{prop:controlled-contours}), and the classification of proper shiftable products and nondeep external returns (Proposition~\ref{prop:proper-ray-products}). Throughout, we work in the ray algebra $A_0=kP$ of Setting~\ref{set:projective-companion} over a field of arbitrary characteristic, fix a positive primitive character $\chi\in L$ as in Setting~\ref{set:fixed-character}, and use the endpoint types of Definition~\ref{def:successor-types} and the primitive return corners $S_x$ of Definition~\ref{def:primitive-return-corner}. Both proofs rely on the finite faithful supports of Definition~\ref{def:faithful-support}, the finite homogeneous categories of Subsection~\ref{subsec:homogeneous-local-categories}, and the quadratic-form criterion of Lemma~\ref{lem:dimension-infinity}.

\subsection{One-sided generating rays}
\label{app:controlled-contours}

We begin by proving Proposition~\ref{prop:controlled-contours}. Recall from Subsection~\ref{subsec:homogeneous-local-categories} that for any finite collection of graded objects $(x_i,w_i)$ with weights $w_i\in L$, we can construct a representation-finite homogeneous category $H$ whose morphism spaces are spanned by rays $r\colon x_i\to x_j$ of weight $\ell(r)=w_j-w_i$. By Lemma~\ref{lem:endpoint-chains}(3), the corner bimodule $A_0(x,z)$ of a generating ray $f\colon x\to z$ is cyclic over at least one endpoint ring, so every ray in $A_0(x,z)$ has weight $\ell(f)$ plus a nonnegative multiple of a positive local weight. Since $\chi$ is positive, $A_0(x,z)$ contains no ray of weight $\ell(f)-n\chi$ for $n>0$, and a ray of weight $\ell(f)+\chi$, when it exists, is the proper successor $f^+$. Moreover, by Lemma~\ref{lem:generating-paths}, every factor of a generating shiftable ray $g$ is itself a generating ray (as the corner bimodule of $g$ has dimension at least~$2$).

\subsubsection{Bimodule retractions from homogeneous categories}

To transfer finite representation type from a finite homogeneous category $H$ to the incidence algebra $kQ$ of a finite poset $Q$, we use a general criterion for when a family of selected rays in $H$ admits a $kQ$-bimodule retraction, even when $H$ contains additional morphisms outside $kQ$ that do not form an ideal. Write $a\lessdot b$ when $a<b$ is a cover relation in $Q$ (so no element of $Q$ lies strictly between $a$ and $b$).

\begin{lemma}
\label{lem:configuration-retraction}
Let $H$ be a finite homogeneous category whose objects are indexed by a finite poset $Q$. Suppose that for each pair $a\leq b$ in $Q$ we have chosen a nonzero basis morphism $r_{ab}\in H(a,b)$ such that $r_{aa}=1_a$ and $r_{ab}r_{bc}=r_{ac}$ for all $a\leq b\leq c$. Assume that for every cover relation $a\lessdot b$ in $Q$ and every homogeneous basis morphism $h$ in $H$ with the appropriate endpoints,
\[
  \begin{split}
  r_{ab}h=r_{ac}\quad(a\leq c)&\quad\Longrightarrow\quad b\leq c,\\
  hr_{ab}=r_{cb}\quad(c\leq b)&\quad\Longrightarrow\quad c\leq a.
  \end{split}
\]
Then the linear projection $H\to kQ$ sending each selected basis morphism $r_{ab}$ to the corresponding poset basis element of $kQ$ and vanishing on all other basis morphisms of $H$ is a $kQ$-bimodule retraction. In particular, if $H$ is representation-finite, then so is the incidence algebra $kQ$.
\end{lemma}
\begin{proof}
The relations $r_{aa}=1_a$ and $r_{ab}r_{bc}=r_{ac}$ embed the incidence algebra $kQ$ as a unital subalgebra of $H$. We first extend the two cover implications to arbitrary pairs $a\leq b$ in $Q$ by induction on the interval $[a,b]$. If $a<b$, choose a cover $a\lessdot a_1\leq b$. Any equality $r_{ab}h=r_{ac}$ with $a\leq c$ can be rewritten as $r_{aa_1}(r_{a_1b}h)=r_{ac}$. Since $r_{ac}\ne0$, the inner product $r_{a_1b}h$ is a nonzero basis morphism, so the cover assumption for $a\lessdot a_1$ gives $a_1\leq c$; because $H$ is thin, $r_{a_1b}h$ must equal $r_{a_1c}$, and the induction hypothesis on $[a_1,b]$ then yields $b\leq c$. The second implication extends in the same way by factoring out the last cover below $b$.

It follows that the composition in $H$ of a selected basis morphism $r_{ab}$ with an unselected basis morphism $h$ (on either side) can never equal a selected basis morphism $r_{ac}$ or $r_{cb}$: otherwise the extended implications would give $b\leq c$ (respectively, $c\leq a$), forcing $r_{ab}h=r_{ab}r_{bc}$ and hence $h=r_{bc}$ (respectively, $h=r_{ca}$) by thinness of $H$. Thus the linear projection $H\to kQ$ onto the span of the selected morphisms $\{r_{ab}\}_{a\leq b}$ is a $kQ$-bimodule homomorphism restricting to the identity on $kQ$. The final assertion then follows from Lemma~\ref{lem:split-representation-transfer}.
\end{proof}

We will apply Lemma~\ref{lem:configuration-retraction} when $Q$ is the reachability poset of an oriented tree and the selected morphisms $r_{ab}$ are the evaluations of the unique oriented paths in the tree. Since the cover relations of $Q$ are precisely the arrows of the tree, it suffices to verify the two implications of Lemma~\ref{lem:configuration-retraction} at each tree arrow (and any candidate factor $h$ whose required weight has negative categorical degree $\deg_\JJ(h)<0$ is automatically absent).

The first tree we will use is the six-vertex tree with two adjacent trivalent vertices:
\begin{equation}
  \sixlifttree
  \label{eq:six-lift-tree}
\end{equation}
In our application below, the two arrows labeled $\alpha$ will be lifts of the same ray $\alpha$, with $(v_4,v_5)$ obtained from $(v_2,v_3)$ by shifting both object weights by $-\chi$ so that the six graded objects are pairwise distinct. Assigning dimension~$2$ to the trivalent vertices $v_2,v_5$ and dimension~$1$ to the four leaves gives $N=D=12$, so Lemma~\ref{lem:dimension-infinity} rules out every orientation of~\eqref{eq:six-lift-tree}. The second tree we will use is the eight-vertex tree of type $\widetilde E_7$ with arm lengths $(3,3,1)$, which is likewise representation-infinite in every orientation by the dimension vector recorded after Lemma~\ref{lem:dimension-infinity}.

\subsubsection{Two distinct first arrows are impossible}

Suppose for a contradiction that a generating left-sided ray $f\colon x\to z$ can be represented by ordinary paths starting with two distinct ordinary arrows $a_1\colon x\to y_1$ and $a_2\colon x\to y_2$. Write $f=a_1w_1=a_2w_2$; since $a_1\ne a_2$ are ordinary arrows, they are incomparable under right factorization ($a_1h=a_2$ with $\deg_\JJ(h)>0$ would contradict the ordinariness of $a_2$, while $\deg_\JJ(h)=0$ would force $a_1=a_2$), and both suffixes $w_1,w_2$ are nonidentity generating rays. Because $\widehat q_xf=\widehat q_xa_iw_i\ne0$, both $\widehat q_xa_1$ and $\widehat q_xa_2$ are nonzero.

We first show that one of $a_1,a_2$ is left-sided and the other is two-sided:
\begin{itemize}
\item \emph{Both $a_1$ and $a_2$ cannot be two-sided.} If both were two-sided, then $\widehat q_xa_i=a_i\widehat q_{y_i}$, so $\widehat q_xf=a_i(\widehat q_{y_i}w_i)$ would be right-divisible by both $a_1$ and $a_2$, as is $f=a_iw_i$. In the faithful projective support at $x$, the four rays $\{a_1,a_2,f,\widehat q_xf\}$ would then form a four-point crown: $a_1$ and $a_2$ are incomparable; $f$ and $\widehat q_xf$ are incomparable because $f$ is generating and left-sided ($f\widehat q_z=0$); and neither $f$ nor $\widehat q_xf$ can precede $a_1$ or $a_2$ for degree reasons. This contradicts Lemma~\ref{lem:dimension-infinity}.

\item \emph{Both $a_1$ and $a_2$ cannot be left-sided.} If both were left-sided, consider the six rays
\[
  1_x,\quad \widehat q_x,\quad a_1,\quad a_2,\quad \widehat q_xa_1,\quad \widehat q_xa_2
\]
in the faithful projective support at $x$. We claim that their induced poset has only the comparisons generated by
\[
  1_x\longrightarrow\widehat q_x\longrightarrow\widehat q_xa_1,\widehat q_xa_2,\qquad
  1_x\longrightarrow a_1,a_2.
\]
Indeed, $\widehat q_x$ cannot precede $a_i$ (as $a_i$ is ordinary) and $a_i$ cannot precede $\widehat q_x$ (as $a_ih=\widehat q_x$ would place $y_i\in S_x$ by Lemma~\ref{lem:primitive-cyclic-corners}(3) and make $a_i$ two-sided); $a_1$ and $a_2$ (and hence $\widehat q_xa_1$ and $\widehat q_xa_2$) are incomparable; and any cross comparison $a_ih=\widehat q_xa_j$ ($i\ne j$) would multiply on the right by $w_j$ to give $a_i(hw_j)=\widehat q_xa_jw_j=\widehat q_xf\ne0$, forcing $hw_j$ to be the unique ray $w_i^+$ in $A_0(y_i,z)$ of weight $\ell(w_i)+\chi$, which contradicts $a_iw_i^+=0$ from the $(\mathsf L,*)$ row of~\eqref{eq:product-type-table}. All reverse comparisons are ruled out because $a_1,a_2$ are generating or by degree. The resulting six-element poset is the orientation of~\eqref{eq:six-lift-tree} with two outgoing leaves at each trivalent vertex, which is excluded by Lemma~\ref{lem:dimension-infinity}.
\end{itemize}

Thus one of the two initial arrows of a generating left-sided ray must have type~$\mathsf B$ and the other type~$\mathsf L$; dually, if a generating right-sided ray had two distinct terminal arrows, one would have type~$\mathsf B$ and the other type~$\mathsf R$. Passing to the opposite category $A_0^{\op}$, it therefore suffices to rule out a generating right-sided ray $f\colon x\to z$ admitting two factorizations
\[
  f=\gamma\alpha=\eta\beta,\qquad
  \alpha\colon y\to z\text{ ordinary of type }\mathsf B,\quad
  \beta\colon t\to z\text{ ordinary of type }\mathsf R.
\]
Here the prefixes $\gamma\colon x\to y$ and $\eta\colon x\to t$ are nonidentity generating rays by Lemma~\ref{lem:generating-paths}, and are incomparable (for instance, $\gamma h=\eta$ would imply $\gamma(h\beta)=\gamma\alpha=f\ne0$, so $h\beta=\alpha$ because both lie in $A_0(y,z)$ with the same weight, contradicting the fact that $\alpha$ and $\beta$ are distinct ordinary arrows). Using the local generators $\widehat q_y$ and $\widehat q_z$, we have $\gamma\widehat q_y\alpha=\gamma\alpha\widehat q_z=f\widehat q_z\ne0$, so $\gamma$ has type~$\mathsf R$ (since $f$ has type~$\mathsf R$, $\widehat q_x\gamma=0$). Moreover, if a local ray $u\in R_x$ has weight $\chi$, then primitivity of $\chi$ gives $u=\widehat q_x$, and hence $uf=\widehat q_xf=0\ne f\widehat q_z$ because $f$ is right-sided; thus
\begin{equation}
  uf\ne f\widehat q_z\qquad\text{for every local ray }u\in R_x.
  \label{eq:absent-source-action}
\end{equation}

We now distinguish two cases according to whether $\alpha$ admits a primitive reverse.

\medskip
\noindent\textbf{Case 1: $\alpha$ has a primitive reverse $\lambda\colon z\to y$.}
Here $\alpha\lambda=\widehat q_y$ and $\lambda\alpha=\widehat q_z$. Set $\zeta=\beta\lambda\colon t\to y$, which is nonzero because $\zeta\alpha=\beta\widehat q_z\ne0$. In a finite homogeneous category $H$, choose the six graded objects
\begin{equation}
  \begin{array}{c|rrrrrr}
  \text{vertex}&v_1&v_2&v_3&v_4&v_5&v_6\\ \hline
  \text{object}&x&y&z&y&z&t\\
  \text{weight}&-\ell(\gamma)&0&\ell(\alpha)&-\chi&\ell(\alpha)-\chi&\ell(\alpha)-\chi-\ell(\beta)
  \end{array}
  \label{eq:six-lift-coordinates}
\end{equation}
and select the tree~\eqref{eq:six-lift-tree}, whose composite oriented paths evaluate to the nonzero rays $\widehat q_y,\zeta,\widehat q_z,\alpha\widehat q_z,\beta\widehat q_z,f$. To verify the cover conditions of Lemma~\ref{lem:configuration-retraction}, note that every left-multiplication check $r_{ab}h=r_{ac}$ with $b\nleq c$ requires $\deg_\JJ(h)<0$, while the right-multiplication checks $hr_{ab}=r_{cb}$ with $c\nleq a$ and $\deg_\JJ(h)\geq0$ reduce to the following forbidden factorizations:
\[
  \begin{array}{c|l}
  \text{forbidden factorization}&\text{reason for exclusion}\\ \hline
  h\gamma=\widehat q_y&(\gamma h)f=\gamma\widehat q_y\alpha=f\widehat q_z\text{, contradicting~\eqref{eq:absent-source-action}}\\
  h\gamma=\lambda&hf=h\gamma\alpha=\widehat q_z\implies(fh)f=f\widehat q_z\text{, contradicting~\eqref{eq:absent-source-action}}\\
  h\gamma=\zeta&(\eta h)f=\eta\beta\lambda\alpha=f\widehat q_z\text{, contradicting~\eqref{eq:absent-source-action}}\\
  h\alpha=\beta\text{ or }h\beta=\alpha&\alpha,\beta\text{ are distinct ordinary arrows}\\
  h\lambda=\gamma&\ell(h)=\ell(\gamma)-\ell(\lambda)=\ell(f)-\chi\text{ lies below the generator }f
  \end{array}
\]
Thus Lemma~\ref{lem:configuration-retraction} provides a bimodule retraction onto the incidence algebra of~\eqref{eq:six-lift-tree}, contradicting Lemma~\ref{lem:dimension-infinity}. Consequently, $\alpha$ cannot admit a primitive reverse.

\medskip
\noindent\textbf{Case 2: $\alpha$ has no primitive reverse.}
We first note that no ray $h\colon t\to y$ can satisfy $\beta\widehat q_z=h\alpha$. Otherwise, in the faithful injective support at $z$, the four rays $\{\widehat q_z,\alpha,\alpha\widehat q_z,\beta\widehat q_z\}$ would form a four-point crown: $\widehat q_z$ and $\alpha$ are incomparable because $\alpha$ is ordinary and has no primitive reverse; $\alpha\widehat q_z$ and $\beta\widehat q_z$ are incomparable because $u\alpha\widehat q_z=\beta\widehat q_z$ or $u\beta\widehat q_z=\alpha\widehat q_z$ cancels $\widehat q_z$ to $u\alpha=\beta$ or $u\beta=\alpha$; and no upper ray $\alpha\widehat q_z,\beta\widehat q_z$ can precede $\widehat q_z$ or $\alpha$ (the only nonnegative-degree possibility $u\beta\widehat q_z=\alpha$ would factor the ordinary arrow $\alpha$ through $\widehat q_z$).

Now choose the eight graded objects in $H$:
\begin{equation}
  \begin{gathered}
  \begin{array}{c|rrrr}
  \text{vertex}&v_1&v_2&v_3&v_4\\ \hline
  \text{object}&x&x&y&y\\
  \text{weight}&-\ell(\gamma)&\chi-\ell(\gamma)&0&\chi
  \end{array}
  \\[1ex]
  \begin{array}{c|rrrr}
  \text{vertex}&v_5&v_6&v_7&v_8\\ \hline
  \text{object}&z&z&t&t\\
  \text{weight}&\ell(\alpha)&\ell(\alpha)-\chi&\ell(\alpha)+\chi-\ell(\beta)&\ell(\alpha)-\chi-\ell(\beta)
  \end{array}
  \end{gathered}
  \label{eq:eight-lift-coordinates}
\end{equation}
and select the oriented tree of type $\widetilde E_7$:
\begin{equation}
\begin{tikzpicture}[>=Latex,x=1.2cm,y=1.1cm,every node/.style={font=\small}]
\node (G) at (0,0) {$v_7$}; \node (B) at (1,0) {$v_2$};
\node (D) at (2,0) {$v_4$}; \node (C) at (3,0) {$v_3$};
\node (E) at (4,0) {$v_5$}; \node (F) at (5,0) {$v_6$};
\node (H) at (6,0) {$v_8$}; \node (A) at (3,1) {$v_1$};
\draw[->] (B)--node[above] {$\eta$}(G);
\draw[->] (B)--node[above] {$\gamma$}(D);
\draw[->] (C)--node[above] {$\widehat q_y$}(D);
\draw[->] (C)--node[above] {$\alpha$}(E);
\draw[->] (F)--node[above] {$\widehat q_z$}(E);
\draw[->] (H)--node[above] {$\beta$}(F);
\draw[->] (A)--node[right] {$\gamma$}(C);
\end{tikzpicture}
\label{eq:eight-lift-tree}
\end{equation}
Its composite oriented paths evaluate to the nonzero rays $\gamma\widehat q_y$, $f=\gamma\alpha$, and $\beta\widehat q_z$. After discarding negative-degree factors, the cover conditions of Lemma~\ref{lem:configuration-retraction} across the seven arrows of~\eqref{eq:eight-lift-tree} reduce to the following forbidden factorizations:
\[
  \begin{array}{c|l}
  \text{forbidden factorization}&\text{reason for exclusion}\\ \hline
  \widehat q_yh=\alpha\text{ or }h\widehat q_z=\alpha&\alpha\text{ is an ordinary arrow}\\
  \alpha h=\widehat q_y\text{ or }h\alpha=\widehat q_z&\alpha\text{ has no primitive reverse}\\
  \gamma h=\eta\text{ or }\eta h=\gamma&\gamma\text{ and }\eta\text{ are incomparable}\\
  h\widehat q_y=\gamma&\ell(h)=\ell(\gamma)-\chi\text{ lies below the generator }\gamma\\
  h\alpha=\beta\widehat q_z&\text{excluded by the injective-support crown above}\\
  h\gamma=\gamma\widehat q_y&hf=h\gamma\alpha=\gamma\widehat q_y\alpha=f\widehat q_z\text{, contradicting~\eqref{eq:absent-source-action}}\\
  h\gamma=\widehat q_y&(\gamma h)f=\gamma\widehat q_y\alpha=f\widehat q_z\text{, contradicting~\eqref{eq:absent-source-action}}\\
  h\widehat q_z=f&\ell(h)=\ell(f)-\chi\text{ lies below the generator }f
  \end{array}
\]
(In the second row, note that $\alpha h=\widehat q_y$ forces $\alpha(h\alpha)=\widehat q_y\alpha=\alpha\widehat q_z\ne0$ and hence $h\alpha=\widehat q_z$, and dually $h\alpha=\widehat q_z$ forces $\alpha h=\widehat q_y$, so either equality would make $h$ a primitive reverse of $\alpha$.) Lemma~\ref{lem:configuration-retraction} therefore yields a bimodule retraction onto the incidence algebra of~\eqref{eq:eight-lift-tree}. Assigning dimension~$4$ to $v_3$, dimensions $3,2,1$ along the two arms of length~$3$, and dimension~$2$ to $v_1$ gives $N=D=48$, contradicting Lemma~\ref{lem:dimension-infinity}. This completes the proof of Proposition~\ref{prop:controlled-contours}.

\subsection{Proper products and external returns}
\label{app:proper-products}

We now prove Proposition~\ref{prop:proper-ray-products}. Recall that the finite homogeneous categories of Subsection~\ref{subsec:homogeneous-local-categories} are representation-finite by Lemma~\ref{lem:homogeneous-projection}, as are their corner subcategories and object quotients, and in particular the incidence algebras of the faithful projective and injective supports of Definition~\ref{def:faithful-support}.

\subsubsection{One-sided lengths and the two-object return}

We begin by noting a mild strengthening of Case~1 in the proof of Lemma~\ref{lem:opposite-controls}: if $f\colon u\to c$ and $g\colon c\to v$ are generating rays whose corner bimodules are generated only at $c$, and if the radical generator $q=q_c$ of $R_c$ satisfies $fq\ne0$, $qg\ne0$, and $q^2\ne0$, then
\begin{equation}
  fg=0.
  \label{eq:strong-star-vanishing}
\end{equation}
Indeed, Case~1 of the proof of Lemma~\ref{lem:opposite-controls} uses only the four nonzero composite rays $fg,fq,qg,q^2$ together with the fact that $f$ and $g$ are generating rays and satisfy the radical inclusions~\eqref{eq:opposite-control-inclusions}; it does not require $fq^2$ or $q^2g$ to be nonzero.

We now prove the dimension bound $\dim A_0(b,y)\leq3$ (with equality only when $N_b=3$) for a generating left-sided ray $f\colon b\to y$ in Proposition~\ref{prop:proper-ray-products}(1). If $\widehat q_b=0$ or $\widehat q_bf=0$, then $\dim A_0(b,y)=1$; if $\widehat q_bf\ne0$, then $A_0(b,y)=R_bf$ by Lemma~\ref{lem:primitive-action-generation}, so $\dim A_0(b,y)\leq N_b$, which already gives the bound when $N_b\leq3$. Suppose for a contradiction that $N_b\geq4$ and $\widehat q_b^2f\ne0$ (so $\dim A_0(b,y)\geq3$). We reach a contradiction in three steps:

\medskip
\noindent\textbf{Step 1: Support points and forward comparisons.}
In the faithful projective support at $b$, retain the six rays
\[
  Q_i=\widehat q_b^i\quad(0\leq i\leq3),\qquad
  Y_j=\widehat q_b^jf\quad(j=1,2).
\]
Clearly $Q_0<Q_1<Q_2<Q_3$, and $Q_i<Y_j$ whenever $i\leq j$. Conversely, if $i>j$, an inequality $Q_i\leq Y_j$ would require a ray $h\in A_0(b,y)$ of weight $\ell(f)-(i-j)\chi$, which is impossible because $A_0(b,y)=R_bf$ has only weights $\ell(f)+a\chi$ with $a\geq0$. Moreover, $Y_1$ and $Y_2$ are incomparable: any comparison between them would produce a local ray at $y$ of weight $\chi$, which by primitivity of $\chi$ would force $f\widehat q_y\ne0$ and make $f$ two-sided.

\medskip
\noindent\textbf{Step 2: Excluding reverse comparisons.}
Suppose that $Y_j<Q_i$ for some $j\in\{1,2\}$ and $i\in\{0,1,2,3\}$. Left cancellation of $\widehat q_b^j$ gives $fh=\widehat q_b^a$ for a ray $h\colon y\to b$ and $a=i-j\in\{1,2\}$. Since $a\leq2$ and $N_b\geq4$, we have $fh\widehat q_b=\widehat q_b^{a+1}\ne0$, so $h\widehat q_b\ne0$. By the dual of Lemma~\ref{lem:primitive-action-generation}, the corner bimodule $A_0(y,b)$ is generated at $b$: writing $h=l\widehat q_b^u$ for the generating ray $l\colon y\to b$ gives $fl=\widehat q_b^m$ with $1\leq m\leq a\leq2$. Note that $y\notin S_b$ because $f$ is one-sided, and both $lf$ and $l\widehat q_b$ are nonzero because $f(lf)=\widehat q_b^mf\ne0$ and $f(l\widehat q_b)=\widehat q_b^{m+1}\ne0$. If $A_0(y,b)$ were also cyclic at $y$, Lemma~\ref{lem:primitive-action-generation} would give $q_yl=l\widehat q_b\ne0$ with $\ell(q_y)=\chi$, and Lemma~\ref{lem:primitive-cyclic-corners}(1) applied to $l$ and $f$ would yield $lf=q_y$ and $fl=\widehat q_b$, placing $y\in S_b$, a contradiction. Thus both $l\colon y\to b$ and $f\colon b\to y$ are generated only at $b$. Since $l\widehat q_b\ne0$, $\widehat q_bf\ne0$, and $\widehat q_b^2\ne0$, the strengthened star identity~\eqref{eq:strong-star-vanishing} gives $lf=0$, contradicting $f(lf)=\widehat q_b^mf\ne0$.

\medskip
\noindent\textbf{Step 3: The dimension contradiction.}
By Steps~1 and~2, the induced poset on $\{Q_0,Q_1,Q_2,Q_3,Y_1,Y_2\}$ has cover diagram
\[
  Q_0\longrightarrow Q_1\longrightarrow Q_2\longrightarrow Q_3,\qquad
  Q_1\longrightarrow Y_1,\qquad Q_2\longrightarrow Y_2,
\]
which is an orientation of the representation-infinite tree~\eqref{eq:six-lift-tree} (with $N=D=12$ for dimensions~$2$ on $Q_1,Q_2$ and~$1$ on the four leaves), contradicting Lemma~\ref{lem:dimension-infinity}. Thus $\widehat q_b^2f=0$ and $\dim A_0(b,y)\leq2$ whenever $N_b\geq4$; the dual argument gives the same bound for a generating right-sided ray.

\medskip
Before proving the second part of Proposition~\ref{prop:proper-ray-products}(1), we isolate the two-object return calculation used in Step~2 above.

\begin{lemma}
\label{lem:nondeep-two-object-return}
Work in Settings~\ref{set:projective-companion} and~\ref{set:fixed-character}. Let $b$ be a supporting vertex, let $y\notin S_b$, and let $e\colon b\to y$ and $l\colon y\to b$ be generating rays satisfying
\begin{equation}
  el=\widehat q_b^a\qquad\text{with}\qquad\widehat q_b^{a+1}\ne0.
  \label{eq:external-return-assumption}
\end{equation}
Then $e$ has type~$\mathsf L$, $l$ has type~$\mathsf R$, and
\begin{equation}
  N_b=4,\qquad el=\widehat q_b^2,\qquad le=0.
  \label{eq:nondeep-two-object-table}
\end{equation}
\end{lemma}
\begin{proof}
Since $el\widehat q_b=\widehat q_b^{a+1}\ne0$, we have $\widehat q_be\ne0$ and $l\widehat q_b\ne0$, so both $A_0(b,y)$ and $A_0(y,b)$ are generated at $b$ by Lemma~\ref{lem:primitive-action-generation}. Neither $A_0(b,y)$ nor $A_0(y,b)$ can be cyclic at $y$: if $A_0(b,y)$ were cyclic at $y$, Lemma~\ref{lem:primitive-action-generation} and Lemma~\ref{lem:primitive-cyclic-corners}(1) would give $el=\widehat q_b$ and $le=q_y$, placing $y\in S_b$; if $A_0(y,b)$ were cyclic at $y$, then $q_yl=l\widehat q_b\ne0$ would imply $(eq_y)l=el\widehat q_b=\widehat q_b^{a+1}\ne0$, so $eq_y=\widehat q_be\ne0$ and $A_0(b,y)$ would be cyclic at $y$ as well. Thus $e$ has type~$\mathsf L$, $l$ has type~$\mathsf R$, and both corner bimodules are generated only at $b$. By~\eqref{eq:strong-star-vanishing} applied to $l\colon y\to b$ and $e\colon b\to y$, we obtain $le=0$. Moreover, $a\geq2$ because $a=1$ would place $y\in S_b$ by Lemma~\ref{lem:primitive-cyclic-corners}(3). Note also that $\widehat q_b^ae=e(le)=0$ and $l\widehat q_b^a=(le)l=0$.

To rule out $a\geq3$, set $E=\ell(e)$ and consider in a finite homogeneous category $H$ the six graded objects
\[
  B_i=(b,i\chi)\quad(0\leq i\leq3),\qquad
  Y_+=(y,E+\chi),\qquad Y_-=(y,E+(2-a)\chi),
\]
together with the selected morphisms along the oriented tree
\begin{equation}
  B_0\xrightarrow{\widehat q_b}B_1\xrightarrow{\widehat q_b}B_2\xrightarrow{\widehat q_b}B_3,\qquad
  B_1\xrightarrow{e}Y_+,\qquad Y_-\xrightarrow{l}B_2.
  \label{eq:proper-six-tree}
\end{equation}
Every oriented path in~\eqref{eq:proper-six-tree} evaluates to a nonzero ray in $H$ because $\widehat q_b^3$ (as $a\geq2$ and $\widehat q_b^{a+1}\ne0$), $\widehat q_be$, and $l\widehat q_b$ are all nonzero. Writing $\ell(l)=a\chi-E$, any ray in $A_0(b,y)=R_be$ has weight $E+m\chi$ with $0\leq m<a$, and any ray in $A_0(y,b)=lR_b$ has weight $\ell(l)+m\chi$ with $0\leq m<a$. Relative to the generating weights $E=\ell(e)$ (for $B_i\to Y_\pm$) and $\ell(l)$ (for $Y_\pm\to B_i$), the required exponent $m$ of $\chi$ in each cross morphism space is:
\[
  \begin{array}{c|cccc}
  \text{morphism space}&B_i\to Y_+&Y_-\to B_i&B_i\to Y_-&Y_+\to B_i\\ \hline
  \text{exponent }m\text{ of }\chi&1-i&i-2&2-a-i&i-a-1
  \end{array}
\]
When $a\geq3$, the exponents in the last two columns are negative for all $0\leq i\leq3$, while the exponents in the first two columns lie in $\{0,\ldots,a-1\}$ only for the selected paths $B_0,B_1\to Y_+$ and $Y_-\to B_2,B_3$. The only remaining unselected morphism space between distinct objects of~\eqref{eq:proper-six-tree} is $H(Y_-,Y_+)$ (of weight $(a-1)\chi$), which cannot compose on either side with any nonidentity selected morphism because $Y_-$ is a source and $Y_+$ is a sink in~\eqref{eq:proper-six-tree}. Thus Lemma~\ref{lem:configuration-retraction} gives a bimodule retraction onto the incidence algebra of~\eqref{eq:proper-six-tree}, contradicting Lemma~\ref{lem:dimension-infinity} (with $N=D=12$). Hence $a=2$, and therefore $\widehat q_b^4=(el)^2=e(le)l=0$ while $\widehat q_b^3\ne0$, so $N_b=4$.
\end{proof}

We now complete the proof of Proposition~\ref{prop:proper-ray-products}(1) by showing that if $rs\ne0$ and $r$ both have type~$\mathsf L$, then $rs$ and $r$ have the same endpoint layer (and dually for type~$\mathsf R$).
\begin{itemize}
\item \emph{Generating first factor.} Suppose that $f\colon b\to y$ is a generating ray of type~$\mathsf L$ and that $fv$ is a proper shiftable ray of type~$\mathsf L$, so that $fv=\widehat q_b^ih$ for a generating ray $h\colon b\to z$ of type~$\mathsf L$ and $i\geq1$. By the dimension bound just proved, $A_0(b,z)$ has dimension at most~$3$ and $fv$ is shiftable, so $i=1$ and $fv=\widehat q_bh$. Since $f$ and $h$ are one-sided, $y,z\notin S_b$, so Lemma~\ref{lem:external-hom-chain} implies that $f$ and $h$ are comparable under factorization:
  \begin{itemize}
  \item If $fc=h$ for a nonidentity ray $c\colon y\to z$, then $fv=\widehat q_bfc=f^+\!c$ shows that $\ell(v)=\ell(c)+\chi$, so $v=c^+$. Applying the $(\mathsf L,*)$ row of~\eqref{eq:product-type-table} to the nonzero product $fc=h$ of type~$\mathsf L$ forces the second substitution $fc^+=fv$ to vanish, contradicting $fv=\widehat q_bh\ne0$.
  \item If $y=z$ and $f=h$, then $\ell(v)=\chi$, so $v=\widehat q_y$ by primitivity, making $f\widehat q_y=fv\ne0$ and contradicting the hypothesis that $f$ has type~$\mathsf L$.
  \item If $hc=f$ for a ray $c\colon z\to y$, then $\widehat q_bh=fv=h(cv)$, so $cv\in R_z$ has weight $\chi$ and therefore equals $\widehat q_z$, making $h\widehat q_z=\widehat q_bh\ne0$ and contradicting the hypothesis that $h$ has type~$\mathsf L$.
  \end{itemize}
  Thus a generating ray $f$ of type~$\mathsf L$ cannot have a proper shiftable product $fv$ of type~$\mathsf L$.

\item \emph{Proper first factor.} Now let $r=\widehat q_b^ir_0$ ($i\geq1$, $r_0$ generating of type~$\mathsf L$) be a proper ray of type~$\mathsf L$ such that $rs=\widehat q_b^i(r_0s)$ has type~$\mathsf L$. Then $\widehat q_b(r_0s)$ is a factor of $\widehat q_b^{i+1}(r_0s)=\widehat q_b(rs)\ne0$, so $r_0s$ is shiftable of type~$\mathsf L$ or~$\mathsf B$. If $r_0s$ had type~$\mathsf B$, then $(rs)\widehat q_z=\widehat q_b^i(r_0s\widehat q_z)=\widehat q_b^{i+1}(r_0s)=\widehat q_b(rs)\ne0$, contradicting the hypothesis that $rs$ has type~$\mathsf L$. Thus $r_0s$ has type~$\mathsf L$, so by the generating case just proved $r_0s$ has endpoint layer~$0$, and hence $rs=\widehat q_b^i(r_0s)$ has endpoint layer $i$, matching the endpoint layer of $r$. The statement for type~$\mathsf R$ follows by duality.
\end{itemize}

\subsubsection{Absence of external two-sided action at the return base}
\label{app:external-two-sided-action}

We next prove Proposition~\ref{prop:proper-ray-products}(4). Let $(e,l)$ be a nondeep external return at $b$ through $y\notin S_b$ (so $el=\widehat q_b^2$, $\widehat q_b^3\ne0$, $N_b=4$, and $le=0$ by Lemma~\ref{lem:nondeep-two-object-return}), and suppose for a contradiction that there exists a generating two-sided ray $f\colon b\to c$ with $c\notin S_b$ and $\widehat q_bf\ne0$.

First observe that no ray $v\colon c\to b$ can satisfy $fv=\widehat q_b^2$. Indeed, since $f$ is generating and two-sided, $A_0(b,c)$ is cyclic at both endpoints by Lemma~\ref{lem:primitive-action-generation}; if $fv=\widehat q_b^2\ne0$, then Lemma~\ref{lem:primitive-cyclic-corners}(1) supplies a generating ray $h'\colon c\to b$ with $fh'=\widehat q_b$ and $h'f=q_c$, which places $c\in S_b$ and contradicts $c\notin S_b$.

Now apply Lemma~\ref{lem:external-hom-chain} at $b$ to compare the outgoing generating rays $e\colon b\to y$ and $f\colon b\to c$. We cannot have $fa=e$ for a ray $a\colon c\to y$, as this would give $f(al)=el=\widehat q_b^2$, which was just ruled out. Hence $eh=f$ for some ray $h\colon y\to c$. In the faithful projective support at $b$, consider the four rays
\[
  e,\quad \widehat q_be;\qquad \widehat q_b^3,\quad \widehat q_bf.
\]
All four upward comparisons from $\{e,\widehat q_be\}$ to $\{\widehat q_b^3,\widehat q_bf\}$ hold via right multiplication by $l\widehat q_b$, $l$, $h\widehat q_c$ (using $\widehat q_bf=f\widehat q_c=eh\widehat q_c$), and $h$ (using $\widehat q_be\cdot h=\widehat q_bf$). On the other hand:
\begin{itemize}
\item $e$ and $\widehat q_be$ are incomparable because $e$ is left-sided;
\item $\widehat q_b^3$ and $\widehat q_bf$ are incomparable: $\widehat q_b^3u=\widehat q_bf$ cancels to $\widehat q_b^2u=f$, contradicting the fact that $f$ is generating, while $(\widehat q_bf)v=\widehat q_b^3$ cancels to $fv=\widehat q_b^2$, which was ruled out above;
\item no upper ray can precede a lower ray: $\widehat q_b^3u=\widehat q_b^te$ ($t\in\{0,1\}$) would require a ray in $A_0(b,y)=R_be$ of weight $\ell(e)-(3-t)\chi<\ell(e)$; $(\widehat q_bf)u=e$ would contradict the fact that $e$ is generating; and $(\widehat q_bf)u=\widehat q_be\ne0$ cancels to $fu=e$, which was already ruled out.
\end{itemize}
Thus $\{e,\widehat q_be,\widehat q_b^3,\widehat q_bf\}$ forms a forbidden four-point crown in the faithful projective support at $b$, contradicting Lemma~\ref{lem:dimension-infinity}. The incoming assertion for $g\colon z\to b$ follows by passing to $A_0^{\op}$.

\subsubsection{Products with one-sided factors}

We now prove Proposition~\ref{prop:proper-ray-products}(2) and begin the proof of Proposition~\ref{prop:proper-ray-products}(3). By Lemma~\ref{lem:successor-substitutions} and the layer-preservation property of Proposition~\ref{prop:proper-ray-products}(1), if $r\colon x\to y$ and $s\colon y\to z$ are generating rays whose product $rs$ is proper and shiftable, the pair of endpoint types of $(r,s)$ must be one of
\[
  (\mathsf B,\mathsf B),\qquad (\mathsf B,\mathsf L),\qquad (\mathsf R,\mathsf B),\qquad (\mathsf L,\mathsf R),
\]
and when $rs$ has type~$\mathsf L$ or~$\mathsf R$, Proposition~\ref{prop:proper-ray-products}(1) forces $rs$ to have endpoint layer~$1$ inside a three-dimensional corner bimodule $A_0(x,z)$ with $N_x=3$ (respectively, $N_z=3$).

For Proposition~\ref{prop:proper-ray-products}(2), let $r\colon x\to y$ of type~$\mathsf B$ and $s\colon y\to z$ of type~$\mathsf L$ be generating rays such that $rs$ is proper and shiftable. By~\eqref{eq:product-type-table}, $rs$ has type~$\mathsf L$, so by the preceding paragraph $rs=\widehat q_xf$ where $f\colon x\to z$ is the generating ray of $A_0(x,z)$ (of type~$\mathsf L$) and $N_x=3$. Suppose for a contradiction that $y\notin S_x$. Since $z\notin S_x$ (as $f$ is one-sided), Lemma~\ref{lem:external-hom-chain} compares $r$ and $f$:
\begin{itemize}
\item If $ra=f$ for a ray $a\colon y\to z$, then $rs=\widehat q_xra=r\widehat q_ya$ cancels to $s=\widehat q_ya$, contradicting the fact that $s$ is generating.
\item If $fa=r$ for a ray $a\colon z\to y$, then $f(as)=rs=\widehat q_xf$, so $as\in R_z$ has weight $\chi$ and equals $\widehat q_z$, making $f\widehat q_z=\widehat q_xf\ne0$ and contradicting the fact that $f$ has type~$\mathsf L$.
\end{itemize}
Thus $y\in S_x$, so $r$ admits a primitive reverse. The case $(\mathsf R,\mathsf B)$ follows by duality, proving Proposition~\ref{prop:proper-ray-products}(2).

Now consider Proposition~\ref{prop:proper-ray-products}(3): let $r\colon x\to y$ of type~$\mathsf L$ and $s\colon y\to z$ of type~$\mathsf R$ be generating rays whose product $rs$ is proper and shiftable. By~\eqref{eq:product-type-table}, $rs$ has type~$\mathsf B$, so $rs=\widehat q_x^jf=f\widehat q_z^j$ for the generating ray $f\colon x\to z$ and some layer $j\geq1$ with $\widehat q_x^{j+1}f=f\widehat q_z^{j+1}\ne0$. If $x=z$, then $f=1_x$ and Lemma~\ref{lem:nondeep-two-object-return} immediately gives $j=2$, $N_x=4$, and that $(r,s)$ is a nondeep external return at $x$, proving~(3)(a).

Assume henceforth that $x\ne z$. Then $f\colon x\to z$ is a generating ray of type~$\mathsf B$, while $y\notin S_x$ and $y\notin S_z$ because $r$ and $s$ are one-sided. We claim that $z\in S_x$ (so $S_x=S_z$). Suppose to the contrary that $z\notin S_x$. Then Lemma~\ref{lem:external-hom-chain} at $x$ compares $r\colon x\to y$ and $f\colon x\to z$. We cannot have $ra=f$ (as $rs=f\widehat q_z^j=ra\widehat q_z^j$ would cancel to $s=a\widehat q_z^j$, contradicting the fact that $s$ is generating), so $fa=r$ for a generating ray $a\colon z\to y$. Canceling $f$ from $f(as)=rs=f\widehat q_z^j$ gives $as=\widehat q_z^j$ with $\widehat q_z^{j+1}\ne0$, so Lemma~\ref{lem:nondeep-two-object-return} implies that $(a,s)$ is a nondeep external return at $z$ (with $j=2$). However, $f\colon x\to z$ is an incoming generating two-sided ray at $z$ with $x\notin S_z$ and $f\widehat q_z\ne0$, which contradicts Proposition~\ref{prop:proper-ray-products}(4)! Therefore $z\in S_x$.

\subsubsection{A complementary path forces an endpoint return}

To finish the proof of Proposition~\ref{prop:proper-ray-products}(3)(b), continue with $x\ne z$ and $z\in S_x$. Since $S_x$ is a primitive return corner, the generating ray $f\colon x\to z$ admits a primitive reverse $g\colon z\to x$ satisfying $fg=\widehat q_x$ and $gf=\widehat q_z$. Note that
\begin{equation}
  r(sg)=rsg=\widehat q_x^jfg=\widehat q_x^{j+1}\ne0,\qquad
  (gr)s=grs=gf\widehat q_z^j=\widehat q_z^{j+1}\ne0,
  \label{eq:complementary-closing-products}
\end{equation}
so both $sg\colon y\to x$ and $gr\colon z\to y$ are nonzero rays, as is $gfg=\widehat q_zg$ (since $\widehat q_z^2=gfgf\ne0$ from $j\geq1$ and $\widehat q_z^{j+1}\ne0$).

Set $F=\ell(f)$ and $E=\ell(r)$. In a finite homogeneous category $H$, retain the six graded objects
\[
  \begin{gathered}
  B_0=(z,F),\quad B_1=(x,\chi),\quad B_2=(z,F+\chi),\quad B_3=(x,2\chi),\\
  Y_-=(y,E+(1-j)\chi),\qquad Y_+=(y,E+\chi),
  \end{gathered}
\]
and select the oriented tree with spine $B_0\xrightarrow{g}B_1\xrightarrow{f}B_2\xrightarrow{g}B_3$ and branches $B_1\xrightarrow{r}Y_+$ and $Y_-\xrightarrow{s}B_2$. By~\eqref{eq:complementary-closing-products}, every oriented path in this six-vertex tree evaluates to a nonzero ray in $H$. After discarding negative-degree factors, checking the cover conditions of Lemma~\ref{lem:configuration-retraction} at the five arrows of this tree reduces to the following six possibilities for an unselected factor $h$:
\[
  \begin{array}{c|l}
  \text{candidate factor or weight}&\text{consequence}\\ \hline
  fh=r&hs=\widehat q_z^j\text{ (canceling }f\text{ from }fhs=f\widehat q_z^j\text{)}\\
  \ell(r)+\ell(h)=\ell(f)&\ell(h)=\ell(s)-j\chi<\ell(s)\text{ (impossible as }s\text{ is generating)}\\
  rh=\widehat q_x&y\in S_x\text{ by Lemma~\ref{lem:primitive-cyclic-corners}(3) (impossible)}\\
  hf=s&rh=\widehat q_x^j\text{ (canceling }f\text{ from }rhf=\widehat q_x^jf\text{)}\\
  \ell(h)+\ell(s)=\ell(f)&\ell(h)=\ell(r)-j\chi<\ell(r)\text{ (impossible as }r\text{ is generating)}\\
  hs=\widehat q_z&y\in S_z\text{ by Lemma~\ref{lem:primitive-cyclic-corners}(3) (impossible)}
  \end{array}
\]
The second, third, fifth, and sixth rows are impossible in all cases. In the first row, any ray $h\colon z\to y$ satisfying $fh=r$ is a factor of the generating shiftable ray $r$ and is therefore generating by Lemma~\ref{lem:generating-paths}, yielding a return $hs=\widehat q_z^j$ at $z$ with $\widehat q_z^{j+1}\ne0$. Similarly, in the fourth row, any ray $h\colon y\to x$ satisfying $hf=s$ is generating and yields a return $rh=\widehat q_x^j$ at $x$ with $\widehat q_x^{j+1}\ne0$.

If neither the first nor the fourth row occurs, then all six rows are excluded, so Lemma~\ref{lem:configuration-retraction} gives a bimodule retraction onto the incidence algebra of the six-vertex tree, contradicting Lemma~\ref{lem:dimension-infinity} ($N=D=12$). Thus at least one of the two factorizations $fh=r$ (with $hs=\widehat q_z^j$) or $hf=s$ (with $rh=\widehat q_x^j$) must hold. Applying Lemma~\ref{lem:nondeep-two-object-return} to the resulting return $(h,s)$ at $z$ or $(r,h)$ at $x$ shows that $j=2$ and that $(h,s)$ (respectively, $(r,h)$) is a nondeep external return at $z$ (respectively, at $x$). This completes the proof of Proposition~\ref{prop:proper-ray-products}(3)(b), and hence of Proposition~\ref{prop:proper-ray-products}.


\begin{thebibliography}{99}

\bibitem{ARS}
M.~Auslander, I.~Reiten, and S.~O.~Smal\o,
\emph{Representation theory of Artin algebras},
Cambridge Studies in Advanced Mathematics, vol.~36,
Cambridge University Press, Cambridge, 1995.

\bibitem{BGRS}
R.~Bautista, P.~Gabriel, A.~V.~Roiter, and L.~Salmer\'on,
\emph{Representation-finite algebras and multiplicative bases},
Invent. Math. \textbf{81} (1985), 217--285.

\bibitem{BautistaSalmeron}
R.~Bautista and L.~Salmer\'on,
\emph{On non-standard algebras and their representations},
Arch. Math. (Basel) \textbf{56} (1991), 16--34.

\bibitem{BES}
J.~Bia\l kowski, K.~Erdmann, and A.~Skowro\'nski,
\emph{Deformed preprojective algebras of generalized Dynkin type},
Trans. Amer. Math. Soc. \textbf{359} (2007), 2625--2650.

\bibitem{BongartzMild}
K.~Bongartz,
\emph{On mild contours in ray categories},
Algebr. Represent. Theory \textbf{17} (2014), 137--159.

\bibitem{BongartzSurvey}
K.~Bongartz,
\emph{On representation-finite algebras and beyond},
in \emph{Advances in representation theory of algebras},
EMS Series of Congress Reports, European Mathematical Society, Z\"urich,
2014, 65--101.

\bibitem{BongartzGabriel}
K.~Bongartz and P.~Gabriel,
\emph{Covering spaces in representation-theory},
Invent. Math. \textbf{65} (1982), 331--378.

\bibitem{BretscherGabriel}
O.~Bretscher and P.~Gabriel,
\emph{The standard form of a representation-finite algebra},
Bull. Soc. Math. France \textbf{111} (1983), 21--40.

\bibitem{Gerstenhaber}
M.~Gerstenhaber,
\emph{On the deformation of rings and algebras},
Ann. of Math. (2) \textbf{79} (1964), 59--103.

\bibitem{GerstenhaberSchack}
M.~Gerstenhaber and S.~D.~Schack,
\emph{Relative Hochschild cohomology, rigid algebras, and the Bockstein},
J. Pure Appl. Algebra \textbf{43} (1986), 53--74.

\bibitem{IyamaRealization}
O.~Iyama,
\emph{A relationship between homological properties and
representation-theoretic realization of artin algebras},
Trans. Amer. Math. Soc. \textbf{357} (2005), 709--734.

\bibitem{IyamaI}
O.~Iyama,
\emph{$\tau$-Categories I: Ladders},
Algebr. Represent. Theory \textbf{8} (2005), 297--321.

\bibitem{Jans}
J.~P.~Jans,
\emph{On the indecomposable representations of algebras},
Ann. of Math. (2) \textbf{66} (1957), 418--429.

\bibitem{Kaplan}
D.~Kaplan,
\emph{Multiplicative preprojective algebras of Dynkin quivers},
J. Pure Appl. Algebra \textbf{227} (2023), article 107146.

\bibitem{Kupisch}
H.~Kupisch,
\emph{Symmetrische Algebren mit endlich vielen unzerlegbaren
Darstellungen. I}, J. Reine Angew. Math. \textbf{219} (1965), 1--25;
\emph{II}, J. Reine Angew. Math. \textbf{245} (1970), 1--14.

\bibitem{Linckelmann}
M.~Linckelmann,
\emph{Integrable derivations and stable equivalences of Morita type},
Proc. Edinb. Math. Soc. (2) \textbf{61} (2018), 343--362.

\bibitem{Lowen}
W.~Lowen,
\emph{Hochschild cohomology, the characteristic morphism and derived deformations},
Compos. Math. \textbf{144} (2008), 1557--1580.

\bibitem{Riedtmann1980}
C.~Riedtmann,
\emph{Algebren, Darstellungsk\"ocher, \"Uberlagerungen und zur\"uck},
Comment. Math. Helv. \textbf{55} (1980), 199--224.

\bibitem{Riedtmann1984}
C.~Riedtmann,
\emph{On stable blocks of Auslander-algebras},
Trans. Amer. Math. Soc. \textbf{283} (1984), 485--505.

\bibitem{RingelVossieck}
C.~M.~Ringel,
\emph{Dieter Vossieck and the development of the representation theory of
Artin algebras}, lecture at the workshop \emph{Discrete Categories in
Representation Theory}, Bielefeld, 2018,
\href{https://arxiv.org/abs/1808.02301}{arXiv:1808.02301}.

\bibitem{Schedler}
T.~Schedler,
\emph{Zeroth Hochschild homology of preprojective algebras over the integers},
Adv. Math. \textbf{299} (2016), 451--542.

\end{thebibliography}
\end{document}
